\documentclass[11pt]{amsart}
\usepackage[margin=2.5cm]{geometry}
\usepackage[T1]{fontenc}
\usepackage{lmodern}
\usepackage{amsmath,amssymb,amsthm,mathtools}
\usepackage[shortlabels]{enumitem}
\usepackage{booktabs}
\usepackage[hidelinks]{hyperref}
\newcommand{\R}{\mathbb R}
\newcommand{\Z}{\mathbb Z}
\newcommand{\N}{\mathbb N}
\newcommand{\eps}{\varepsilon}
\newcommand{\cP}{\mathcal P}
\newcommand{\cD}{\mathcal D}
\newcommand{\cM}{\mathcal M}
\newcommand{\cE}{\mathcal E}
\newcommand{\cR}{\mathcal R}
\newcommand{\cZ}{\mathcal Z}
\newcommand{\cT}{\mathcal T}
\newcommand{\cX}{\mathcal X}
\newcommand{\relint}{\operatorname{relint}}
\newcommand{\Bot}{\operatorname{bot}}
\newcommand{\Top}{\operatorname{top}}
\newcommand{\dist}{\operatorname{dist}}
\newcommand{\Sh}{\operatorname{Sh}}
\newcommand{\Ov}{\operatorname{Ov}}
\newcommand{\Ovgap}{\operatorname{Ov}_{\mathrm{gap}}}
\newcommand{\Fgap}{F_{\mathrm{gap}}}
\newcommand{\Exc}{\operatorname{Exc}}
\newcommand{\invF}{\operatorname{inv}F}
\newcommand{\amax}{\alpha_{\max}}
\newcommand{\hmax}{h_{\max}}
\newcommand{\bbar}{\bar\beta}
\newcommand{\Lam}{\Lambda}
\newcommand{\Fz}{F^{0}}
\newcommand{\Fs}{F_{s}}
\newcommand{\Tstar}{T^{*}}
\newcommand{\frc}[1]{\{#1\}}
\newcommand{\dd}{\,\mathrm d}
\newtheorem{theorem}{Theorem}[section]
\newtheorem{proposition}[theorem]{Proposition}
\newtheorem{lemma}[theorem]{Lemma}
\newtheorem{corollary}[theorem]{Corollary}
\newtheorem{claim}[theorem]{Claim}
\theoremstyle{definition}
\newtheorem{definition}[theorem]{Definition}
\newtheorem{example}[theorem]{Example}
\newtheorem{notation}[theorem]{Notation}
\theoremstyle{remark}
\newtheorem{remark}[theorem]{Remark}
\numberwithin{equation}{section}

\title[Deficiency of order $k^{1/4}$]{Packing $k^2-c$ unit squares:\\ a lower bound of order $k^{1/4}$ for the deficiency}
\author{Sungjoon Ryu}
\date{Preprint, version 1.0, 7 October 2026}

\raggedbottom
\begin{document}
\begin{abstract}
Let $s(n)$ be the side of the smallest square into which $n$ non-overlapping unit squares can be packed, and let $M(k)$
be the largest number of unit squares that can be placed in $[0,k]^2$ so that they are pairwise disjoint as closed sets;
then $s(k^2-c)=k$ exactly when $c<k^2-M(k)$. Building on a previous paper [R], where $k^2-M(k)\ge0.0353\log k$ is shown
for every integer $k\ge2$, we prove that $k^2-M(k)\ge0.1\,k^{1/4}$ for every integer $k\ge4.62\cdot10^{12}$ and that
$\liminf_{k\to\infty}(k^2-M(k))/k^{1/4}\ge0.1696$. Consequently $s(k^2-c)=k$ for every integer
$k\ge\max(4.62\cdot10^{12},\,10^4c^4+1)$, so the largest $c$ with $s(k^2-c)=k$ grows at least like $k^{1/4}$; by recent
work on the wasted area it is $O(k^{3/5})$. The proof follows upward paths from the bottom side (and, by reflection,
downward paths from the top side) that adapt to the orientation of the squares they cross and enter squares only through
their bottom sides (top sides for the downward paths), organized over a continuum of scales, and it uses the lemmas of
[R] on the cost of inclinations. The constants $0.1$, $4.62\cdot10^{12}$ and $0.1696$ rely on the computer-assisted
Lemma~4.10 of [R], whose box computations have not yet been reproduced by an independent implementation. Without it,
using the analytic Lemma~4.9 of [R] instead, the same argument gives $k^2-M(k)\ge0.1\,k^{1/4}$ for every integer
$k\ge2.17\cdot10^{13}$ and the constant $0.1411$. This paper has not been peer reviewed.
\end{abstract}
\maketitle
\tableofcontents

\section{Introduction and main results}\label{sec:intro}

\newcounter{IntroCount}
\providecommand{\theHIntroCount}{\arabic{IntroCount}}
\newenvironment{IntroStatement}[2]{\par\addvspace{\medskipamount}%
  \renewcommand{\theIntroCount}{#2}\refstepcounter{IntroCount}%
  \noindent\textbf{#1~#2.}\ \itshape\ignorespaces}%
 {\par\addvspace{\medskipamount}}

For $n\ge1$ let $s(n)$ be the side of the smallest square that contains $n$ unit squares with pairwise disjoint
interiors (rotations allowed); trivially $s(k^2)=k$, and $s(n)\le\lceil\sqrt n\,\rceil$. For $n=k^2-c$ with $c$ small
compared with $k$ it is natural to expect that one cannot do better than the $k\times k$ grid with $c$ squares removed.
Friedman's survey \cite{Fri} contains the conjecture that $s(n^2-c)=n$ implies $s((n+1)^2-c)=n+1$, and Chung and Graham
\cite[concluding remarks]{CG09} wrote that the same should hold for $n^2-c$ unit squares when $c$ is fixed and $n$ is
sufficiently large. Following the notes \cite[file \texttt{s12/search/FRIEDMAN.md}]{Dan}, put
\[
c^*(k):=\max\{c\in\Z_{\ge0}:\ s(k^2-c)=k\}.
\]
Since $s$ is non-decreasing, $s(k^2-c)=k$ holds exactly for $0\le c\le c^*(k)$.

The classical lower bound for the wasted area is that of Roth and Vaughan \cite{RV}: their theorem
\cite[p.~170, Theorem]{RV} states that the wasted area $w(\alpha)$ in a square of side $\alpha$ satisfies
$w(\alpha)\gg(\alpha\|\alpha\|)^{1/2}$ whenever $\alpha(\alpha-[\alpha])>1/6$, where $\|\alpha\|$ is the distance from
$\alpha$ to the nearest integer; this vanishes as $\alpha\to k^-$, and it gives no information on $c^*(k)$. In the
previous paper [R] it is shown that $c^*(k)\to\infty$, with a logarithmic rate (see below). In the other direction,
Erd\H{o}s and Graham \cite{EG} packed a square of side $x$ with wasted area $O(x^{7/11})$, and this exponent was lowered
several times (see the introductions of \cite{McC,Bui}); the exponent $3/5$ was claimed by Chung and Graham \cite{CG}, an
error in their argument was pointed out by Arslanov and Bui \cite{AB}, and the bound $O(x^{3/5})$ was re-established,
independently and by different methods, in the preprints of Bui \cite{Bui} and of McClenagan \cite{McC}. Applied to $x$
slightly below $k$, this gives $s(k^2-Ck^{3/5})<k$ for large $k$, that is, $c^*(k)=O(k^{3/5})$ (the refereed bound
$O(x^{(3+\sqrt2)/7}\log x)$ of \cite{CG09} gives $c^*(k)=O(k^{(3+\sqrt2)/7}\log k)$). In this paper we show that $c^*(k)$
grows at least like $k^{1/4}$.

Following [R], a \emph{square} is a closed unit square, and a
\emph{closed packing} of $[0,k]^2$ is a finite family of squares contained in $[0,k]^2$ that are pairwise disjoint as
closed sets. For a closed packing with $N$ squares, $W:=k^2-N$ is the area of $[0,k]^2$ not covered by the squares; it
is an integer. Let $M(k)$ be the largest size of a closed packing of $[0,k]^2$ and $W_{\min}(k):=k^2-M(k)$ the least
value of $W$.
By [R, Lemma~\ref{imp:L2.1}] and its proof, for integers $k\ge2$ and $0\le c<k^2$ one has $s(k^2-c)=k$ if and only if
$k^2-c>M(k)$, that is, if and only if $c<W_{\min}(k)$; hence $c^*(k)=W_{\min}(k)-1$. Lower bounds for $W$ that hold for
all closed packings of $[0,k]^2$ therefore give $s(k^2-c)=k$ for all $c$ below them.

In [R] it is shown that $W\ge\kappa\log k$ for every integer $k\ge2$ and every closed packing, with $\kappa=0.0353$
($\kappa=0.0319$ without the computer-assisted [R, Lemma~\ref{imp:L4.10}]), and that $W\ge1.99954\,(\log k-13.06675)/30.418$;
in particular $c^*(k)\to\infty$. In this paper we prove a bound of order $k^{1/4}$ for large $k$. Throughout, $\log$ is
the natural logarithm and $\delta:=10^{-5}$.

\begin{IntroStatement}{Theorem}{A}\label{thm:A}
Let $k\ge4.62\cdot10^{12}$ be an integer. Then every closed packing of $[0,k]^2$ by $N$ unit squares satisfies
\[
W=k^2-N\ \ge\ 0.1\,k^{1/4};
\]
in particular $W\ge\lceil0.1\,k^{1/4}\rceil\ge147$. Moreover
\[
\liminf_{k\to\infty}\frac{W_{\min}(k)}{k^{1/4}}\ \ge\ c_\infty:=\frac{16\cdot2^{-1/4}\,(1-\lambda)^{5/4}}{5\cdot5.0001^{1/4}\sqrt{AA'}}
=0.1696524\ldots,
\]
where $A:=4\sqrt{9/0.32}/(2-3\cdot10^{-6})$, $A':=A/\cos(1.5\cdot10^{-6})+10^{-8}$ and $\lambda:=2\delta+2\cdot10^{-12}$.
\end{IntroStatement}

The digits of $c_\infty$ shown are those of its decimal expansion, i.e.\ the value is rounded down; $A$ is the constant
of [R, Lemma~\ref{imp:L4.13}] ($A=4\sqrt{K_w^*/Q_*}/(2-3\cdot10^{-6})$ with $K_w^*=9$, $Q_*=0.32$).

The threshold $4.62\cdot10^{12}$ has little margin. With the parameters used in the proof, the last inequality of the
argument (condition~(i) of Proposition~\ref{prop:monotone}) holds exactly for $k\ge k_*$, where $k_*$ is about
$4.6167\cdot10^{12}$, and $4.62\cdot10^{12}$ is $k_*$ rounded up to three significant digits
($4.62\cdot10^{12}/k_*-1<7.3\cdot10^{-4}$; Section~\ref{subsec:const-sens}). If one of the constants entering the proof were
larger than the tolerance stated in Section~\ref{subsec:const-sens} (for instance $A'\le10.60905$ or
$C_\Lambda\le1.0150$), this inequality would fail at $4.62\cdot10^{12}$ and the threshold would have to be raised.

\begin{IntroStatement}{Theorem}{A$_{13}$}\label{thm:A13}
Let $k\ge2.17\cdot10^{13}$ be an integer. Then every closed packing of $[0,k]^2$ satisfies $W\ge0.1\,k^{1/4}$, in
particular $W\ge216$. Moreover
\[
\liminf_{k\to\infty}\frac{W_{\min}(k)}{k^{1/4}}\ \ge\ c^{(13)}_\infty:=\frac{16\cdot2^{-1/4}\,(1-\lambda)^{5/4}}{5\cdot5.0001^{1/4}\sqrt{A_{13}A'_{13}}}
=0.1411593\ldots,
\]
where $A_{13}:=4\sqrt{13/0.32}/(2-3\cdot10^{-6})$, $A'_{13}:=A_{13}/\cos(1.5\cdot10^{-6})+10^{-8}$, and $\lambda$ is as in
Theorem~\ref{thm:A}.
\end{IntroStatement}

The proof of Theorem~\ref{thm:A13} does not use [R, Lemma~\ref{imp:L4.10}]; the value of $c^{(13)}_\infty$ is again rounded
down. Its threshold has little margin as well: it is about $2.1696\cdot10^{13}$ rounded up (Section~\ref{subsec:const-sens}).

\begin{IntroStatement}{Corollary}{A}\label{cor:A}
Let $c\ge0$ be an integer. Then $s(k^2-c)=k$ for every integer $k\ge\max(4.62\cdot10^{12},\,10^4c^4+1)$. In particular
$c^*(k)\ge\lceil0.1\,k^{1/4}\rceil-1$ for every integer $k\ge4.62\cdot10^{12}$, and $s(k^2-c)=k$ for every integer
$k\ge4.62\cdot10^{12}$ if $c\le146$. Without [R, Lemma~\ref{imp:L4.10}] (by Theorem~\ref{thm:A13}), the same statements
hold with $2.17\cdot10^{13}$ in place of $4.62\cdot10^{12}$ and, in the last one, with $c\le215$ in place of $c\le146$.
\end{IntroStatement}

Together with the upper bound recalled above, $c^*(k)$ thus lies between $\lceil0.1\,k^{1/4}\rceil-1$ and $O(k^{3/5})$
for $k\ge4.62\cdot10^{12}$ (the upper bound relying on the preprints \cite{Bui,McC}).

The proofs of Theorems~\ref{thm:A}, \ref{thm:A13} and Corollary~\ref{cor:A} are given in Section~\ref{sec:assembly};
the numerical constants and thresholds are collected in Section~\ref{sec:constants}, where the direction of rounding of
every displayed value is stated. We do not prove bounds of the form $W\ge c_*k^{1/4}$ for all $k<k_0$, where
$k_0=4.62\cdot10^{12}$ in Theorem~\ref{thm:A} and $k_0=2.17\cdot10^{13}$ in Theorem~\ref{thm:A13}; Section~\ref{sec:constants}
gives lower bounds for $W$ at single values of $k$ (Table~\ref{tab:const-k}).

\emph{Dependence on [R].} The proofs use the results of [R] listed in Section~\ref{sec:imported}, where their statements
are quoted with the numbering of [R]; their proofs are in [R]. A few facts about these proofs are also used; they are
listed in Remark~\ref{rem:imp-proofs}. Among the results used, [R, Lemma~\ref{imp:L4.10}] is computer-assisted: its proof
is a branch-and-bound computation over boxes, whose verdicts were replayed with the same programs \cite{LevR} but have
not yet been reproduced by an independent implementation. It enters only through the multiplicity bound $K_w^*=9$,
and through it the constants $A$, $A'$ and $C_\Lambda$ (Definition~\ref{def:flows-consts}). Thus Theorem~\ref{thm:A}, the
value of $c_\infty$, the threshold $4.62\cdot10^{12}$ and the form of Corollary~\ref{cor:A} with this threshold depend on
it. Theorem~\ref{thm:A13} and the form of Corollary~\ref{cor:A} with the threshold $2.17\cdot10^{13}$ do not: they use
instead the analytic [R, Lemma~\ref{imp:L4.9}] ($K_w=13$) everywhere (Remark~\ref{rem:imp-dep}), whose proof contains only
a short finite check of angle conditions, carried out in [R] by a small program.

\emph{Comparison with [R].} The bound of [R] holds for every $k\ge2$ and grows like $\log k$, with asymptotic slope
$1.99954/30.418>0.0657$; Theorems~\ref{thm:A} and~\ref{thm:A13} give the order $k^{1/4}$, but only for $k\ge k_0$. For a
fixed integer $c\ge0$, Corollary~\ref{cor:A} gives $s(k^2-c)=k$ as soon as $k\ge\max(4.62\cdot10^{12},10^4c^4+1)$ (note
$\log(4.62\cdot10^{12})<29.17$), whereas for $c\ge1$ the bound of [R] gives it for $\log k>15.22\,c+13.07$
[R, Remark~7.4]. Compared with [R], the quantization in columns of [R, Section~5] is replaced by
Sections~\ref{sec:P12}--\ref{sec:P4}, and [R, Lemma~\ref{imp:L4.14}] by Section~\ref{sec:P6}. The lemmas of
[R, Section~4] quoted in Section~\ref{sec:imported} are used as stated there, with two exceptions explained in that
section: the multiplicity bounds of [R, Lemmas~\ref{imp:L4.9} and~\ref{imp:L4.10}] are applied to families of pairs that
need not be edges of the visibility graphs of [R], as their proofs allow (Proposition~\ref{imp:pairs}), and for
Theorem~\ref{thm:A13} [R, Lemmas~\ref{imp:L4.13}--\ref{imp:L4.15}] are used with $A_{13}$ in place of $A$
(Remark~\ref{rem:imp-dep}).

\emph{Idea of the proof.} As in [R], we consider horizontal lines $y$ with $y_0\le d(y)\le(1-\varepsilon)(\frac k2-3)$,
where $d(y)=\min(y,k-y)$, whose height is not close to an integer (the set $H$ of Definition~\ref{def:flows-lines}), and
we integrate over them with the measure $d(y)^{-3/4}dy$. On such a line with little waste, either the line meets a square
of inclination at least $\beta(y)=c_1d(y)^{-3/4}$, and such lines are paid for by [R, Lemma~\ref{imp:L4.13}]; or all
squares on the line are nearly axis-parallel, and then the short chords cut from the ramps of squares near their bottom
and top add up to almost a full unit ([R, Lemmas~\ref{imp:L3.4}, \ref{imp:L3.5} and~\ref{imp:L4.15}]). The squares that
carry these short chords form the sets $\mathcal Z(s)$ of Definition~\ref{def:flows-lines}. In [R] such squares are
controlled through vertical columns below them; here we use paths. A path starts at a point $(x,0)$ of the bottom side of
$[0,k]^2$ and moves upwards, following the orientation of the squares it meets: in a gap it keeps the direction $n_X$ of
the last square $X$ it passed, it enters a square only through the relative interior of the bottom side, and it crosses
the square parallel to its sides (Section~\ref{sec:flows}). A path dies when its total length in gaps reaches $\delta$;
of several paths that arrive at the same entry point only one continues (merge); a path that meets a square elsewhere
than on its bottom side stops (end collision); and in the flow of scale $s$ a path stops when it enters a square of
inclination at least $\alpha(s)=c_0s^{-1/2}$ (tilt). The height identity of Section~\ref{sec:P12} shows that a path that has
crossed $m$ squares of inclination less than $\alpha(s)$ with total gap less than $\delta$ is within
$0.50001\,\alpha(s)^2q_y+\delta\approx0.50001\,c_0^2+\delta$ of the integer $m$ when it reaches a bottom side at height
$q_y\approx s$; hence the ramps of a square whose bottom side is reached alive lie close to integers, and no path of
the flow of scale $s$ reaches the bottom side of a square of $\mathcal Z(s)$ alive. The shadow inequality of Section~\ref{sec:P4} turns this into a lower bound
for the losses: for each scale $s$, the measure $N(s)$ of the set of tilt-terminated starting points is at least
$|\mathcal Z(s)|/(\varepsilon s+b_W)$ minus the other losses (Section~\ref{sec:P7}). Deaths, merges, end collisions and
overlaps of paths in gaps cost $O(W/\delta)$ in total (Sections~\ref{sec:P3} and~\ref{sec:P5}), and exits through the
side walls cost $O(\sqrt s)$ per scale. Tilt terminations are paid for by a version of [R, Lemma~\ref{imp:L4.14}] for
paths (Section~\ref{sec:P6}): $\int\alpha(\Tstar(x))\,dx+\int\alpha(\tilde T^*(x))\,dx\le A'W$, where $\Tstar(x)$ is the
least scale at which the path from $(x,0)$ is tilt-terminated, and $\tilde T^*$ is the same for paths from the top side.
Integrating over the scales with the weight $|\alpha'(s)|$ and over the lines with the measure $d(y)^{-3/4}dy$
(Section~\ref{sec:P8}) gives a left-hand side of order $k^{1/4}$ and a right-hand side $BW+O(\log k)$. The tilt
threshold $c_0s^{-1/2}$ is chosen so that the accumulated loss of height
$\sum_i(1-\cos\varphi_i)\le\alpha(s)^2q_y/(2\cos\alpha(s))$ over about $s$ crossed squares stays of order $c_0^2/2$. This is possible because paths enter squares only through
bottom sides: a contact with another side of a square ends the path once, as an end collision, and the end collisions cost
$O(\delta W)$ in total (Section~\ref{sec:P3}). (The columns of [R] use the threshold $1/(2s)$.)

\emph{Organisation.} Section~\ref{sec:imported} collects the definitions and results of [R] that are used.
Section~\ref{sec:flows} defines the parameters, the flows of paths and the quantities attached to them, and proves their
basic properties. Section~\ref{sec:P12} proves the height identity and the quantization of the bottom-side entries.
Section~\ref{sec:P3} bounds the end collisions. Section~\ref{sec:P4} proves the shadow inequality, with its exceptional
heights, and the bound for the wall term. Section~\ref{sec:P5} bounds deaths, merges and the overlap of paths in gaps.
Section~\ref{sec:P6} proves the path version of [R, Lemma~\ref{imp:L4.14}]. Section~\ref{sec:P7} proves the per-scale
count and the lower bound for $|\mathcal Z(s)|$. Section~\ref{sec:P8} proves the integration inequality $(\bigstar)$
(Theorem~\ref{thm:star}).
Section~\ref{sec:assembly} deduces Theorems~\ref{thm:A}, \ref{thm:A13} and Corollary~\ref{cor:A}, and
Section~\ref{sec:constants} collects the constants and thresholds.

\emph{Status.} This paper has not been peer reviewed, and neither has [R]. How the arguments have been checked so far,
by AI-based reviews and numerical tests only, is described at the end of the paper (Verification status).

\section{Results used from [R]}\label{sec:imported}

\newcounter{ImpCount}
\providecommand{\theHImpCount}{\arabic{ImpCount}}
\newenvironment{ImpStatement}[2]{\par\addvspace{\medskipamount}%
  \renewcommand{\theImpCount}{#2}\refstepcounter{ImpCount}%
  \noindent\textbf{#1~#2 of [R].}\ \itshape\ignorespaces}%
 {\par\addvspace{\medskipamount}}

In this section we collect the definitions and results of
[R] (S.~Ryu, \emph{Packing $k^2-c$ unit squares: $s(k^2-c)=k$ for all large $k$}, Preprint, version~1.2, 2026)
that are used in this paper. Each statement is quoted verbatim and carries the number it has in [R]; only
cross-references are adapted (a reference such as \ref{imp:L4.13} points to the quotation in this section).
Letters that are bound inside a quoted statement (for instance $\alpha$, $\beta$, $\delta$, $\lambda$, $\varphi$,
$\Lambda$, $X$, $Y$, $I$, $J$) have the meaning given there; they are unrelated to the objects with the same
names in Sections~\ref{sec:flows}--\ref{sec:constants} unless this is said explicitly. Proofs are in [R].
Throughout, $k\ge4$ is an integer.

\subsection{Packings and the reduction}

The following definitions are those of [R, Sections~1 and~2].
For $n\ge 1$ let $s(n)$ denote the side of the smallest square that contains $n$ unit squares with pairwise
disjoint interiors (rotations allowed), and
\[
c^*(k):=\max\{c\in\Z_{\ge0}:\ s(k^2-c)=k\}.
\]
Throughout, a \emph{square} is a closed unit square. A \emph{closed packing} of $[0,k]^2$ is a finite family of squares
contained in $[0,k]^2$ that are pairwise disjoint as closed sets. Put $M(k)$ for the maximum size of a closed packing
of $[0,k]^2$ and, for a given closed packing, $W:=k^2-N$, the area of $[0,k]^2$ not covered by the squares.

\begin{ImpStatement}{Lemma}{2.1}\label{imp:L2.1}
For integers $k\ge 2$ and $n\ge 1$: $s(n)<k$ if and only if $n\le M(k)$. Moreover $M(k)\le k^2-1$
\textup{(}Corollary~\ref{imp:C3.3} below\textup{)}, so $c^*(k)=k^2-M(k)-1$.
\end{ImpStatement}

The last sentence of the proof of [R, Lemma~\ref{imp:L2.1}] reads: ``Finally, $s(n)\le k$ always holds for
$n\le k^2$, so for integers $0\le c<k^2$ we have $s(k^2-c)=k$ iff $k^2-c>M(k)$.'' We use the lemma in this
form. Note that $W$ is an integer.

\subsection{Elementary geometry}

The following definitions are those of [R, Section~3].
The \emph{inclination} $a(S)\in[0,\pi/4]$ of a square $S$ is the least angle between a side of $S$ and a side of
$[0,k]^2$; the inclination $\theta(S_1,S_2)$ between two squares is the least angle between their sides
(this is the quotient metric on $\R/(\pi/2)\Z$, with values in $[0,\pi/4]$), and we write $S_0:=\partial[0,k]^2$ with
$\theta(S,S_0):=a(S)$. Likewise, for a single side $\sigma$ of $[0,k]^2$ we put
$\theta(S,\sigma)=\theta(\sigma,S):=a(S)$. (These are the definitions of Roth and Vaughan.)

\begin{ImpStatement}{Lemma}{3.1}\label{imp:L3.1}
Let $S$ be a square with inclination $a$ and lowest point at height $v$.
\begin{enumerate}[label=(\alph*)]
\item The vertical extent of $S$ is $[v,v+\cos a+\sin a]$.
\item For $y\in\mathcal M(S):=[v+\sin a,\ v+\cos a]$ (the \emph{middle band}, of length $\cos a-\sin a$) the line
$y=\mathrm{const}$ meets $S$ in a segment of length $\sec a\ge 1$; for $y$ in the \emph{ramp set}
$\mathcal R(S):=[v,v+\sin a)\cup(v+\cos a,\ v+\cos a+\sin a]$ the length is $<\sec a$. If $a=0$, then $\mathcal R(S)=\emptyset$.
\item Symmetrically, for $x$ in the \emph{vertical middle band} $\mathcal M'(S)$, an interval of length
$\cos a-\sin a$ \textup{(}defined in the proof; for $a=0$ it is the horizontal extent of $S$\textup{)}, the vertical line through $x$ meets $S$ in a segment of length $\sec a$, entering $S$ through a
side of slope $\pm\tan a$; the lower endpoint of that segment has height in $[v,v+\sin a]$. The complement of
$\mathcal M'(S)$ in the horizontal extent of $S$ (the \emph{vertical ramp set}) has total length $2\sin a$.
\end{enumerate}
\end{ImpStatement}

(In the proof of [R, Lemma~\ref{imp:L3.1}], for $0<a\le\pi/4$ and up to a reflection, the vertices of $S$ are
$P_0=(x_0,v)$, $P_1=P_0+(\cos a,\sin a)$, $P_3=P_0+(-\sin a,\cos a)$, $P_2=P_1+P_3-P_0$, and
$\mathcal M'(S):=[x_0,x_0+\cos a-\sin a]$. Part~(c) is not used in this paper.)

For $y\in[0,k]$ let $\omega(y)$ be the length of the part of the segment $[0,k]\times\{y\}$ not covered by squares
(the \emph{line waste}); by Fubini $\int_0^k\omega(y)\,dy=W$. Similarly, for a vertical segment $\sigma$ we call the
length of $\sigma$ not covered by squares its \emph{vertical waste}; integrating the vertical waste of
$\{x\}\times[0,k]$ over $x$ also gives $W$. Since a packing is finite, all functions of $y$ (or of $x$) considered
in [R] are piecewise continuous with finitely many pieces, and all sets of $y$ (or of $x$) defined there are finite
unions of intervals.

\begin{ImpStatement}{Lemma}{3.2}\label{imp:L3.2}
Let $k\ge 2$. On every horizontal line, the squares of a closed packing of $[0,k]^2$ meet the segment
$[0,k]\times\{y\}$ in pairwise disjoint closed intervals, and at most $k-1$ of these have length $\ge1$.
\end{ImpStatement}

\begin{ImpStatement}{Corollary}{3.3}\label{imp:C3.3}
For $k\ge2$ every closed packing of $[0,k]^2$ has $W\ge1$; that is, $M(k)\le k^2-1$.
\end{ImpStatement}

For a square $S$ and $y\in\R$ let $c_S(y)$ be the length of the horizontal chord $S\cap(\R\times\{y\})$ (zero if the line
misses $S$).

\begin{ImpStatement}{Lemma}{3.4}\label{imp:L3.4}
Let $k\ge 2$, $y\in(0,k)$ and $0<\alpha\le\frac\pi4$, and suppose that every square meeting the line $y$ has inclination
$<\alpha$. Then
\[
\sum_{S:\ 0<c_S(y)<1}c_S(y)\ >\ 1-\omega(y)-(k-1)(\sec\alpha-1),
\]
the sum being over the squares $S$ of the packing. Moreover, for a square $S$ let $\Phi(S):=\{y:0<c_S(y)<1\}$. If
$a:=a(S)>0$, then $\Phi(S)\subset\mathcal R(S)$, $|\Phi(S)|=2\sin a\cos a$ and $\int_{\Phi(S)}c_S(y)\,dy=\sin a\cos a$; if $a(S)=0$, then
$\Phi(S)=\emptyset$.
\end{ImpStatement}

Only the second part (``Moreover \dots'') of [R, Lemma~\ref{imp:L3.4}] is used in this paper.

\begin{ImpStatement}{Lemma}{3.5}\label{imp:L3.5}
Let $k\ge2$ and $y\in(0,k)$, and put
\[
E(y):=\sum_{S:\ c_S(y)\ge1}\big(c_S(y)-1\big),
\]
the sum being over the squares $S$ of the packing. Then
\[
\sum_{S:\ 0<c_S(y)<1}c_S(y)\ \ge\ 1-\omega(y)-E(y).
\]
\end{ImpStatement}

The proof of [R, Lemma~\ref{imp:L3.5}] uses [R, Lemma~\ref{imp:L3.2}].

\subsection{The sharpened fundamental lemma and the cost of inclinations}

The following conventions are those of [R, Section~4]:
throughout [R, Section~4], $k\ge4$ (this is used, for instance, in [R, Lemma~\ref{imp:L4.8}(b)] and in the proof of
[R, Lemma~\ref{imp:L4.10}]; in this paper $k\ge4$ always holds), a closed packing of $[0,k]^2$ is fixed,
$S_0:=\partial[0,k]^2$, and the \emph{waste} is the set of points of $(0,k)^2$ that lie in no square; its area is $W$.

\begin{ImpStatement}{Lemma}{4.1}\label{imp:L4.1}
Let $I$ be an interval of length $\lambda$, let $0\le\varphi\le1$, and let $g\ge0$ be a continuous, convex, piecewise linear
function on $I$ with finitely many break points, such that $|g'|\ge\varphi$ off the break points and the slope of $g$
increases by at least $2$ at every break point. Then $\int_Ig\ge\varphi(2-\varphi)\lambda^2/4$.
\end{ImpStatement}

Let $(s,t)$ be Cartesian coordinates in the plane. For a square $Z$ let $[\tau^-(Z),\tau^+(Z)]$ be its projection to the
$t$-axis, and for $\tau^-(Z)\le t\le\tau^+(Z)$ let $l_Z(t)$ and $r_Z(t)$ be the least and the largest $s$ with $(s,t)\in Z$.
The vertices of $Z$ at which $t=\tau^+(Z)$ or $t=\tau^-(Z)$ are its \emph{$t$-extreme} vertices.

\begin{ImpStatement}{Lemma}{4.2}\label{imp:L4.2}
The function $l_Z$ is convex and $r_Z$ is concave. Each of them is piecewise linear with at most one break point, at which
the slope of $l_Z$ increases, and that of $r_Z$ decreases, by at least $2$. Off the break points $l_Z'(t)=\tan\psi$, where
$\psi\in(-\frac\pi2,\frac\pi2)$ is the angle between the $t$-axis and the side of $Z$ containing $(l_Z(t),t)$, and similarly
for $r_Z$. Moreover $\tau^+(Z)-\tau^-(Z)\ge1$.
\end{ImpStatement}

\emph{The rectangles $R(S_1,S_2)$.}
For a square $S_1$ and $S_2$ either another square or $S_0$, let
$L$ be the segment from the centre $c_1$ of $S_1$ to the centre $c_2$ of $S_2$, respectively to a point $c_2$ of $S_0$ nearest to
$c_1$ (any choice will do), and let $R=R(S_1,S_2)$ be the closed rectangle of points at distance at most $\frac12$ from the line through $L$ whose
orthogonal projection lies on $L$. This is the rectangle used by Roth and Vaughan. Let $u$ be the unit vector along $L$
and $n$ a unit normal, and use the coordinates $(s,t)$ of $c_1+su+tn$; thus $R=\{0\le s\le|L|,\ |t|\le\frac12\}$. (Here, and in the quoted
statements below, $s$ and $t$ are these coordinates; they are unrelated to the scale $s$ and the ray parameter $t$ of
Section~\ref{sec:flows}.)

\begin{ImpStatement}{Proposition}{4.3}\label{imp:P4.3}
Let $\theta:=\theta(S_1,S_2)$, and let $m'$ be the number of squares $Z\notin\{S_1,S_2\}$ of the packing that have a
$t$-extreme vertex $v\in R$ with $|t(v)|<\frac12$. Let $\mathcal T\subset(-\frac12,\frac12)$ be the set of the numbers $t(v)$
for these vertices, and let $\lambda_1,\dots,\lambda_q$ be the lengths of the intervals into which the points of $\mathcal T$
cut $(-\frac12,\frac12)$. Then the waste in $R$ has area at least
$\frac{\theta(2-\theta)}4\sum_r\lambda_r^2\ge\theta(2-\theta)/(4(m'+1))$.
\end{ImpStatement}

\begin{ImpStatement}{Lemma}{4.4}\label{imp:L4.4}
Let $Z$ be a square, $e$ a unit vector, and $v$ a vertex of $Z$ at which $p\mapsto p\cdot e$ is smallest on $Z$. For
$0\le\nu\le\frac34$ the line $\{p:p\cdot e=v\cdot e+\nu\}$ meets $Z$ in a segment of length at least $\min(2\nu,1)$.
\end{ImpStatement}

[R, Lemma~\ref{imp:L4.4}] is used only inside the proof of [R, Proposition~\ref{imp:P4.5}].

\begin{ImpStatement}{Proposition}{4.5}\label{imp:P4.5}
Let $S_1$ be a square and $S_2$ a square or $S_0$, at distance less than $d$, with $\theta:=\theta(S_1,S_2)\le\theta_1$, where
$d,\theta_1\le10^{-5}$. Then $\sum_r\lambda_r^2\ge Q_*:=0.32$ in Proposition~\ref{imp:P4.3}, so the waste in $R(S_1,S_2)$ has
area at least $Q_*\theta(2-\theta)/4$.
\end{ImpStatement}

\emph{Visibility graphs.}
Let $\mathcal V_h$ consist of the squares of the packing together with the left and the right side of $[0,k]^2$. Two
distinct members $P,Q$ of $\mathcal V_h$ are \emph{consecutive on the line $y$} ($0<y<k$) if both meet the segment
$[0,k]\times\{y\}$ and the open segment between them on that line meets no square. Let $G_h$ be the graph on
$\mathcal V_h$ in which $P$ and $Q$ are adjacent if they are consecutive on a set of lines $y$ of positive measure.
Define $\mathcal V_v$ and $G_v$ in the same way, with vertical lines and the bottom and the top side of $[0,k]^2$.
For $0<d\le10^{-4}$ put $\Lambda=\Lambda(d):=d+\sqrt2$ and $R_0=R_0(d):=\sqrt{\Lambda^2+\frac14}<1.5001$, and
\[
K:=27,\qquad q_0:=10 .
\]
(The symbols $\Lambda(d)$, $R_0(d)$, $K=27$ and $q_0$ are used only in this section, as in [R]; they are unrelated to the
loss terms $\Lambda_0$, $\Lambda^*(s)$ of Section~\ref{sec:flows} and to the multiplicity bound denoted $K$ in
Sections~\ref{sec:P3} and~\ref{sec:P5}.)

\begin{ImpStatement}{Lemma}{4.8}\label{imp:L4.8}
Let $0<d\le10^{-4}$.
\begin{enumerate}[label=(\alph*)]
\item If $S_1$ is a square and $S_2$ is a square or $S_0$, at distance $<d$, then $m'\le q_0-1=9$ in
Proposition~\ref{imp:P4.3}; hence the waste in $R(S_1,S_2)$ has area at least $\theta(2-\theta)/(4q_0)$, $\theta=\theta(S_1,S_2)$.
\item Let $k\ge4$ and let $G$ be $G_h$ or $G_v$. For an edge $e=\{P,Q\}$ of $G$ whose members are at distance $<d$ put $R_e:=R(P,Q)$ if
both are squares, and $R_e:=R(P,S_0)$ if $P$ is a square and $Q$ a side. Then every point lies in $R_e$ for at most $K$
such edges $e$ \textup{(}the two sides of $[0,k]^2$ that are vertices of $G$ are at distance $k>d$\textup{)}.
\end{enumerate}
\end{ImpStatement}

Part~(a) is used in this paper; part~(b) only supplies the definition of $R_e$. For a pair $e$ we write $f_e$ for the
area of the waste in the rectangle $R_e$ attached to $e$ (as in the proof of [R, Lemma~\ref{imp:L4.13}]).

\begin{ImpStatement}{Lemma}{4.9}\label{imp:L4.9}
Let $0<d\le10^{-4}$, let $G$ be $G_h$ or $G_v$, and define $R_e$ for the edges $e$ of $G$ whose members are at distance $<d$ as
in Lemma~\ref{imp:L4.8}(b). Then every point of the waste lies in $R_e$ for at most $K_w:=13$ such edges.
\end{ImpStatement}

The paragraph of [R] that follows the proof of [R, Lemma~\ref{imp:L4.9}] reads:
``The proof of Lemma~4.9 uses no property of $G$ except that the only sides involved are two opposite ones; the bound
holds for any family of such pairs at distance $<d$. A computer-assisted refinement of the same counting gives the bound $9$, which
is sharp for the number of all pairs at distance $<d$ (see the remark after the proof), though not known to be sharp for the edges
of a single graph $G$.''

\begin{ImpStatement}{Lemma}{4.10}\label{imp:L4.10}
Under the assumptions of Lemma~\ref{imp:L4.9}, every point of the waste lies in $R_e$ for at most $K_w^*:=9$ of the edges $e$
considered there.
\end{ImpStatement}

\begin{ImpStatement}{Remark}{4.11}\label{imp:R4.11}
The proof bounds the number of \emph{all} pairs of squares (or of a square and one of two opposite sides) at distance less than $d$ whose
rectangles contain $p$, without using the graph $G$, and for this number the bound $9$ is sharp [\dots]
\end{ImpStatement}

(We quote the first sentence of [R, Remark~\ref{imp:R4.11}]; the rest of the remark describes a configuration
showing the sharpness and is not used. In the remark, $p$ is the point of the waste considered in the proof of
[R, Lemma~\ref{imp:L4.10}].)

\begin{ImpStatement}{Lemma}{4.12}\label{imp:L4.12}
Let $A,B$ be distinct squares, or let one of them be the left or the right side of $[0,k]^2$, and let $E_{AB}$ be the set
of $y\in[0,k]$ such that on the line $y$ both $A$ and $B$ are met, $A$ lies to the left of $B$, and no square is met
strictly between them. For $y\in E_{AB}$ let $g(y)$ be the length of the gap between them, and let $\theta:=\theta(A,B)$.
\begin{enumerate}[label=(\alph*)]
\item All $y\in E_{AB}$ with $g(y)<1$ lie in one component of $E_{AB}$.
\item On every component $J$ of $E_{AB}$ the function $g$ is convex, and $\int_Ig\ge\theta(2-\theta)|I|^2/4$ for every interval
$I\subset J$. In particular, for $\eta>0$ the set $\{y\in J:g(y)<\eta\}$ is an interval.
\end{enumerate}
The same holds with the roles of the axes exchanged (vertical lines, $A$ below $B$, the bottom and the top side).
\end{ImpStatement}

\begin{ImpStatement}{Lemma}{4.13}\label{imp:L4.13}
Let $0<\bar\beta\le1.5\cdot10^{-6}$, and let $Y\subset[0,k]$ and $\beta:Y\to[0,\bar\beta]$ be such that for every $y\in Y$:
$\omega(y)<\frac12$ and the line $y$ meets a square of inclination $\ge\beta(y)$. Then
\[
\int_Y2\beta(y)\,dy\ \le\ A\,W,\qquad A:=\frac{4\sqrt{K_w^*/Q_*}}{2-3\cdot10^{-6}}\ <\ 10.6067 .
\]
\end{ImpStatement}

\begin{ImpStatement}{Lemma}{4.14}\label{imp:L4.14}
Let $0<\delta\le10^{-5}$ and $0<\bar\beta\le3\cdot10^{-6}$. Let $X,X'\subset[0,k]$. For $x\in X$ let $\bar y_x\in(0,\frac k2-1]$ and
$\beta(x)\in[0,\bar\beta]$ be such that the segment $\sigma_x=\{x\}\times[0,\bar y_x]$ has vertical waste $<\delta$ and meets a square of
inclination $\ge\beta(x)$; for $x\in X'$ let $\bar y'_x$, $\beta'(x)$ satisfy the same with the segment
$\sigma_x'=\{x\}\times[k-\bar y'_x,k]$. Then
\[
\int_X\beta(x)\,dx+\int_{X'}\beta'(x)\,dx\ \le\ A\,W .
\]
\end{ImpStatement}

In this paper [R, Lemma~\ref{imp:L4.14}] is used only inside the proof of [R, Lemma~\ref{imp:L4.15}], where it is
applied with $\bar\beta:=\alpha$ and $\bar y_x:=y=d(y)$ (or, for $y>\frac k2$, with $X:=\emptyset$ and
$\bar y'_x:=k-y=d(y)$). Section~\ref{sec:P6} proves a version of it for the paths of Section~\ref{sec:flows}.

For $y\in(0,k)$ let $d(y):=\min(y,k-y)$ be the distance to the nearer of the bottom and the top side.

\begin{ImpStatement}{Lemma}{4.15}\label{imp:L4.15}
Let $0<\delta\le10^{-5}$ and $0<\alpha\le1.5\cdot10^{-6}$, let $y\in(0,k)$ with $d(y)\le\frac k2-1$, and suppose that every square
meeting the line $y$ has inclination $<\alpha$. Then, with $E(y)$ as in Lemma~\ref{imp:L3.5},
\[
E(y)\ \le\ 0.50001\,\alpha\Big(A+\frac\alpha\delta\Big)W .
\]
\end{ImpStatement}

The proof of [R, Lemma~\ref{imp:L4.15}] begins with the inequality
\begin{equation}\label{imp:sec}
\sec x-1\ \le\ 0.50001\,x^2\qquad(0\le x\le10^{-3}).
\end{equation}
For completeness we include a proof: for $0<x\le10^{-3}$ we have $\cos x\ge1-x^2/2>0$, hence
$\sec x-1=(1-\cos x)/\cos x\le(x^2/2)/(1-x^2/2)=x^2/(2-x^2)\le x^2/(2-10^{-6})<0.5000003\,x^2$; for $x=0$ there is
nothing to prove. The proof of [R, Lemma~\ref{imp:L4.15}] also uses [R, Lemma~\ref{imp:L4.14}] and, through it,
the pair estimates [R, Proposition~\ref{imp:P4.5}], [R, Lemma~\ref{imp:L4.8}(a)] and [R, Lemma~\ref{imp:L4.10}]
with $d=\delta$, exactly as the proof of [R, Lemma~\ref{imp:L4.14}] does (the proof of [R, Lemma~\ref{imp:L4.13}]
uses them with $d=d'=2\bar\beta$).

\subsection{The form of the multiplicity bound used in this paper}

The formal statements of [R, Lemmas~\ref{imp:L4.9} and~\ref{imp:L4.10}] concern edges of the visibility graphs
$G_h$ and $G_v$. In Sections~\ref{sec:P3}, \ref{sec:P5} and~\ref{sec:P6} the pairs to which the bound is applied are
pairs of squares that arise from the paths of Section~\ref{sec:flows} (for instance consecutive squares entered by a
path, or the two squares from which two paths reach a common point), or a square and the bottom side (respectively the
top side) of $[0,k]^2$; they need not be edges of $G_h$ or $G_v$. We therefore use the bound in the following form, which is what
the proofs of [R, Lemmas~\ref{imp:L4.9} and~\ref{imp:L4.10}] establish according to the paragraph of [R] following
the proof of [R, Lemma~\ref{imp:L4.9}] and to [R, Remark~\ref{imp:R4.11}], both quoted above.

\begin{proposition}[multiplicity for families of pairs]\label{imp:pairs}
Let $0<d\le10^{-4}$, and let $\sigma,\sigma'$ be two opposite sides of $[0,k]^2$. Let $\mathfrak F$ be a family of
distinct unordered pairs $\{P,Q\}$, where $P$ is a square of the packing and $Q$ is either another square of the
packing or one of $\sigma,\sigma'$, such that the distance between $P$ and $Q$ is less than $d$. For
$e=\{P,Q\}\in\mathfrak F$ put $R_e:=R(P,Q)$ if $Q$ is a square and $R_e:=R(P,S_0)$ if $Q$ is a side, and let $f_e$
be the area of the waste in $R_e$. Then every point of the waste lies in $R_e$ for at most $K_w^*=9$ pairs
$e\in\mathfrak F$; by the analytic argument of [R, Lemma~\ref{imp:L4.9}] it lies in $R_e$ for at most $K_w=13$ pairs
$e\in\mathfrak F$. Consequently
\[
\sum_{e\in\mathfrak F}f_e\ \le\ K_w^*\,W=9W,\qquad\text{and also}\qquad \sum_{e\in\mathfrak F}f_e\ \le\ K_w\,W=13W .
\]
\end{proposition}

\begin{proof}
The pointwise bounds are the statements quoted above: the paragraph after the proof of [R, Lemma~\ref{imp:L4.9}]
states that the bound $13$ ``holds for any family of such pairs at distance $<d$'', where the sides involved are two
opposite ones, and [R, Remark~\ref{imp:R4.11}] states that the proof of [R, Lemma~\ref{imp:L4.10}] bounds ``the number
of all pairs of squares (or of a square and one of two opposite sides) at distance less than $d$ whose rectangles contain
$p$'' by $9$. (In that proof the count runs over the set $I$ of all centres within distance
$\rho_1:=\sqrt{(d+\frac{\sqrt2}2)^2+\frac14}$ of $p$, and every pair counted has a member whose centre lies in $I$; the
proof bounds separately the pairs with both centres in $I$, the pairs of a square with centre in $I$ and a side, and
the pairs of a square with centre in $I$ and a square with centre outside $I$, using only distances and angles, never
the graph.)
For the sum, the multiplicity function $p\mapsto\#\{e\in\mathfrak F:p\in R_e\}$ is a finite sum of indicator functions
of closed rectangles, hence Borel, and by Tonelli's theorem
\[
\sum_{e\in\mathfrak F}f_e=\int_{\text{waste}}\#\{e\in\mathfrak F:\ p\in R_e\}\,dp\ \le\ 9\,|\text{waste}|=9W,
\]
and in the same way with $13$; here $|\text{waste}|=W$ because the squares have total area $N$ and the waste differs from
$[0,k]^2\setminus\bigcup\mathcal P$ by a null set. ($\mathfrak F$ is finite, since the packing is finite.)
\end{proof}

When Proposition~\ref{imp:pairs} is applied, the pairs are always unordered and each is counted once; in the
applications, $\sigma$ and $\sigma'$ are the bottom and the top side of $[0,k]^2$. For a square $P$ near a corner of
$[0,k]^2$ the rectangle $R(P,S_0)$ may stand on a side other than $\sigma,\sigma'$ (the side of $[0,k]^2$ nearest to the
centre of $P$); this is the situation of [R, Lemmas~\ref{imp:L4.8}(b) and~\ref{imp:L4.9}], whose proofs allow it.

\subsection{Dependence on a computer-assisted result}

\begin{remark}[dependence on {[R, Lemma~\ref{imp:L4.10}]}]\label{rem:imp-dep}
The proof of [R, Lemma~\ref{imp:L4.10}] is computer-assisted: it consists of a branch-and-bound computation over a
cover of a parameter domain by boxes, with rigorous ball arithmetic (see [R, proof of Lemma~\ref{imp:L4.10}]).
In this paper [R, Lemma~\ref{imp:L4.10}] enters only through the number $K_w^*=9$, namely
\begin{enumerate}[label=(\roman*)]
\item through the constant $A=4\sqrt{K_w^*/Q_*}/(2-3\cdot10^{-6})$ of [R, Lemmas~\ref{imp:L4.13}--\ref{imp:L4.15}];
\item through Proposition~\ref{imp:pairs} with the bound $9$, in Sections~\ref{sec:P3}, \ref{sec:P5} and~\ref{sec:P6}
(hence in the constants $C_\Lambda$ and $A'$ of Definition~\ref{def:flows-consts}).
\end{enumerate}
Theorem~\ref{thm:A} and its constants $c_\infty$ and $k_0=4.62\cdot10^{12}$ use [R, Lemma~\ref{imp:L4.10}].
Theorem~\ref{thm:A13} does not: it uses the analytic [R, Lemma~\ref{imp:L4.9}] ($K_w=13$) in all places where
$K_w^*=9$ is used. In (ii) this is the second bound of Proposition~\ref{imp:pairs}. In (i) it rests on the following
facts about the proofs in [R]. In the proof of [R, Lemma~\ref{imp:L4.13}], [R, Lemma~\ref{imp:L4.10}] is used only
in the sentence ``By Lemma~4.10 (with $d=d'$), every point of the waste lies in $R_e$ for at most $K_w^*$ of the pairs
with $|Y_e^c|>0$, so $\sum f_e\le K_w^*W$, the sum over these pairs'', and these pairs are edges of $G_h$ (``they are
edges of $G_h$''); in the
proof of [R, Lemma~\ref{imp:L4.14}] it is used only in ``$\sum_ef_e\le K_w^*W$ (Lemma~4.10 with $d=\delta$)'', for
pairs that are ``edges of $G_v$ whose members are at distance $<\delta$''; and the proof of
[R, Lemma~\ref{imp:L4.15}] uses [R, Lemma~\ref{imp:L4.10}] only through [R, Lemma~\ref{imp:L4.14}]. Replacing $9$ by
$13$ in these sentences ([R, Lemma~\ref{imp:L4.9}] applies to these edges), $K_w^*$ by $K_w$ in the formulas that follow
them, and leaving the rest of the proofs unchanged
(the last step of the proof of [R, Lemma~\ref{imp:L4.13}] uses only $A\ge2$), one obtains [R,
Lemmas~\ref{imp:L4.13}--\ref{imp:L4.15}] with $A$ replaced by
\[
A_{13}:=\frac{4\sqrt{K_w/Q_*}}{2-3\cdot10^{-6}}=\frac{4\sqrt{13/0.32}}{2-3\cdot10^{-6}}\ <\ 12.7476 ,
\]
as recorded in [R, Remark~7.5]: ``with the analytic bound $K_w=13$ of Lemma~4.9 in place of $K_w^*$, the proof of
Section~6 gives $A<12.7476$''. Numerically $A=10.60661762\ldots<10.6067$ and $A_{13}=12.74756790\ldots<12.7476$
(digits truncated; the bounds are rounded up).
\end{remark}

\begin{remark}[where the results of {[R]} are used]\label{rem:imp-uses}
[R, Lemma~\ref{imp:L2.1}] (with [R, Corollary~\ref{imp:C3.3}]) gives Corollary~\ref{cor:A} from Theorems~\ref{thm:A}
and~\ref{thm:A13} (Section~\ref{sec:assembly}). [R, Lemma~\ref{imp:L3.1}(a)(b)] gives the vertical extent of a square,
the positions of the ramps and the middle chords used in Sections~\ref{sec:flows}, \ref{sec:P12}, \ref{sec:P3},
\ref{sec:P6} and~\ref{sec:P7}. The second part of
[R, Lemma~\ref{imp:L3.4}] ($\Phi(S)\subset\mathcal R(S)$, $\int_{\Phi(S)}c_S=\sin a\cos a$, $\Phi(S)=\emptyset$ if $a(S)=0$)
is used in Sections~\ref{sec:P12}, \ref{sec:P7} and~\ref{sec:P8}, and [R, Lemma~\ref{imp:L3.5}] gives the short chords
on lines of kind~(iii) in Sections~\ref{sec:P7} and~\ref{sec:P8}; [R, Lemma~\ref{imp:L3.2}] enters through the proof
of [R, Lemma~\ref{imp:L3.5}].
[R, Proposition~\ref{imp:P4.5}] and [R, Lemma~\ref{imp:L4.8}(a)] give the lower bounds for the waste in the
rectangles $R_e$ (Sections~\ref{sec:P3}, \ref{sec:P5}, \ref{sec:P6}). [R, Lemmas~\ref{imp:L4.1} and~\ref{imp:L4.12}]
are mentioned only for comparison, in Section~\ref{sec:P6}, whose slit inequalities are proved directly; like
[R, Lemma~\ref{imp:L4.2}, Proposition~\ref{imp:P4.3} and Lemma~\ref{imp:L4.4}], they are needed only for the
statements and proofs in [R] of the results quoted here. [R, Lemmas~\ref{imp:L4.9}, \ref{imp:L4.10} and Remark~\ref{imp:R4.11}] are used through
Proposition~\ref{imp:pairs}. [R, Lemmas~\ref{imp:L4.13} and~\ref{imp:L4.15}] bound the lines of kinds~(ii) and~(iii) in
Section~\ref{sec:P8}.
\end{remark}

\begin{remark}[facts taken from the proofs in {[R]}]\label{rem:imp-proofs}
Besides the statements quoted above, the following facts about the \emph{proofs} in [R] (version~1.2) are used. They
are not part of the quoted statements; each is stated where it is used, and can be checked in [R] at the place
indicated.
\begin{enumerate}[label=(\alph*)]
\item The proofs of [R, Lemmas~\ref{imp:L4.9} and~\ref{imp:L4.10}] use only distances and angles, the disjointness of
the squares, the fact that $p$ is a point of the waste, and the fact that the sides involved are two opposite ones;
they do not use the visibility graphs ([R, the paragraph after the proof of Lemma~\ref{imp:L4.9}, and
Remark~\ref{imp:R4.11}]). This is used in Proposition~\ref{imp:pairs} and its proof.
\item In the proofs of [R, Lemmas~\ref{imp:L4.13}--\ref{imp:L4.15}], [R, Lemma~\ref{imp:L4.10}] enters only through
the two sentences quoted in Remark~\ref{rem:imp-dep}. Theorem~\ref{thm:A13} rests on this.
\item The proofs of [R, Lemmas~\ref{imp:L4.8}(b), \ref{imp:L4.9} and~\ref{imp:L4.10}] allow the rectangle $R(P,S_0)$ of a
square $P$ near a corner to stand on any side of $[0,k]^2$, and the proof of [R, Lemma~\ref{imp:L4.10}] shows that all
such rectangles that contain a given point of the waste stand on one and the same side. This is used after
Proposition~\ref{imp:pairs} and in Remark~\ref{rem:P3-corner}.
\item The distance from a square of inclination $a$ to a side of $[0,k]^2$ is the distance from its centre to that side
minus $\frac12(\cos a+\sin a)$ ([R, proof of Proposition~\ref{imp:P4.5}(b)]; used in Remark~\ref{rem:P3-corner}).
\item The proof of [R, Lemma~\ref{imp:L4.15}] begins with the inequality \eqref{imp:sec} and uses
[R, Lemma~\ref{imp:L4.14}] as described after [R, Lemma~\ref{imp:L4.15}] above.
\end{enumerate}
\end{remark}

\section{Path flows}\label{sec:flows}

Throughout this section $k\ge4$ is an integer and $\mathcal P$ is a closed packing of $[0,k]^2$ with $N$ squares;
$W=k^2-N$. As in [R], the \emph{waste} is the set of points of $(0,k)^2$ that lie in no square of $\mathcal P$; its
area is $W$. We write $|\cdot|$ for Lebesgue measure on $\R$ (for sets of starting abscissas $x\in(0,k)$ we also call
it the \emph{floor measure}), $e_1=(1,0)$, $e_2=(0,1)$, and $\ell_y:=\R\times\{y\}$ for the horizontal line at height
$y$. The letter $\pi$ denotes the parameter vector of Definition~\ref{def:flows-params}, $\pi_x$ a path, and $\pi$ also
the number; the meaning is always clear from the context.

\subsection{Parameters and constraints}\label{ssec:flows-params}

\begin{definition}[parameters]\label{def:flows-params}
\emph{Fixed constants.} $\delta:=10^{-5}$, $\bar\beta:=1.5\cdot10^{-6}$, $b_W:=2.0002$, and, from [R],
$Q_*=0.32$, $K_w^*=9$, $K_w=13$ (see Section~\ref{sec:imported}). Some results of Sections~\ref{sec:P3}--\ref{sec:P6}
are stated for every $\delta\in(0,10^{-5}]$; everywhere else $\delta=10^{-5}$.

\emph{Free parameters.} $\pi=(c_0,c_1,y_0,\varepsilon,\omega_0)$.

\emph{Derived quantities.} For $s>0$, $t>0$ and $0<y<k$ put
\begin{gather*}
\alpha(s):=c_0s^{-1/2},\qquad \amax:=\alpha(y_0),\qquad \hat\beta(t):=c_1t^{-3/4},\\
d(y):=\min(y,k-y),\qquad \beta(y):=\hat\beta(d(y)),\qquad b^*:=\hat\beta\big((1-\varepsilon)y_0\big),\\
y_1:=\tfrac k2-3,\qquad \hmax:=\tfrac k2-1,\\
w_0:=0.50001\,c_0^2\Big(1+\frac{1.0001}{y_0}\Big)+\delta+1.0001\,b^*+10^{-12},\qquad h_0:=1-2w_0 .
\end{gather*}
$\alpha(s)$ is the \emph{tilt threshold} at scale $s$ and $\beta(y)$ the \emph{line threshold} at height $y$.
\end{definition}

The parameter vectors used for the main theorems are
\[
\pi^*:=(0.3078,\ 0.2754,\ 42\,108\,000\,000,\ 5.8\cdot10^{-6},\ 1.7\cdot10^{-5})\quad\text{(Theorem~\ref{thm:A})},
\]
\[
\pi^*_{13}:=(0.3674,\ 0.3012,\ 60\,000\,000\,000,\ 5\cdot10^{-6},\ 1.4\cdot10^{-5})\quad\text{(Theorem~\ref{thm:A13})}.
\]

\begin{definition}[constraints]\label{def:flows-constraints}
The parameters are subject to the following conditions.
\begin{enumerate}
\item[(C0)] $c_0,c_1>0$, $y_0\in\Z_{>0}$, $0<\varepsilon<1$, $\omega_0>0$.
\item[(C1)] $\alpha(y_0)\le\bar\beta=1.5\cdot10^{-6}$.
\item[(C2)] $\hat\beta(y_0)\le1.5\cdot10^{-6}$.
\item[(C3)] $c_1(1-\varepsilon)^{-3/4}<c_0y_0^{1/4}$ (not used in the proofs; see below).
\item[(C4)] $\omega_0<\frac12$.
\item[(C5)] $w_0<\frac12$.
\item[(C6)] $\delta\le\frac12\cos2\amax$.
\item[(C7)] $y_0+1\le(1-\varepsilon)y_1$.
\item[(Q1)] $\alpha(y_0)\le10^{-3}$.
\item[(Q2)] $b^*\le10^{-4}$.
\end{enumerate}
\end{definition}

The role of each condition is explained at the end of Section~\ref{ssec:flows-lines}, after the sets of heights to which
some of them refer have been defined.

\begin{lemma}[elementary consequences]\label{lem:flows-constraints}
\begin{enumerate}[label=(\alph*)]
\item \textup{(C1)} holds if and only if $y_0\ge(c_0/\bar\beta)^2$. Under \textup{(C1)}, $\amax\le1.5\cdot10^{-6}$,
\textup{(Q1)} holds, and \textup{(C6)} holds for every $\delta\in(0,10^{-5}]$.
\item The functions $\alpha$ and $\hat\beta$ are decreasing. If \textup{(C2)} holds and $\varepsilon\le0.9963$, then
\textup{(Q2)} holds.
\item \textup{(C3)} holds if and only if $\hat\beta((1-\varepsilon)s)<\alpha(s)$ for every $s\ge y_0$.
\item Under \textup{(C0)}, \textup{(C1)} and \textup{(C7)}, for every $s\in[y_0,y_1]$:
$0<\alpha(s)\le\amax\le\bar\beta<\pi/4$, $s+2\le\hmax$, and $s<k/2$; moreover $0<\hmax<k$.
\item If \textup{(C7)} holds for some $k$, it holds, with the same $\pi$, for every larger $k$.
\end{enumerate}
\end{lemma}

\begin{proof}
(a) $c_0y_0^{-1/2}\le\bar\beta$ is equivalent to $y_0\ge(c_0/\bar\beta)^2$. If (C1) holds, then
$\amax=\alpha(y_0)\le1.5\cdot10^{-6}\le10^{-3}$, and $2\amax\le3\cdot10^{-6}$ gives
$\frac12\cos2\amax\ge\frac12(1-\frac12(3\cdot10^{-6})^2)>0.49>10^{-5}\ge\delta$.

(b) The first statement is clear. We have $b^*=(1-\varepsilon)^{-3/4}\hat\beta(y_0)\le1.5\cdot10^{-6}(1-\varepsilon)^{-3/4}$.
If $\varepsilon\le0.9963$, then $(1-\varepsilon)^{3/4}\ge0.0037^{3/4}>0.015$, because
$0.0037^3=5.0653\cdot10^{-8}>5.0625\cdot10^{-8}=0.015^4$; hence $b^*<1.5\cdot10^{-6}/0.015=10^{-4}$.

(c) For $s>0$, $\hat\beta((1-\varepsilon)s)/\alpha(s)=\frac{c_1}{c_0}(1-\varepsilon)^{-3/4}s^{-1/4}$ is strictly decreasing
in $s$, so it is $<1$ for all $s\ge y_0$ if and only if it is $<1$ at $s=y_0$, which is (C3).

(d) As $\alpha$ is decreasing and $s\ge y_0$, $0<\alpha(s)\le\alpha(y_0)=\amax\le\bar\beta<\pi/4$. From $s\le y_1=\frac k2-3$
we get $s+2\le\frac k2-1=\hmax<k$ and $s<\frac k2$; and $\hmax\ge1$ as $k\ge4$.

(e) $(1-\varepsilon)y_1=(1-\varepsilon)(\frac k2-3)$ increases with $k$.
\end{proof}

\begin{lemma}[the parameter vectors]\label{lem:flows-pistar}
\begin{enumerate}[label=(\alph*)]
\item $\pi^*$ satisfies \textup{(C0)--(C7)}, \textup{(Q1)} and \textup{(Q2)} for every integer $k\ge4.62\cdot10^{12}$.
At $\pi^*$ one has $(c_0/\bar\beta)^2=205\,200^2=42\,107\,040\,000\le y_0$, $b^*<10^{-8}$, $0.0473813<w_0<0.0473814$ and
$0.9052372<h_0<0.9052374$.
\item $\pi^*_{13}$ satisfies \textup{(C0)--(C7)}, \textup{(Q1)} and \textup{(Q2)} for every integer
$k\ge2.17\cdot10^{13}$, and $w_0<0.0676$ there.
\end{enumerate}
\end{lemma}

\begin{proof}
(a) (C0) is clear. (C1): $c_0^2=0.3078^2=0.09474084$ and $\bar\beta^2y_0=2.25\cdot10^{-12}\cdot42\,108\,000\,000=0.094743$, so
$c_0^2\le\bar\beta^2y_0$; also $c_0/\bar\beta=205\,200$ exactly. (C2), (Q2): $y_0\ge10^{10}$ and
$(1-\varepsilon)y_0\ge10^{10}$ give $\hat\beta(y_0),b^*\le0.2754\cdot10^{-7.5}<0.2754/(3\cdot10^7)<10^{-8}$. (C3):
$c_1(1-\varepsilon)^{-3/4}\le0.2754/(1-5.8\cdot10^{-6})<0.2755$ and $c_0y_0^{1/4}\ge0.3078\cdot10^{2.5}>97$. (C4) is
clear. (C5): $0.50001\,c_0^2=0.0473713674084$, $1.0001/y_0<10^{-10}$ and $1.0001\,b^*<1.0001\cdot10^{-8}$, so
\begin{gather*}
w_0\ >\ 0.0473713674+10^{-5}\ >\ 0.0473813,\\
w_0\ <\ 0.0473713674085\,(1+10^{-10})+10^{-5}+1.0001\cdot10^{-8}+10^{-12}\ <\ 0.0473814,
\end{gather*}
which gives the stated bounds for $w_0$ and $h_0=1-2w_0$, and (C5). (C6) and (Q1) follow from (C1) by
Lemma~\ref{lem:flows-constraints}(a). (C7): for $k\ge4.62\cdot10^{12}$,
$(1-\varepsilon)y_1\ge(1-10^{-5})(2.31\cdot10^{12}-3)>2\cdot10^{12}>y_0+1$.

(b) In the same way: $c_0^2=0.3674^2=0.13498276\le0.135=2.25\cdot10^{-12}\cdot6\cdot10^{10}$ gives (C1);
$\hat\beta(y_0),b^*<0.3012/(3\cdot10^7)<1.1\cdot10^{-8}$ gives (C2) and (Q2); $c_1(1-\varepsilon)^{-3/4}<0.302<97<c_0y_0^{1/4}$
gives (C3); $w_0<0.50001\cdot0.13498276\,(1+10^{-10})+10^{-5}+1.0001\cdot1.1\cdot10^{-8}+10^{-12}<0.0676$ gives (C5); and for
$k\ge2.17\cdot10^{13}$, $(1-\varepsilon)y_1\ge(1-10^{-5})(1.085\cdot10^{13}-3)>10^{13}>y_0+1$.
\end{proof}

\subsection{Lines, windows and constants}\label{ssec:flows-lines}

\begin{definition}[lines and windows]\label{def:flows-lines}
A square $S$ \emph{meets the line $y$} if $S\cap\ell_y\ne\emptyset$. For $y\in\R$, $\frc y$ is the fractional part
and $\dist(y,\Z)=\min(\frc y,1-\frc y)$. Put
\[
H:=\{y\in(0,k):\ y_0\le d(y)\le(1-\varepsilon)y_1,\ \frc y\in[w_0,1-w_0]\},\qquad H_b:=H\cap(0,\tfrac k2),\quad
H_t:=H\cap(\tfrac k2,k).
\]
Every $y\in H$ is of exactly one of the following \emph{kinds}:
(i) $\omega(y)\ge\omega_0$;
(ii) $\omega(y)<\omega_0$ and the line $y$ meets a square $S$ with $a(S)\ge\beta(y)$;
(iii) $\omega(y)<\omega_0$ and every square meeting the line $y$ has $a(S)<\beta(y)$.
Let $Y_i,Y_{ii},Y_{iii}\subset H$ be the corresponding sets, $Y_b:=Y_{iii}\cap(0,\frac k2)$ and $Y_t:=Y_{iii}\cap(\frac k2,k)$.
For $s\in[y_0,y_1]$ let
\[
\mathcal W_s:=[(1-\varepsilon)s,\ s],\qquad V_s:=[(1-\varepsilon)s-1.0001,\ s+1.0001]\quad(\text{so }|V_s|=\varepsilon s+b_W),
\]
\[
\mathcal Z(s):=\{Z\in\mathcal P:\ \Phi(Z)\cap Y_b\cap\mathcal W_s\ne\emptyset\},\qquad
\tilde{\mathcal Z}(s):=\{Z\in\mathcal P:\ \Phi(Z)\cap Y_t\cap(k-\mathcal W_s)\ne\emptyset\},
\]
where $k-Y:=\{k-y:\ y\in Y\}$ for a set $Y\subset\R$, and $\Phi(Z)=\{y:0<c_Z(y)<1\}$ ([R, Lemma~\ref{imp:L3.4}]). For a Borel set $Y\subset H$,
$\ell(Y):=\int_Yd(y)^{-3/4}\,dy$ (the \emph{line measure}).
\end{definition}

Under (C7), $d(y)\le(1-\varepsilon)y_1<\frac k2-1$ on $H$; in particular $\frac k2\notin H$, so $Y_{iii}=Y_b\sqcup Y_t$,
and $d(y)=y$ on $H_b$.

\begin{definition}[constants]\label{def:flows-consts}
With $Q_*,K_w^*,K_w$ from [R] put
\begin{gather*}
A:=\frac{4\sqrt{K_w^*/Q_*}}{2-3\cdot10^{-6}},\qquad A':=\frac A{\cos\bar\beta}+10^{-8},\qquad C_\Lambda:=1+1185.6\,\delta^2;\\
A_{13}:=\frac{4\sqrt{K_w/Q_*}}{2-3\cdot10^{-6}},\qquad A'_{13}:=\frac{A_{13}}{\cos\bar\beta}+10^{-8},\qquad
C^{(13)}_\Lambda:=1+1712.6\,\delta^2 .
\end{gather*}
($A$ is the constant of [R, Lemma~\ref{imp:L4.13}].) For the integration of Section~\ref{sec:P8} put
\begin{gather*}
\invF:=\frac{4(1-\varepsilon)^{3/4}\big(1-(1-\varepsilon)^{3/4}\big)}{3\varepsilon\,\big(1+b_W/(\varepsilon y_0)\big)},\qquad
\Gamma:=1.00002\Big(2c_1Ay_0^{-1/2}+\tfrac45c_1^2\delta^{-1}y_0^{-5/4}\Big),\\
\Omega_W:=\int_{y_0}^{y_1}2(s+2)\tan\alpha(s)\,\frac{c_0}2s^{-3/2}\,ds,\qquad
K_W:=\frac{4c_1}{c_0\,\invF}\cdot1.000001\,c_0^2,\\
B:=\frac1{\omega_0y_0^{3/4}}+\frac A{2c_1}+\frac{2c_1}{c_0\,\invF}\Big(A'+\frac{\alpha(y_0)\,C_\Lambda}\delta\Big)+\Gamma .
\end{gather*}
In the analytic variant (Theorem~\ref{thm:A13}), $A,A',C_\Lambda$ are replaced by $A_{13},A'_{13},C^{(13)}_\Lambda$ in
$\Gamma$ and $B$.
\end{definition}

$\Gamma$ denotes the explicit expression above; Section~\ref{sec:P8} shows that it is an upper bound for the integral
$\Gamma_H:=\int_H0.50001\,\beta(y)\big(A+\beta(y)/\delta\big)\,d(y)^{-3/4}dy$. (The numbers $\Gamma$ and $\Gamma_H$ are
not to be confused with the trajectories of Definition~\ref{def:flows-traj}, which always carry a floor point as an
index, as in $\Gamma_x$, $\Gamma^0_x$.) Numerically (rounded up)
$A<10.6067$, $A'<10.6068$, $C_\Lambda<1.0000002$, $A_{13}<12.7476$, $A'_{13}<12.7476$ and $C^{(13)}_\Lambda<1.0000002$;
see Section~\ref{sec:constants}. The constants $C_\Lambda$ and $A'$ are those of the bounds of
Sections~\ref{sec:P5} and~\ref{sec:P6}. The subscript $W$ in $b_W$, $\Omega_W$ and $K_W$ is only a label (window,
wall), as in the wall time $t_W$ of Definition~\ref{def:flow}; it does not refer to the waste area $W$, and $K_W$ is
unrelated to the multiplicity bounds $K_w=13$ and $K_w^*=9$ of [R].

The roles of the conditions of Definition~\ref{def:flows-constraints} are as follows.
(C0) is standing. (C1) gives $\amax\le1.5\cdot10^{-6}$ and $\amax\le\bar\beta$, the hypotheses on $\amax$ in
Sections~\ref{sec:P3}, \ref{sec:P5} and~\ref{sec:P6}, and the bound $\tan x\le1.000001\,x$ for $0\le x\le\amax$ used
for the wall term in Section~\ref{sec:P8}; it implies (Q1) and (C6) (Lemma~\ref{lem:flows-constraints}). (C2) makes
all line thresholds $\beta(y)$, $y\in H$, at most $1.5\cdot10^{-6}$; this is the
hypothesis of [R, Lemmas~\ref{imp:L4.13} and~\ref{imp:L4.15}] in Section~\ref{sec:P8}, and it bounds the inclinations of
the squares of $\mathcal Z(s)$ (Sections~\ref{sec:P12} and~\ref{sec:P7}). (C3) is equivalent to
$\hat\beta((1-\varepsilon)s)<\alpha(s)$ for all $s\ge y_0$ (Lemma~\ref{lem:flows-constraints}); it is not used in any
proof of this paper and is kept in the list only for reference (it holds at $\pi^*$, at $\pi^*_{13}$ and at all other
parameters used in Sections~\ref{sec:assembly} and~\ref{sec:constants}). The results below whose hypotheses include
(C0)--(C6) or (C0)--(C7) hold without (C3). (C4) is the hypothesis $\omega(y)<\frac12$ of
[R, Lemma~\ref{imp:L4.13}] for the lines of kind~(ii). (C5) gives $h_0>0$, used in the lower bound for the measure of
$H$ in Section~\ref{sec:P8}. (C6) is the hypothesis of the two-source lemma of Section~\ref{sec:P5}, on which the bounds
for deaths, merges and gap overlaps there rest; it is not used in Section~\ref{sec:P3}. (C7) depends on $k$; it gives
$y_0<y_1$ and is used in this section (Lemma~\ref{lem:flows-constraints}(d) and the scale flows) and in
Sections~\ref{sec:P8} and~\ref{sec:assembly}. (Q1) gives $\sec\alpha(s)-1\le0.50001\,\alpha(s)^2$ for $s\ge y_0$, by
the inequality \eqref{imp:sec} (Section~\ref{sec:P12}). (Q2) gives $a(Z)<b^*\le10^{-4}$ and
$\cos a(Z)+\sin a(Z)<1.0001$ for the squares $Z\in\mathcal Z(s)$ (Section~\ref{sec:P12}). Finally $0<\delta\le10^{-5}$
is the hypothesis of [R, Proposition~\ref{imp:P4.5} and Lemmas~\ref{imp:L4.14} and~\ref{imp:L4.15}] and of
Sections~\ref{sec:P3}--\ref{sec:P6}.

\subsection{Squares, frames and the floor}\label{ssec:flows-frames}

\begin{definition}[frames and sides]\label{def:flows-frames}
Let $S$ be a square. Its \emph{phase} $\varphi_S\in(-\pi/4,\pi/4]$ is the unique number in this interval that is
congruent modulo $\pi/2$ to the angle between the positive $x$-axis and a side of $S$. Put
\[
u_S:=(\cos\varphi_S,\sin\varphi_S),\qquad n_S:=(-\sin\varphi_S,\cos\varphi_S),
\]
let $c_S$ be the centre, so that $S=\{c_S+\xi u_S+\zeta n_S:\ |\xi|,|\zeta|\le\frac12\}$, and $a(S)=|\varphi_S|$ is the
inclination of [R]; for squares $S_1,S_2$, $\theta(S_1,S_2)=\min_{m\in\Z}|\varphi_{S_1}-\varphi_{S_2}-m\pi/2|$. The
\emph{sides} of $S$ are
\begin{gather*}
\Bot(S):=\{c_S-\tfrac12n_S+\xi u_S:|\xi|\le\tfrac12\},\qquad \Top(S):=\{c_S+\tfrac12n_S+\xi u_S:|\xi|\le\tfrac12\},\\
\{c_S+\tfrac12u_S+\zeta n_S:|\zeta|\le\tfrac12\}\ \text{(right side)},\qquad \{c_S-\tfrac12u_S+\zeta n_S:|\zeta|\le\tfrac12\}\ \text{(left side)},
\end{gather*}
with outer normals $-n_S$, $n_S$, $u_S$, $-u_S$; their \emph{relative interiors} ($\relint$) are given by
$|\xi|<\frac12$, resp.\ $|\zeta|<\frac12$, and the vertices are $c_S\pm\frac12u_S\pm\frac12n_S$. The four relative
interiors and the four vertices partition $\partial S$. The \emph{line of} $\Bot(S)$ is
$\{v:(v-c_S)\cdot n_S=-\frac12\}$, and similarly for the other sides. (The centre $c_S$ is a point, with coordinates
$c_{S,x}$, $c_{S,y}$; the chord length $c_S(y)$ of Section~\ref{sec:imported} is a function of the height $y$ and always
carries its argument.)

\emph{The $45^\circ$ convention.} If the sides of $S$ make the angle $\pi/4$ with the axes, then $\varphi_S=+\pi/4$;
thus $n_S=(-1,1)/\sqrt2$, $\Bot(S)$ is the lower right side (outer normal $(1,-1)/\sqrt2$), and the lower left side
is the left side.

\emph{The floor.} The floor is the pseudo-square $X_0$ with $\varphi_{X_0}:=0$, $u_{X_0}:=e_1$, $n_{X_0}:=e_2$, whose
``top side'' is $[0,k]\times\{0\}$ with relative interior $(0,k)\times\{0\}$. It is not a member of $\mathcal P$.
\end{definition}

For $X,Y$ squares or the floor we have
\begin{equation}\label{eq:flows-frame}
\begin{gathered}
n_X\cdot e_2=\cos\varphi_X,\qquad e_2=\cos\varphi_X\,n_X+\sin\varphi_X\,u_X,\\
n_X\cdot n_Y=u_X\cdot u_Y=\cos(\varphi_Y-\varphi_X),\qquad u_X\cdot n_Y=-\sin(\varphi_Y-\varphi_X),
\end{gathered}
\end{equation}
and the point $c_S+\xi u_S+\zeta n_S$ has height $c_{S,y}+\xi\sin\varphi_S+\zeta\cos\varphi_S$. In particular the
lowest point of $S$ has height $v_S=c_{S,y}-\frac12(\cos a(S)+\sin a(S))$, and the points of $\Top(S)$ have heights
$\ge c_{S,y}+\frac12(\cos a(S)-\sin a(S))=v_S+\cos a(S)>v_S$. Since $S\subset[0,k]^2$ and the abscissas
$\xi\mapsto c_{S,x}\mp\frac12\sin\varphi_S+\xi\cos\varphi_S$ of the points of $\Top(S)$, resp.\ $\Bot(S)$, are strictly
increasing in $\xi$, the points of $\relint\Top(S)$ and of $\relint\Bot(S)$ have abscissas in $(0,k)$.

\begin{lemma}[reflection]\label{lem:flows-refl}
Let $\rho(x,y):=(x,k-y)$.
\begin{enumerate}[label=(\alph*)]
\item $\rho\mathcal P:=\{\rho S:S\in\mathcal P\}$ is a closed packing of $[0,k]^2$ with $N$ squares and the same $W$, and
$\rho$ maps the waste of $\mathcal P$ onto the waste of $\rho\mathcal P$. For every square $S$:
$a(\rho S)=a(S)$, $\varphi_{\rho S}\equiv-\varphi_S\pmod{\pi/2}$, $v_{\rho S}=k-(v_S+\cos a(S)+\sin a(S))$ and
$c_{\rho S}(k-y)=c_S(y)$ for all $y$.
\item $\mathcal R(\rho S)=k-\mathcal R(S)$ and $\Phi(\rho S)=k-\Phi(S)$; the line waste and the excess $E$ of
[R, Lemma~\ref{imp:L3.5}] satisfy $\omega_{\rho\mathcal P}(y)=\omega_{\mathcal P}(k-y)$ and
$E_{\rho\mathcal P}(y)=E_{\mathcal P}(k-y)$; $d(k-y)=d(y)$ and $\beta(k-y)=\beta(y)$; the line $k-y$ meets $\rho S$ if and
only if the line $y$ meets $S$.
\item Since $k\in\Z$: $\dist(k-y,\Z)=\dist(y,\Z)$ for all $y$, hence $k-H=H$ and $H_t=k-H_b$. The kind of $y\in H$ for
$\rho\mathcal P$ equals the kind of $k-y$ for $\mathcal P$; thus $Y_{iii}[\rho\mathcal P]=k-Y_{iii}[\mathcal P]$,
$Y_b[\rho\mathcal P]=k-Y_t[\mathcal P]$, and $\rho Z\in\mathcal Z(s)[\rho\mathcal P]$ if and only if
$Z\in\tilde{\mathcal Z}(s)[\mathcal P]$.
\end{enumerate}
\end{lemma}

\begin{proof}
(a) $\rho$ is an isometry of $\R^2$ mapping $[0,k]^2$ onto itself, so it maps squares to squares, preserves
disjointness and areas, and maps $(0,k)^2$ onto itself. It maps the direction $(\cos\psi,\sin\psi)$ to
$(\cos\psi,-\sin\psi)$, hence the side directions of $\rho S$ have angles congruent to $-\varphi_S$, and the least angle
between a side and the axes is unchanged. The highest point of $S$ has height $v_S+\cos a+\sin a$
([R, Lemma~\ref{imp:L3.1}(a)]) and is mapped to the lowest point of $\rho S$. Finally
$\rho S\cap\ell_{k-y}=\rho(S\cap\ell_y)$.

(b) With $v':=v_{\rho S}=k-v-\cos a-\sin a$ ($v:=v_S$, $a:=a(S)$), [R, Lemma~\ref{imp:L3.1}(b)] gives
$\mathcal R(\rho S)=[v',v'+\sin a)\cup(v'+\cos a,v'+\cos a+\sin a]=[k-v-\cos a-\sin a,\,k-v-\cos a)\cup(k-v-\sin a,\,k-v]$,
which is $k-\big([v,v+\sin a)\cup(v+\cos a,v+\cos a+\sin a]\big)=k-\mathcal R(S)$, including the open and closed ends. The
statements on $\Phi$, $\omega$, $E$ and on meeting lines follow from $\rho S\cap\ell_{k-y}=\rho(S\cap\ell_y)$ and
$\rho\big(([0,k]\times\{y\})\setminus\bigcup\mathcal P\big)=([0,k]\times\{k-y\})\setminus\bigcup\rho\mathcal P$; $d(k-y)=d(y)$
is clear.

(c) $k-\Z=\Z$. If $y\in\Z$, both fractional parts vanish and $0\notin[w_0,1-w_0]$ (as $w_0>0$); otherwise
$\frc{k-y}=1-\frc y$ and $[w_0,1-w_0]$ is invariant under $t\mapsto1-t$. Together with $d(k-y)=d(y)$ this gives $k-H=H$.
By (b), $\omega$, the inclinations of the squares meeting a line and $\beta$ are transported by $y\mapsto k-y$, so the
kinds are; and $k-(\frac k2,k)=(0,\frac k2)$. The last statement follows from $\Phi(\rho Z)=k-\Phi(Z)$:
$(k-\Phi(Z))\cap(k-Y_t)\cap\mathcal W_s\ne\emptyset$ if and only if $\Phi(Z)\cap Y_t\cap(k-\mathcal W_s)\ne\emptyset$.
\end{proof}

\subsection{Rules of motion}\label{ssec:flows-rules}

\begin{definition}[flows: ray steps and events]\label{def:flow}
Let $0<\alpha_F<\pi/4$, $\delta>0$ and $0<h_F<k$. The \emph{flow} $F=F(\alpha_F,\delta,h_F)$ assigns to every floor
point $x\in(0,k)$ a \emph{path} $\pi_x$. A path moves through \emph{states} $(p,d,X,g)$: $p\in[0,k]^2$ is the current
point; $X$ is the floor or a square, the \emph{source} (the square last passed); $d=n_X$ is the direction; and $g\ge0$
is the \emph{cumulative gap}, the sum of the Euclidean lengths of the gap segments travelled so far. The initial state
is $((x,0),e_2,X_0,0)$.

A \emph{ray step} from the state $(p,n_X,X,g)$ uses $r(t):=p+tn_X$, $t\ge0$, and
\begin{gather*}
t_S:=\min\{t\ge0:\ r(t)\in Y\text{ for some square }Y\in\mathcal P,\ Y\ne X\}\quad(+\infty\text{ if there is none}),\\
t_W:=\min\{t\ge0:\ r(t)_x\in\{0,k\}\}\quad(+\infty\text{ if there is none}),\qquad t_D:=\delta-g,\qquad
t_H:=\frac{h_F-p_y}{\cos\varphi_X},\\
t^*:=\min(t_S,t_W,t_D,t_H).
\end{gather*}
(The minimum defining $t_S$ exists because the squares are closed and finite in number.) The first of the following
rules that applies is applied; this is the \emph{priority} D $>$ W $>$ contact $>$ H.
\begin{enumerate}
\item[(D)] If $t_D=t^*$, the path terminates at $r(t_D)$ (\emph{death}, D).
\item[(W)] Otherwise, if $t_W=t^*$, the path terminates at $r(t_W)$ (\emph{wall}, W).
\item[(C)] Otherwise, if $t_S=t^*$, put $q:=r(t_S)$ and $g\leftarrow g+t_S$; let $Y$ be the square containing $q$
(it is unique, as the squares are disjoint). If $q\in\relint\Bot(Y)$, then $(Y,q)$ is an \emph{entry candidate} and is
processed by rule~R1 (Definition~\ref{def:R1}) and then rule~R2 (Definition~\ref{def:R2}). Otherwise the path
terminates at $q$ (\emph{end collision}, E).
\item[(H)] Otherwise ($t_H<\min(t_S,t_W,t_D)$) the path terminates at $r(t_H)$; it is \emph{completed} (H).
\end{enumerate}
A point $q=r(t_S)$ with $t^*=t_S$ is called a \emph{first contact point}; it is a contact point in case (C), and the
end point of a death or wall termination if $t_D=t_S$ or $t_W=t_S$.
\end{definition}

Since contact has priority over H, a contact at height exactly $h_F$ is processed; since D has priority over contact,
a processed contact has $g+t_S<\delta$ (Lemma~\ref{lem:flows-traj}(f)). The rules R1 and R2 are defined below in the
order in which they can be stated: first rule R2 and the passing rule, together with the \emph{free paths}, which ignore
R1 (Definition~\ref{def:R2}); then, by a recursion over the squares that uses the free paths, rule R1 and the paths of
$F$ (Definition~\ref{def:R1}). The terminations M (merge) and T (tilt) are introduced there. The terminations D, W, E,
M, T are \emph{losses}; H is not. (The letter $d$ for the direction of a ray, with components $d_x$, $d_y$, is not to
be confused with the function $d(y)=\min(y,k-y)$ of Definition~\ref{def:flows-params}.)

\begin{definition}[rule R2 and passing; free paths]\label{def:R2}
\emph{Rule R2 (tilt).} If a path is the winner at an entry candidate $(Y,q)$ (Definition~\ref{def:R1}) and
$a(Y)\ge\alpha_F$, it terminates at $q$ (\emph{tilt termination}, T). A contact processed by rule~(C) at the relative
interior of a left or right side, or at a vertex, is an end collision whatever $a(Y)$ is; tilt terminations occur only
at entries through $\relint\Bot(Y)$.

\emph{Passing.} If the winner has $a(Y)<\alpha_F$, it \emph{passes} $Y$: it traverses the \emph{pass segment}
$[q,q+n_Y]$, its state becomes $(q+n_Y,n_Y,Y,g)$ (the length travelled inside $Y$ is not added to $g$), and if
$(q+n_Y)_y\ge h_F$ it terminates at the \emph{exit point} $q+n_Y$ with H; otherwise a new ray step begins.

\emph{Free paths.} The \emph{free path} $\pi^f_x$ is defined by Definition~\ref{def:flow} and the two rules above,
every entry candidate being treated as won (rule~R1 is ignored).
\end{definition}

\begin{lemma}[top sides are not hit first]\label{lem:flows-top}
In every ray step of a free path or of a path of a flow, the first contact point is not in $\relint\Top(Y)$ for any
square $Y$. Consequently an end collision occurs at a point of the relative interior of the left or the right side of
a square, or at a vertex.
\end{lemma}

\begin{proof}
Let the step start at $p$ with source $X$, so $d=n_X$ with $|\varphi_X|<\pi/4$ ($X$ is the floor, or a square passed
by the path, so $a(X)<\alpha_F<\pi/4$). Suppose $q=r(t_S)=c_Y+\frac12n_Y+\xi u_Y$ with $|\xi|<\frac12$. If $t_S=0$, then
$q=p$; if $X$ is a square, $p\in X$ and $X\cap Y=\emptyset$, which is impossible; if $X$ is the floor, $q=(x,0)$ has
height $0$, so it is a lowest point of $Y\subset[0,k]^2$, whereas the points of $\Top(Y)$ have heights $>v_Y$. So
$t_S>0$. As $|\varphi_Y-\varphi_X|<\pi/4+\pi/4=\pi/2$, we have $n_Y\cdot d=\cos(\varphi_Y-\varphi_X)>0$ by
\eqref{eq:flows-frame}. For $0<t_S-t$ small, $r(t)=q-(t_S-t)d$ satisfies
$(r(t)-c_Y)\cdot n_Y=\frac12-(t_S-t)\cos(\varphi_Y-\varphi_X)\in(-\frac12,\frac12)$ and
$|(r(t)-c_Y)\cdot u_Y|=|\xi-(t_S-t)\,d\cdot u_Y|<\frac12$, so $r(t)\in Y$ with $t<t_S$, contradicting the minimality of
$t_S$. The second statement follows since $\partial Y$ is the disjoint union of the relative interiors of the four
sides and the four vertices.
\end{proof}

\begin{lemma}[DAG lemma]\label{lem:DAG}
Let a path (free, or of a flow) pass a square $X$, entering it at $q_X\in\relint\Bot(X)$, and let its next contact be a
point $q_Y\in\relint\Bot(Y)$ of a square $Y$. Then
\[
c_{Y,y}-c_{X,y}\ \ge\ \tfrac12\big(\cos a(X)-\sin a(X)\big)+\tfrac12\big(\cos a(Y)-\sin a(Y)\big)\ >\ 0 .
\]
Consequently the squares of the successive entry candidates of a path have strictly increasing centre heights. In
particular they are distinct, a path has at most one entry candidate at each square, and two squares with the same
centre height never both carry entry candidates of one path.
\end{lemma}

\begin{proof}
Write $q_X=c_X-\frac12n_X+\tau_Xu_X$ with $|\tau_X|<\frac12$. Then
$c_{X,y}=q_{X,y}+\frac12\cos\varphi_X-\tau_X\sin\varphi_X\le q_{X,y}+\frac12(\cos a(X)+\sin a(X))$. In the same way
$c_{Y,y}=q_{Y,y}+\frac12\cos\varphi_Y-\tau_Y\sin\varphi_Y\ge q_{Y,y}+\frac12(\cos a(Y)-\sin a(Y))$. The path goes from
$q_X$ to the exit point $q_X+n_X$ and then a distance $t\ge0$ in the direction $n_X$, so
$q_{Y,y}=q_{X,y}+(1+t)\cos\varphi_X\ge q_{X,y}+\cos a(X)$. Combining the three relations gives the inequality. Its first
term is positive because $a(X)<\alpha_F<\pi/4$, and the second is $\ge0$ because $a(Y)\le\pi/4$. The consequences follow
by induction along the path (the first entry candidate comes from the floor, and every later one is the next contact
after a pass).
\end{proof}

\begin{lemma}[branch maps]\label{lem:flows-branch}
Let $\sigma=(Z_1,\dots,Z_m)$ be a sequence of distinct squares with $a(Z_i)<\pi/4$ for $i<m$, and put $Z_0:=X_0$. Let
$\tau^\sigma_1(x)$ be the $u_{Z_1}$-coordinate (relative to $c_{Z_1}$) of the intersection of the vertical line
$\{x\}\times\R$ with the line of $\Bot(Z_1)$, and, for $2\le j\le m$, let $\tau^\sigma_j(x)$ be the $u_{Z_j}$-coordinate of
the intersection of the line of $\Bot(Z_j)$ with the line through $c_{Z_{j-1}}+\frac12n_{Z_{j-1}}+\tau^\sigma_{j-1}(x)u_{Z_{j-1}}$
in the direction $n_{Z_{j-1}}$. Then each $\tau^\sigma_j:\R\to\R$ is an affine bijection with slope
$\prod_{i=1}^{j}\sec(\varphi_{Z_i}-\varphi_{Z_{i-1}})\ge1$. If the free path of $x$ passes $Z_1,\dots,Z_{m-1}$ as its
first $m-1$ passes and its next contact is a point $q\in\relint\Bot(Z_m)$, then
$q=c_{Z_m}-\frac12n_{Z_m}+\tau^\sigma_m(x)\,u_{Z_m}$.
\end{lemma}

\begin{proof}
For $i\le m$ we have $|\varphi_{Z_{i-1}}|<\pi/4$ and $\varphi_{Z_i}\in(-\pi/4,\pi/4]$, so
$|\varphi_{Z_i}-\varphi_{Z_{i-1}}|<\pi/2$. Let $L_1=\{p_0+su_1:s\in\R\}$ and $L_2=\{c_2+ru_2:r\in\R\}$ be lines with unit
directions $u_i=(\cos\varphi_i,\sin\varphi_i)$, $|\varphi_2-\varphi_1|<\pi/2$, and put $n_i:=(-\sin\varphi_i,\cos\varphi_i)$.
Since $n_1\cdot n_2=\cos(\varphi_2-\varphi_1)\ne0$, the line through $p_0+su_1$ in the direction $n_1$ meets $L_2$ in exactly
one point $c_2+ru_2$; taking the inner product of $p_0+su_1+tn_1=c_2+ru_2$ with $u_1$ gives
$(p_0-c_2)\cdot u_1+s=r\cos(\varphi_2-\varphi_1)$, so $r$ is an affine function of $s$ with slope
$\sec(\varphi_2-\varphi_1)\ge1$. The step from the floor is the case $L_1=\R\times\{0\}$, $s=x$, $n_1=e_2$; the passage
from the line of $\Bot(Z)$ to the line of $\Top(Z)$ by $v\mapsto v+n_Z$ preserves the $u_Z$-coordinate relative to
$c_Z$. Composing gives the slope. For the last statement: the first ray step of the free path is vertical, so its first
entry point lies on $\{x\}\times\R$ and on the line of $\Bot(Z_1)$; the exit point is the entry point plus $n_{Z_1}$,
with the same $u_{Z_1}$-coordinate; the next contact lies on the line through the exit point in the direction
$n_{Z_1}$; and so on by induction.
\end{proof}

\begin{definition}[rule R1 and the paths of a flow]\label{def:R1}
Number the squares of $\mathcal P$ as $Y_1,\dots,Y_N$ so that $c_{Y_1,y}\le c_{Y_2,y}\le\dots\le c_{Y_N,y}$ (squares with
equal centre heights in a fixed arbitrary order). For $i=1,\dots,N$ in turn and for every $q\in\relint\Bot(Y_i)$ let
$\mathcal A_{Y_i}(q)$ be the set of $x\in(0,k)$ such that $(Y_i,q)$ is an entry candidate of the free path
$\pi^f_x$ and $x$ is the winner at every entry candidate of $\pi^f_x$ preceding $(Y_i,q)$; if
$\mathcal A_{Y_i}(q)\ne\emptyset$, the \emph{winner} $w(Y_i,q)$ is the element $x\in\mathcal A_{Y_i}(q)$ for which the pair
$(g_x(q),x)$ is lexicographically smallest, where $g_x(q)$ is the cumulative gap of $\pi^f_x$ at $q$ (after the update
in rule~(C)). We say that $x$ is the winner at $(Y,q)$ if $x=w(Y,q)$.

The path $\pi_x$ of $F$ is $\pi^f_x$ if $x$ is the winner at every entry candidate of $\pi^f_x$; otherwise it is the
initial part of $\pi^f_x$ up to the first entry candidate $(Y,q)$ at which $x$ is not the winner, and it terminates
there at $q$ (\emph{merge}, M). In words, \emph{rule R1 (merge)}: among the paths that reach an entry candidate $(Y,q)$
alive, the one with the lexicographically smallest $(g,x)$ continues, and the others terminate at $q$ with M; then
rule~R2 is applied to the winner.
\end{definition}

\begin{lemma}[R1 is well defined]\label{lem:flows-R1}
\begin{enumerate}[label=(\alph*)]
\item For every square $Y$ and every $q\in\relint\Bot(Y)$ the set $\mathcal A_Y(q)$ is finite.
\item The recursion of Definition~\ref{def:R1} is well defined: at stage $i$ it uses only winners at squares $Y_l$ with
$l<i$. A path of $F$ reaches an entry candidate $(Y,q)$ without having terminated if and only if $x\in\mathcal A_Y(q)$, and
up to its termination $\pi_x$ coincides with $\pi^f_x$; thus the paths of $F$ obey Definition~\ref{def:flow} with
entry processing R1 followed by R2.
\item At most one path of $F$ passes, or is T-terminated at, a given entry point $(Y,q)$.
\end{enumerate}
\end{lemma}

\begin{proof}
(a) Let $x\in\mathcal A_Y(q)$. Before $q$, the free path $\pi^f_x$ passes squares $Z_1,\dots,Z_{m-1}$ with
$a(Z_i)<\alpha_F<\pi/4$, which are distinct by Lemma~\ref{lem:DAG}, and $q$ is its next contact. By
Lemma~\ref{lem:flows-branch}, $q=c_Y-\frac12n_Y+\tau^\sigma_m(x)u_Y$ with $\sigma=(Z_1,\dots,Z_{m-1},Y)$, and as
$\tau^\sigma_m$ is a bijection, $x$ is determined by $\sigma$ and $q$. There are only finitely many sequences of
distinct squares, so $\mathcal A_Y(q)$ is finite and the lexicographic minimum exists (distinct $x$ give distinct
pairs).

(b) By Lemma~\ref{lem:DAG}, an entry candidate of $\pi^f_x$ preceding $(Y_i,q)$ lies on a square $Y_l$ with
$c_{Y_l,y}<c_{Y_i,y}$, hence $l<i$; so the winners used at stage $i$ have been defined before. The rules D, W, E, H, R2
and the passing rule depend only on the state of the path itself; rule R1 stops losers and does not change the motion
of winners. Hence, by induction along the path, $\pi_x$ coincides with $\pi^f_x$ until it terminates, and it reaches
$(Y,q)$ without having terminated exactly when $x$ won all earlier candidates, i.e.\ when $x\in\mathcal A_Y(q)$.

(c) A path that passes, or is T-terminated at, $(Y,q)$ is the winner $w(Y,q)$, which is unique.
\end{proof}

The order in which R1 and R2 are applied at an entry candidate is discussed in Remark~\ref{rem:flows-order} at the
end of this section.

\subsection{Trajectories}\label{ssec:flows-traj}

\begin{definition}[trajectories]\label{def:flows-traj}
Let $\pi_x$ be a path of a flow $F$ (or a free path). Its \emph{trajectory} $\Gamma_x$ is the polygonal line formed,
in order, by its \emph{gap segments} $[p,p+t^*n_X]$, one for each ray step (of length $t^*\ge0$), and its pass
segments $[q,q+n_Y]$, one for each passed square. The last point of $\Gamma_x$ is the \emph{termination point}
$e(x)$ (for a path completed after a pass, the exit point). The \emph{termination height} is $\tau(x):=e(x)_y$, except
that $\tau(x):=h_F$ if $\pi_x$ is completed (H). The endpoints of the segments, including $(x,0)$ and $e(x)$, are the
\emph{event points}. The path is \emph{alive at height $y$} if $\tau(x)>y$. Each path has exactly one
\emph{termination type} $\star\in\{\mathrm T,\mathrm D,\mathrm E,\mathrm M,\mathrm W,\mathrm H\}$.
\end{definition}

\begin{lemma}[trajectories]\label{lem:flows-traj}
Let $\pi_x$ be a path of a flow $F=F(\alpha_F,\delta,h_F)$ or a free path, $x\in(0,k)$.
\begin{enumerate}[label=(\alph*)]
\item Every segment of $\Gamma_x$ has a unit direction $d$ with $d_y\ge\cos\alpha_F>0$. Hence, if $z\ne z'$ are points
of $\Gamma_x$ and $z$ comes before $z'$, then $z_y<z'_y$.
\item Every ray step starts at a point $p$ with $p_x\in(0,k)$, $p_y<h_F$ and cumulative gap $g<\delta$. If its source
$X$ is a square, then $p\in\relint\Top(X)$, $t_S>0$ and $t^*>0$.
\item The open part $\{p+tn_X:\ 0<t<t^*\}$ of a gap segment lies in the waste.
\item If $t_S<\infty$, the point $r(t_S)$ lies in exactly one square $Y$, and $r(t_S)\in\partial Y$.
\item A pass segment $[q,q+n_Y]$ lies in $Y$ and meets no other square; its open part lies in $\operatorname{int}Y$;
$q\in\relint\Bot(Y)$ and $q+n_Y\in\relint\Top(Y)$.
\item At a processed contact the cumulative gap is $g+t_S<\delta$. If a step ends with a death at a first contact
point ($t_D=t_S$), the cumulative gap there equals $\delta$.
\item No event point lies in the interior of a square. The points of $\Gamma_x$ that lie in a square $Z$ are first
contact points on $\partial Z$, and the points of the pass segment of $Z$ if $\pi_x$ passes $Z$; in the latter case the
initial point of that pass segment is a first contact point.
\item $\pi_x$ has at most $h_F/\cos\alpha_F+1$ passes and at most $h_F/\cos\alpha_F+2$ ray steps.
\end{enumerate}
\end{lemma}

\begin{proof}
(a) The gap segments have directions $n_X$, where $X$ is the floor ($\cos\varphi_X=1$) or a passed square
($a(X)<\alpha_F$), and the pass segments have directions $n_Y$ with $Y$ passed; by \eqref{eq:flows-frame},
$n\cdot e_2=\cos\varphi>\cos\alpha_F$. Along a segment of positive length the height is therefore strictly increasing.

(f) At the start of every ray step $g<\delta$: initially $g=0$, and a new step begins only after a pass, which
follows a processed contact. A processed contact has $t_S=t^*$ and $t_D\ne t^*$ (otherwise D applies), so
$t_D>t_S$, i.e.\ $g+t_S<\delta$. If $t_D=t_S=t^*$, rule D applies at $r(t_D)$, where the cumulative gap is $g+t_D=\delta$.

(b) The first step starts at $(x,0)$ with $x\in(0,k)$ and $g=0$. Every later step starts at an exit point $p=q+n_Y$
of a pass that did not end with H, so $p_y<h_F$, $p\in\relint\Top(Y)$ (by (e)), whose abscissas lie in $(0,k)$, and
$g<\delta$ by (f). If $X$ is a square, then $p\in X$ and $p\notin Y'$ for every square $Y'\ne X$; as the set of $t\ge0$
with $r(t)\in\bigcup_{Y'\ne X}Y'$ is closed, $t_S>0$. Moreover $t_D=\delta-g>0$, $t_W>0$ (as $p_x\in(0,k)$) and
$t_H=(h_F-p_y)/\cos\varphi_X>0$, so $t^*>0$.

(c) Let $0<t<t^*$. As $t<t^*\le t_S$, $r(t)$ lies in no square $Y\ne X$; if $X$ is a square,
$(r(t)-c_X)\cdot n_X=\frac12+t>\frac12$, so $r(t)\notin X$. By (b), $p_x\in(0,k)$, and $t<t^*\le t_W$ gives
$r(t)_x\in(0,k)$. Finally $r(t)_y>p_y\ge0$ and $r(t)_y<p_y+t^*\cos\varphi_X\le p_y+t_H\cos\varphi_X=h_F<k$. Hence $r(t)$
lies in $(0,k)^2$ and in no square.

(d) $r(t_S)$ lies in a square by the definition of $t_S$, and in only one since the squares are disjoint. If $t_S>0$,
then $r(t)\notin Y$ for $0\le t<t_S$, so $r(t_S)$ is a limit of points outside $Y$ and lies on $\partial Y$. If $t_S=0$,
the source is the floor by (b), $r(0)=(x,0)$ has height $0$, and no point of height $0$ is interior to a square in
$[0,k]^2$.

(e) $q=c_Y-\frac12n_Y+\xi u_Y$ with $|\xi|<\frac12$, so $q+tn_Y=c_Y+\xi u_Y+(t-\frac12)n_Y$, which lies in $Y$ for
$0\le t\le1$ and in $\operatorname{int}Y$ for $0<t<1$; $q+n_Y\in\relint\Top(Y)$. Other squares are disjoint from $Y$.

(g) The event points are: $(x,0)$, of height $0$, not interior to any square; entry points and exit points, on
$\partial Y$ by (e); end collision points, first contact points on $\partial Y$ by (d); death points $r(t_D)$, which lie
in no square if $t_D<t_S$ (as in (c)) and are first contact points if $t_D=t_S$; wall points, whose abscissa is $0$ or
$k$; and the end points of completions, which are either gap points $r(t_H)$ with $t_H<t_S$ (in no square, as in (c))
or exit points. For the second statement: points of open gap parts lie in the waste by (c); points of the pass segment
of a square $Y\ne Z$ lie in $Y$, hence not in $Z$; the start point of a ray step is $(x,0)$ (covered above, if in $Z$ it
is $r(0)$ with $t_S=0$, a first contact point) or an exit point of a passed square $X$ (if it lies in $Z$, then $Z=X$
and it is a point of the pass segment of $Z$); the end point $r(t^*)$ of a ray step lies in no square if $t^*<t_S$ and
is a first contact point if $t^*=t_S$. Finally, a pass segment of $Z$ starts at the entry point, which is the first
contact point of the preceding ray step.

(h) An entry candidate is processed only if $t_S\le t_H$, i.e.\ at height $\le h_F$. Each pass raises the height by
$\cos\varphi_Y\ge\cos\alpha_F$ and the gap segments do not decrease it, so the $j$-th entry point has height
$\ge(j-1)\cos\alpha_F$. Hence $(j-1)\cos\alpha_F\le h_F$ for every entry, which gives the bound on passes; every ray step
except the first follows a pass.
\end{proof}

\subsection{Piecewise-affine structure}\label{ssec:flows-structure}

\begin{lemma}[piecewise-affine structure]\label{lem:structure}
Let $F=F(\alpha_F,\delta,h_F)$ be a flow. There is a finite set $\Sigma_F\subset[0,k]$ containing $0$ and $k$ such
that on every component $I$ of $[0,k]\setminus\Sigma_F$ (a \emph{piece} of $F$):
\begin{enumerate}[label=(\alph*)]
\item the combinatorial type of $\pi_x$ is the same for all $x\in I$: the number of ray steps, the source of each step,
the event that ends each step (D, W, contact, H), for each contact the square met and the side (relative interior of
the bottom, left or right side) that contains the contact point, for each entry candidate whether $x$ is the winner and
whether it then passes or is T-terminated, and the termination type;
\item the start point $p_j(x)$ and the start gap $g_j(x)$ of the $j$-th ray step, its length $t^*_j(x)$, its contact
point $q_j(x)$ (if it ends with a contact) and the termination point $e(x)$ are affine functions of $x\in I$;
\item for every $x\in I$, the line that carries a ray step of $\pi_x$ contains no vertex of a square. In particular no
first contact point of $\pi_x$ (Definition~\ref{def:flow}) is a vertex of a square: for $x\in I$ no entry candidate, no
end collision and no D- or W-termination at a first contact point ($t_D=t_S$ or $t_W=t_S$) lies at a vertex.
Consequently, first contact points at vertices of squares occur only for the finitely many $x\in\Sigma_F$.
\end{enumerate}
The same holds for the free paths, with a finite set $\Sigma^f_F\subset\Sigma_F$ and without the R1 outcomes in (a).
\end{lemma}

\begin{lemma}[expansion]\label{lem:expansion}
On a piece $I$ of Lemma~\ref{lem:structure} (for $F$ or for the free paths) let $X_j$ be the source of the $j$-th ray
step ($X_0$ the floor).
\begin{enumerate}[label=(\alph*)]
\item $p_j'(x)=\lambda_ju_{X_j}$ with $\lambda_0=1$ and $\lambda_j=\prod_{i=1}^j\sec(\varphi_{X_i}-\varphi_{X_{i-1}})\ge1$.
\item If the $j$-th step ends with a contact at the relative interior of $\Bot(Y)$, then $q_j'(x)=\mu_ju_Y$ with
$\mu_j=\lambda_j/\cos(\varphi_Y-\varphi_{X_j})\ge\lambda_j\ge1$; if the path passes $Y$, then $X_{j+1}=Y$ and
$\lambda_{j+1}=\mu_j$.
\item If the $j$-th step ends with a contact at the relative interior of the left or the right side of $Y$, then
$q_j'(x)=\nu n_Y$ with $\nu\ne0$. (The derivative factors $\mu_j$ and $\nu$ are unrelated to the measures $\mu_Z(y)$
and $\nu_y$ of Definition~\ref{def:flows-line} and Lemma~\ref{lem:flows-density}.)
\item For $y\in\R$, on the set of $x\in I$ with $p_{j,y}(x)<y<p_{j,y}(x)+t^*_j(x)\cos\varphi_{X_j}$, the abscissa of the
point of height $y$ on the $j$-th gap segment is an affine function of $x$ with derivative
$\lambda_j/\cos\varphi_{X_j}\ge1$. If the path passes $Y$ after step $j$, then on the set of $x\in I$ with
$q_{j,y}(x)<y<q_{j,y}(x)+\cos\varphi_Y$ the abscissa of the point of height $y$ on the pass segment of $Y$ is affine
with derivative $\mu_j/\cos\varphi_Y\ge1$.
\end{enumerate}
\end{lemma}

\begin{proof}[Proof of Lemmas~\ref{lem:structure} and~\ref{lem:expansion}]
\emph{Step 1: one ray step of the free paths.} Let $I$ be an open interval and $j\ge0$, and suppose that for every
$x\in I$ the free path $\pi^f_x$ has a $j$-th ray step, with the same source $X$, with start point $p(x)$ and start gap
$g(x)$ affine in $x$, and with $p'(x)=\lambda u_X$, $\lambda\ge1$. Put $d:=n_X$; since $p'\perp d$, the lines
$\{p(x)+td:t\in\R\}$, $x\in I$, are parallel and are indexed by $\xi(x):=p(x)\cdot u_X$, an affine function with slope
$\lambda>0$.

(i) For each square $Y\ne X$ the lines that meet $Y$ are those with $\xi$ between the least and the largest of the four
numbers $v\cdot u_X$, $v$ a vertex of $Y$. Let $C_1$ be the finite set of $x\in I$ with $\xi(x)=v\cdot u_X$ for some
vertex $v$ of some square. On a component $I_1$ of $I\setminus C_1$ and for each $Y\ne X$, either no line meets $Y$, or
every line meets the interior of $Y$ in a non-degenerate segment that contains no vertex of $Y$. In the latter case the
initial point of the segment (in the direction $d$) lies in the relative interior of a side $\sigma_Y$ of $Y$; it moves
continuously with $x$ and never meets a vertex, so $\sigma_Y$ is the same for all $x\in I_1$. The line is not parallel
to $\sigma_Y$ (a segment along $\sigma_Y$ would begin at a vertex), so the parameters $t_Y(x)\le t'_Y(x)$ of the initial
and final points are affine functions of $x$.

(ii) If $X$ is a square, $p(x)\in X$ lies in no other square, so $0\notin[t_Y(x),t'_Y(x)]$; the open sets
$\{t_Y>0\}$ and $\{t'_Y<0\}$ cover $I_1$, which is connected, so one of them is $I_1$. If $X$ is the floor, then
$d=e_2$, $p(x)=(x,0)$, and $t_Y(x)\ge0$ is the height of the lowest point of $Y$ on the vertical line through $(x,0)$.
In both cases the set $\mathcal Y$ of squares met at some $t\ge0$ is the same for all $x\in I_1$, and
$t_S(x)=\min_{Y\in\mathcal Y}t_Y(x)$ ($+\infty$ if $\mathcal Y=\emptyset$).

(iii) For $x\in I_1$ and distinct $Y,Y'\in\mathcal Y$, $t_Y(x)\ne t_{Y'}(x)$, since otherwise the point
$p(x)+t_Y(x)d$ would lie in $Y\cap Y'$. Hence the sets $\{x\in I_1:t_Y(x)<t_{Y'}(x)\ \forall Y'\in\mathcal Y\setminus\{Y\}\}$,
$Y\in\mathcal Y$, are open, disjoint and cover $I_1$; one of them, for $Y=Y_S$, is $I_1$. Thus the first square met and
the side $\sigma:=\sigma_{Y_S}$ are constant on $I_1$, and $t_S=t_{Y_S}$ is affine.

(iv) $t_D=\delta-g$ and $t_H=(h_F-p_y)/\cos\varphi_X$ are affine. If $d_x=0$ then $t_W\equiv+\infty$ (as $p_x\in(0,k)$
by Lemma~\ref{lem:flows-traj}(b)); if $d_x>0$ then $t_W=(k-p_x)/d_x$, and if $d_x<0$ then $t_W=-p_x/d_x$; in each case
$t_W$ is affine or $+\infty$.

(v) Let $C_2$ be the set of zeros in $I_1$ of those differences of two of the functions $t_S,t_W,t_D,t_H$ that are finite
and not identically zero; it is finite. On a component $I_2$ of $I_1\setminus C_2$ all order relations ($<$, $=$) among
the four functions are constant, so the rule applied (D, W, contact, H) is the same for all $x\in I_2$, and $t^*$ is
affine.

(vi) By (i), for $x\in I_2$ the line $\{p(x)+td:t\in\R\}$ contains no vertex of a square, and if $t_S(x)<\infty$, the
first contact point $p(x)+t_S(x)d$ is the initial point of the segment in which this line meets $Y_S$, so it lies in
$\relint\sigma$ and is not a vertex; this holds whichever rule (D, W or contact) is applied at it. If the rule is
contact, $q(x)=p(x)+t_S(x)d$ is affine and lies in $\relint\sigma$; by Lemma~\ref{lem:flows-top}, $\sigma$ is the
bottom, the left or the right side of $Y_S$, the same for all $x\in I_2$. Whether $a(Y_S)\ge\alpha_F$ is fixed. In the case of a pass, let $C_3$ be the
zero (if any) of the affine function $q_y(x)+\cos\varphi_{Y_S}-h_F$ if it is not identically zero; on each component of
$I_2\setminus C_3$ it is fixed whether the pass ends with H, and otherwise the next step has source $Y_S$, start point
$q(x)+n_{Y_S}$ and start gap $g(x)+t_S(x)$, all affine.

(vii) \emph{Derivatives.} If $\sigma=\Bot(Y)$ ($Y=Y_S$), its line has direction $u_Y$, so $q'=\mu u_Y$ for some $\mu$;
on the other hand $q'=p'+t_S'd=\lambda u_X+t_S'n_X$. Taking the inner product with $u_X$ and using
\eqref{eq:flows-frame} gives $\mu\cos(\varphi_Y-\varphi_X)=\lambda$. As $|\varphi_X|<\pi/4$ and
$\varphi_Y\in(-\pi/4,\pi/4]$, $\cos(\varphi_Y-\varphi_X)\in(0,1]$, so $\mu=\lambda/\cos(\varphi_Y-\varphi_X)\ge\lambda$.
After a pass the next start point is $q+n_Y$, with derivative $\mu u_Y$. If $\sigma$ is a left or right side, its line
has direction $n_Y$, so $q'=\nu n_Y$, and $\nu\ne0$ because $q'\cdot u_X=\lambda\ne0$. For the crossings: the point of
height $y$ on the gap segment is $p+tn_X$ with $t=(y-p_y)/\cos\varphi_X$, of abscissa $p_x-(y-p_y)\tan\varphi_X$, whose
derivative is $\lambda\cos\varphi_X+\lambda\sin\varphi_X\tan\varphi_X=\lambda/\cos\varphi_X\ge1$; the point of height $y$ on
the pass segment of $Y$ has abscissa $q_x-(y-q_y)\tan\varphi_Y$, with derivative
$\mu\cos\varphi_Y+\mu\sin\varphi_Y\tan\varphi_Y=\mu/\cos\varphi_Y\ge1$.

\emph{Step 2: the free paths.} Start with $I=(0,k)$, $j=0$, $X=X_0$, $p(x)=(x,0)$, $g=0$, $\lambda=1$, and apply Step~1
repeatedly. On each resulting interval the paths either terminate in the $j$-th step or pass and satisfy the hypotheses
of Step~1 for $j+1$ with $\lambda_{j+1}=\mu$. By Lemma~\ref{lem:flows-traj}(h) the number of steps is bounded, so after
finitely many subdivisions we obtain a finite set $\Sigma^f_F\ni0,k$ with properties (a)--(c) for the free paths, and
Step~1(vii) gives Lemma~\ref{lem:expansion} for them (with $\lambda_j$ as stated, by induction).

\emph{Step 3: the outcomes of R1.} For a free piece $I$ (a component of $[0,k]\setminus\Sigma^f_F$) and a square $Y$
carrying an entry candidate of the common type, write the contact point as $q^I(x)=c_Y-\frac12n_Y+\tau^I(x)u_Y$, where
$\tau^I$ is affine with slope $\ge1$ (Step~1(vii)), and let $g^I(x)$ be the cumulative gap at $q^I(x)$, also affine.
We show by induction over the stages $i=1,\dots,N$ of Definition~\ref{def:R1} that for every free piece $I$ the sets
$A_{Y_i}(I):=\{x\in I:\ x\in\mathcal A_{Y_i}(q^I(x))\}$ and $L_{Y_i}(I):=\{x\in A_{Y_i}(I):\ x\ne w(Y_i,q^I(x))\}$ are finite
unions of intervals and points (empty if $Y_i$ carries no candidate of the type of $I$); these two families of sets
are used only in this proof. Assume this for all stages
$l<i$. The candidates of the type of $I$ that precede the one at $Y_i$ lie on squares $Y_l$ with $l<i$
(Lemma~\ref{lem:DAG}), so $A_{Y_i}(I)=I\setminus\bigcup_lL_{Y_l}(I)$, a finite union of intervals and points. A point
$x\in A_{Y_i}(I)$ is in $L_{Y_i}(I)$ if and only if some $x'\ne x$ in $\mathcal A_{Y_i}(q^I(x))$ has
$(g_{x'},x')<_{\rm lex}(g^I(x),x)$. Such an $x'$ either lies in the finite set $\Sigma^f_F$, and since $\tau^I$ is
injective and $x'$ has at most one candidate at $Y_i$, it affects at most one $x\in I$; or it lies in a free piece $I'$
(possibly $I'=I$) with $x'\in A_{Y_i}(I')$ and $q^{I'}(x')=q^I(x)$, i.e.\ $x'=\psi(x):=(\tau^{I'})^{-1}(\tau^I(x))$, an
affine map defined on the interval $\{x\in I:\tau^I(x)\in\tau^{I'}(I')\}$. Therefore $L_{Y_i}(I)$ is the union of
finitely many points and of the finitely many sets
\[
\big\{x\in A_{Y_i}(I)\cap\operatorname{dom}\psi:\ \psi(x)\in A_{Y_i}(I'),\ \big(g^{I'}(\psi(x)),\psi(x)\big)<_{\rm lex}\big(g^I(x),x\big)\big\},
\]
each defined by finitely many affine equalities and inequalities in $x$, hence a finite union of intervals and points.
(For $I'=I$ we have $\psi=\mathrm{id}$ and the strict inequality fails.)

\emph{Step 4: the flow.} Let $\Sigma_F$ consist of $\Sigma^f_F$ and the endpoints of the intervals and the isolated
points of all the sets $L_Y(I)$ and $A_Y(I)$. On each component of $[0,k]\setminus\Sigma_F$ all outcomes of R1 are constant, so by
Definition~\ref{def:R1} and Lemma~\ref{lem:flows-R1}(b) the path $\pi_x$ is the free path truncated at the same
candidate (or not truncated) for all $x$; (a)--(c) and Lemma~\ref{lem:expansion} follow from the free case (the
termination point of an M-terminated path is the affine contact point).
\end{proof}

\subsection{The master flow, the scale flows and the ceiling flows}\label{ssec:flows-family}

\begin{definition}[master flow]\label{def:F0}
Assume \textup{(C0)} and \textup{(C1)}. The \emph{master flow} is $\Fz:=F(\amax,\delta,\hmax)$ with $\amax=\alpha(y_0)$
and $\hmax=\frac k2-1$; by (C0) and Lemma~\ref{lem:flows-constraints}(a), $0<\amax\le1.5\cdot10^{-6}<\pi/4$, and
$0<\hmax<k$ as $k\ge4$, so this is a flow. Its
paths are $\pi^0_x$, its trajectories $\Gamma^0_x$, its termination heights $\tau_0(x)$. Rule R1 is decided only in
$\Fz$. The \emph{R1-passed entries} of $\pi^0_x$ are the entry candidates of $\pi^0_x$ at which $x$ is the winner; we list
them in path order as $(Z_1,q_1),\dots,(Z_J,q_J)$, $J=J(x)\ge0$, and put $\eta_j:=(q_j)_y$. For $j<J$, $\pi^0_x$ passes
$Z_j$; $Z_J$ is passed or T-terminates $\pi^0_x$.
\end{definition}

\begin{definition}[scale flows]\label{def:Fs}
Assume \textup{(C0)}, \textup{(C1)} and \textup{(C7)}, and let $s\in[y_0,y_1]$. The path $\pi^s_x$ of the \emph{scale
flow} $\Fs$ is $\pi^0_x$ truncated at the first, in path order, of the following events.
\begin{enumerate}
\item[(a)] $T_s$: an R1-passed entry $(Z_j,q_j)$ with $a(Z_j)\ge\alpha(s)$; $\pi^s_x$ terminates at $q_j$ with T.
\item[(b)] $H_s$: either a ray step of $\pi^0_x$, with start point $p$, direction $d$ and parameters
$t_S,t_W,t_D$ (in $\Fz$), such that $t^{(s)}_H:=(s+2-p_y)/d_y<\min(t_S,t_W,t_D)$; then $\pi^s_x$ terminates at
$p+t^{(s)}_Hd$ with H; or a pass of $\pi^0_x$ whose exit point has height $\ge s+2$; then $\pi^s_x$ terminates at the
exit point with H.
\item[(c)] The termination of $\pi^0_x$ (D, W, E, M or T), if no event (a) or (b) precedes it; $\pi^s_x$ then
terminates at the same point with the same type.
\end{enumerate}
If $\pi^0_x$ is T-terminated at an R1-passed entry $(Z_j,q_j)$ that is not preceded by an event (a) or (b), then (a)
holds there as well, because $a(Z_j)\ge\amax\ge\alpha(s)$; such an \emph{inherited} T termination is a $T_s$
termination. Thus ``T-terminated in $\Fs$'' and ``$T_s$-terminated in $\Fs$'' mean the same. We put $h_{\Fs}:=s+2$,
write $\tau_s(x)$ for the termination height of $\pi^s_x$ ($\tau_s:=s+2$ for H), and
\[
\mathcal T_s:=\{x\in(0,k):\ \pi^s_x\text{ is T-terminated}\},\qquad N(s):=|\mathcal T_s| .
\]
\end{definition}

If a ray step of $\pi^0_x$ ends with H in $\Fz$, then $t^{(s)}_H\le t_H<\min(t_S,t_W,t_D)$ (as $s+2\le\hmax$), so (b)
applies; if $\pi^0_x$ is completed at an exit point, that exit point has height $\ge\hmax\ge s+2$, so (b) applies there
or earlier. Hence (c) never concerns a completion. A contact at height exactly $s+2$ has $t_S=t^{(s)}_H$, so (b) does not
apply in its step: as in $\Fz$, contact precedes H.

\begin{definition}[the tilt time]\label{def:flows-Tstar}
Let $(Z_j,q_j)$, $j\le J(x)$, be the R1-passed entries of $\pi^0_x$ and $\eta_j=(q_j)_y$. Put
\[
s_j:=\max\big(y_0,\ \eta_j-2,\ (c_0/a(Z_j))^2\big)\ \text{ if }a(Z_j)>0,\qquad s_j:=+\infty\ \text{ if }a(Z_j)=0,
\]
and $\Tstar(x):=\min_js_j$ if this minimum is $\le y_1$; otherwise $\Tstar(x)$ is undefined. Let $\operatorname{dom}\Tstar$
be the set of $x$ where $\Tstar(x)$ is defined. Only R1-passed entries enter: deaths, merges and end collisions do not.
\end{definition}

\begin{lemma}[scale flows]\label{lem:flows-Ts}
Assume \textup{(C0)}, \textup{(C1)}, \textup{(C7)}, let $x\in(0,k)$ and $s\in[y_0,y_1]$.
\begin{enumerate}[label=(\alph*)]
\item $\eta_1<\eta_2<\dots<\eta_J$. For $j\ge2$, every point of $\Gamma^0_x$ before $q_j$ has height $<\eta_j$; every point
before $q_1$ has height $\le\eta_1$.
\item $\pi^s_x$ is T-terminated if and only if there is $j$ with $a(Z_j)\ge\alpha(s)$ and $\eta_j\le s+2$; its termination
point is then $q_j$ for the least such $j$.
\item The trajectory of $\pi^s_x$ is an initial part of $\Gamma^0_x$. Every square passed by $\pi^s_x$ is passed by
$\pi^0_x$ and has inclination $<\alpha(s)$. Every D, W, E or M termination of $\Fs$ is a termination of $\Fz$ of the
same type at the same point. Every first contact point on the trajectory of $\pi^s_x$ has height $\le s+2$.
\item If $y_0\le s\le s'\le y_1$, then $\mathcal T_s\subset\mathcal T_{s'}$.
\item $\{x\in\operatorname{dom}\Tstar:\ \Tstar(x)\le s\}=\mathcal T_s$.
\item There is a finite set $\Sigma_s\supset\Sigma_{\Fz}$ such that on each component of $[0,k]\setminus\Sigma_s$ the
truncation event of $\pi^s_x$ (its type and the step at which it occurs) is constant and the termination point of
$\pi^s_x$ is affine in $x$; Lemma~\ref{lem:structure}(a)--(c) holds for $\Fs$ with these pieces. In particular
$\mathcal T_s$ and, for every $y$ and $\star$, the set $\{x:\pi^s_x\text{ is }\star\text{-terminated},\ \tau_s(x)<y\}$ are finite
unions of intervals and points.
\end{enumerate}
\end{lemma}

\begin{proof}
(a) For $j\ge2$, $q_j$ is the first contact of the ray step that starts at the exit point $p$ of $Z_{j-1}$; this step
has a square source, so $t_S>0$ (Lemma~\ref{lem:flows-traj}(b)) and $p_y<\eta_j$; all earlier points have height
$\le p_y$ by Lemma~\ref{lem:flows-traj}(a), and $\eta_{j-1}<p_y$. Before $q_1$ there is no entry candidate (a lost
candidate would have terminated the path), so the points before $q_1$ lie on the first, vertical, gap segment and have
height $\le\eta_1$.

(c) The first statement is the definition. A square passed by $\pi^s_x$ is an R1-passed entry $(Z_j,q_j)$ of $\pi^0_x$
at which (a) of Definition~\ref{def:Fs} does not hold, so $a(Z_j)<\alpha(s)$. D, W, E, M terminations of $\Fs$ come from
(c) of Definition~\ref{def:Fs}. Let $q$ be a first contact point on the trajectory of $\pi^s_x$, ending a ray step that
starts at $p$. If $p_y\ge s+2$, then $p$ (which has positive height) is an exit point of height $\ge s+2$, and $\pi^s_x$
would have terminated at $p$ by (b); so $p_y<s+2$. If $q_y>s+2$, then $t^{(s)}_H<t_S$, and $t_S=t^*\le\min(t_D,t_W)$, so
$t^{(s)}_H<\min(t_S,t_W,t_D)$ and $\pi^s_x$ would have terminated before $q$ by (b). Hence $q_y\le s+2$.

(b) ($\Leftarrow$) Let $j$ be least with $a(Z_j)\ge\alpha(s)$ and $\eta_j\le s+2$. Up to its first truncation, $\pi^s_x$
coincides with $\pi^0_x$, and $\pi^0_x$ reaches $q_j$ and wins there. No event (a) precedes $q_j$: an earlier R1-passed
entry $Z_i$ ($i<j$) has $\eta_i<\eta_j\le s+2$ by (a), hence $a(Z_i)<\alpha(s)$ by the minimality of $j$. No event (b)
precedes $q_j$: the exit points before $q_j$ have height $<\eta_j\le s+2$ by (a); every ray step before $q_j$ ends with a
contact at an entry $q_i$ of height $\eta_i\le\eta_j\le s+2$, so $t_S\le t^{(s)}_H$ in that step and (b) does not apply.
Hence $\pi^s_x$ reaches $q_j$ with the state of $\pi^0_x$, is the R1 winner there (R1 is decided in $\Fz$), and is
T-terminated at $q_j$ by (a). ($\Rightarrow$) If $\pi^s_x$ is T-terminated, its termination point is an R1-passed entry
$(Z_j,q_j)$ with $a(Z_j)\ge\alpha(s)$ (for an inherited termination, $a(Z_j)\ge\amax\ge\alpha(s)$), and $\eta_j\le s+2$ by
(c). By the first part, the termination point is $q_j$ for the least such $j$.

(d) As $\alpha(s')\le\alpha(s)$ and $s+2\le s'+2$, the condition in (b) for $s$ implies the condition for $s'$.

(e) For every $j$: if $a(Z_j)=0$, then $s_j=+\infty>s$ and $a(Z_j)<\alpha(s)$. If $a(Z_j)>0$, then, as $y_0\le s$,
$s_j\le s$ holds if and only if $\eta_j\le s+2$ and $(c_0/a(Z_j))^2\le s$, i.e.\ $a(Z_j)\ge c_0s^{-1/2}=\alpha(s)$. Hence
$\min_js_j\le s$ if and only if some $j$ satisfies the condition of (b), i.e.\ if and only if $x\in\mathcal T_s$. Finally,
$\min_js_j\le s\le y_1$ means that $x\in\operatorname{dom}\Tstar$ and $\Tstar(x)\le s$.

(f) On a piece $I$ of $\Fz$ the R1-passed entries $Z_j$ are fixed, and $\eta_j(x)$, the exit heights, $t^*_j(x)$ and
$t^{(s)}_H(x)$ (for each ray step) are affine. Let $\Sigma_s$ consist of $\Sigma_{\Fz}$ and the zeros in each piece of the
finitely many affine functions $\eta_j-(s+2)$, (exit height)$-(s+2)$ and $t^{(s)}_H-\min(t_S,t_W,t_D)$ that are not
identically zero. On each component of $[0,k]\setminus\Sigma_s$ the truth value of each of the conditions (a), (b), (c) of
Definition~\ref{def:Fs} at each event of the common type is constant, hence so is the first truncation; the termination
point is $q_j(x)$, $p_j(x)+t^{(s)}_H(x)d_j$, an exit point, or the termination point of $\pi^0_x$, all affine. Thus
Lemma~\ref{lem:structure}(a),(b) for $\Fs$ follow from the same statements for $\Fz$; and since the ray steps of
$\pi^s_x$ are initial parts of ray steps of $\pi^0_x$ (with the same lines) and each component of $[0,k]\setminus\Sigma_s$
lies in a piece of $\Fz$, Lemma~\ref{lem:structure}(c) for $\Fs$ follows from Lemma~\ref{lem:structure}(c) for $\Fz$.
The last statement follows since on each such component the type is fixed and $\tau_s$ is affine or the constant $s+2$.
\end{proof}

By Lemma~\ref{lem:flows-Ts}(c), Lemma~\ref{lem:flows-traj}(a) and (c)--(g) hold for the paths of $\Fs$ as well
(with $h_F=s+2$ and $\alpha_F:=\alpha(s)$; the length $t^*$ of a ray step of $\pi^s_x$ is $t^{(s)}_H$ at an $H_s$
truncation in a gap, and the length $t^*$ of the ray step of $\pi^0_x$ otherwise): the trajectory of $\pi^s_x$ is an initial part of $\Gamma^0_x$, and the additional end point
$p+t^{(s)}_Hd$ of an H termination in a gap has $t^{(s)}_H<\min(t_S,t_W,t_D)$, so, as in the proof of
Lemma~\ref{lem:flows-traj}(c), it lies in the waste.

\begin{definition}[ceiling flows]\label{def:flows-ceiling}
For any object or quantity $Q$ defined from the floor flows of the packing $\mathcal P$ (paths, trajectories,
termination heights, $N(s)$, $\Tstar$, and the losses and line quantities defined in Section~\ref{ssec:flows-quantities}
below, such as $\mathcal D$, $\mathcal M$, $\cE$, $\Lam_0$), its \emph{ceiling version} $\tilde Q$ is the same quantity computed
for the reflected packing $\rho\mathcal P$ and transported back by $\rho$: floor points $x$ are unchanged and heights are
mapped by $y\mapsto k-y$. Thus the ceiling master flow is $\tilde F^0:=\rho\circ\Fz[\rho\mathcal P]\circ\rho$; its path
$\tilde\pi^0_x:=\rho(\pi^0_x[\rho\mathcal P])$ starts at $(x,k)$ and moves downwards; similarly $\tilde F_s$,
$\tilde N(s):=N(s)[\rho\mathcal P]$, $\tilde T^*:=\Tstar[\rho\mathcal P]$, $\tilde\Lam_0:=\Lam_0[\rho\mathcal P]$. A
ceiling path is alive at height $y$ if its termination height $\tilde\tau(x):=k-\tau[\rho\mathcal P](x)$ is $<y$.
\end{definition}

By Lemma~\ref{lem:flows-refl}, $\rho\mathcal P$ is a closed packing of $[0,k]^2$ with the same $N$ and $W$; hence every
statement of this paper about the floor flows of an arbitrary closed packing applies to the ceiling flows, with heights
reflected.

\begin{remark}[one R1 for all scales]\label{rem:flows-joint}
All scale flows are truncations of the single flow $\Fz$, in which R1 is decided once. Therefore an entry point is
passed or T-terminated by at most one path (Lemma~\ref{lem:flows-R1}(c)), simultaneously for $\Fz$ and all $\Fs$; and
$\mathcal T_s$ increases with $s$ and equals $\{\Tstar\le s\}$ (Lemma~\ref{lem:flows-Ts}(d)(e)). These two facts are used
in Sections~\ref{sec:P6} and~\ref{sec:P8}. They are properties of this construction: if R1 were decided separately in
each $\Fs$, the winners could depend on $s$, and neither fact would follow.
\end{remark}

\subsection{Losses and line quantities}\label{ssec:flows-quantities}

In this subsection $F\in\{\Fz\}\cup\{\Fs:\ s\in[y_0,y_1]\}$, with $h_F=\hmax$, resp.\ $h_F=s+2$, paths $\pi_x$ and
termination heights $\tau=\tau_F$. For $\Fs$, ``T-terminated'' means $T_s$-terminated (Definition~\ref{def:Fs}).

\begin{definition}[losses]\label{def:flows-losses}
For $\star\in\{\mathrm T,\mathrm D,\mathrm E,\mathrm M,\mathrm W\}$ and $y\in\R$ put
\[
L^F_\star(y):=\big|\{x\in(0,k):\ \pi_x\text{ is }\star\text{-terminated in }F,\ \tau_F(x)<y\}\big|,\qquad
L^F(y):=\sum_\star L^F_\star(y).
\]
For the master flow let $\mathcal D$, $\mathcal M$, $\cE$ be the sets of $x$ for which $\pi^0_x$ is D-, M-, E-terminated,
and $L_E:=|\cE|$; their ceiling versions are $\tilde{\mathcal D}$, $\tilde{\mathcal M}$, $\tilde{\cE}$, $\tilde L_E$. For a
square $Z$ let $m^F(Z)$ be the floor measure of the set of $x$ such that $\pi_x$ has an R1-passed entry at $Z$ in $F$
(i.e.\ passes $Z$ or is T-terminated at $\Bot(Z)$ in $F$).
\end{definition}

\begin{corollary}[tilt losses of a scale flow]\label{cor:flows-LT}
For $s\in[y_0,y_1]$ and every $y$, $L^{\Fs}_T(y)\le N(s)$.
\end{corollary}

\begin{proof}
By Definition~\ref{def:Fs} and Lemma~\ref{lem:flows-Ts}(b), the set of $x$ that are T-terminated in $\Fs$ is $\mathcal T_s$,
so $L^{\Fs}_T(y)=|\{x\in\mathcal T_s:\tau_s(x)<y\}|\le|\mathcal T_s|=N(s)$.
\end{proof}

\begin{definition}[line quantities]\label{def:flows-line}
Let $y\in(0,h_F)$. If $\tau(x)\ge y$, the trajectory $\Gamma_x$ meets $\ell_y$ in exactly one point $P_y(x)$, with
abscissa $X_y(x)$: indeed $\Gamma_x$ is a connected compact polygonal line that starts at height $0$ and ends at height
$\ge\tau(x)\ge y$ (for a completed path the end has height $\ge h_F$), and heights increase strictly along it
(Lemma~\ref{lem:flows-traj}(a)). Put $\mathrm{Wa}_y:=\ell_y\cap\text{waste}$, identified with the set of its abscissas;
$|\mathrm{Wa}_y|=\omega(y)$. For a square $Z$ let
\[
A_Z(y):=\{x:\ \tau(x)>y,\ P_y(x)\in\operatorname{int}Z\},\qquad \mu_Z(y):=|A_Z(y)|,
\]
\[
A_{\rm gap}(y):=\{x:\ \tau(x)>y,\ P_y(x)\in\mathrm{Wa}_y\},\qquad \Fgap(y):=|A_{\rm gap}(y)|,
\]
\begin{gather*}
\Sh(y):=\sum_{Z\in\mathcal P}\big(c_Z(y)-\mu_Z(y)\big)_+,\qquad \Ov(y):=\sum_{Z\in\mathcal P}\big(\mu_Z(y)-c_Z(y)\big)_+,\\
\mathrm{Bd}(y):=k-L(y)-\sum_{Z\in\mathcal P}\mu_Z(y)-\Fgap(y).
\end{gather*}
Convention for $\mu_Z$: a path is counted if it is alive at $y$ in the strict sense $\tau(x)>y$ and crosses $\ell_y$ in
the interior $\operatorname{int}Z$ of $Z$ in $\R^2$. (The set $\operatorname{int}Z\cap\ell_y$ is contained in the relative
interior of the chord $Z\cap\ell_y$, and it is empty when $\ell_y$ contains a side of $Z$, although then $c_Z(y)=1$.)
Let $\nu_y$ be the image of the floor measure restricted to $A_{\rm gap}(y)$ under $x\mapsto X_y(x)$. Finally, for a
point $q$ of the waste, $M_g(q)$ is the number of $x$ such that $q$ lies in the open part of a gap segment of
$\Gamma^0_x$, and $M_g:=\sup_qM_g(q)$ (the \emph{gap multiplicity}).
\end{definition}

For a square $Z$, $\mu_Z(y)=0$ for all $y$ if no path of $F$ passes $Z$ (Lemma~\ref{lem:flows-density}(a) below).
More generally, by Lemma~\ref{lem:flows-density}(a) every $x\in A_Z(y)$ passes $Z$ and so has an R1-passed entry at $Z$;
hence $\mu_Z(y)\le m^F(Z)$ for every $y$, and $\mu_Z\equiv0$ if $m^F(Z)=0$.

\begin{lemma}[crossings and the gap density]\label{lem:flows-density}
Let $y\in(0,h_F)$ and let $x$ satisfy $\tau(x)>y$.
\begin{enumerate}[label=(\alph*)]
\item $P_y(x)\in\operatorname{int}Z$ if and only if $\pi_x$ passes $Z$ in $F$ and $P_y(x)$ is a point of the open part of
the pass segment of $Z$, i.e.\ $q_{Z,y}(x)<y<q_{Z,y}(x)+\cos\varphi_Z$ for its entry point $q_Z(x)$.
\item $P_y(x)\in\mathrm{Wa}_y$ if and only if $P_y(x)$ lies on the open part of a gap segment of $\Gamma_x$.
\item $A_{\rm gap}(y)$ is the union of finitely many points and finitely many disjoint open intervals $J_c$; on each
$J_c$, $P_y(x)$ lies on the open part of the gap segment of a fixed ray step with a fixed source $X_c$, and $X_y$ is affine
on $J_c$ with slope $\kappa_c\ge1$. Consequently $\nu_y$ is absolutely continuous, with density
\[
\rho_y=\sum_c\kappa_c^{-1}\,\mathbf 1_{X_y(J_c)}\qquad\text{(almost everywhere)},
\]
which vanishes almost everywhere outside $\mathrm{Wa}_y$, and $\Fgap(y)=\nu_y(\R)=\int_{\mathrm{Wa}_y}\rho_y$. (The
density always carries the height as a subscript, $\rho_y$; the letter $\rho$ alone denotes the reflection of
Lemma~\ref{lem:flows-refl}.)
\item $A_Z(y)$ is the union of finitely many points and finitely many open intervals on each of which $X_y$ is affine with
slope $\ge1$.
\end{enumerate}
\end{lemma}

\begin{proof}
Since $\tau(x)>y$, the point $P_y(x)$ is not the termination point (whose height is $\ge\tau(x)$, also for completed
paths). The other event points are $(x,0)$, of height $0<y$, and entry and exit points, which lie on the boundaries of
squares (Lemma~\ref{lem:flows-traj}(e)). Hence $P_y(x)$ is either an entry or exit point, or a point of the open part of
a gap segment (which lies in the waste, Lemma~\ref{lem:flows-traj}(c)), or a point of the open part of a pass segment
(which lies in $\operatorname{int}$ of the passed square, Lemma~\ref{lem:flows-traj}(e)). As the interiors of the squares,
their boundaries and the waste are pairwise disjoint, (a) and (b) follow.

(c) Let $I$ be a piece of $F$ (Lemma~\ref{lem:structure}, Lemma~\ref{lem:flows-Ts}(f)). For each ray step $j$ of the
common type (for $\Fs$, $t^*_j$ is the length of the step in $\Fs$, as stated after Lemma~\ref{lem:flows-Ts}), the set of $x\in I$ with $p_{j,y}(x)<y<p_{j,y}(x)+t^*_j(x)\cos\varphi_{X_j}$ is an open interval, given by two
affine strict inequalities; for such $x$, the path continues above height $y$, so $\tau(x)>y$, and by (b)
$x\in A_{\rm gap}(y)$. By (b), $A_{\rm gap}(y)\cap I$ is the union of these intervals, which are disjoint because the
height is strictly increasing along $\Gamma_x$. There are finitely many pieces, and $\Sigma_F$ (resp.\ $\Sigma_s$) is
finite. On each interval the slope of $X_y$ is $\kappa_c=\lambda_j/\cos\varphi_{X_j}\ge1$ by
Lemma~\ref{lem:expansion}(d). The image of Lebesgue measure on an interval $J$ under an injective affine map with slope
$\kappa$ is $|\kappa|^{-1}$ times Lebesgue measure on the image; summing over $c$ gives the density. The images lie in
$\mathrm{Wa}_y$ by (b).

(d) In the same way, by (a), on a piece $I$ the set $A_Z(y)\cap I$ is empty unless $Z$ is passed in the common type, and
then it is the open interval given by $q_{Z,y}(x)<y<q_{Z,y}(x)+\cos\varphi_Z$; the slope is $\mu/\cos\varphi_Z\ge1$ by
Lemma~\ref{lem:expansion}(d).
\end{proof}

\begin{definition}[gap overlap and loss terms]\label{def:flows-ovgap}
With $\rho_y$ as in Lemma~\ref{lem:flows-density} put
\[
\Ovgap(y):=\int_{\mathrm{Wa}_y}(\rho_y-1)_+\,dq,\qquad U(y):=\int_{\mathrm{Wa}_y}(1-\rho_y)_+\,dq
\]
(these do not depend on the version of $\rho_y$),
\[
\Lam^F(y):=\big[L^F_D+L^F_E+L^F_M+L^F_W\big](y)+\Ovgap^F(y),
\]
\[
\Lam_0:=|\mathcal D|+|\mathcal M|+L_E+\sup_{0<y<\hmax}\Ovgap^{\Fz}(y),\qquad \Lam^*(s):=\Lam_0+2(s+2)\tan\alpha(s),
\]
and, for the ceiling flows, $\tilde\Lam_0:=\Lam_0[\rho\mathcal P]$ and $\tilde\Lam^*(s):=\tilde\Lam_0+2(s+2)\tan\alpha(s)$.
\end{definition}

Since $\Ovgap(y)\le\int_{\mathrm{Wa}_y}\rho_y=\Fgap(y)\le k$ and $\mathcal D,\mathcal M,\cE$ are disjoint subsets of
$(0,k)$, we have $\Lam_0\le2k<\infty$. The term $2(s+2)\tan\alpha(s)$ in $\Lam^*(s)$ accounts for wall terminations of
$\Fs$; it is not of the form $O(W/\delta)$ and is integrated separately in Section~\ref{sec:P8}.

\begin{lemma}[measurability]\label{lem:flows-meas}
\begin{enumerate}[label=(\alph*)]
\item The sets $\mathcal D$, $\mathcal M$, $\cE$, $\mathcal T_s$, $\operatorname{dom}\Tstar$, and, for every $y$ and every
square $Z$, the sets $\{x:\pi_x\text{ is }\star\text{-terminated in }F,\ \tau_F(x)<y\}$, $A_Z(y)$ and $A_{\rm gap}(y)$
are finite unions of intervals and points.
\item $\Tstar$ is a Borel function on the Borel set $\operatorname{dom}\Tstar$, and
$\{(x,s):\ x\in\operatorname{dom}\Tstar,\ \Tstar(x)\le s\}$ is a Borel subset of $\R^2$.
\item $s\mapsto N(s)$ is non-decreasing on $[y_0,y_1]$, hence Borel.
\item For fixed $F$, the functions $y\mapsto L^F_\star(y)$, $\mu_Z(y)$, $\Fgap(y)$, $c_Z(y)$, $\omega(y)$, $\Sh(y)$, $\Ov(y)$
and $\mathrm{Bd}(y)$ are Borel on $(0,h_F)$.
\end{enumerate}
\end{lemma}

\begin{proof}
(a) On each piece of $\Fz$ the termination type is constant, and on each piece of $\Fs$ as well
(Lemma~\ref{lem:flows-Ts}(f)); $\tau$ is affine or constant on pieces; there are finitely many pieces and finitely many
exceptional points. This gives the statements for $\mathcal D,\mathcal M,\cE$ and the loss sets; for $\mathcal T_s$ see
Lemma~\ref{lem:flows-Ts}(f); for $A_Z(y)$ and $A_{\rm gap}(y)$ see Lemma~\ref{lem:flows-density}. On a piece $I$ of $\Fz$
the R1-passed entries $Z_j$ are fixed and $\eta_j$ is affine, so each $s_j$ is a continuous piecewise affine function on
$I$ with at most three pieces, and $\min_js_j$ is a continuous piecewise affine function with finitely many pieces; hence $\operatorname{dom}\Tstar\cap I=\{x\in I:\min_js_j(x)\le y_1\}$
is a finite union of intervals.

(b) By the above, $\Tstar$ is continuous on $\operatorname{dom}\Tstar\cap I$ for each piece $I$, and the finitely many
remaining points form a Borel set; so $\Tstar$ is Borel. The set in question is the preimage of the closed set
$\{(t,s):t\le s\}$ under the Borel map $(x,s)\mapsto(\Tstar(x),s)$ on $\operatorname{dom}\Tstar\times\R$.

(c) This is Lemma~\ref{lem:flows-Ts}(d).

(d) Fix a piece $I$ of $F$ (Lemma~\ref{lem:structure}, resp.\ Lemma~\ref{lem:flows-Ts}(f)) and consider the planar sets
\begin{gather*}
\{(x,y):\ x\in I,\ 0<y<h_F,\ \pi_x\ \star\text{-terminated},\ \tau(x)<y\},\\
\{(x,y):\ x\in I,\ 0<y<h_F,\ x\in A_Z(y)\},\qquad \{(x,y):\ x\in I,\ 0<y<h_F,\ x\in A_{\rm gap}(y)\}.
\end{gather*}
On $I$ the termination type is constant and $\tau$ is affine or constant. By Lemma~\ref{lem:flows-density}(a),(b), for
$x\in I$ the condition $x\in A_Z(y)$ holds if and only if $Z$ is passed in the common type and
$q_{Z,y}(x)<y<q_{Z,y}(x)+\cos\varphi_Z$, and the condition $x\in A_{\rm gap}(y)$ holds if and only if
$p_{j,y}(x)<y<p_{j,y}(x)+t^*_j(x)\cos\varphi_{X_j}$ for one of the finitely many ray steps $j$ of the common type; the
functions of $x$ occurring here are affine on $I$ (Lemma~\ref{lem:structure}(b)). Hence each of the three sets is a finite
union of sets of the form $\{(x,y):\ x\in I,\ l_1(x,y)\diamond_10,\dots,l_m(x,y)\diamond_m0\}$ with affine functions $l_i$
and relations $\diamond_i\in\{<,\le,=\}$. Each such set is the intersection of finitely many open or closed half-planes
(the strip $I\times\R$ is the intersection of two open half-planes, and an equation is the intersection of two closed
half-planes); it is a convex polyhedral set, and it is Borel because half-planes are open or closed. For each of the
finitely many floor points $x\in\Sigma_F\cap(0,k)$ (resp.\ $\Sigma_s\cap(0,k)$) the set of $y\in(0,h_F)$ that satisfy
the defining condition is an open half-line, an open interval or a finite union of open intervals (the height ranges of
the open parts of the segments of $\Gamma_x$), intersected with $(0,h_F)$; adding these finitely many Borel subsets of
vertical lines keeps the three planar sets, taken over all of $(0,k)$, Borel. By Tonelli's theorem the measures of
their sections at height $y$, which are $L^F_\star(y)$, $\mu_Z(y)$ and $\Fgap(y)$, are Borel functions of $y$. Next, if
$a(Z)>0$, then $c_Z(y)$ is the length of the horizontal section of a convex polygon, a continuous piecewise linear
function of $y$ with finitely many pieces; if $a(Z)=0$, then $c_Z$ is the indicator function of the closed interval
$[v_Z,v_Z+1]$ (the chord has length $1$ also at the heights of the two horizontal sides). In both cases $c_Z$ is Borel.
Further $\omega(y)=k-\sum_Zc_Z(y)$ for $y\in[0,k]$, since the chords are pairwise disjoint subsets of
$[0,k]\times\{y\}$; and $\Sh$, $\Ov$ and $\mathrm{Bd}$ are finite sums and positive parts of the functions above.
\end{proof}

\begin{definition}[exceptional heights]\label{def:flows-exc}
For $F\in\{\Fz\}\cup\{\Fs\}$ let $\Exc_F$ be the set of $y\in(0,h_F)$ for which there are a piece $I$ of $F$ (a component
of $[0,k]\setminus\Sigma_{\Fz}$, resp.\ of $[0,k]\setminus\Sigma_s$) and an event point of the common type of the paths
$\pi_x$, $x\in I$, whose height is the constant function $y$ on $I$. Put $\Exc:=\Exc_{\Fz}$.
\end{definition}

\begin{lemma}[exceptional heights]\label{lem:flows-exc}
\begin{enumerate}[label=(\alph*)]
\item $\Exc_F$ is finite.
\item If $y\in(0,h_F)\setminus\Exc_F$, then only finitely many $x$ outside the finite exceptional set of the pieces
($\Sigma_{\Fz}$, resp.\ $\Sigma_s$) have an event point at height $y$.
\item $\Exc_{\Fs}\subset\Exc\cap(0,s+2)$ for every $s\in[y_0,y_1]$. Thus $\Exc$ is one finite set of heights that serves for
all scale flows at once.
\end{enumerate}
\end{lemma}

\begin{proof}
(a) There are finitely many pieces, and the paths of a piece have finitely many event points
(Lemma~\ref{lem:flows-traj}(h)). (b) On each piece the height of each event point is affine in $x$ and not the constant
$y$, so it equals $y$ for at most one $x$. (c) The pieces of $\Fs$ are subintervals of the pieces of $\Fz$. An event
point of $\pi^s_x$ is an event point of $\pi^0_x$, or the end point $p+t^{(s)}_Hd$ of an H termination in a gap, whose
height is $s+2\notin(0,s+2)$. An affine function that is constant on a subinterval of a piece of $\Fz$ is constant on the
whole piece. Finally $\Exc_{\Fs}\subset(0,s+2)$ by definition.
\end{proof}

Section~\ref{sec:P4} describes the heights in $\Exc$ and shows that they cannot be omitted from the shadow inequality.

\begin{remark}[the order ``M before T'']\label{rem:flows-order}
At an entry candidate $(Y,q)$ with $a(Y)\ge\alpha_F$, rule R1 is applied before rule R2: the winner is T-terminated
and the others are M-terminated. Consider the opposite order, in which every path that reaches such a candidate alive
is T-terminated and R1 is applied only at candidates with $a(Y)<\alpha_F$. At a candidate with $a(Y)\ge\alpha_F$ all
arriving paths terminate at $q$ in both orders, and at a candidate with $a(Y)<\alpha_F$ the two orders agree; the
rules D, W, E, H do not depend on labels. By induction over the centre heights, the two orders produce the same
trajectories, termination points and termination heights; only the labels T and M differ. Hence every quantity that is
defined from trajectories and termination heights alone (such as those of Definition~\ref{def:flows-line}) is the same
for both orders; for the shadow inequality of Section~\ref{sec:P4} see Theorem~\ref{thm:P4}(g). The injectivity used in
Section~\ref{sec:P6} comes from rule R1 at the \emph{lower} member of each consecutive pair of entries (a passed square,
or the floor), not from this order. The order is used in Section~\ref{sec:P8}, for the inclusion
$\{x:\pi^s_x\text{ is T-terminated in }F_s\}\subset\{\Tstar\le s\}$ (Lemma~\ref{lem:flows-Ts}(e)): with the order
``M before T'', a path that is T-terminated in $\Fz$ is the R1 winner at its terminating entry, so that entry is one of
the R1-passed entries from which $\Tstar$ is defined (Definition~\ref{def:flows-Tstar}).
\end{remark}

\begin{remark}[degenerate cases]\label{rem:flows-degenerate}
\emph{Vertices and grazing.} Let a ray first meet a square at a vertex; this happens in particular for a ray running
along the line of a side. By the priority D $>$ W $>$ contact $>$ H the path terminates there with D if $t_D=t_S$,
otherwise with W if $t_W=t_S$ (the vertex then lies on a side wall of the container), and otherwise with an end
collision E, since a vertex does not lie in $\relint\Bot(Y)$. By Lemma~\ref{lem:structure}(c), which holds for the
scale flows by Lemma~\ref{lem:flows-Ts}(f), all this happens only for the finitely many $x\in\Sigma_F$ of each flow $F$
(resp.\ $x\in\Sigma_s$ for $\Fs$).
\emph{Ties.} Coincidences such as $t_D=t_S$ can occur for sets of $x$ of positive measure (for example for
paths crossing a horizontal gap of constant width); they are resolved by the priority D $>$ W $>$ contact $>$ H, and a
tie in $g$ at a merge is resolved by $x$. \emph{Squares on the floor.} If an axis-parallel square $Y$ rests on the floor,
then for $(x,0)\in\relint\Bot(Y)$ the first step has $t_S=0$ and an entry candidate $(Y,(x,0))$ with $g=0$; no other path
reaches a point of height $0$, so $x$ wins and passes $Y$ (as $a(Y)=0<\alpha_F$). The two abscissas of the end points of
$\Bot(Y)$, and the abscissa of a vertex of a tilted square on the floor, give end collisions with $\tau=0$; these are
finitely many $x$. \emph{Walls.} The floor points $0$ and $k$ are excluded; a path reaches a wall only through rule W, and
open gap parts and pass segments stay in $(0,k)\times\R$. \emph{Heights $h_F$ and $s+2$.} A contact at height exactly
$h_F$ (in $\Fz$) or $s+2$ (in $\Fs$) is processed. \emph{Inclination $\pi/4$.} By the $45^\circ$ convention such a
square has a bottom side; as $\alpha_F<\pi/4$ it is never passed: at an entry candidate on its bottom side the R1
winner is T-terminated and the other arriving paths are M-terminated, and a first contact elsewhere gives E (or D or W at
a tie). \emph{Exceptional heights} are treated by Lemma~\ref{lem:flows-exc}.
\end{remark}

\section{P1 and P2: the height identity and quantization at bottom sides}\label{sec:P12}

This section proves the height identity P1 (Theorem~\ref{thm:P1} and Corollary~\ref{cor:P1}) and the quantization P2 at bottom sides (Theorem~\ref{thm:P2}), and derives from them the form of P2 that is used in Section~\ref{sec:P7} (Theorem~\ref{thm:P2Z}): if $Z\in\cZ(s)$, then $Z\subset\R\times V_s$, no trajectory of the scale flow $\Fs$ meets the bottom side of $Z$, every trajectory of $\Fs$ meets $Z$ at most in its own termination point, and $\mu^{\Fs}_Z\equiv0$ (under the conditions stated there). The ceiling versions follow by reflection. Besides the definitions of Section~\ref{sec:flows} and the well-definedness of R1 (Lemmas~\ref{lem:DAG} and~\ref{lem:flows-R1}), the proofs use only [R, Lemma~\ref{imp:L3.1}], the second part of [R, Lemma~\ref{imp:L3.4}], and the inequality $\sec x-1\le0.50001\,x^2$ for $0\le x\le10^{-3}$ from the proof of [R, Lemma~\ref{imp:L4.15}], quoted as \eqref{imp:sec} in Section~\ref{sec:imported}.

\subsection{Notation and standing assumptions}\label{ssec:P12-notation}

Throughout this section $k\ge4$ is an integer, $\cP$ is a closed packing of $[0,k]^2$ with $N$ squares, and the parameters of Section~\ref{sec:flows} satisfy $\delta>0$, $c_0,c_1>0$, $y_0>0$, $0<\eps<1$ and $\omega_0>0$. The letter $s$ always denotes a number in $[y_0,y_1]$, where $y_1=k/2-3$. Further conditions on the parameters are stated in each result that uses them. We recall the notions of Section~\ref{sec:flows} on which the proofs rely.

\paragraph{Squares.}
For a square $S$ let $\varphi_S\in(-\frac\pi4,\frac\pi4]$ be its phase angle (a square whose sides make the angle $\frac\pi4$ with the axes has $\varphi_S=\frac\pi4$), $u_S=(\cos\varphi_S,\sin\varphi_S)$, $n_S=(-\sin\varphi_S,\cos\varphi_S)$, and $c_S$ its centre, so that
\[
S=\{c_S+\xi u_S+\zeta n_S:\ |\xi|,|\zeta|\le\tfrac12\};
\]
$a(S)=|\varphi_S|\in[0,\frac\pi4]$ is the inclination of $S$. The bottom and top sides of $S$ are the closed segments
\[
\Bot(S)=\{c_S-\tfrac12n_S+\xi u_S:\ |\xi|\le\tfrac12\},\qquad
\Top(S)=\{c_S+\tfrac12n_S+\xi u_S:\ |\xi|\le\tfrac12\},
\]
and their relative interiors $\relint\Bot(S)$, $\relint\Top(S)$ are given by $|\xi|<\frac12$. The floor is the pseudo-square $X_0$ with $\varphi_{X_0}=0$, $n_{X_0}=e_2:=(0,1)$ and ``top side'' $[0,k]\times\{0\}$; it is not a square. The height of a point $z$ is its second coordinate $z_y$. Since $|\varphi_S|\le\frac\pi4$,
\begin{equation}\label{eq:P12-ny}
n_S\cdot e_2=\cos\varphi_S\ \ge\ \cos\tfrac\pi4\ >\ 0
\end{equation}
for every square $S$ and for $S=X_0$. We write $\ell_y=\R\times\{y\}$; $c_S(y)$ is the length of $S\cap\ell_y$; $v_S$ is the height of the lowest point of $S$; $\Phi(S)=\{y:\ 0<c_S(y)<1\}$; and $\cR(S)$ is the ramp set of [R, Lemma~\ref{imp:L3.1}]: with $a=a(S)$ and $v=v_S$,
\[
\cR(S)=[v,\ v+\sin a)\cup(v+\cos a,\ v+\cos a+\sin a],\qquad \cR(S)=\emptyset\ \text{ if } a=0 .
\]
For $y\in\R$, $\frc{y}\in[0,1)$ is the fractional part of $y$ and $\dist(y,\Z)=\min(\frc{y},1-\frc{y})$.

\paragraph{Flows.}
(Definitions~\ref{def:flow}, \ref{def:R1} and~\ref{def:R2}.) A flow $F=F(\alpha_F,\delta,h_F)$ moves a path $\pi_x$ from each floor point $(x,0)$, $x\in(0,k)$. The state $(p,d,X,g)$ consists of the current point $p$, the direction $d=n_X$, the last traversed square $X$ (initially the floor $X_0$) and the cumulative gap $g$; the initial state is $((x,0),e_2,X_0,0)$. In this section a pass of Section~\ref{sec:flows} is called a \emph{traversal} and a pass segment a \emph{traversal segment}. We number the ray steps of $\pi_x$ by $j=0,1,2,\dots$. Step $j$ starts at the point $p_j$, with last traversed square $X_j$ and with cumulative gap $G_j$ (the start gap $g_j$ of Lemma~\ref{lem:structure}; here the letters $g_i$ are reserved for the lengths of the gap segments, see Theorem~\ref{thm:P1}); its ray is $r_j(t)=p_j+t\,n_{X_j}$ ($t\ge0$), and
\begin{gather*}
t_{S,j}=\min\{t\ge0:\ r_j(t)\in Y\ \text{for a square}\ Y\ne X_j\}\quad(=+\infty\ \text{if there is no such }t),\\
t_{W,j}=\min\{t\ge0:\ (r_j(t))_x\in\{0,k\}\},\qquad t_{D,j}=\delta-G_j,\qquad t_{H,j}=\frac{h_F-p_{j,y}}{\cos\varphi_{X_j}},\\
t^*_j=\min(t_{S,j},t_{W,j},t_{D,j},t_{H,j}).
\end{gather*}
At $t^*_j$ the first applicable rule in the following list is applied: D, if $t_{D,j}=t^*_j$ (termination at $r_j(t^*_j)$); W, if $t_{W,j}=t^*_j$ (termination at $r_j(t^*_j)$); contact, if $t_{S,j}=t^*_j$; H otherwise (the path is completed at $r_j(t^*_j)$). At a contact, the point $q=r_j(t_{S,j})$ lies in a unique square $Y$, the cumulative gap becomes $G_j+t_{S,j}$, and $(Y,q)$ is an \emph{entry candidate} if $q\in\relint\Bot(Y)$; otherwise the path terminates at $q$ with E. An entry candidate is processed by R1 (among all paths that reach $(Y,q)$ alive, the one with the lexicographically least pair $(g,x)$ wins and the others terminate at $q$ with M; R1 is well defined, see Definition~\ref{def:R1} and Lemmas~\ref{lem:DAG} and~\ref{lem:flows-R1}), then by R2 (the winner terminates at $q$ with T if $a(Y)\ge\alpha_F$); otherwise the winner \emph{traverses} $Y$: it moves along the segment $[q,q+n_Y]$, its state becomes $(q+n_Y,n_Y,Y,G_j+t_{S,j})$, and it is completed (H) if $(q+n_Y)_y\ge h_F$ and starts its next ray step otherwise. If step $j$ ends with the traversal of a square, we call that square $X_{j+1}$ and put
\[
q_j:=r_j(t_{S,j})\ \ (\text{entry point}),\qquad p_{j+1}:=q_j+n_{X_{j+1}}\ \ (\text{exit point}),\qquad G_{j+1}:=G_j+t_{S,j}.
\]
The \emph{trajectory} $\Gamma_x$ of $\pi_x$ is the polygonal path formed, in this order, by the gap segments $[p_j,r_j(t^*_j)]$ and the traversal segments $[q_j,p_{j+1}]$; its last point is the termination point $e(x)$. For $z,z'\in\Gamma_x$ we say that $z$ \emph{precedes} $z'$ if $z$ comes before $z'$ along $\Gamma_x$; by Lemma~\ref{lem:P12-traj}\,(F1) below this holds if and only if $z_y<z'_y$.

\paragraph{First-contact points.}
If $t^*_j=t_{S,j}<\infty$, the point $r_j(t_{S,j})$ is called the \emph{first-contact point of step $j$}, whatever rule (D, W or contact) is applied there. It is the last point of the gap segment of step $j$ and therefore lies on $\Gamma_x$. The number $G_j+t_{S,j}$ is called the \emph{cumulative gap at} $r_j(t_{S,j})$; at a contact it is the new value of the state variable $g$.

\paragraph{Master flow and scale flows.}
(Definitions~\ref{def:F0} and~\ref{def:Fs}.) Put $\alpha(s)=c_0s^{-1/2}$. The master flow is $\Fz=F(\amax,\delta,\hmax)$ with $\amax=\alpha(y_0)$ and $\hmax=k/2-1$; R1 is decided in $\Fz$ only. For $\Fz$ to be a flow we need $0<\amax<\pi/4$ and $0<\hmax<k$; the latter holds as $k\ge4$, and the former is assumed whenever $\Fz$ or $\Fs$ is used in this section (it holds under (Q1) below, since then $\amax\le10^{-3}$, and hence under the constraint (C1) of Section~\ref{sec:flows}). We write $\pi^0_x$ and $\Gamma^0_x$ for the path and the trajectory of $x$ in $\Fz$, and use the notation $p_j,X_j,G_j,q_j,\dots$ for $\pi^0_x$. For $s\in[y_0,y_1]$ the path of $x$ in the scale flow $\Fs$ is $\pi^0_x$ truncated at the first point of $\Gamma^0_x$ that is a point of type $(T_s)$ or $(H_s)$ below, or, if there is no such point, at the termination point of $\pi^0_x$:
\begin{itemize}
\item[($T_s$)] the point $q$ of an entry candidate $(Y,q)$ at which $\pi^0_x$ wins R1 and $a(Y)\ge\alpha(s)$;
\item[($H_s$)] a point $r_j(t)$ of the gap segment of a step $j$ with $t=t^{(s)}_{H,j}:=(s+2-p_{j,y})/\cos\varphi_{X_j}<\min(t_{S,j},t_{W,j},t_{D,j})$; or an exit point $p_{j+1}$ with $p_{j+1,y}\ge s+2$.
\end{itemize}
We write $\Gamma^s_x$ for the trajectory of $x$ in $\Fs$; it is the part of $\Gamma^0_x$ from $(x,0)$ up to the truncation point. The $\Fs$ path of $x$ \emph{traverses} a square $X$ if the traversal segment of $X$ is contained in $\Gamma^s_x$, and it \emph{enters} $X$ if it reaches an entry candidate $(X,q)$ alive and wins R1 there.

\paragraph{Ceiling flows.}
Let $\rho(x,y)=(x,k-y)$. Then $\rho\cP=\{\rho S:\ S\in\cP\}$ is a closed packing of $[0,k]^2$ (Lemma~\ref{lem:P12-refl}). The ceiling flows are $\tilde F^0=\rho\circ\Fz[\rho\cP]\circ\rho$ and $\tilde F_s=\rho\circ\Fs[\rho\cP]\circ\rho$: the trajectory of $x$ in $\tilde F_s$ is the image under $\rho$ of the trajectory of $x$ in $\Fs[\rho\cP]$, and the first-contact points, entry candidates, entries, termination points and termination types of $\tilde F_s$ are the images under $\rho$ of those of $\Fs[\rho\cP]$. The ceiling quantities are defined through $\rho$; for instance $\mu^{\tilde F_s}_Z(y)=\mu^{\Fs[\rho\cP]}_{\rho Z}(k-y)$ and $\Sh^{\tilde F_s}(y)=\Sh^{\Fs[\rho\cP]}(k-y)$.

\paragraph{Quantities.}
(Section~\ref{sec:flows}.) For $F\in\{\Fz,\Fs\}$, $\tau_F(x)$ is the termination height of the path of $x$ (with the convention of Section~\ref{sec:flows} for completed paths); for an E, D or W termination it is the height of the termination point. For $\tau_F(x)\ge y$, $P_y(x)$ is the point of $\Gamma_x\cap\ell_y$, which is unique by Lemma~\ref{lem:P12-traj}\,(F1). With $|\cdot|$ the Lebesgue measure,
\[
\mu^F_Z(y)=\bigl|\{x\in(0,k):\ \tau_F(x)>y,\ P_y(x)\in\operatorname{int}Z\}\bigr|,\qquad
\Sh^F(y)=\sum_{S\in\cP}\bigl(c_S(y)-\mu^F_S(y)\bigr)_+ .
\]

\paragraph{Parameters and classes of lines.}
(Section~\ref{sec:flows}.) Put $\hat\beta(t)=c_1t^{-3/4}$ for $t>0$ (a decreasing function), $d(y)=\min(y,k-y)$, $\beta(y)=\hat\beta(d(y))$, and
\[
b^*:=\hat\beta((1-\eps)y_0)=c_1\bigl((1-\eps)y_0\bigr)^{-3/4},\qquad
w_0:=0.50001\,c_0^2\Bigl(1+\frac{1.0001}{y_0}\Bigr)+\delta+1.0001\,b^*+10^{-12}.
\]
Let $\omega(y)$ be the line waste, and
\begin{align*}
H&:=\{y\in(0,k):\ y_0\le d(y)\le(1-\eps)y_1,\ \frc{y}\in[w_0,1-w_0]\},\\
Y_{iii}&:=\{y\in H:\ \omega(y)<\omega_0\ \text{and}\ a(S)<\beta(y)\ \text{for every}\ S\in\cP\ \text{meeting}\ \ell_y\},
\end{align*}
$Y_b:=Y_{iii}\cap(0,k/2)$, $Y_t:=Y_{iii}\cap(k/2,k)$, $\mathcal W_s:=[(1-\eps)s,\ s]$, $V_s:=[(1-\eps)s-1.0001,\ s+1.0001]$,
\[
\cZ(s):=\{Z\in\cP:\ \Phi(Z)\cap Y_b\cap\mathcal W_s\ne\emptyset\},\qquad
\tilde\cZ(s):=\{Z\in\cP:\ \Phi(Z)\cap Y_t\cap(k-\mathcal W_s)\ne\emptyset\},
\]
where $k-Y:=\{k-y:\ y\in Y\}$ for $Y\subset\R$. Since $w_0\ge\delta>0$,
\begin{equation}\label{eq:P12-frac}
\frc{y}\in[w_0,1-w_0]\iff\dist(y,\Z)\ge w_0\qquad(y\in\R);
\end{equation}
indeed both conditions fail if $y\in\Z$, and if $y\notin\Z$ both say $\min(\frc y,1-\frc y)\ge w_0$.

\paragraph{Conditions.}
The following conditions on the parameters are used; each result states which of them it needs.
\begin{itemize}
\item[(Q1)] $\alpha(y_0)=c_0y_0^{-1/2}\le10^{-3}$.
\item[(Q2)] $b^*=c_1\bigl((1-\eps)y_0\bigr)^{-3/4}\le10^{-4}$.
\item[(C2)] $\hat\beta(y_0)=c_1y_0^{-3/4}\le1.5\cdot10^{-6}$.
\end{itemize}
(C2) is one of the constraints of Section~\ref{sec:flows}, and (Q1) follows from the constraint (C1), $\alpha(y_0)\le1.5\cdot10^{-6}$, of Section~\ref{sec:flows}.

\subsection{Trajectories}\label{ssec:P12-traj}

The next lemma restates, in the notation of this section, parts of Lemma~\ref{lem:flows-traj} and of its proof: (F1)
is part (a) there, (F2) follows from the proofs of parts (b) and (c), (F3) is part (d), (F4) is part (e), (F5) is part
(f) together with part (b), and (F6) is part (g). We repeat the short proofs in the form used below.

\begin{lemma}\label{lem:P12-traj}
Let $\pi_x$ be a path of a flow $F$, with the notation of Subsection~\ref{ssec:P12-notation}.
\begin{enumerate}[label=\textup{(F\arabic*)}]
\item Every gap or traversal segment of $\Gamma_x$ has the form $\{z+t\,n:\ 0\le t\le t_0\}$ with $t_0\ge0$ and $n=n_X$ for a square $X$ or for the floor, and $n$ has positive second component. Hence the height is strictly increasing along $\Gamma_x$, distinct points of $\Gamma_x$ have distinct heights, and $\Gamma_x$ meets every horizontal line in at most one point.
\item If $0<t<t_{S,j}$, then $r_j(t)$ lies in no square. If $j\ge1$, then $p_j\in\Top(X_j)$ and $t_{S,j}>0$.
\item If $t_{S,j}<\infty$, then $r_j(t_{S,j})$ lies in exactly one square, and on the boundary of that square.
\item If step $j$ ends with the traversal of $X_{j+1}$, then $q_j\in\relint\Bot(X_{j+1})$, $p_{j+1}\in\relint\Top(X_{j+1})$, and the segment $[q_j,p_{j+1}]$ is contained in $X_{j+1}$ and meets no other square.
\item $G_j<\delta$ for every step $j$ of $\pi_x$. If $t^*_j=t_{S,j}$, then the cumulative gap at the first-contact point $r_j(t_{S,j})$ satisfies $G_j+t_{S,j}\le\delta$, with equality only if rule D is applied there.
\item Let $Z$ be a square and $z\in\Gamma_x\cap Z$. Then
\begin{enumerate}[label=\textup{(\roman*)}]
\item $z$ is the first-contact point of some step $j$ \textup{(}so $t^*_j=t_{S,j}$\textup{)} and $Z$ is the square containing it, or
\item $Z=X_{j+1}$ for some $j$ and $z\in[q_j,p_{j+1}]$.
\end{enumerate}
In case \textup{(ii)}, the entry point $q_j$ is a point of type \textup{(i)} and lies in $\relint\Bot(Z)$.
\end{enumerate}
\end{lemma}

\begin{proof}
(F4) Let $Y=X_{j+1}$, with centre $c$ and frame $u,n$. Since $(Y,q_j)$ is an entry candidate, $q_j=c-\frac12n+\xi u$ with $|\xi|<\frac12$. For $0\le\lambda\le1$ we get $q_j+\lambda n=c+(\lambda-\frac12)n+\xi u\in Y$, and for $\lambda=1$ this is $p_{j+1}=c+\frac12n+\xi u\in\relint\Top(Y)$. Since the squares of $\cP$ are pairwise disjoint, the segment $[q_j,p_{j+1}]\subset Y$ meets no other square.

(F1) The gap segment of step $j$ is $\{p_j+t\,n_{X_j}:\ 0\le t\le t^*_j\}$ and the traversal segment of $X_{j+1}$ is $\{q_j+t\,n_{X_{j+1}}:\ 0\le t\le1\}$; the vectors $n_{X_j}$, $n_{X_{j+1}}$ are unit vectors with positive second component by \eqref{eq:P12-ny}. Hence the height increases strictly along every segment of positive length, and since consecutive segments share their endpoints, it increases strictly along $\Gamma_x$. Two distinct points of $\Gamma_x$ are separated along $\Gamma_x$ by a part of positive length, so they have distinct heights.

(F2) Let $0<t<t_{S,j}$. For a square $Y\ne X_j$ we have $r_j(t)\notin Y$ by the minimality of $t_{S,j}$. If $j\ge1$, then $p_j\in\Top(X_j)$ by (F4) (applied to step $j-1$), that is $(p_j-c_{X_j})\cdot n_{X_j}=\frac12$; hence $(r_j(t)-c_{X_j})\cdot n_{X_j}=\frac12+t>\frac12$ and $r_j(t)\notin X_j$. Moreover, for $j\ge1$ the point $p_j$ lies in $X_j$, and the finitely many squares $Y\ne X_j$ are compact and disjoint from $X_j$; so the distance from $p_j$ to their union is positive, and $t_{S,j}>0$.

(F3) The minimum defining $t_{S,j}$ is attained, since there are finitely many squares and they are closed; so $r_j(t_{S,j})$ lies in a square $Y\ne X_j$, and in no other square because the squares are disjoint. If $t_{S,j}>0$, then $r_j(t)\notin Y$ for $0<t<t_{S,j}$ by (F2), so $r_j(t_{S,j})$ is a limit of points outside $Y$ and lies on $\partial Y$. If $t_{S,j}=0$, then $j=0$ by (F2) and $r_0(0)=(x,0)$; the points $(x,-\eta)$, $\eta>0$, are not in $Y\subset[0,k]^2$, so again $(x,0)\in\partial Y$.

(F5) $G_0=0<\delta$. If step $j$ ends with a traversal, the rule applied at $t^*_j=t_{S,j}$ is the contact rule, so rule D was not applicable: $t_{D,j}\ne t^*_j$, that is $t_{S,j}<t_{D,j}=\delta-G_j$, and $G_{j+1}=G_j+t_{S,j}<\delta$. By induction $G_j<\delta$ for every step. If $t^*_j=t_{S,j}$, then $t_{S,j}\le t_{D,j}=\delta-G_j$; equality means $t_{D,j}=t^*_j$, and then rule D is applied.

(F6) The points of $\Gamma_x$ are the points $r_j(t)$, $0\le t\le t^*_j$, of the gap segments and the points of the traversal segments. Let $z\in\Gamma_x\cap Z$.
\begin{itemize}
\item $z=r_j(t)$ with $t>0$. Since $t\le t^*_j\le t_{S,j}$, (F2) shows that $z$ lies in no square unless $t=t^*_j=t_{S,j}$; then $z$ is the first-contact point of step $j$, and $Z$ is the square containing it by (F3). This is case (i).
\item $z=r_j(0)=p_j$ with $j\ge1$. Then $p_j\in X_j$ by (F4), so $Z=X_j$ and $z\in[q_{j-1},p_j]$. This is case (ii).
\item $z=r_0(0)=(x,0)$. Then $t_{S,0}=0$, hence $t^*_0=0=t_{S,0}$, and $z$ is the first-contact point of step $0$. This is case (i).
\item $z$ lies on a traversal segment $[q_j,p_{j+1}]$. By (F4) it lies in $X_{j+1}$ and in no other square, so $Z=X_{j+1}$. This is case (ii).
\end{itemize}
In case (ii), step $j$ ended with the contact rule at $q_j=r_j(t_{S,j})$, so $t^*_j=t_{S,j}$ and $q_j$ is the first-contact point of step $j$; it lies in $Z$, and in $\relint\Bot(Z)$ by (F4).
\end{proof}

\subsection{P1: the height identity}\label{ssec:P12-P1}

\begin{theorem}[P1, height identity]\label{thm:P1}
Let $\pi_x$ be a path of a flow $F$ that traverses the squares $X_1,\dots,X_m$ \textup{(}$m\ge0$\textup{)}, in this order, and let $q:=r_m(t_{S,m})$ be the first-contact point of step $m$ \textup{(}so $t^*_m=t_{S,m}$\textup{)}. Put $\varphi_i:=\varphi_{X_i}$ for $1\le i\le m$, $\varphi_0:=0$, and $g_i:=t_{S,i}$ for $0\le i\le m$. Then $g_i$ is the length of the gap segment of step $i$,
\begin{equation}\label{eq:P12-P1}
q=(x,0)+\sum_{i=0}^mg_i\,n_{X_i}+\sum_{i=1}^mn_{X_i},
\end{equation}
that is,
\begin{equation}\label{eq:P12-P1y}
q_y=\sum_{i=1}^m\cos\varphi_i+\sum_{i=0}^mg_i\cos\varphi_i,\qquad
q_x=x-\sum_{i=1}^m\sin\varphi_i-\sum_{i=0}^mg_i\sin\varphi_i,
\end{equation}
and the cumulative gap at $q$ is $G_m+t_{S,m}=\sum_{i=0}^mg_i$.
\end{theorem}

\begin{proof}
For $i<m$, step $i$ ends with the traversal of $X_{i+1}$. Hence $t^*_i=t_{S,i}=g_i$, the gap segment of step $i$ is $[p_i,q_i]$ with $q_i=p_i+g_i\,n_{X_i}$, and
\[
p_{i+1}=q_i+n_{X_{i+1}},\qquad G_{i+1}=G_i+g_i .
\]
Moreover $q=p_m+g_m\,n_{X_m}$. Starting from $p_0=(x,0)$ and $G_0=0$ and adding these relations for $i=0,\dots,m-1$ gives \eqref{eq:P12-P1} and $G_m+t_{S,m}=\sum_{i=0}^mg_i$. Since $|n_{X_i}|=1$, $g_i$ is the length of the segment from $p_i$ to $p_i+g_i\,n_{X_i}$. Inserting $n_{X_i}=(-\sin\varphi_i,\cos\varphi_i)$, which also holds for $i=0$ because $n_{X_0}=(0,1)$ and $\varphi_0=0$, gives \eqref{eq:P12-P1y}.
\end{proof}

\begin{remark}\label{rem:P12-points}
The identity is a restatement of the definition of the trajectory: neither R1, nor R2, nor the height bound $h_F$ is used. It describes every point of $\Gamma_x$, but in two different forms. Suppose that the path has traversed $X_1,\dots,X_m$ before the point $z\in\Gamma_x$ considered.
\begin{enumerate}[label=\textup{(\alph*)}]
\item If $z=r_m(t)$ with $0\le t\le t^*_m$ is a point of the gap segment of step $m$, then $z$ satisfies \eqref{eq:P12-P1} with $g_m$ replaced by $t$.
\item If $z=q_m+\lambda\,n_{X_{m+1}}$ with $0\le\lambda\le1$ is a point of the traversal segment of $X_{m+1}$, then
\[
z=(x,0)+\sum_{i=0}^mg_i\,n_{X_i}+\sum_{i=1}^mn_{X_i}+\lambda\,n_{X_{m+1}}.
\]
\end{enumerate}
Thus \eqref{eq:P12-P1} holds verbatim, with the partial gap length, at the points of gap segments; at the points of a traversal segment the additional term $\lambda\,n_{X_{m+1}}$ appears. Both formulas follow from the relations in the proof of Theorem~\ref{thm:P1}.
\end{remark}

\begin{corollary}[P1, corollary]\label{cor:P1}
In the situation of Theorem~\ref{thm:P1}, let $0<\alpha<\pi/2$ and assume $a(X_i)<\alpha$ for $1\le i\le m$. Put $g:=\sum_{i=0}^mg_i$. Then
\begin{equation}\label{eq:P12-cor}
-(\sec\alpha-1)\,q_y\ \le\ -\sum_{i=1}^m(1-\cos\varphi_i)\ \le\ q_y-m\ \le\ \sum_{i=0}^mg_i\cos\varphi_i\ \le\ g\ \le\ \delta .
\end{equation}
If moreover $\alpha\le10^{-3}$, then $(\sec\alpha-1)\,q_y\le0.50001\,\alpha^2q_y$. Consequently $q_y-m\in[-0.50001\,\alpha^2q_y,\ \delta]$, and $q_y-m\in[-0.50001\,\alpha^2q_y,\ \delta)$ if rule D is not applied at $q$ \textup{(}in particular, if $q$ is the point of an entry candidate\textup{)}.
\end{corollary}

\begin{proof}
By \eqref{eq:P12-P1y},
\[
q_y-m=-\sum_{i=1}^m(1-\cos\varphi_i)+\sum_{i=0}^mg_i\cos\varphi_i .
\]
Since $g_i\ge0$ and $0<\cos\varphi_i\le1$, the second sum lies in $[0,g]$, and the first sum is nonnegative. This gives the second, third and fourth inequalities in \eqref{eq:P12-cor}; the fifth is (F5) of Lemma~\ref{lem:P12-traj}, because $g$ is the cumulative gap at $q$ by Theorem~\ref{thm:P1}. For the first inequality let $1\le i\le m$. Since $|\varphi_i|=a(X_i)<\alpha<\pi/2$ and $\sec$ is even and increasing on $[0,\pi/2)$,
\[
1-\cos\varphi_i=\cos\varphi_i\,(\sec\varphi_i-1)\ \le\ \cos\varphi_i\,(\sec\alpha-1).
\]
Summing over $i$ and using $\sum_{i=1}^m\cos\varphi_i\le q_y$, which holds by \eqref{eq:P12-P1y} because the remaining terms are nonnegative, we get $\sum_{i=1}^m(1-\cos\varphi_i)\le(\sec\alpha-1)\,q_y$. (For $m=0$ the sums over $i\ge1$ are empty and $q_y=g_0\in[0,g]$.) If $\alpha\le10^{-3}$, then $\sec\alpha-1\le0.50001\,\alpha^2$ by the inequality \eqref{imp:sec}, and $q_y\ge0$, so $(\sec\alpha-1)\,q_y\le0.50001\,\alpha^2q_y$. Finally, by (F5) the cumulative gap $g$ at $q$ equals $\delta$ only if rule D is applied at $q$; at the point of an entry candidate the contact rule is applied.
\end{proof}

\subsection{Three geometric lemmas}\label{ssec:P12-geom}

\begin{lemma}[bottom side]\label{lem:P12-bot}
Let $Z$ be a square, $\varphi=\varphi_Z$ and $a=a(Z)$.
\begin{enumerate}[label=\textup{(\alph*)}]
\item $v_Z=c_{Z,y}-\frac12(\cos a+\sin a)$, and the heights of the points of $\Bot(Z)$ fill exactly the interval $[v_Z,\ v_Z+\sin a]$. In particular the lowest point of $Z$ lies on $\Bot(Z)$, and
\[
q_y-\sin a\ \le\ v_Z\ \le\ q_y\qquad\text{for every } q\in\Bot(Z)\ \text{(endpoints included)}.
\]
\item If $a>0$, the closure of $\cR(Z)$ is $[v_Z,\ v_Z+\sin a]\cup[v_Z+\cos a,\ v_Z+\cos a+\sin a]$. If $a=0$, then $\cR(Z)=\emptyset$.
\end{enumerate}
\end{lemma}

\begin{proof}
(a) The point $c_Z+\xi u_Z+\zeta n_Z$ has height $c_{Z,y}+\xi\sin\varphi+\zeta\cos\varphi$. Since $\cos\varphi>0$, the minimum over $|\xi|,|\zeta|\le\frac12$ is attained at $\zeta=-\frac12$ and $\xi=-\frac12\operatorname{sgn}\varphi$ (any $\xi$ if $\varphi=0$), and equals $c_{Z,y}-\frac12\cos\varphi-\frac12|\sin\varphi|=c_{Z,y}-\frac12(\cos a+\sin a)$. On $\Bot(Z)$, where $\zeta=-\frac12$, the height $c_{Z,y}-\frac12\cos a+\xi\sin\varphi$ is an affine function of $\xi\in[-\frac12,\frac12]$ whose range is $[c_{Z,y}-\frac12\cos a-\frac12\sin a,\ c_{Z,y}-\frac12\cos a+\frac12\sin a]=[v_Z,\ v_Z+\sin a]$. The remaining assertions follow. Only $\cos\varphi>0$ and $|\sin\varphi|=\sin a$ were used, so the convention $\varphi=\frac\pi4$ for squares at $45^\circ$ plays no role.

(b) This follows from the formula for $\cR(Z)$ in [R, Lemma~\ref{imp:L3.1}(b)], with $v=v_Z$: for $a>0$ both intervals have the positive length $\sin a$, so their closures are the closed intervals displayed.
\end{proof}

\begin{lemma}[the auxiliary inequality]\label{lem:P12-aux}
Let $f(a):=\sin a+1-\cos a$ for $a\in[0,\frac\pi4]$.
\begin{enumerate}[label=\textup{(\alph*)}]
\item $f$ is strictly increasing on $[0,\frac\pi4]$, $f(0)=0$ and $f(\frac\pi4)=1$.
\item There is a unique $a^*\in(0,\frac\pi4]$ with $f(a^*)=1.0001\,a^*$, and for $a\in[0,\frac\pi4]$
\[
f(a)\le1.0001\,a\iff a\le a^* .
\]
\item $2.000133\cdot10^{-4}<a^*<2.000134\cdot10^{-4}$, that is, $a^*=2.000133\ldots\cdot10^{-4}$. In particular $f(a)\le1.0001\,a$ for $0\le a\le2\cdot10^{-4}$; for these $a$ this also follows directly from $f(a)\le a+a^2/2$.
\end{enumerate}
\end{lemma}

\begin{proof}
(a) $f'(a)=\cos a+\sin a>0$ on $[0,\frac\pi4]$, and $f(\frac\pi4)=\frac{\sqrt2}2+1-\frac{\sqrt2}2=1$.

(b) Let $h(a):=f(a)-1.0001\,a$. Then $h(0)=0$, $h'(a)=(\cos a+\sin a)-1.0001$, $h'(0)=-10^{-4}<0$, and $h''(a)=\cos a-\sin a>0$ for $0\le a<\frac\pi4$, so $h$ is strictly convex on $[0,\frac\pi4]$. Also $h(\frac\pi4)=1-1.0001\cdot\frac\pi4>1-1.0001\cdot0.7854>0.2$. Since $h'(0)<0$, we have $h(a)<0$ for all sufficiently small $a>0$, so by the intermediate value theorem $h$ has a zero $a^*\in(0,\frac\pi4)$. If $0<a<a^*$, strict convexity gives $h(a)<(1-\frac a{a^*})\,h(0)+\frac a{a^*}\,h(a^*)=0$. If $a^*<a\le\frac\pi4$, put $\lambda=a^*/a\in(0,1)$; then $0=h(a^*)=h((1-\lambda)\cdot0+\lambda a)<(1-\lambda)\,h(0)+\lambda\,h(a)=\lambda\,h(a)$, so $h(a)>0$. Hence $\{a\in[0,\frac\pi4]:\ h(a)\le0\}=[0,a^*]$, and $a^*$ is the only zero of $h$ in $(0,\frac\pi4]$.

(c) For $t>0$, integrating the inequality $\cos t\le1$ repeatedly gives
\[
t-\frac{t^3}6\le\sin t\le t-\frac{t^3}6+\frac{t^5}{120},\qquad \frac{t^2}2-\frac{t^4}{24}\le1-\cos t\le\frac{t^2}2,
\]
and therefore
\[
\frac t2-\frac{t^2}6-\frac{t^3}{24}-10^{-4}\ \le\ \frac{h(t)}t\ \le\ \frac t2-\frac{t^2}6+\frac{t^4}{120}-10^{-4}.
\]
For $t_1=2.000133\cdot10^{-4}$ the right-hand side is less than $-1.7\cdot10^{-11}$, and for $t_2=2.000134\cdot10^{-4}$ the left-hand side is greater than $3.2\cdot10^{-11}$ (direct evaluation; both bounds are rounded towards zero). Hence $h(t_1)<0<h(t_2)$, and by (b), $t_1<a^*<t_2$. Finally, $\sin a\le a$ and $1-\cos a\le a^2/2$ give $f(a)\le a+a^2/2\le1.0001\,a$ for $0\le a\le2\cdot10^{-4}$.
\end{proof}

Parts (b)--(d) of the next lemma, and the first sentence of part (a), are also contained in Lemma~\ref{lem:flows-refl};
the second sentence of part (a), which describes $\rho$ on the bottom and the top side, is new. We include the proofs.

\begin{lemma}[reflection]\label{lem:P12-refl}
Let $\rho(x,y)=(x,k-y)$.
\begin{enumerate}[label=\textup{(\alph*)}]
\item For every square $Z\subset[0,k]^2$, $\rho Z$ is a square in $[0,k]^2$ with $a(\rho Z)=a(Z)=:a$, $v_{\rho Z}=k-(v_Z+\cos a+\sin a)$ and $c_{\rho Z}(k-y)=c_Z(y)$ for all $y$. If $a<\frac\pi4$, then $\varphi_{\rho Z}=-\varphi_Z$ and $\rho(\Bot(\rho Z))=\Top(Z)$, $\rho(\relint\Bot(\rho Z))=\relint\Top(Z)$.
\item $\cR(\rho Z)=k-\cR(Z)$ and $\Phi(\rho Z)=k-\Phi(Z)$.
\item $\dist(k-y,\Z)=\dist(y,\Z)$ for all $y\in\R$; hence $\frc{k-y}\in[w_0,1-w_0]\iff\frc{y}\in[w_0,1-w_0]$.
\item $\rho\cP$ is a closed packing of $[0,k]^2$ with $N$ squares; $\omega_{\rho\cP}(k-y)=\omega_{\cP}(y)$; a square $\rho S$ of $\rho\cP$ meets $\ell_{k-y}$ if and only if $S$ meets $\ell_y$, and $a(\rho S)=a(S)$; and $\beta(k-y)=\beta(y)$. Consequently
\[
k-H=H,\qquad Y_{iii}[\rho\cP]=k-Y_{iii}[\cP],\qquad Y_b[\rho\cP]=k-Y_t[\cP],
\]
and, for every square $Z\in\cP$ and every $s$,
\[
Z\in\tilde\cZ(s)[\cP]\iff\rho Z\in\cZ(s)[\rho\cP].
\]
\end{enumerate}
Here $[\cP]$ and $[\rho\cP]$ indicate the packing with respect to which a set is formed. The integrality of $k$ is used only in \textup{(c)} and, through \textup{(c)}, in \textup{(d)}.
\end{lemma}

\begin{proof}
(a) $\rho$ is an isometry of $\R^2$ that maps $[0,k]^2$ onto itself and maps horizontal and vertical lines to horizontal and vertical lines. Hence $\rho Z$ is a square in $[0,k]^2$, and the least angle between a side of $\rho Z$ and a side of $[0,k]^2$ equals that of $Z$, so $a(\rho Z)=a(Z)$. By [R, Lemma~\ref{imp:L3.1}(a)] the highest point of $Z$ has height $v_Z+\cos a+\sin a$; $\rho$ maps it to the lowest point of $\rho Z$, so $v_{\rho Z}=k-(v_Z+\cos a+\sin a)$. Since $\rho Z\cap\ell_{k-y}=\rho(Z\cap\ell_y)$, the chords have equal lengths. For the last assertion write $c=c_Z$, $\varphi=\varphi_Z$, $u=u_Z$, $n=n_Z$; then
\[
\rho(c+\xi u+\zeta n)=\rho(c)+\xi(\cos\varphi,-\sin\varphi)+\zeta(-\sin\varphi,-\cos\varphi)=\rho(c)+\xi u'-\zeta n',
\]
with $u'=(\cos(-\varphi),\sin(-\varphi))$ and $n'=(-\sin(-\varphi),\cos(-\varphi))$. If $|\varphi|<\frac\pi4$, then $-\varphi\in(-\frac\pi4,\frac\pi4)$ is the phase angle of $\rho Z$ and $(u',n')$ its frame; the point of $Z$ with frame coordinates $(\xi,\zeta)$ is mapped to the point of $\rho Z$ with frame coordinates $(\xi,-\zeta)$. Hence $\rho$ maps $\Top(Z)$ ($\zeta=\frac12$) onto $\Bot(\rho Z)$ ($\zeta=-\frac12$), preserving relative interiors; as $\rho$ is an involution, $\rho(\Bot(\rho Z))=\Top(Z)$.

(b) If $a=0$, both ramp sets are empty. Otherwise, by (a) and [R, Lemma~\ref{imp:L3.1}(b)], with $v=v_Z$ and $v'=v_{\rho Z}=k-v-\cos a-\sin a$,
\begin{align*}
\cR(\rho Z)&=[v',\ v'+\sin a)\cup(v'+\cos a,\ v'+\cos a+\sin a]\\
&=[k-v-\cos a-\sin a,\ k-v-\cos a)\cup(k-v-\sin a,\ k-v]\\
&=k-\bigl((v+\cos a,\ v+\cos a+\sin a]\cup[v,\ v+\sin a)\bigr)=k-\cR(Z),
\end{align*}
including the open and closed ends. The identity for $\Phi$ follows from $c_{\rho Z}(k-y)=c_Z(y)$.

(c) Since $k\in\Z$, $k-\Z=\Z$. If $y\in\Z$, both distances are $0$. Otherwise $\frc{k-y}=1-\frc{y}$, and $\min(\frc y,1-\frc y)$ is invariant under $\frc y\mapsto1-\frc y$. The second assertion follows from the first and \eqref{eq:P12-frac}.

(d) $\rho$ maps $[0,k]^2$ bijectively onto itself and preserves disjointness, so $\rho\cP$ is a closed packing with $N$ squares. Since $([0,k]\times\{k-y\})\setminus\bigcup\rho\cP=\rho\bigl(([0,k]\times\{y\})\setminus\bigcup\cP\bigr)$, the line wastes agree. Also $\rho S\cap\ell_{k-y}=\rho(S\cap\ell_y)$ and $a(\rho S)=a(S)$ by (a), and $d(k-y)=\min(k-y,y)=d(y)$, so $\beta(k-y)=\beta(y)$. Each of the three conditions defining $H$ is invariant under $y\mapsto k-y$ (the last one by (c)), so $k-H=H$. Comparing the definitions, $k-y\in Y_{iii}[\rho\cP]$ if and only if $y\in Y_{iii}[\cP]$. Hence
\[
Y_b[\rho\cP]=Y_{iii}[\rho\cP]\cap(0,k/2)=\bigl(k-Y_{iii}[\cP]\bigr)\cap(0,k/2)=k-\bigl(Y_{iii}[\cP]\cap(k/2,k)\bigr)=k-Y_t[\cP].
\]
Finally, by (b),
\[
\Phi(\rho Z)\cap Y_b[\rho\cP]\cap\mathcal W_s=\bigl(k-\Phi(Z)\bigr)\cap\bigl(k-Y_t[\cP]\bigr)\cap\mathcal W_s=k-\bigl(\Phi(Z)\cap Y_t[\cP]\cap(k-\mathcal W_s)\bigr),
\]
so one of these sets is empty if and only if the other one is.
\end{proof}

\subsection{P2: quantization at bottom sides}\label{ssec:P12-P2}

\begin{lemma}[quantization, general form]\label{lem:P12-core}
Let $0<\alpha<\pi/2$. Let $\pi_x$ be a path of a flow $F$ that traverses the squares $X_1,\dots,X_m$ \textup{(}$m\ge0$\textup{)}, all with $a(X_i)<\alpha$, and suppose that the first-contact point $q$ of step $m$ lies on the closed side $\Bot(Z)$ of a square $Z$ \textup{(}whatever rule is applied at $q$\textup{)}. Let $g$ be the cumulative gap at $q$, so that $g\le\delta$ by \textup{(F5)}; put $a:=a(Z)$ and $K:=(\sec\alpha-1)\,q_y$ \textup{(}a local abbreviation, used only in this lemma and its proof\textup{)}. Then
\begin{align}
[v_Z,\ v_Z+\sin a]&\subset\bigl[\,m-K-\sin a,\ \ m+g+\sin a\,\bigr],\label{eq:P12-low}\\
[v_Z+\cos a,\ v_Z+\cos a+\sin a]&\subset\bigl[\,m+\cos a-K-\sin a,\ \ m+\cos a+g+\sin a\,\bigr],\label{eq:P12-up}
\end{align}
and every $y$ in the closure of $\cR(Z)$ satisfies
\[
\dist(y,\Z)\ <\ K+\delta+\sin a+(1-\cos a).
\]
\end{lemma}

\begin{proof}
By Theorem~\ref{thm:P1} and Corollary~\ref{cor:P1}, $m-K\le q_y\le m+g$, and by Lemma~\ref{lem:P12-bot}(a), $q_y-\sin a\le v_Z\le q_y$. Hence
\begin{gather*}
v_Z\ge q_y-\sin a\ge m-K-\sin a,\qquad v_Z+\sin a\le q_y+\sin a\le m+g+\sin a,\\
v_Z+\cos a\ge q_y+\cos a-\sin a\ge m+\cos a-K-\sin a,\\
v_Z+\cos a+\sin a\le q_y+\cos a+\sin a\le m+\cos a+g+\sin a,
\end{gather*}
which proves \eqref{eq:P12-low} and \eqref{eq:P12-up}. If $a=0$, then $\cR(Z)=\emptyset$ and there is nothing more to prove. Let $a>0$ and put $w:=K+\delta+\sin a+(1-\cos a)$. As $K\ge0$, $\delta>0$, $g\le\delta$ and $1-\cos a>0$,
\[
K+\sin a<w,\quad g+\sin a<w,\quad K+\sin a+(1-\cos a)<w,\quad g+\sin a-(1-\cos a)<w .
\]
By Lemma~\ref{lem:P12-bot}(b) the closure of $\cR(Z)$ is the union of the two intervals on the left of \eqref{eq:P12-low} and \eqref{eq:P12-up}. For $y$ in the first one, \eqref{eq:P12-low} gives $-(K+\sin a)\le y-m\le g+\sin a$, so $|y-m|<w$. For $y$ in the second one, \eqref{eq:P12-up} and $m+\cos a=(m+1)-(1-\cos a)$ give
\[
-\bigl(K+\sin a+(1-\cos a)\bigr)\ \le\ y-(m+1)\ \le\ g+\sin a-(1-\cos a),
\]
so $|y-(m+1)|<w$. Since $m,m+1\in\Z$, $\dist(y,\Z)<w$ in both cases.
\end{proof}

\begin{remark}\label{rem:P12-mechanism}
The integer near which the ramps of $Z$ lie is the number $m$ of squares traversed before $q$ (or $m+1$ for the upper ramp). The quantization comes from the integer height $0$ of the floor and from the fact that each traversal adds $\cos\varphi_i$, which is close to $1$, to the height.
\end{remark}

\begin{lemma}[scale flows]\label{lem:P12-Fs}
Let $s\in[y_0,y_1]$ and $x\in(0,k)$.
\begin{enumerate}[label=\textup{(\alph*)}]
\item $\Gamma^s_x$ is an initial part of $\Gamma^0_x$; up to its termination point the $\Fs$ path coincides with $\pi^0_x$ \textup{(}same steps, first-contact points and cumulative gaps\textup{)}. A point of $\Gamma^0_x$ that precedes a point of $\Gamma^s_x$ lies on $\Gamma^s_x$.
\item Every square $X$ traversed by the $\Fs$ path of $x$ satisfies $a(X)<\alpha(s)$.
\item If $p_j\in\Gamma^s_x$ and the $\Fs$ path does not terminate at $p_j$, then $p_{j,y}<s+2$. Every first-contact point of $\pi^0_x$ that lies on $\Gamma^s_x$ has height $\le s+2$.
\item Let $Z$ be a square and $z\in\Gamma^s_x\cap Z$. Then $z$ is of type \textup{(i)} or \textup{(ii)} in Lemma~\ref{lem:P12-traj}\,\textup{(F6)} for $\pi^0_x$, and in case \textup{(ii)} the entry point $q_j\in\relint\Bot(Z)$ lies on $\Gamma^s_x$. At every first-contact point on $\Gamma^s_x$ the cumulative gap is $\le\delta$, with equality only if rule D is applied there.
\item A truncation point of type $(T_s)$ lies in $\relint\Bot(Y)$ for the square $Y$ of its entry candidate, and in no other square. A truncation point of type $(H_s)$ lies in no square, or it is the exit point of a square traversed by the $\Fs$ path and lies in no other square.
\end{enumerate}
\end{lemma}

\begin{proof}
(a) This is the definition of $\Fs$ (Subsection~\ref{ssec:P12-notation}).

(b) The path $\pi^0_x$ traverses $X$, so it won R1 at the entry candidate $(X,q)$, where $q\in\relint\Bot(X)$ is the entry point. If $a(X)\ge\alpha(s)$, then $q$ is a point of type $(T_s)$, so the truncation point precedes or equals $q$, and $\Gamma^s_x$ contains no point of the traversal segment of $X$ other than $q$. This contradicts the assumption that $X$ is traversed by the $\Fs$ path.

(c) For $j=0$, $p_{0,y}=0<s+2$. For $j\ge1$, $p_j$ is an exit point; if $p_{j,y}\ge s+2$, it is a point of type $(H_s)$, and the $\Fs$ path would terminate at $p_j$ or before. This proves the first assertion. Now let $q=r_j(t_{S,j})$, with $t^*_j=t_{S,j}$, lie on $\Gamma^s_x$. If $q=p_j$, then $t_{S,j}=0$, so $j=0$ by (F2) and $q_y=0$. Otherwise $p_j$ precedes $q$, so the $\Fs$ path does not terminate at $p_j$ and $p_{j,y}<s+2$; thus $t^{(s)}_{H,j}>0$. If we had $t^{(s)}_{H,j}<\min(t_{S,j},t_{W,j},t_{D,j})$, the point $r_j(t^{(s)}_{H,j})$ of the gap segment of step $j$ would be a point of type $(H_s)$ preceding $q$, and $q$ would not lie on $\Gamma^s_x$. Since $t_{S,j}=t^*_j\le\min(t_{W,j},t_{D,j})$, we get $t^{(s)}_{H,j}\ge t_{S,j}$, and hence
\[
q_y=p_{j,y}+t_{S,j}\cos\varphi_{X_j}\ \le\ p_{j,y}+t^{(s)}_{H,j}\cos\varphi_{X_j}=s+2 .
\]

(d) By (a), $\Gamma^s_x\subset\Gamma^0_x$, so Lemma~\ref{lem:P12-traj}\,(F6) applies to $z$ and $\pi^0_x$. In case (ii), $q_j$ precedes or equals $z$, so $q_j\in\Gamma^s_x$ by (a). The statement on the cumulative gap is (F5) for $\pi^0_x$.

(e) A point of type $(T_s)$ is the point $q\in\relint\Bot(Y)$ of an entry candidate $(Y,q)$; it lies in no other square by disjointness. Let the truncation point be of type $(H_s)$. If it is an exit point $p_{j+1}$, then the traversal segment $[q_j,p_{j+1}]$ precedes it, so $X_{j+1}$ is traversed by the $\Fs$ path, and $p_{j+1}$ lies in $X_{j+1}$ only (F4). Otherwise it is $r_j(t)$ with $t=t^{(s)}_{H,j}<\min(t_{S,j},t_{W,j},t_{D,j})\le t_{S,j}$. If $t>0$, this point lies in no square by (F2). If $t=0$ \textup{(}so $j\ge1$, because $p_{0,y}=0<s+2$\textup{)}, then $p_{j,y}=s+2$ and the point is the exit point $p_j$, which is covered by the previous case. \textup{(}$t<0$ is impossible, because $r_j(t)$ lies on the gap segment.\textup{)}
\end{proof}

For a square $Z$, $s\in[y_0,y_1]$ and a point $q$ we put
\begin{equation}\label{eq:P12-w}
w(s,Z;q):=0.50001\,\alpha(s)^2q_y+\delta+\sin a(Z)+\bigl(1-\cos a(Z)\bigr)=0.50001\,\alpha(s)^2q_y+\delta+f\bigl(a(Z)\bigr),
\end{equation}
with $f$ as in Lemma~\ref{lem:P12-aux}. When $q$ is clear from the context we also write $w(s,Z)$.

\begin{theorem}[P2, quantization at bottom sides]\label{thm:P2}
Assume \textup{(Q1)}. Let $s\in[y_0,y_1]$ and let $Z$ be a square.
\begin{enumerate}[label=\textup{(\alph*)}]
\item Let $q\in\Bot(Z)$ \textup{(}the closed side, endpoints included\textup{)} lie on the trajectory $\Gamma^s_x$ of an $\Fs$ path. Then $q$ is a first-contact point of $\pi^0_x$, $q_y\le s+2$, and every $y$ in the closure of $\cR(Z)$ satisfies
\[
\dist(y,\Z)\ <\ w(s,Z;q)=0.50001\,\alpha(s)^2q_y+\delta+\sin a(Z)+\bigl(1-\cos a(Z)\bigr).
\]
Here $0.50001\,\alpha(s)^2q_y=0.50001\,c_0^2\bigl(1+\frac{q_y-s}s\bigr)\le0.50001\,c_0^2\bigl(1+\frac2s\bigr)$.
\item If moreover $a(Z)\le a^*$ \textup{(}Lemma~\ref{lem:P12-aux}\textup{)}, then
\[
w(s,Z;q)\ \le\ 0.50001\,c_0^2\Bigl(1+\frac{q_y-s}s\Bigr)+\delta+1.0001\,a(Z).
\]
For $a(Z)\in(a^*,\frac\pi4]$ this inequality is false.
\item \textup{(Ceiling.)} Let $q$ lie on $\rho(\Bot(\rho Z))$ and on the trajectory of an $\tilde F_s$ path. Then $k-q_y\le s+2$, and every $y$ in the closure of $\cR(Z)$ satisfies
\[
\dist(y,\Z)\ <\ 0.50001\,\alpha(s)^2(k-q_y)+\delta+\sin a(Z)+\bigl(1-\cos a(Z)\bigr).
\]
If $a(Z)<\frac\pi4$, then $\rho(\Bot(\rho Z))=\Top(Z)$.
\end{enumerate}
Part \textup{(a)} applies in particular to every entry candidate $(Z,q)$ reached by an $\Fs$ path that has not terminated before $q$, before R1 is applied at $(Z,q)$; to E-contacts at the two endpoints of $\Bot(Z)$; and to D- and W-terminations at points of $\Bot(Z)$.
\end{theorem}

\begin{proof}
(a) By Lemma~\ref{lem:P12-Fs}(d), $q$ is of type (i) or (ii) for $\pi^0_x$. In case (ii), $q$ lies on the traversal segment $[q_j,q_j+n_Z]$ of $Z$, whose points $q_j+\lambda n_Z$ have frame coordinate $\zeta=-\frac12+\lambda$ with respect to $Z$; since $q\in\Bot(Z)$ means $\zeta=-\frac12$, we get $\lambda=0$ and $q=q_j$, which is of type (i). So $q$ is the first-contact point of some step $m$ of $\pi^0_x$. The squares $X_1,\dots,X_m$ traversed by $\pi^0_x$ before $q$ have their traversal segments before $q$ on $\Gamma^0_x$, hence on $\Gamma^s_x$ by Lemma~\ref{lem:P12-Fs}(a); so they are traversed by the $\Fs$ path and $a(X_i)<\alpha(s)$ by Lemma~\ref{lem:P12-Fs}(b). Since $\alpha$ is decreasing and $s\ge y_0$, (Q1) gives $0<\alpha(s)\le\alpha(y_0)\le10^{-3}<\pi/2$. Lemma~\ref{lem:P12-core}, applied with $\alpha:=\alpha(s)$, gives $\dist(y,\Z)<(\sec\alpha(s)-1)\,q_y+\delta+f(a(Z))$ for $y$ in the closure of $\cR(Z)$, and $(\sec\alpha(s)-1)\,q_y\le0.50001\,\alpha(s)^2q_y$ by Corollary~\ref{cor:P1}. The bound $q_y\le s+2$ is Lemma~\ref{lem:P12-Fs}(c). Finally $\alpha(s)^2q_y=c_0^2q_y/s=c_0^2\bigl(1+\frac{q_y-s}s\bigr)$, and $q_y-s\le2$.

(b) By \eqref{eq:P12-w} and (a), $w(s,Z;q)-0.50001\,c_0^2\bigl(1+\frac{q_y-s}s\bigr)-\delta=f(a(Z))$. By Lemma~\ref{lem:P12-aux}(b), $f(a(Z))\le1.0001\,a(Z)$ holds if and only if $a(Z)\le a^*$.

(c) By the definition of the ceiling flow, $\rho q$ lies on $\Bot(\rho Z)$ and on the trajectory of an $\Fs[\rho\cP]$ path, and $\rho Z$ is a square of the closed packing $\rho\cP$ (Lemma~\ref{lem:P12-refl}(d)). Condition (Q1) does not involve the packing. By (a), applied to $\rho\cP$, $\rho Z$ and $\rho q$, whose height is $k-q_y$: $k-q_y\le s+2$, and $\dist(y',\Z)<0.50001\,\alpha(s)^2(k-q_y)+\delta+f(a(\rho Z))$ for every $y'$ in the closure of $\cR(\rho Z)$. By Lemma~\ref{lem:P12-refl}(a)--(c), $a(\rho Z)=a(Z)$, the closure of $\cR(\rho Z)$ is $k$ minus the closure of $\cR(Z)$, and $\dist(k-y,\Z)=\dist(y,\Z)$; taking $y'=k-y$ gives the assertion. The last sentence is Lemma~\ref{lem:P12-refl}(a).

The final assertion of the theorem holds because each of the listed points is a point of $\Bot(Z)$ on $\Gamma^s_x$: the point of an entry candidate reached by the $\Fs$ path is a first-contact point on $\Gamma^s_x$, and so are the E-, D- and W-termination points mentioned.
\end{proof}

\subsection{Squares of \texorpdfstring{$\cZ(s)$}{Z(s)}: the form used later}\label{ssec:P12-Z}

\begin{lemma}[window]\label{lem:P12-window}
Assume \textup{(Q2)} or \textup{(C2)}. Let $s\in[y_0,y_1]$ and $Z\in\cZ(s)$, and let $y^*\in\Phi(Z)\cap Y_b\cap\mathcal W_s$. Put $b_s:=\hat\beta\bigl(\max\{(1-\eps)s,\ y_0\}\bigr)$. Then:
\begin{enumerate}[label=\textup{(\alph*)}]
\item $0<a(Z)<\beta(y^*)\le b_s\le\min\{b^*,\ \hat\beta(y_0)\}$, and $b_s\le10^{-4}$.
\item $Z\subset\R\times\bigl((1-\eps)s-1.0001,\ s+1.0001\bigr)\subset\R\times V_s$. In particular every point of $Z$ has height $<s+1.0001<s+2$.
\item $y^*\in\cR(Z)$ and $\dist(y^*,\Z)\ge w_0$.
\item $w(s,Z;q)<w_0-10^{-12}$ for every point $q\in Z$.
\end{enumerate}
\end{lemma}

\begin{proof}
(a) Since $y^*\in\Phi(Z)$, the set $\Phi(Z)$ is nonempty, so $a(Z)>0$ by the second part of [R, Lemma~\ref{imp:L3.4}]. Since $c_Z(y^*)>0$, the square $Z$ meets $\ell_{y^*}$, and $y^*\in Y_{iii}$ gives $a(Z)<\beta(y^*)$. As $y^*\in Y_b\subset H$ and $y^*<k/2$, we have $d(y^*)=y^*\ge y_0$; and $y^*\in\mathcal W_s$ gives $y^*\ge(1-\eps)s>0$. Since $\hat\beta$ is decreasing, $\beta(y^*)=\hat\beta(y^*)\le\hat\beta(\max\{(1-\eps)s,y_0\})=b_s$. Moreover $b_s\le\hat\beta(y_0)$, and $b_s\le\hat\beta((1-\eps)s)\le\hat\beta((1-\eps)y_0)=b^*$ because $s\ge y_0$. Under (Q2), $b_s\le b^*\le10^{-4}$; under (C2), $b_s\le\hat\beta(y_0)\le1.5\cdot10^{-6}<10^{-4}$.

(b) Let $a=a(Z)$. By [R, Lemma~\ref{imp:L3.1}(a)] the vertical extent of $Z$ is $[v_Z,\ v_Z+\cos a+\sin a]$, and it contains $y^*$ because $c_Z(y^*)>0$. Since $\cos a\le1$ and $\sin a\le a$, (a) gives $\cos a+\sin a\le1+a<1.0001$. Hence $v_Z\ge y^*-(\cos a+\sin a)>y^*-1.0001\ge(1-\eps)s-1.0001$ and $v_Z+\cos a+\sin a\le y^*+(\cos a+\sin a)<y^*+1.0001\le s+1.0001$.

(c) As $a(Z)>0$, the second part of [R, Lemma~\ref{imp:L3.4}] gives $\Phi(Z)\subset\cR(Z)$, so $y^*\in\cR(Z)$. Since $y^*\in H$, $\frc{y^*}\in[w_0,1-w_0]$, and \eqref{eq:P12-frac} gives $\dist(y^*,\Z)\ge w_0$.

(d) Let $q\in Z$. By (b), $0\le q_y<s+1.0001$, so
\[
0.50001\,\alpha(s)^2q_y=0.50001\,c_0^2\,\frac{q_y}s\ \le\ 0.50001\,c_0^2\Bigl(1+\frac{1.0001}s\Bigr)\ \le\ 0.50001\,c_0^2\Bigl(1+\frac{1.0001}{y_0}\Bigr).
\]
By (a), $a(Z)<b_s\le10^{-4}<a^*$, so Lemma~\ref{lem:P12-aux} gives $f(a(Z))\le1.0001\,a(Z)<1.0001\,b_s\le1.0001\,b^*$. Adding, $w(s,Z;q)<0.50001\,c_0^2(1+1.0001/y_0)+\delta+1.0001\,b^*=w_0-10^{-12}$.
\end{proof}

\begin{theorem}[P2 for the squares of $\cZ(s)$]\label{thm:P2Z}
Assume \textup{(Q1)}, and \textup{(Q2)} or \textup{(C2)}. Let $s\in[y_0,y_1]$ and $Z\in\cZ(s)$.
\begin{enumerate}[label=\textup{(\alph*)}]
\item \textup{(Inclination and window.)} $0<a(Z)<\hat\beta\bigl(\max\{(1-\eps)s,\,y_0\}\bigr)\le\min\{b^*,\hat\beta(y_0)\}$, $a(Z)<10^{-4}$, and $Z\subset\R\times V_s$.
\item \textup{(No contact with the bottom side.)} No point of the closed side $\Bot(Z)$ lies on the trajectory of any $\Fs$ path. In particular, no $\Fs$ path reaches an entry candidate $(Z,q)$; hence no $\Fs$ path enters or traverses $Z$, the set of $x$ whose $\Fs$ path enters $Z$ is empty, and no M- or T-termination of $\Fs$ occurs at a point of $Z$. No $\Fs$ path terminates with D or W at a point of $\Bot(Z)$, and no $\Fs$ path makes an E-contact at one of the two endpoints of $\Bot(Z)$.
\item \textup{(Trajectories meet $Z$ only at their ends.)} Every $\Fs$ trajectory $\Gamma^s_x$ meets $Z$ in at most one point. Such a point is the termination point $e(x)$ of the $\Fs$ path; this termination is of type E, D or W, it is also the termination of $\pi^0_x$, and $e(x)\in\partial Z\setminus\Bot(Z)$.
\item \textup{(No through-paths.)} $\mu^{\Fs}_Z(y)=0$ for every $y\in\R$. The same holds if $\operatorname{int}Z$ is replaced by $Z$ in the definition of $\mu^{\Fs}_Z$, and also if the condition $\tau_{\Fs}(x)>y$ is replaced by $\tau_{\Fs}(x)\ge y$; in the last variant with $Z$ in place of $\operatorname{int}Z$, the set $\{x:\ \tau_{\Fs}(x)\ge y,\ P_y(x)\in Z\}$ has at most $2(N+1)$ elements.
\item \textup{(Shadow.)} $\Sh^{\Fs}(y)\ge\sum_{Z'\in\cZ(s)}c_{Z'}(y)$ for every $y\in\R$; no height is excluded.
\item \textup{(Ceiling.)} If $Z\in\tilde\cZ(s)$, then $\rho Z\in\cZ(s)[\rho\cP]$ and \textup{(a)--(e)} hold for the packing $\rho\cP$, the square $\rho Z$ and the flow $\Fs[\rho\cP]$. Equivalently: $0<a(Z)<10^{-4}$ and $Z\subset\R\times(k-V_s)$; no point of $\Top(Z)=\rho(\Bot(\rho Z))$ lies on the trajectory of any $\tilde F_s$ path, and no $\tilde F_s$ path enters $Z$; every $\tilde F_s$ trajectory meets $Z$ in at most one point, which is its termination point, of type E, D or W, and lies in $\partial Z\setminus\Top(Z)$; $\mu^{\tilde F_s}_Z\equiv0$ in all the readings of \textup{(d)}; and $\Sh^{\tilde F_s}(y)\ge\sum_{Z'\in\tilde\cZ(s)}c_{Z'}(y)$ for every $y$.
\end{enumerate}
\end{theorem}

\begin{proof}
Fix $y^*\in\Phi(Z)\cap Y_b\cap\mathcal W_s$.

(a) This is Lemma~\ref{lem:P12-window}(a),(b).

(b) Suppose that a point $z\in\Bot(Z)$ lies on $\Gamma^s_x$. By Theorem~\ref{thm:P2}(a) (which uses (Q1)), $z$ is a first-contact point of $\pi^0_x$, and every $y$ in the closure of $\cR(Z)$ satisfies $\dist(y,\Z)<w(s,Z;z)$. By Lemma~\ref{lem:P12-window}(d), $w(s,Z;z)<w_0$. But $y^*\in\cR(Z)$ and $\dist(y^*,\Z)\ge w_0$ by Lemma~\ref{lem:P12-window}(c), a contradiction. The remaining assertions of (b) are special cases. Indeed, the point $q$ of an entry candidate $(Z,q)$ reached by an $\Fs$ path is a point of $\Bot(Z)$ on its trajectory; entering $Z$, traversing $Z$, and M- or T-terminations at a point of $Z$ all require such an entry candidate (M and T occur only at points of entry candidates, and a point of $Z$ lies in no other square); and the termination points and E-contact points mentioned are points of $\Bot(Z)$ on the trajectory.

(c) Let $z\in\Gamma^s_x\cap Z$. By Lemma~\ref{lem:P12-Fs}(d), $z$ is of type (i) or (ii) for $\pi^0_x$. In case (ii) the entry point $q_j\in\relint\Bot(Z)$ would lie on $\Gamma^s_x$, which (b) excludes. So $z=r_j(t_{S,j})$ is the first-contact point of a step $j$ of $\pi^0_x$, $Z$ is the square containing it, and $z\notin\Bot(Z)$ by (b). Consider the rule applied by $\pi^0_x$ at $z$: it is D if $t_{D,j}=t^*_j$, else W if $t_{W,j}=t^*_j$, else the contact rule; in the last case $(Z,z)$ is not an entry candidate because $z\notin\relint\Bot(Z)$, so the rule is E. In all three cases $\pi^0_x$ terminates at $z$, so $z$ is the last point of $\Gamma^0_x$; as $z\in\Gamma^s_x$ and $\Gamma^s_x$ is an initial part of $\Gamma^0_x$, $z$ is also the last point of $\Gamma^s_x$, that is $z=e(x)$, and $\Gamma^s_x\cap Z=\{z\}$. The $\Fs$ path is not truncated at $z$ by $(T_s)$ or $(H_s)$: by Lemma~\ref{lem:P12-Fs}(e), a $(T_s)$ point lies in $\relint\Bot(Y)$ for the only square $Y$ containing it, which would be $Z$, while $z\notin\Bot(Z)$; and an $(H_s)$ point lies in no square or in a square traversed by the $\Fs$ path only, which would again be $Z$, while $Z$ is not traversed by (b). Hence the termination of the $\Fs$ path at $z$ is the E, D or W termination of $\pi^0_x$. Finally $z\in\partial Z$ by (F3), and $z\notin\Bot(Z)$.

(d) For every $x$, by (c), $\Gamma^s_x\cap Z$ is empty or consists of the point $e(x)\in\partial Z$; in particular $\Gamma^s_x\cap\operatorname{int}Z=\emptyset$. Hence no $x$ satisfies $P_y(x)\in\operatorname{int}Z$, whatever $y$ and whichever aliveness condition is used; so $\mu^{\Fs}_Z(y)=0$, also in the variant with $\tau_{\Fs}(x)\ge y$. Now replace $\operatorname{int}Z$ by $Z$. If $\tau_{\Fs}(x)>y$ and $P_y(x)\in Z$, then $P_y(x)=e(x)$ by (c), and the termination is of type E, D or W, so $\tau_{\Fs}(x)$ is the height of $e(x)$, which is $y$; a contradiction. So the set is empty in the variant with $\tau_{\Fs}(x)>y$. In the variant with $\tau_{\Fs}(x)\ge y$, every $x$ in the set satisfies $e(x)=P_y(x)\in\partial Z\cap\ell_y$. Since $a(Z)>0$ by (a), no side of $Z$ is horizontal, and so $\ell_y$ meets $\partial Z$ in at most two points: indeed $Z\cap\ell_y$ is empty or a segment $[A,B]$ (possibly a single point), and if a point strictly between $A$ and $B$ lay on a side $\sigma$ of $Z$, then $A$ and $B$, which lie in the closed half-plane bounded by the line of $\sigma$ that contains $Z$, would both lie on that line, so $\ell_y$ would be the line of $\sigma$ and $\sigma$ would be horizontal. It remains to show that at most $N+1$ paths of $\Fs$ terminate at a given point $z\in Z$. By the proof of (c), such a path terminates at the first-contact point $z$ of its last step $j$, so $z=p_j+t\,n_{X_j}$ with $t=t_{S,j}\ge0$, where $X_j$ is the floor or the last square traversed. Let $x\ne x'$ be two such paths with the same $X:=X_j$ (with their own $p,t$ and $p',t'$). If $X$ is the floor, then $p=(x,0)$ and $p'=(x',0)$, and $z=(x,t)=(x',t')$ gives $x=x'$. If $X$ is a square, then $p,p'\in\Top(X)$ by (F2), so $p-p'$ is parallel to $u_X$, while $p-p'=(t'-t)\,n_X$ is parallel to $n_X$; hence $p=p'$. Both paths then left $X$ at $p$, so both traversed $X$ after the entry candidate $(X,p-n_X)$, and both won R1 there (R1 is decided in $\Fz$, and the $\Fs$ paths are initial parts of the $\Fz$ paths). This is impossible, because R1 has exactly one winner at each entry candidate (Definition~\ref{def:R1}, Lemma~\ref{lem:flows-R1}). So distinct paths terminating at $z$ have distinct $X$, and there are at most $N+1$ choices of $X$. Altogether the set has at most $2(N+1)$ elements, and it has measure $0$.

(e) By (d), $\mu^{\Fs}_{Z'}(y)=0$ for every $Z'\in\cZ(s)$ and every $y$. As all terms of the sum defining $\Sh^{\Fs}(y)$ are nonnegative,
\[
\Sh^{\Fs}(y)=\sum_{S\in\cP}\bigl(c_S(y)-\mu^{\Fs}_S(y)\bigr)_+\ \ge\ \sum_{Z'\in\cZ(s)}\bigl(c_{Z'}(y)-0\bigr)_+=\sum_{Z'\in\cZ(s)}c_{Z'}(y).
\]

(f) By Lemma~\ref{lem:P12-refl}(d), $\rho Z\in\cZ(s)[\rho\cP]$, and $\rho\cP$ is a closed packing of $[0,k]^2$ with $N$ squares. The conditions (Q1), (Q2), (C2) do not involve the packing, so (a)--(e) hold for $\rho\cP$, $\rho Z$ and $\Fs[\rho\cP]$. By definition, the trajectories, entries and termination points of $\tilde F_s$ are the images under $\rho$ of those of $\Fs[\rho\cP]$. By Lemma~\ref{lem:P12-refl}(a), $a(Z)=a(\rho Z)<10^{-4}<\frac\pi4$, so $\rho(\Bot(\rho Z))=\Top(Z)$ and $\rho\bigl(\partial(\rho Z)\setminus\Bot(\rho Z)\bigr)=\partial Z\setminus\Top(Z)$; and $\rho Z\subset\R\times V_s$ gives $Z\subset\R\times(k-V_s)$. The measure statements follow from $\mu^{\tilde F_s}_Z(y)=\mu^{\Fs[\rho\cP]}_{\rho Z}(k-y)$, and similarly for the variants of (d). For the shadow, using $c_{\rho S}(k-y)=c_S(y)$ and $\cZ(s)[\rho\cP]=\{\rho Z':\ Z'\in\tilde\cZ(s)\}$ (Lemma~\ref{lem:P12-refl}(d)),
\[
\Sh^{\tilde F_s}(y)=\Sh^{\Fs[\rho\cP]}(k-y)\ \ge\sum_{Z''\in\cZ(s)[\rho\cP]}c_{Z''}(k-y)=\sum_{Z'\in\tilde\cZ(s)}c_{Z'}(y).
\qedhere
\]
\end{proof}

\subsection{Remarks on the hypotheses and on degenerate cases}\label{ssec:P12-remarks}

\begin{remark}[what is used]\label{rem:P12-used}
\begin{enumerate}[label=\textup{(\arabic*)}]
\item Of the definition of $\cZ(s)$, Lemma~\ref{lem:P12-window} and Theorem~\ref{thm:P2Z} use only the existence of $y^*\in\Phi(Z)$ with $(1-\eps)s\le y^*\le s$, $y_0\le y^*<k/2$, $\frc{y^*}\in[w_0,1-w_0]$, and $a(S)<\beta(y^*)$ for every $S\in\cP$ meeting $\ell_{y^*}$. The conditions $\omega(y^*)<\omega_0$ and $d(y^*)\le(1-\eps)y_1$ are not used. Under (Q2), the bound $y^*\ge y_0$ is used only for the inequality $\beta(y^*)\le b_s$ in Lemma~\ref{lem:P12-window}(a); the other conclusions hold without it, since $y^*\ge(1-\eps)s$ gives $a(Z)<\hat\beta((1-\eps)s)\le b^*\le10^{-4}$.
\item No relation between the inclination of $Z\in\cZ(s)$ and $\alpha(s)$ is needed. A square of $\cZ(s)$ may or may not satisfy $a(Z)\ge\alpha(s)$, that is, it may or may not be a square at which an entering $\Fs$ path would be terminated with T; by Theorem~\ref{thm:P2Z}(b), no $\Fs$ path reaches an entry candidate of $Z$ at all. In particular, no inequality between $\hat\beta((1-\eps)s)$ and $\alpha(s)$ is assumed anywhere in this section.
\item Condition (Q1) is used only in Theorem~\ref{thm:P2}(a),(c), through the inequality~\eqref{imp:sec}, and hence in Theorem~\ref{thm:P2Z}(b)--(f). If (Q1) is dropped but $\alpha(s)<\pi/2$, Lemma~\ref{lem:P12-core} still gives the conclusion of Theorem~\ref{thm:P2}(a) with the term $(\sec\alpha(s)-1)\,q_y$ in place of $0.50001\,\alpha(s)^2q_y$.
\item Conditions (Q2) and (C2) are used only in Lemma~\ref{lem:P12-window}, to obtain $a(Z)<10^{-4}$. This gives the window (b) and, through Lemma~\ref{lem:P12-aux}, the bound $f(a(Z))\le1.0001\,a(Z)$ used in (d). Either of the two conditions suffices. Since (C2) is a constraint of Section~\ref{sec:flows} and (Q1) follows from (C1), Theorem~\ref{thm:P2Z} holds under the constraints (C1) and (C2) of that section, without (Q2).
\item The value of $\delta>0$ is not used; it enters only through $w_0$. The integrality of $k$ is used only in the ceiling statements, through Lemma~\ref{lem:P12-refl}(c),(d).
\item The constant $1.0001$ in $V_s$ and in the term $0.50001\,c_0^2(1+1.0001/y_0)$ of $w_0$ has a single source: the vertical extent $\cos a(Z)+\sin a(Z)$ of a square $Z\in\cZ(s)$ is less than $1.0001$.
\end{enumerate}
\end{remark}

\begin{remark}[degenerate cases]\label{rem:P12-degenerate}
\begin{enumerate}[label=\textup{(\arabic*)}]
\item \emph{Vertices and grazing contacts.} A contact at a vertex of a square, in particular the contact of a ray running along the line of a side, whose first point in the square is a vertex, is an E-contact and not an entry candidate (if instead rule D or W applies at that first-contact point, it is a D- or W-termination, see (2) and (3)); by Lemma~\ref{lem:structure}(c), first-contact points at vertices occur only for finitely many $x$. Lemma~\ref{lem:P12-core} and Theorem~\ref{thm:P2}(a) hold at all points of the closed side $\Bot(Z)$, including its endpoints; so for $Z\in\cZ(s)$ the endpoints of $\Bot(Z)$ lie on no $\Fs$ trajectory (Theorem~\ref{thm:P2Z}(b)), and E-contacts at the other two vertices of $Z$ are termination points (Theorem~\ref{thm:P2Z}(c),(d)).
\item \emph{Death ties} ($t_{D,j}=t_{S,j}=t^*_j$). The cumulative gap at the first-contact point is then exactly $\delta$ (F5). Lemma~\ref{lem:P12-core} allows $g\le\delta$ and its conclusion is strict, so death ties on $\Bot(Z)$ are excluded for $Z\in\cZ(s)$; death ties at other points of $\partial Z$ are termination points.
\item \emph{Wall ties} ($t_{W,j}=t_{S,j}=t^*_j$, rule W). The point is a first-contact point and is treated exactly like a death tie.
\item \emph{Contacts at height exactly $s+2$.} Contacts have priority over H, so such a contact is processed; Theorem~\ref{thm:P2}(a) allows $q_y=s+2$. For $Z\in\cZ(s)$ every point of $Z$ has height $<s+1.0001$ (Lemma~\ref{lem:P12-window}(b)).
\item \emph{Ties in R1} (equal cumulative gaps). Theorem~\ref{thm:P2} concerns every path that reaches a point of $\Bot(Z)$, in particular every path that reaches an entry candidate alive, before R1 is applied; so the outcome of R1 plays no role.
\item \emph{Axis-parallel squares} ($a(Z)=0$). Then $\cR(Z)=\emptyset$, so Theorem~\ref{thm:P2}(a) is vacuous, and $\Phi(Z)=\emptyset$, so $Z\notin\cZ(s)$. For such squares a set of paths of positive measure can terminate on $\Bot(Z)$: if the bottom side of an axis-parallel square $Z$ has height $\delta$ and no other square meets the strip below it, every floor path starting below $Z$ reaches $\Bot(Z)$ with cumulative gap $\delta$ and terminates there with D. This is why the count in Theorem~\ref{thm:P2Z}(d) uses $a(Z)>0$.
\item \emph{Contacts reached directly from the floor} ($m=0$). Then $q_y=g_0\le\delta$, and Lemma~\ref{lem:P12-core} applies with $m=0$.
\item \emph{Squares touching the floor.} If $(x,0)\in Z$, the point $(x,0)$ is the first-contact point of step $0$, with $t_{S,0}=0$ (Lemma~\ref{lem:P12-traj}\,(F6)). For $Z\in\cZ(s)$ this point would be the lowest point of $Z$, which lies on $\Bot(Z)$ (Lemma~\ref{lem:P12-bot}(a)), and it is excluded by Theorem~\ref{thm:P2Z}(b). No lower bound on $(1-\eps)s-1.0001$ is needed for this.
\item \emph{Squares at $45^\circ$.} Lemma~\ref{lem:P12-bot} does not depend on the convention $\varphi=\frac\pi4$. If $a(Z)=\frac\pi4$, then $f(a(Z))=1$, so $w(s,Z;q)>\frac12\ge\dist(y,\Z)$ and Theorem~\ref{thm:P2}(a) holds trivially. Squares of $\cZ(s)$ have $a(Z)<10^{-4}$.
\item \emph{Exceptional heights.} Theorem~\ref{thm:P2Z}(d),(e) hold at every height $y$; no exceptional set is removed.
\item \emph{Ceiling.} The ceiling statements follow by reflection (Lemma~\ref{lem:P12-refl}), using that $k$ is an integer.
\end{enumerate}
\end{remark}

\section{End collisions (P3)}\label{sec:P3}

This section bounds the measure of the set of floor points whose path ends with an end collision (event~E). It also proves facts that are used again in Section~\ref{sec:P5}: the basic properties of sources and delays (Lemma~\ref{lem:P3-reach}), the fact that a source carries at most one path through a given point (Lemma~\ref{lem:P3-samesource}), the separation of the floor flow from the ceiling flow (Lemma~\ref{lem:P3-separation}), and the form of the multiplicity bound of~[R] that is used in both sections (Lemma~\ref{lem:P3-mult}). Results of~[R] are quoted in Section~\ref{sec:imported}; a reference such as [R, Lemma~\ref{imp:L4.10}] points to that quotation.

\subsection{Setting}\label{ssec:P3-setting}

\paragraph{Standing assumptions.} Throughout this section and Section~\ref{sec:P5}:
\begin{itemize}
\item $k\ge4$ is an integer, $\cP$ is a closed packing of $[0,k]^2$, and $W=k^2-|\cP|$. The \emph{waste} is the set of points of $(0,k)^2$ that lie in no square of $\cP$; its area is $W$.
\item $0<\delta\le10^{-5}$, $0<\amax\le1.5\cdot10^{-6}$ and $\hmax=k/2-1$. (For the parameters of Section~\ref{sec:flows}, $\amax=\alpha(y_0)$, and $\amax\le1.5\cdot10^{-6}$ is constraint~(C1).)
\item $\Fz=F(\amax,\delta,\hmax)$ is the master flow (Definition~\ref{def:F0}) and $\tilde F^0$ the ceiling master flow (Definition~\ref{def:flows-ceiling}). For $x\in(0,k)$ we write $\pi_x$ for the path of $\Fz$ that starts at $(x,0)$ (denoted $\pi^0_x$ in Section~\ref{sec:flows}).
\end{itemize}

\paragraph{Squares and sides.} We use the frames of Definition~\ref{def:flows-frames}: a square $S$ has the phase $\varphi_S\in(-\pi/4,\pi/4]$, the frame $u_S=(\cos\varphi_S,\sin\varphi_S)$, $n_S=(-\sin\varphi_S,\cos\varphi_S)$ and the centre $c_S$, and a point $z$ lies in $S$ if and only if $|(z-c_S)\cdot u_S|\le\frac12$ and $|(z-c_S)\cdot n_S|\le\frac12$. The sides $\Bot(S)$, $\Top(S)$, the right side and the left side are given by $(z-c_S)\cdot n_S=-\frac12$, $(z-c_S)\cdot n_S=\frac12$, $(z-c_S)\cdot u_S=\frac12$ and $(z-c_S)\cdot u_S=-\frac12$, with outer normals $-n_S$, $n_S$, $u_S$, $-u_S$; the right and the left side are the \emph{lateral sides}. The relative interior $\relint\sigma$ of a side $\sigma$ is $\sigma$ without its end points; the four relative interiors and the four vertices partition $\partial S$. The inclination is $a(S)=|\varphi_S|$, and $\theta(S_1,S_2)\in[0,\pi/4]$ is the distance of $\varphi_{S_1}$ and $\varphi_{S_2}$ in $\R/(\pi/2)\Z$; if $|\varphi_{S_1}-\varphi_{S_2}|\le\pi/2$, then $\theta(S_1,S_2)=\min\big(|\varphi_{S_1}-\varphi_{S_2}|,\ \pi/2-|\varphi_{S_1}-\varphi_{S_2}|\big)$. The floor is the pseudo-square $X_0$ with $\varphi_{X_0}=0$, $u_{X_0}=e_1$, $n_{X_0}=e_2$. We write $\sigma_{\rm bot}:=[0,k]\times\{0\}$ and $\sigma_{\rm top}:=[0,k]\times\{k\}$ for the bottom and the top side of the container, $S_0:=\partial[0,k]^2$, and $\theta(S,\sigma):=a(S)$ for a square $S$ and a side $\sigma$ of $[0,k]^2$, as in~[R]. For $X,Y\in\cP\cup\{X_0\}$ and $\Delta:=\varphi_Y-\varphi_X$, by~\eqref{eq:flows-frame},
\begin{equation}\label{eq:P3-frames}
u_X\cdot u_Y=n_X\cdot n_Y=\cos\Delta,\qquad u_Y\cdot n_X=\sin\Delta,\qquad n_Y\cdot u_X=-\sin\Delta .
\end{equation}

\paragraph{Paths.} We recall the rules of Definitions~\ref{def:flow}, \ref{def:R1} and~\ref{def:R2} in the form used below. A ray step of a path starts in a state $(p,n_X,X,g)$, where either $X=X_0$, $p=(x,0)$ and $g=0$, or $X$ is the last square passed by the path (the \emph{source}) and $p\in\relint\Top(X)$; $g<\delta$ is the cumulative gap. With $r(t)=p+t\,n_X$ the step has the parameters $t_S$, $t_W$, $t_D=\delta-g$, $t_H=(\hmax-p_y)/\cos\varphi_X$ and the length $t^*=\min(t_S,t_W,t_D,t_H)$, and it ends with the first applicable of the rules D, W, contact and~H. A contact at $q=r(t_S)\in Y$ is an \emph{entry candidate} $(Y,q)$ if $q\in\relint\Bot(Y)$, and an end collision~E otherwise. At an entry candidate, R1 selects one winner among the paths that reach it alive, and the other ones stop there with~M. The winner stops with~T if $a(Y)\ge\amax$; otherwise it passes $Y$ along the pass segment $[q,q+n_Y]$ and, unless it is completed there (if $q+n_Y$ has height $\ge\hmax$), starts its next ray step at $q+n_Y\in\relint\Top(Y)$.

\paragraph{Facts from Section~\ref{sec:flows}.} We use the following results of Section~\ref{sec:flows}.
\begin{enumerate}[label=(S\arabic*)]
\item\label{it:P3-S1} (Lemmas~\ref{lem:DAG} and~\ref{lem:flows-traj}(a)) Heights increase strictly along every trajectory, and the squares of the successive entry candidates of a path are distinct; in particular a path passes each square at most once.
\item\label{it:P3-S2} (Lemma~\ref{lem:structure}) There is a finite set $\Sigma:=\Sigma_{\Fz}\subset[0,k]$, containing $0$ and $k$, such that on every component $I$ of $[0,k]\setminus\Sigma$ (a \emph{piece}) the combinatorial type of $\pi_x$ does not depend on $x\in I$ (the sources and the events of all steps, the squares and the sides met by the contacts, the outcomes of R1, the termination type), the start point $p_j(x)$, the start gap $g_j(x)$, the length $t^*_j(x)$ and the contact point $q_j(x)$ of the $j$-th ray step are affine functions of $x\in I$, and for no $x\in I$ is a first contact point of $\pi_x$ (a contact, or the end point of a D or W termination with $t_D=t_S$ or $t_W=t_S$) a vertex of a square.
\item\label{it:P3-S3} (Lemma~\ref{lem:expansion}(a),(b)) On a piece, if $X$ is the source of the $j$-th step, then $p_j'(x)=\lambda_j u_X$ with $\lambda_j\ge1$; if the step ends with a contact in $\relint\Bot(Y)$, then $q_j'(x)=\mu_j u_Y$ with $\mu_j=\lambda_j/\cos(\varphi_Y-\varphi_X)\ge\lambda_j\ge1$.
\item\label{it:P3-S4} (Lemma~\ref{lem:flows-traj}(b),(c),(f)) Every ray step starts at a point $p$ with $p_x\in(0,k)$ and cumulative gap $g<\delta$; if its source $X$ is a square, then $p\in\relint\Top(X)$ and $t_S>0$. The open part $\{p+t\,n_X:0<t<t^*\}$ of a gap segment lies in the waste. At a processed contact, $g+t_S<\delta$.
\item\label{it:P3-S5} (Lemma~\ref{lem:flows-R1}(c)) At most one path passes a given entry point $(Y,q)$.
\end{enumerate}
By~\ref{it:P3-S2}, every set of floor points that is defined below through the type of the path is, up to a finite set, a finite union of pieces; in particular it is Lebesgue measurable. We write $|\cdot|$ for Lebesgue measure (on $(0,k)$, on a line or in the plane, according to the context) and $\mathcal H^1$ for length.

\paragraph{The ceiling flow.} Let $\rho(x,y):=(x,k-y)$. By Lemma~\ref{lem:flows-refl}, $\rho\cP$ is a closed packing of $[0,k]^2$ with the same $W$; $\tilde F^0$ is the master flow of $\rho\cP$ transported by $\rho$, and the ceiling versions $\tilde\cE,\tilde\cD,\dots$ of all quantities are computed in this way (Definition~\ref{def:flows-ceiling}). Every statement below about $\Fz$ holds for an arbitrary closed packing; applied to $\rho\cP$ and transported by $\rho$, it gives the corresponding statement for $\tilde F^0$. The reflection $\rho$ is an isometry, maps the waste of $\rho\cP$ onto the waste of $\cP$ and $S_0$ onto itself, and preserves inclinations, the quotient distance $\theta$, distances and areas.

\subsection{Sources and delays}\label{ssec:P3-sources}

\begin{definition}[reaching a point]\label{def:P3-reach}
Let $x\in(0,k)$ and consider a ray step of $\pi_x$ with start state $(p,n_X,X,g)$ and length $t^*$. For $t\ge0$ we say that $\pi_x$ \emph{reaches} the point $q:=p+t\,n_X$ \emph{from the source $X$ with delay $t$}
\begin{itemize}
\item \emph{in the gap case}, if $0<t<t^*$, that is, if $q$ lies in the open part of the gap segment of the step;
\item \emph{in the contact case}, if $t=t^*=t_S$ and the step ends with a contact, that is, $t_D>t^*$ and $t_W>t^*$.
\end{itemize}
A square $X$ is a \emph{source of $q$} if some path reaches $q$ from the source $X$.
\end{definition}

\begin{lemma}[sources and delays]\label{lem:P3-reach}
Let $\pi_x$ reach $q$ from the source $X$ with delay $t$, in a ray step with start state $(p,n_X,X,g)$.
\begin{enumerate}[label=(\alph*)]
\item $g+t<\delta$; in particular $0\le t<\delta$, and $t>0$ if $X$ is a square.
\item If $X$ is a square, then $a(X)<\amax$; in particular $|\varphi_X|<\amax$.
\item $p\in X$ (for $X=X_0$: $p\in\sigma_{\rm bot}$), $p_y\le q_y\le\hmax$, and $q_y<\hmax$ in the gap case.
\item In the gap case, $q$ is a point of the waste in $(0,k)\times(0,\hmax)$.
\item If $X$ is a square, then $(q-c_X)\cdot n_X=\frac12+t$ and $(q-c_X)\cdot u_X=(p-c_X)\cdot u_X\in(-\frac12,\frac12)$. If $X=X_0$, then $q=(x,t)$.
\item The path $\pi_x$ reaches $q$ in at most one of its steps and with at most one delay; and $\pi_x$ has at most one step with source $X$.
\end{enumerate}
\end{lemma}

\begin{proof}
(a) In the gap case $t<t^*\le t_D=\delta-g$. In the contact case $g+t=g+t_S<\delta$ by~\ref{it:P3-S4}. As $g\ge0$, $t<\delta$. If $X$ is a square, then $t_S>0$ by~\ref{it:P3-S4}, so $t=t_S>0$ in the contact case, and $t>0$ in the gap case by definition.

(b) A square is passed only by the winner of R1 at an entry candidate on it, and only if R2 does not stop that winner, that is, only if its inclination is $<\amax$.

(c) If $X$ is a square, $p\in\relint\Top(X)\subset X$; if $X=X_0$, $p=(x,0)\in\sigma_{\rm bot}$. The second coordinate of $n_X$ is $\cos\varphi_X>0$, so $p_y\le q_y=p_y+t\cos\varphi_X$. Since $t\le t^*\le t_H$, we get $q_y\le p_y+t_H\cos\varphi_X=\hmax$, and $q_y<\hmax$ if $t<t^*$.

(d) By~\ref{it:P3-S4} and (c).

(e) For a square $X$, $p\in\relint\Top(X)$ means $(p-c_X)\cdot n_X=\frac12$ and $(p-c_X)\cdot u_X\in(-\frac12,\frac12)$, and $q-p=t\,n_X$. For $X=X_0$, $p=(x,0)$ and $n_{X_0}=e_2$.

(f) By~\ref{it:P3-S1} two different points of the trajectory have different heights, and within one step $t\mapsto p+t\,n_X$ is injective; so $q$ is reached at most once. The source of a step is the floor (for the first step only) or the square passed just before the step, and each square is passed at most once by~\ref{it:P3-S1}; so at most one step has source $X$.
\end{proof}

By Definition~\ref{def:P3-reach}, the gap multiplicity $M_g(q)$ of Definition~\ref{def:flows-line}, the number of $x$ such that $q$ lies in the open part of a gap segment of the trajectory of $\pi_x$, is the number of $x\in(0,k)$ such that $\pi_x$ reaches $q$ in the gap case.

\begin{lemma}[one source, one path]\label{lem:P3-samesource}
\begin{enumerate}[label=(\alph*)]
\item If two paths $\pi_x$ and $\pi_{x'}$ start ray steps at the same point $p\in\relint\Top(X)$ of a square $X$ after passing $X$, then $x=x'$.
\item For every point $q$ and every $X\in\cP\cup\{X_0\}$ there is at most one $x\in(0,k)$ such that $\pi_x$ reaches $q$ from the source $X$.
\end{enumerate}
\end{lemma}

\begin{proof}
(a) A pass of $X$ moves a path from its entry point $z\in\relint\Bot(X)$ to $z+n_X$. So both paths passed $X$ from the entry point $(X,p-n_X)$, and $x=x'$ by~\ref{it:P3-S5}.

(b) Let $\pi_x$ and $\pi_{x'}$ reach $q$ from the source $X$ with delays $t,t'$ in steps with start points $p,p'$. If $X=X_0$, then $p=(x,0)$, $p'=(x',0)$ and $q=p+t\,e_2=p'+t'e_2$, so $x=x'$. If $X$ is a square, then $p,p'\in\Top(X)$, so $p-p'$ is a multiple of $u_X$; and $p-p'=(t'-t)\,n_X$ is a multiple of $n_X$. Hence $p=p'$, and (a) gives $x=x'$.
\end{proof}

\begin{lemma}[floor squares and ceiling squares]\label{lem:P3-separation}
Call a square a \emph{floor square} if it is the source of a point reached by a path of $\Fz$, or if it contains the contact point of a ray step of $\Fz$ that ends with a contact. Define \emph{ceiling squares} in the same way with $\tilde F^0$. Then every floor square meets the half-plane $\{y\le\hmax\}$, every ceiling square meets $\{y\ge k-\hmax\}$, and no square is both a floor square and a ceiling square.
\end{lemma}

\begin{proof}
If $X$ is the source of a point $q$ reached by $\pi_x$ in a step with start point $p$, then $p\in X$ and $p_y\le\hmax$ by Lemma~\ref{lem:P3-reach}(c). A contact point $q$ of a step ending with a contact is reached in the contact case, so $q_y\le\hmax$ by the same lemma. Applying this to $\rho\cP$ and reflecting, every ceiling square meets $\{y\ge k-\hmax\}$. By~[R, Lemma~\ref{imp:L3.1}(a)] a square of inclination $a$ has vertical extent $\cos a+\sin a\le\sqrt2$, which is less than $2=(k-\hmax)-\hmax$. So no square meets both half-planes.
\end{proof}

\subsection{Pair waste and multiplicity}\label{ssec:P3-mult}

For a square $S_1$ and for $S_2$ a square or $S_0$, $R(S_1,S_2)$ is the rectangle of Roth and Vaughan defined in Section~\ref{sec:imported} (before [R, Proposition~\ref{imp:P4.3}]). As in~[R, Lemma~\ref{imp:L4.8}(b)], for a pair $e$ of two squares $S_1,S_2$ we put $R_e:=R(S_1,S_2)$, and for a pair $e=\{S,\sigma\}$ of a square $S$ and a side $\sigma$ of $[0,k]^2$ we put $R_e:=R(S,S_0)$; $f_e$ is the area of the waste in $R_e$. We use~[R, Lemma~\ref{imp:L4.8}(a)]: \emph{if $0<d\le10^{-4}$ and $S_1$ is a square and $S_2$ a square or $S_0$, at distance $<d$, then $f\ge\theta(2-\theta)/40$ for the waste $f$ in $R(S_1,S_2)$, $\theta=\theta(S_1,S_2)$}; and the multiplicity bound in the following form.

\begin{lemma}[multiplicity at waste points, for families of pairs]\label{lem:P3-mult}
Let $0<d\le10^{-4}$, and let $\mathcal F$ be a set of unordered pairs, each of which consists either of two squares of $\cP$ at distance $<d$, or of a square of $\cP$ and one of the two sides $\sigma_{\rm bot},\sigma_{\rm top}$ at distance $<d$. Then every point of the waste lies in $R_e$ for at most $K$ pairs $e\in\mathcal F$, where $K=9$ by~[R, Lemma~\ref{imp:L4.10} and Remark~\ref{imp:R4.11}], and $K=13$ by~[R, Lemma~\ref{imp:L4.9}] and the paragraph following its proof. (In this section and in Section~\ref{sec:P5}, $K$ denotes this multiplicity bound, $K=K_w^*=9$ or $K=K_w=13$; it is not the constant $K=27$ of [R, Lemma~\ref{imp:L4.8}(b)].) Consequently
\begin{equation}\label{eq:P3-sumf}
\sum_{e\in\mathcal F}f_e\ \le\ K\,W .
\end{equation}
\end{lemma}

\begin{proof}
This is Proposition~\ref{imp:pairs} with $\sigma=\sigma_{\rm bot}$, $\sigma'=\sigma_{\rm top}$ and $\mathfrak F=\mathcal F$. The formal statements of [R, Lemmas~\ref{imp:L4.9} and~\ref{imp:L4.10}] concern the edges of a visibility graph $G_h$ or $G_v$, whereas the pairs of $\mathcal F$ need not be such edges. The form used here is the one that [R] states for the proofs of these lemmas: after the proof of [R, Lemma~\ref{imp:L4.9}], that the bound $13$ ``holds for any family of such pairs at distance $<d$'', the only sides involved being two opposite ones; and in [R, Remark~\ref{imp:R4.11}], that the proof of [R, Lemma~\ref{imp:L4.10}] bounds ``the number of \emph{all} pairs of squares (or of a square and one of two opposite sides) at distance less than $d$ whose rectangles contain $p$'' by $9$, $p$ being a point of the waste. In $\mathcal F$ the only sides are $\sigma_{\rm bot}$ and $\sigma_{\rm top}$, which are opposite. Since $f_e=\int_{\rm waste}\mathbf 1_{R_e}$, summation gives $\sum_{e\in\mathcal F}f_e=\int_{\rm waste}\#\{e\in\mathcal F:\ p\in R_e\}\,dp\le K\,W$.
\end{proof}

\begin{remark}[rectangles of a square and a side near a corner]\label{rem:P3-corner}
Let $e=\{S,\sigma_{\rm bot}\}$ with $S$ at distance $<d$ from $\sigma_{\rm bot}$. Then $S$ is at distance $<d$ from $S_0\supset\sigma_{\rm bot}$, so~[R, Lemma~\ref{imp:L4.8}(a)] applies to $(S,S_0)$ and gives $f_e\ge a(S)(2-a(S))/40$, because $\theta(S,S_0)=a(S)$. The rectangle $R_e=R(S,S_0)$ stands on a side of $[0,k]^2$ nearest to the centre of $S$. If $S$ is close to a corner, this may be the left or the right side rather than $\sigma_{\rm bot}$. Nothing changes in this case. The distance from $S$ to a side of $[0,k]^2$ is the distance from its centre to that side minus $\frac12(\cos a(S)+\sin a(S))$ (proof of [R, Proposition~\ref{imp:P4.5}], part~(b)); so $S$ is at distance $<d$ from the side on which $R_e$ stands as well. [R, Lemma~\ref{imp:L4.8}(a)] makes no assumption on the side on which $R(S,S_0)$ stands, and the proofs of [R, Lemmas~\ref{imp:L4.9} and~\ref{imp:L4.10}] treat the rectangle $R(P,S_0)$ of a pair of a square $P$ and a side whatever side of $[0,k]^2$ it stands on (the proof of [R, Lemma~\ref{imp:L4.10}] shows that all such rectangles containing a given waste point stand on one and the same side); see also the paragraph after Proposition~\ref{imp:pairs}. The same applies to pairs $\{S,\sigma_{\rm top}\}$.
\end{remark}

\begin{remark}[dependence on {[R, Lemma~\ref{imp:L4.10}]}]\label{rem:P3-dependence}
[R, Lemma~\ref{imp:L4.10}] is computer-assisted in~[R] (Remark~\ref{rem:imp-dep}). Every constant below, in this section and in Section~\ref{sec:P5}, that is computed with $K=9$ depends on it; the constants computed with $K=13$ use only [R, Lemma~\ref{imp:L4.9}]. Both values are given.
\end{remark}

\subsection{End collisions}\label{ssec:P3-E}

By Lemma~\ref{lem:flows-top}, a first contact point never lies in the relative interior of a top side, and an end collision occurs at a vertex or at a point of the relative interior of a lateral side. The next lemma determines the lateral side.

\begin{lemma}[the lateral side that is hit]\label{lem:P3-sides}
Consider a ray step with direction $d=n_X$ that ends with a contact at a point $q=r(t_S)$ of the relative interior of a lateral side $\sigma$ of a square $Y$, with outer normal $o$. Then $o\cdot d<0$. Consequently $\varphi_Y\ne\varphi_X$, and $\sigma$ is the left side of $Y$ if $\varphi_Y>\varphi_X$ and the right side of $Y$ if $\varphi_Y<\varphi_X$; we call this side $\sigma_{X,Y}$. An end collision of a step with source $X$ at a square $Y$ occurs at a vertex of $Y$ or in $\relint\sigma_{X,Y}$.
\end{lemma}

\begin{proof}
Write $Y=\{z:|(z-c_Y)\cdot u_Y|\le\frac12,\ |(z-c_Y)\cdot n_Y|\le\frac12\}$. At $q\in\relint\sigma$ the defining inequality $(z-c_Y)\cdot o\le\frac12$ is an equality and the other three are strict; they remain strict in a small disc $B(q,r_0)$, so that
\[
Y\cap B(q,r_0)=\{z\in B(q,r_0):\ (z-q)\cdot o\le0\}.
\]
If $t_S=0$, the source is the floor by~\ref{it:P3-S4} and $q=(x,0)$; the points $q-\varepsilon n_Y$, $0<\varepsilon$ small, lie in $\sigma\subset Y$ (as $\sigma$ has the direction $n_Y$ and $q\in\relint\sigma$) and have negative height, which contradicts $Y\subset[0,k]^2$. So $t_S>0$. For $0<\varepsilon<\min(t_S,r_0)$ the point $q-\varepsilon d=r(t_S-\varepsilon)$ lies in $B(q,r_0)$ and not in $Y$, because $t_S$ is the least parameter at which the ray meets a square other than $X$ (and $Y\ne X$). So $(q-\varepsilon d-q)\cdot o>0$, that is, $o\cdot d<0$.

The outer normals of the lateral sides are $\pm u_Y$, and $u_Y\cdot d=\sin(\varphi_Y-\varphi_X)$ by~\eqref{eq:P3-frames}. If $\varphi_Y=\varphi_X$ both products vanish, which is impossible. Otherwise exactly one of $\pm u_Y$ has negative product with $d$: $-u_Y$ (left side) if $\varphi_Y>\varphi_X$, and $u_Y$ (right side) if $\varphi_Y<\varphi_X$. The last statement follows with Lemma~\ref{lem:flows-top}.
\end{proof}

Let $\cE\subset(0,k)$ be the set of $x$ such that $\pi_x$ terminates with~E, and $L_E:=|\cE|$; $\tilde\cE$ and $\tilde L_E$ are their ceiling versions (Definitions~\ref{def:flows-losses} and~\ref{def:flows-ceiling}). Let $V\subset(0,k)$ be the set of $x$ such that $\pi_x$ has a contact at a vertex of a square; by~\ref{it:P3-S2}, $V\subset\Sigma$ is finite. For $X\in\cP\cup\{X_0\}$ and a square $Y$ with $\varphi_Y\ne\varphi_X$, let
\[
E_{X,Y}:=\{x\in(0,k):\ \pi_x\ \text{terminates with E at a point of}\ \relint\sigma_{X,Y},\ \text{in a step with source}\ X\}.
\]
By Lemma~\ref{lem:P3-sides},
\begin{equation}\label{eq:P3-Ecover}
\cE\ \subset\ V\cup\bigcup_{(X,Y)}E_{X,Y},
\end{equation}
the union being over the finitely many ordered pairs $(X,Y)$ with $E_{X,Y}\ne\emptyset$. For such a pair, $X$ is the floor or a square with $|\varphi_X|<\amax$ (Lemma~\ref{lem:P3-reach}(b)).

\begin{lemma}[one source and one target]\label{lem:P3-pair}
Let $X\in\cP\cup\{X_0\}$ and $Y\in\cP$ with $E_{X,Y}\ne\emptyset$, and put $\Delta:=\varphi_Y-\varphi_X$. Then:
\begin{enumerate}[label=(\alph*)]
\item if $X$ is a square, $\dist(X,Y)<\delta$; if $X=X_0$, $\dist(Y,\sigma_{\rm bot})<\delta$;
\item $|E_{X,Y}|\le\delta\tan|\Delta|$.
\end{enumerate}
\end{lemma}

\begin{proof}
For $x\in E_{X,Y}$ let $p(x)$, $q(x)$ and $t(x)$ be the start point, the contact point and the delay of the step of $\pi_x$ with source $X$; the step is unique by Lemma~\ref{lem:P3-reach}(f), and $\pi_x$ reaches $q(x)$ from $X$ with delay $t(x)$ in the contact case.

(a) $q(x)\in Y$, $p(x)\in X$ (resp.\ $p(x)\in\sigma_{\rm bot}$), and $|q(x)-p(x)|=t(x)<\delta$ by Lemma~\ref{lem:P3-reach}(a),(c).

(b) Put $\xi(x):=(p(x)-c_X)\cdot u_X$ if $X$ is a square and $\xi(x):=x$ if $X=X_0$. Since $q(x)-p(x)$ is parallel to $n_X$, $\xi(x)=(q(x)-c_X)\cdot u_X$, respectively $\xi(x)=q(x)_x$.

\emph{$\xi$ is injective on $E_{X,Y}$.} For $X=X_0$ this is clear. If $X$ is a square and $\xi(x)=\xi(x')$, then $p(x)$ and $p(x')$ are points of $\Top(X)$ with the same $u_X$-coordinate, so $p(x)=p(x')$, and Lemma~\ref{lem:P3-samesource}(a) gives $x=x'$.

\emph{$\xi$ expands.} By~\ref{it:P3-S2}, $E_{X,Y}\setminus\Sigma$ is the union of the pieces $I$ that it contains. On such an $I$ the step with source $X$ has a fixed index $j$, and by~\ref{it:P3-S3} $\xi$ is affine with $\xi'=\lambda_j\ge1$ (for $X=X_0$, $\xi'=1$). The intervals $\xi(I)$ are pairwise disjoint by injectivity, so
\[
|E_{X,Y}|=\sum_I|I|\le\sum_I|\xi(I)|=\mathcal H^1\Big(\bigcup_I\xi(I)\Big)\le\mathcal H^1\big(\xi(E_{X,Y})\big).
\]

\emph{Location of the contact points.} By Lemma~\ref{lem:P3-reach}(a),(e), $q(x)$ lies on the line $\ell$ of $\sigma_{X,Y}$ and satisfies $\frac12<(q(x)-c_X)\cdot n_X<\frac12+\delta$ if $X$ is a square, respectively $0\le q(x)_y<\delta$ if $X=X_0$. The line $\ell$ has the direction $n_Y$. Moving along $\ell$ at unit speed, $(z-c_X)\cdot n_X$ changes at the rate $n_Y\cdot n_X=\cos\Delta>0$ and $(z-c_X)\cdot u_X$ at the rate $n_Y\cdot u_X=-\sin\Delta$, by~\eqref{eq:P3-frames} (for $X=X_0$: $z_y$ and $z_x$, at the rates $\cos\varphi_Y$ and $-\sin\varphi_Y$; here $\Delta=\varphi_Y$). Here $\cos\Delta>0$ because $|\varphi_X|<\amax$ and $|\varphi_Y|\le\pi/4$. Hence the points of $\ell$ satisfying the above condition form a segment of length $\delta/\cos\Delta$, and the values of $(z-c_X)\cdot u_X$ (resp.\ $z_x$) on this segment form an interval of length
\[
\frac{\delta}{\cos\Delta}\,|\sin\Delta|=\delta\tan|\Delta| .
\]
As $\xi(E_{X,Y})$ is contained in this interval, $|E_{X,Y}|\le\mathcal H^1(\xi(E_{X,Y}))\le\delta\tan|\Delta|$.
\end{proof}

The bound of Lemma~\ref{lem:P3-pair}(b) is written with $|\varphi_Y-\varphi_X|$ and not with $\theta(X,Y)$: the phase difference can exceed $\pi/4$ when $Y$ is inclined by nearly $\pi/4$ (then $\theta(X,Y)=\pi/2-|\Delta|$). This includes squares with $\varphi_Y=\pi/4$, whose bottom side is, by the $45^\circ$ convention of Definition~\ref{def:flows-frames}, the lower right side; such squares are never sources, since $\amax<\pi/4$.

\subsection{The bound for end collisions}\label{ssec:P3-prop}

\begin{proposition}[P3]\label{prop:P3}
Under the standing assumptions of Section~\ref{ssec:P3-setting}, let $K$ be the constant of Lemma~\ref{lem:P3-mult} and $c_{\mathrm E}:=41.933$. Then
\[
L_E+\tilde L_E\ \le\ 2\,c_{\mathrm E}\,K\,\delta\,W .
\]
With $K=9$ (using [R, Lemma~\ref{imp:L4.10}]) this gives $L_E+\tilde L_E\le754.8\,\delta W$; with $K=13$ (using only [R, Lemma~\ref{imp:L4.9}]) it gives $L_E+\tilde L_E\le1090.3\,\delta W$. Both numbers are rounded up.
\end{proposition}

\begin{proof}
\emph{Step 1: the ratio $c_{\mathrm E}$.} Let $X\in\cP\cup\{X_0\}$ and $Y\in\cP$ with $E_{X,Y}\ne\emptyset$, $\Delta=\varphi_Y-\varphi_X$ and $\theta:=\theta(X,Y)$ (for $X=X_0$, $\theta=\theta(Y,S_0)=a(Y)$). We claim that $\theta>0$ and
\begin{equation}\label{eq:P3-ratio}
40\tan|\Delta|\ \le\ c_{\mathrm E}\,\theta(2-\theta).
\end{equation}
Indeed $\varphi_X\ne\varphi_Y$ (Lemma~\ref{lem:P3-sides}), $|\varphi_X|<\amax$ and $\varphi_Y\in(-\pi/4,\pi/4]$, so $0<|\Delta|<\pi/4+\amax<\pi/2$ and $\theta=\min(|\Delta|,\pi/2-|\Delta|)>0$.
\begin{itemize}
\item If $|\Delta|\le\pi/4$, then $\theta=|\Delta|$. On $(0,\pi/4]$ the function $\tan\theta/(\theta(2-\theta))$ is the product of the positive increasing functions $\tan\theta/\theta$ and $1/(2-\theta)$, so
\[
\frac{40\tan|\Delta|}{\theta(2-\theta)}\ \le\ \frac{40}{(\pi/4)(2-\pi/4)}=41.93109\ldots
\]
\item If $\pi/4<|\Delta|<\pi/4+\amax$, then $\theta=\pi/2-|\Delta|\in(\pi/4-\amax,\pi/4)$. As $x(2-x)$ is increasing on $[0,1]$ and $\tan$ is increasing,
\[
\frac{40\tan|\Delta|}{\theta(2-\theta)}\ <\ \frac{40\tan(\pi/4+\amax)}{(\pi/4-\amax)(2-\pi/4+\amax)}\ \le\ 41.9313
\]
for $\amax\le1.5\cdot10^{-6}$; the middle term increases with $\amax$ (its numerator increases, and the derivative of its denominator with respect to $\amax$ is $\pi/2-2-2\amax<0$), and its value at $\amax=1.5\cdot10^{-6}$ is $41.93124\ldots$.
\end{itemize}
In both cases the ratio is at most $41.933=c_{\mathrm E}$.

\emph{Step 2: one pair.} Let $(X,Y)$ be as in Step~1, and let $e:=\{X,Y\}$ if $X$ is a square and $e:=\{Y,\sigma_{\rm bot}\}$ if $X=X_0$. By Lemma~\ref{lem:P3-pair}(a) the members of $e$ are at distance $<\delta\le10^{-4}$. By~[R, Lemma~\ref{imp:L4.8}(a)] with $d=10^{-4}$, applied to $(X,Y)$, respectively to $(Y,S_0)$ (Remark~\ref{rem:P3-corner}), $f_e\ge\theta(2-\theta)/40$. With Lemma~\ref{lem:P3-pair}(b) and~\eqref{eq:P3-ratio},
\begin{equation}\label{eq:P3-onepair}
|E_{X,Y}|\ \le\ \delta\tan|\Delta|\ \le\ c_{\mathrm E}\,\delta\,\frac{\theta(2-\theta)}{40}\ \le\ c_{\mathrm E}\,\delta\,f_e .
\end{equation}

\emph{Step 3: the floor flow.} Let $\mathcal Q^{\rm E}_b$ be the set of the pairs $e$ of Step~2, over all ordered pairs $(X,Y)$ with $E_{X,Y}\ne\emptyset$. A pair of two squares arises from at most two ordered pairs, $(X,Y)$ and $(Y,X)$, and a pair $\{Y,\sigma_{\rm bot}\}$ from the single ordered pair $(X_0,Y)$. Since $V$ is finite, \eqref{eq:P3-Ecover} and~\eqref{eq:P3-onepair} give
\[
L_E\ \le\ \sum_{(X,Y)}|E_{X,Y}|\ \le\ 2\,c_{\mathrm E}\,\delta\sum_{e\in\mathcal Q^{\rm E}_b}f_e .
\]

\emph{Step 4: the ceiling flow.} Apply Step~3 to $\rho\cP$ and transport the pairs to $\cP$ by $\rho$. The pair $\{\rho Y,\sigma_{\rm bot}\}$ of $\rho\cP$ becomes $\{Y,\sigma_{\rm top}\}$. Since $\rho$ is an isometry that maps $S_0$ onto itself, it maps $R(\rho S_1,\rho S_2)$ onto $R(S_1,S_2)$ (for a pair with a side, choosing as nearest point of $S_0$ the image of the chosen one), and it maps the waste of $\rho\cP$ onto the waste of $\cP$; so the values $f_e$ and the distances are unchanged. Thus $\tilde L_E\le2c_{\mathrm E}\delta\sum_{e\in\mathcal Q^{\rm E}_c}f_e$, where $\mathcal Q^{\rm E}_c$ is a set of pairs of squares of $\cP$ at distance $<\delta$ and of pairs $\{Y,\sigma_{\rm top}\}$ with $\dist(Y,\sigma_{\rm top})<\delta$.

\emph{Step 5: summation.} The two sets $\mathcal Q^{\rm E}_b$ and $\mathcal Q^{\rm E}_c$ are disjoint. A pair of two squares in $\mathcal Q^{\rm E}_b$ consists of floor squares (each is a source or contains a contact point of $\Fz$), and a pair of two squares in $\mathcal Q^{\rm E}_c$ consists of ceiling squares, so these pairs differ by Lemma~\ref{lem:P3-separation}. The pairs with a side differ by the side. Hence $\mathcal F:=\mathcal Q^{\rm E}_b\cup\mathcal Q^{\rm E}_c$ is a set of pairs as in Lemma~\ref{lem:P3-mult} (with $d=10^{-4}$), and by~\eqref{eq:P3-sumf}
\[
L_E+\tilde L_E\ \le\ 2\,c_{\mathrm E}\,\delta\sum_{e\in\mathcal F}f_e\ \le\ 2\,c_{\mathrm E}\,K\,\delta\,W .
\]

\emph{Step 6: numbers.} For $K=9$, $18\cdot41.933=754.794\le754.8$. For $K=13$, $26\cdot41.933=1090.258\le1090.3$.
\end{proof}

\begin{remark}[what the proof uses]\label{rem:P3-uses}
The proof of Proposition~\ref{prop:P3} uses R1 only through Lemma~\ref{lem:P3-samesource}, the expansion $\lambda_j\ge1$ of~\ref{it:P3-S3}, the classification of end collisions (Lemmas~\ref{lem:flows-top} and~\ref{lem:P3-sides}), [R, Lemma~\ref{imp:L4.8}(a)] and Lemma~\ref{lem:P3-mult}. It does not use the hypothesis $\delta\le\frac12\cos2\amax$ (constraint~(C6)) and none of the results of Section~\ref{sec:P5}. The degenerate cases are covered as follows (see also Remark~\ref{rem:flows-degenerate}). A ray that runs along a lateral side of $Y$ ($\varphi_Y=\varphi_X$) meets $Y$ first at a vertex, so the corresponding $x$ lie in the finite set $V$. A square resting on the floor gives contacts with delay $0$ from the floor; these are entry candidates or vertex contacts (a contact with delay $0$ is never in the relative interior of a lateral side, by the proof of Lemma~\ref{lem:P3-sides}). Contacts at height exactly $\hmax$ are processed (contact precedes~H) and are included. A ray step never starts on a side wall of the container (\ref{it:P3-S4}). Ties between $t_S$, $t_W$, $t_D$, $t_H$ (which may occur for sets of $x$ of positive measure) are decided by the priority rule and do not affect Lemma~\ref{lem:P3-reach}.
\end{remark}

\section{The shadow inequality (\texorpdfstring{P4$^*$}{P4*})}\label{sec:P4}

This section works with one horizontal line $\ell_y$ and one flow $F$, which is either the master flow $\Fz$ or a scale flow $\Fs$. On $\ell_y$ we compare the chord of each square with the floor measure of the paths that cross the square there. The main result, Theorem~\ref{thm:P4}, says that the total shadow $\Sh^F(y)$ is at most the tilt loss below $y$ plus the other losses $\Lam^F(y)$, for every height $y$ outside a finite set $\Exc$ that does not depend on $s$. Proposition~\ref{prop:wall} bounds the wall part of the losses of $\Fs$. The remaining part $\Lam_0$ is bounded in Section~\ref{sec:P5}. Section~\ref{sec:P7} uses the result in the form of Corollary~\ref{cor:P4-scale}.

The arguments use the definitions and the structure lemmas of Section~\ref{sec:flows}. They use no result of~[R].

\subsection{Setting and notation}\label{ssec:P4-notation}

Throughout this section $\cP$ is a closed packing in $[0,k]^2$, where $k\ge4$ is an integer, and the parameters are those of Section~\ref{sec:flows}. The letter $F$ denotes the master flow $\Fz$ (Definition~\ref{def:F0}) or a scale flow $\Fs$ with $s\in[y_0,y_1]$ (Definition~\ref{def:Fs}). We put
\[
\alpha_{\Fz}:=\amax,\quad h_{\Fz}:=\hmax=\tfrac k2-1,\qquad \alpha_{\Fs}:=\alpha(s),\quad h_{\Fs}:=s+2 .
\]
Since $y_0\le s\le y_1=\frac k2-3$ and $\amax=\alpha(y_0)$, we have
\begin{equation}\label{eq:P4-param}
0<\alpha(s)\le\amax<\tfrac\pi4,\qquad \delta>0,\qquad 2<h_{\Fs}\le h_{\Fz}=\hmax<k .
\end{equation}
We recall the notions of Section~\ref{sec:flows} that are used below, in the form in which we use them.

\begin{enumerate}[(N1)]
\item \emph{Frames.} A square $S$ has centre $c_S$ and phase angle $\varphi_S\in(-\frac\pi4,\frac\pi4]$. Put $u_S=(\cos\varphi_S,\sin\varphi_S)$ and $n_S=(-\sin\varphi_S,\cos\varphi_S)$. Then $S$ is the set of points $c_S+\xi u_S+\zeta n_S$ with $|\xi|,|\zeta|\le\frac12$. Its bottom side, top side and lateral sides are
\[
\Bot(S)=\{c_S-\tfrac12n_S+\xi u_S\},\qquad \Top(S)=\{c_S+\tfrac12n_S+\xi u_S\},\qquad \{c_S\pm\tfrac12u_S+\zeta n_S\},
\]
with $|\xi|\le\frac12$, $|\zeta|\le\frac12$. The relative interiors of $\Bot(S)$ and $\Top(S)$ are given by $|\xi|<\frac12$. The inclination is $a(S)=|\varphi_S|\in[0,\frac\pi4]$. The floor is treated as a pseudo-square $X_0$ with $\varphi_{X_0}=0$, $u_{X_0}=e_1$, $n_{X_0}=e_2$.
\item \emph{Ray steps and events} (Definition~\ref{def:flow}). The path $\pi_x$ of a floor point $x\in(0,k)$ in $\Fz$ starts in the state $((x,0),e_2,X_0,0)$. In a state $(p,n_X,X,g)$ it follows the ray $r(t)=p+tn_X$, $t\ge0$. Here $t_S$ is the least $t\ge0$ with $r(t)$ in a square $Y\ne X$ ($+\infty$ if there is none), $t_W$ the least $t\ge0$ with $r(t)_x\in\{0,k\}$ ($+\infty$ if there is none), $t_D=\delta-g$, $t_H=(\hmax-p_y)/(n_X)_y$, and $t^*=\min(t_S,t_W,t_D,t_H)$. At $t^*$ the first applicable rule in the order D ($t_D=t^*$), W ($t_W=t^*$), contact ($t_S=t^*$), H is applied. A contact at $q=r(t_S)$ belongs to exactly one square $Y$. It is an \emph{entry candidate} $(Y,q)$ if $q\in\relint\Bot(Y)$; otherwise the path terminates there with label E.
\item \emph{Entry processing} (Definitions~\ref{def:R1} and~\ref{def:R2}). At an entry candidate $(Y,q)$ of $\Fz$, let $\mathcal A_Y(q)$ be the set of floor points whose paths reach $(Y,q)$ alive. Rule R1 keeps the element of $\mathcal A_Y(q)$ for which $(g_x(q),x)$ is lexicographically smallest, the \emph{R1 winner}. Here $g_x(q)$ is the cumulative gap at $q$. Every other element of $\mathcal A_Y(q)$ terminates at $q$ with label M. Rule R2: if $a(Y)\ge\amax$, the winner terminates at $q$ with label T (order ``M before T''). Otherwise the winner \emph{traverses} $Y$ (it passes $Y$, in the terminology of Section~\ref{sec:flows}) along the pass segment $[q,q+n_Y]$ and continues from the exit point $q+n_Y\in\relint\Top(Y)$ in the state $(q+n_Y,n_Y,Y,g)$; it is completed (label H) if the exit point has height $\ge\hmax$. R1 is decided in $\Fz$ only. Since pairs $(g,x)$ with different $x$ are different and the lexicographic order is total, \emph{an entry candidate has at most one R1 winner}.
\item \emph{Scale flows} (Definition~\ref{def:Fs}). For $s\in[y_0,y_1]$ the path of $x$ in $\Fs$ is its path in $\Fz$, cut at the first of the following points:
 \begin{itemize}
 \item[($T_s$)] an entry candidate at which the path is the R1 winner and $a(Y)\ge\alpha(s)$ (label T);
 \item[($H_s$)] a point of a ray step at height $s+2$ with $t^{(s)}_H:=(s+2-p_y)/(n_X)_y<\min(t_S,t_W,t_D)$, or an exit point of height $\ge s+2$ (label H);
 \end{itemize}
 if the $\Fz$ path terminates before reaching such a point, the $\Fs$ path terminates there as well, with the same label. An $\Fz$ termination with label T that is reached before any cut is an entry candidate with an R1 winner and $a(Y)\ge\amax\ge\alpha(s)$, so it is a $T_s$ point. We write $\cT_s$ for the set of $x$ that terminate with label T in $\Fs$, and $N(s):=|\cT_s|$.
\item \emph{Trajectories} (described in detail in \S\ref{ssec:P4-traj}). $\Gamma_x=\Gamma^F_x$ is the trajectory of $\pi_x$ in $F$, $e(x)$ its termination point, and $\tau(x)=\tau_F(x)$ the height of $e(x)$, except that $\tau(x):=h_F$ for label H. The path is \emph{alive at $y$} if $\tau(x)>y$. Every $x$ has exactly one termination type $\star\in\{T,D,E,M,W,H\}$, which we call its \emph{label}.
\item \emph{Losses.} $L_\star(y):=|\{x:\ x\text{ has label }\star,\ \tau(x)<y\}|$ for $\star\in\{T,D,E,M,W\}$, and $L(y):=\sum_{\star\in\{T,D,E,M,W\}}L_\star(y)$. For $\Fz$, $\cD,\cM,\cE$ are the sets of $x$ with label D, M, E, and $L_E:=|\cE|$. Here $|\cdot|$ is Lebesgue measure.
\item \emph{Lines.} $\ell_y:=\R\times\{y\}$, identified with $\R$ by $(q,y)\mapsto q$. The waste is $(0,k)^2\setminus\bigcup_{Z\in\cP}Z$, and $\mathrm{Wa}_y:=\ell_y\cap\text{waste}$. The line waste is $\omega(y):=|([0,k]\times\{y\})\setminus\bigcup_{Z\in\cP}Z|$, and $c_Z(y):=|Z\cap\ell_y|$ is the length of the closed chord.
\item \emph{Line quantities.} Let $y\in(0,h_F)$. If $\tau(x)\ge y$, then $P_y(x)$ denotes the unique point of $\Gamma_x$ on $\ell_y$ (Theorem~\ref{thm:P4}(a)) and $X_y(x)$ its abscissa. Put
 \begin{align*}
 A_Z(y)&:=\{x:\ \tau(x)>y,\ P_y(x)\in\operatorname{int}Z\},& \mu_Z(y)&:=|A_Z(y)|,\\
 A_{\rm gap}(y)&:=\{x:\ \tau(x)>y,\ P_y(x)\in\mathrm{Wa}_y\},& \Fgap(y)&:=|A_{\rm gap}(y)|.
 \end{align*}
 $\nu_y$ is the image of Lebesgue measure on $A_{\rm gap}(y)$ under $X_y$, and $\rho_y$ is its density (it exists by Theorem~\ref{thm:P4}(e)). Then
 \begin{align*}
 \Sh(y)&:=\sum_{Z\in\cP}\big(c_Z(y)-\mu_Z(y)\big)_+, & \Ov(y)&:=\sum_{Z\in\cP}\big(\mu_Z(y)-c_Z(y)\big)_+,\\
 \Ovgap(y)&:=\int_{\mathrm{Wa}_y}\big(\rho_y(q)-1\big)_+\dd q, & \Lam^F(y)&:=\big[L_D+L_E+L_M+L_W\big](y)+\Ovgap(y),
 \end{align*}
 and $\Lam_0:=|\cD|+|\cM|+L_E+\sup_{0<y<\hmax}\Ovgap^{\Fz}(y)$, $\Lam^*(s):=\Lam_0+2(s+2)\tan\alpha(s)$.
\end{enumerate}
We also use two quantities that appear only in this section:
\begin{equation}\label{eq:P4-BdU}
\mathrm{Bd}(y):=k-L(y)-\sum_{Z\in\cP}\mu_Z(y)-\Fgap(y),\qquad U(y):=\int_{\mathrm{Wa}_y}\big(1-\rho_y(q)\big)_+\dd q .
\end{equation}
$U(y)$ is the part of the line waste that the gap crossers do not use. A superscript $F$ (as in $\Sh^F$, $L^F_T$, $\tau_F$) indicates the flow when necessary. All sums over $Z$ are finite. The sets in (N6)--(N8) are finite unions of intervals and points (Lemma~\ref{lem:P4-meas}), so all these measures are defined. By Theorem~\ref{thm:P4}(e), $\Ovgap(y)\le\int_{\mathrm{Wa}_y}\rho_y=\Fgap(y)\le k$, hence $\Lam_0<\infty$.

\subsection{Trajectories and event points}\label{ssec:P4-traj}

Section~\ref{sec:flows} also records most of the facts of this subsection and of Sections~\ref{ssec:P4-exc}--\ref{ssec:P4-e}. We restate them in the form used here, with proofs, so that the proof of Theorem~\ref{thm:P4} can be read on its own; the statements agree with those of Section~\ref{sec:flows}, under the same hypotheses and conventions. The correspondence is as follows. Lemma~\ref{lem:P4-refine} restates Lemma~\ref{lem:flows-Ts}(f). Definition~\ref{def:P4-traj} is Definition~\ref{def:flows-traj}. Lemma~\ref{lem:P4-traj}(i),(iii) are contained in Lemmas~\ref{lem:flows-traj}(a), \ref{lem:flows-Ts}(c) and~\ref{lem:structure}(b) (with Lemma~\ref{lem:flows-Ts}(f) for $\Fs$), and Lemma~\ref{lem:P4-traj}(ii) is Lemma~\ref{lem:flows-traj}(h), proved here also for $\Fs$ with $h_{\Fs}=s+2$. The first assertion of Lemma~\ref{lem:P4-position}(0) and Lemma~\ref{lem:P4-position}(i)--(iii) are contained in Lemma~\ref{lem:flows-traj}(b),(c),(e),(g) and the remark after Lemma~\ref{lem:flows-Ts}. Definition~\ref{def:P4-exc} defines the sets $\Exc_F$ and $\Exc$ of Definition~\ref{def:flows-exc}, and Lemma~\ref{lem:P4-exc}(i)--(iii) is Lemma~\ref{lem:flows-exc}. Lemma~\ref{lem:P4-crossing}(i),(ii),(iv) are contained in Definition~\ref{def:flows-line} and in Lemma~\ref{lem:flows-density}(a),(b) and its proof. Lemma~\ref{lem:P4-meas} extends Lemma~\ref{lem:flows-meas}(a). Theorem~\ref{thm:P4}(e) contains Lemma~\ref{lem:flows-density}(c); the bound $L^{\Fs}_T\le N(s)$ in Theorem~\ref{thm:P4}(f) is Corollary~\ref{cor:flows-LT}; Theorem~\ref{thm:P4}(g) proves in detail the statement of Remark~\ref{rem:flows-order}; and Remark~\ref{rem:P4-borel} is Lemma~\ref{lem:flows-meas}(d), extended to $\Ovgap$ and $U$. We first collect the structural facts from Section~\ref{sec:flows} that are used.

\begin{lemma}[structure]\label{lem:P4-recall}
The following facts are contained in Lemmas~\ref{lem:DAG}, \ref{lem:flows-R1}, \ref{lem:structure} and~\ref{lem:expansion}.
\begin{enumerate}[(i)]
\item Along every path of $\Fz$, the centre heights of the squares at which the path has entry candidates increase strictly (Lemma~\ref{lem:DAG}). Hence a path has at most one entry candidate on each square, and it traverses each square at most once. Each set $\mathcal A_Y(q)$ is finite, and R1 is well defined (Section~\ref{sec:flows}).
\item There is a finite set $\Sigma_0\subset[0,k]$ containing $0$ and $k$ with the following property on every component $I$ of $[0,k]\setminus\Sigma_0$. The combinatorial type of the $\Fz$ path of $x$ is the same for all $x\in I$. The type consists of the number of ray steps, the source of each step (the floor or the last traversed square), the rule ending each step (D, W, contact or H), for a contact the square and the side or vertex contacted, for an entry candidate the outcomes of R1 and R2, and the label. Moreover the start $p_j(x)$, the cumulative gap $g_j(x)$ at the start, the length $t^*_j(x)$ and the end point $q_j(x)=p_j(x)+t^*_j(x)\,n_{X_j}$ of step $j$ are affine functions of $x\in I$. Here $X_j$ is the source of step $j$.
\item On such an $I$, if step $j$ has source $X$, then $p_j'=\lambda_ju_X$ with $|\lambda_j|\ge1$ (for the floor $\lambda_0=1$). If step $j$ ends at an entry candidate on $\Bot(Y)$, then $q_j'=\mu_ju_Y$ with $\mu_j=\lambda_j/\cos(\varphi_Y-\varphi_X)$, so $|\mu_j|\ge|\lambda_j|\ge1$. If $Y$ is then traversed, the next step has $\lambda_{j+1}=\mu_j$.
\end{enumerate}
\end{lemma}

The components of $[0,k]\setminus\Sigma_0$ are the \emph{pieces} of $\Fz$. For the scale flows we need the following refinement.

\begin{lemma}[pieces of $\Fs$]\label{lem:P4-refine}
Let $s\in[y_0,y_1]$. There is a finite set $\Sigma_s\supseteq\Sigma_0$ with the following property on every component $I'$ of $[0,k]\setminus\Sigma_s$ (a \emph{piece} of $\Fs$). For all $x\in I'$, the $\Fs$ path of $x$ is the $\Fz$ path cut at the same step and by the same rule: $T_s$, $H_s$ inside a ray step, $H_s$ at an exit point, or the termination of the $\Fz$ path. Consequently the combinatorial type of the $\Fs$ path is constant on $I'$. The start, the length and the end point of every ray step of the $\Fs$ path, all entry and exit points, and the termination point are affine functions of $x\in I'$.
\end{lemma}

\begin{proof}
Fix a piece $I$ of $\Fz$ and use the notation of Lemma~\ref{lem:P4-recall} on $I$. We list the points of the $\Fz$ path at which $\Fs$ may cut, in the order of the path.
\begin{itemize}
\item \emph{$H_s$ inside step $j$.} The function $t^{(s)}_{H,j}(x):=(s+2-p_{j,y}(x))/(n_{X_j})_y$ is affine on $I$. Suppose the rule ending step $j$ in $\Fz$ is D, W or contact. Then one of $t_S,t_W,t_D$ equals $t^*_j$ and none is smaller, so $\min(t_S,t_W,t_D)=t^*_j$. The condition for $H_s$ therefore reads $t^{(s)}_{H,j}(x)<t^*_j(x)$, an affine strict inequality. If instead the rule ending step $j$ is H, then $t_H<\min(t_S,t_W,t_D)$ and $t^{(s)}_{H,j}\le t_H$ because $s+2\le\hmax$. So the condition for $H_s$ holds for all $x\in I$.
\item \emph{$T_s$ at the end of step $j$.} The condition is that step $j$ ends at an entry candidate on a square $Y_j$, the path is the R1 winner there, and $a(Y_j)\ge\alpha(s)$. By Lemma~\ref{lem:P4-recall}(ii) it holds for all $x\in I$ or for none.
\item \emph{$H_s$ at the exit point of $Y_j$.} The exit point $q_j+n_{Y_j}$ has height $q_{j,y}(x)+\cos\varphi_{Y_j}$, so the condition is the affine inequality $q_{j,y}(x)+\cos\varphi_{Y_j}\ge s+2$.
\end{itemize}
Let $\Sigma_s$ consist of $\Sigma_0$ together with, for every piece $I$, the zeros in $I$ of those of the finitely many affine functions $t^{(s)}_{H,j}-t^*_j$ and $q_{j,y}+\cos\varphi_{Y_j}-(s+2)$ that do not vanish identically on $I$. On a component $I'$ of $[0,k]\setminus\Sigma_s$ each of these functions vanishes identically or has constant sign, so every condition has the same truth value for all $x\in I'$. The $\Fs$ path is cut at the first point, in the order of the path, at which a condition holds. If no condition holds, it follows the $\Fz$ path to its termination. Hence the step and the rule of the cut are the same for all $x\in I'$. The termination point of the $\Fs$ path is $q_j(x)$ for $T_s$, $p_j(x)+t^{(s)}_{H,j}(x)\,n_{X_j}$ for $H_s$ inside step $j$, the exit point $q_j(x)+n_{Y_j}$ for $H_s$ at an exit point, and the termination point of the $\Fz$ path otherwise. All of these are affine on $I'$. The remaining data are those of $\Fz$.
\end{proof}

We put $\Sigma_F:=\Sigma_0$ if $F=\Fz$ and $\Sigma_F:=\Sigma_s$ if $F=\Fs$. (The set $\Sigma_s$ constructed here may differ by finitely many points from the set of Lemma~\ref{lem:flows-Ts}(f); only its finiteness and the stated properties are used, and the exceptional sets below do not depend on this choice.)

\begin{definition}[trajectory, event points]\label{def:P4-traj}
Let $x\in(0,k)$. The trajectory $\Gamma_x$ of $\pi_x$ in $F$ is the polygonal line consisting, in order, of the following segments.
\begin{itemize}
\item For each ray step of the path in $F$, the \emph{gap segment} $[p,p+t^*d]$. Here $p$ is the start of the step, $d=n_X$ with $X$ the source ($d=e_2$ for the floor), and $t^*\ge0$ is the parameter at which the step ends in $F$. For a step of $\Fs$ cut by $H_s$, $t^*=t^{(s)}_H$.
\item For each square $Y$ traversed in $F$, the \emph{pass segment} $[q,q+n_Y]$, where $q\in\relint\Bot(Y)$ is the entry point.
\end{itemize}
The last point of $\Gamma_x$ is $e(x)$. The \emph{event points} of $\Gamma_x$ are the end points of its segments, including $(x,0)$ and $e(x)$. The \emph{open part} of a segment $[a,b]$ is $[a,b]\setminus\{a,b\}$; it is empty if $a=b$.
\end{definition}

The segments alternate: gap segment, pass segment, gap segment, and so on. The last segment is a gap segment, or a pass segment if the path is completed at an exit point. This is the trajectory of Section~\ref{sec:flows}, described in the form used below.

\begin{lemma}\label{lem:P4-traj}
Let $x\in(0,k)$.
\begin{enumerate}[(i)]
\item Every square traversed by $\pi_x$ in $F$ has $a(Y)<\alpha_F$. Every segment of $\Gamma_x$ has a direction $d$ equal to $e_2$ or to $n_Y$ for such a square, so $d_y>\cos\alpha_F>0$. The height is strictly increasing along $\Gamma_x$.
\item $\Gamma_x$ has at most $2h_F/\cos\alpha_F+3$ segments. It is a compact connected polygonal line from $(x,0)$ to $e(x)$.
\item On each piece of $F$, the number and the kinds of the segments of $\Gamma_x$ are the same for all $x$, and every event point is an affine function of $x$.
\end{enumerate}
\end{lemma}

\begin{proof}
(i) In $\Fz$ a square is traversed only by an R1 winner with $a(Y)<\amax$ (rule R2). If a square is traversed in $\Fs$, then it is traversed in $\Fz$ by the R1 winner, and the entry is not a $T_s$ point; hence $a(Y)<\alpha(s)$. Since $|\varphi_Y|<\alpha_F$, we get $(n_Y)_y=\cos\varphi_Y>\cos\alpha_F$, and $(e_2)_y=1>\cos\alpha_F$. A segment of positive length therefore has strictly increasing height, and a segment of length zero is a single point.

(ii) Let $Y_1,\dots,Y_m$ be the squares traversed in $F$, with entry heights $\eta_1,\dots,\eta_m$. The exit height of $Y_i$ is $\eta_i+\cos\varphi_{Y_i}>\eta_i+\cos\alpha_F$, and by (i) $\eta_{i+1}$ is at least this exit height. Since $\eta_1\ge0$, this gives $\eta_m\ge(m-1)\cos\alpha_F$. On the other hand $\eta_m\le h_F$. In $\Fz$ a contact is processed only if $t_S=t^*\le t_H$, so its height is at most $\hmax$. In $\Fs$, suppose an entry point $q_j$ of the $\Fz$ path has height $>s+2$. In the step ending at $q_j$ we have $\min(t_S,t_W,t_D)=t_S$ and $t^{(s)}_H<t_S$, so the $\Fs$ path is cut by $H_s$ inside this step or earlier. Thus $m\le h_F/\cos\alpha_F+1$, and there are at most $2m+1$ segments.

(iii) This follows from Lemmas~\ref{lem:P4-recall}(ii) and~\ref{lem:P4-refine}. Note that an exit point is the corresponding entry point plus the constant vector $n_Y$.
\end{proof}

\begin{lemma}[position of trajectories]\label{lem:P4-position}
Let $x\in(0,k)$, and let $\Gamma_x$ be its trajectory in $F$.
\begin{enumerate}
\item[(0)] The start $p$ of every ray step satisfies $0<p_x<k$. If the source $X$ is a square, then $p_y\ge\cos\varphi_X\ge\cos\frac\pi4>0$.
\item[(i)] The open part of every gap segment lies in the waste.
\item[(ii)] The open part of the pass segment of a square $Y$ lies in $\operatorname{int}Y$.
\item[(iii)] No event point lies in the interior of a square. Every event point other than $(x,0)$ and $e(x)$ is an entry point or an exit point of a traversed square $Y$, and it lies on $\partial Y$.
\end{enumerate}
\end{lemma}

\begin{proof}
(0) For the floor, $p=(x,0)$ with $0<x<k$. Otherwise $p\in\relint\Top(X)$ and $X\subset[0,k]^2$. Suppose $p_x=0$. The linear function $z\mapsto z_x$ is $\ge0$ on $X$ and vanishes at $p$, so it attains its minimum over $X$ at $p$. Since $p$ is an interior point of the segment $\Top(X)$, the function is constant on $\Top(X)$. Then $(u_X)_x=\cos\varphi_X=0$, which is impossible because $|\varphi_X|\le\frac\pi4$. In the same way $p_x\ne k$. Moreover $p-n_X\in\Bot(X)\subset X\subset[0,k]^2$, so $p_y\ge(n_X)_y=\cos\varphi_X\ge\cos\frac\pi4$.

(i) Let $[p,p+t^*d]$ be a gap segment with $t^*>0$ and source $X$, and let $0<t<t^*$. Every ray step ends in $F$ at a parameter $t^*\le\min(t_S,t_W,t_D)$. In $\Fz$ this holds because $t^*=\min(t_S,t_W,t_D,t_H)$. In $\Fs$, a step cut by $H_s$ has $t^*=t^{(s)}_H<\min(t_S,t_W,t_D)$, and any other step is a step of $\Fz$. Consequently:
\begin{itemize}
\item $r(t)$ lies in no square $Y\ne X$, because $t<t_S$;
\item if $X$ is a square, then $r(t)\notin X$, because $(r(t)-c_X)\cdot n_X=\frac12+t>\frac12$;
\item $0<r(t)_x<k$: by (0), $0<p_x<k$; $r(t')_x\notin\{0,k\}$ for $0\le t'<t_W$; and $t<t_W$; so the claim follows by continuity;
\item $0<r(t)_y<k$: $r(t)_y>p_y\ge0$ because $d_y>0$. Also $r(t)_y$ is less than the end height of the step, which is at most $h_F<k$. In $\Fz$ this is because $t^*\le t_H$. In $\Fs$ the end height is $s+2$ if the step is cut by $H_s$. Otherwise it is at most $s+2$, because a step of $\Fz$ whose end height exceeds $s+2$ satisfies $t^{(s)}_H<t^*\le\min(t_S,t_W,t_D)$ and is cut by $H_s$.
\end{itemize}
Hence $r(t)\in(0,k)^2\setminus\bigcup_ZZ$, i.e.\ $r(t)$ lies in the waste.

(ii) Write $q=c_Y-\frac12n_Y+\xi u_Y$ with $|\xi|<\frac12$. For $0<t<1$ we have $q+tn_Y=c_Y+\xi u_Y+(t-\frac12)n_Y$ with $|\xi|<\frac12$ and $|t-\frac12|<\frac12$, so $q+tn_Y\in\operatorname{int}Y$.

(iii) First, two observations about a ray step with start $p$ and source $X$. (1) For $0<t<t_S$, the point $r(t)$ lies in no square: not in $X$ by the computation in (i), and not in any other square because $t<t_S$. (2) The contact point $r(t_S)$, if it is reached, lies in exactly one square $Y$, and it lies on $\partial Y$. Indeed, if $t_S>0$, it is a limit of the points $r(t)$, $t<t_S$, which are outside $Y$. If $t_S=0$, the source is the floor and the point has height $0$, so it is not interior to $Y$.

Since the segments alternate, every event point other than $(x,0)$ and $e(x)$ is the common end point of a gap segment and the next pass segment (an entry point) or of a pass segment and the next gap segment (an exit point). Entry points lie in $\relint\Bot(Y)$ and exit points in $\relint\Top(Y)$, so both lie on $\partial Y$. The squares are disjoint, so such points lie in no other square. The point $(x,0)$ has height $0$, while $\operatorname{int}Z\subset(0,k)^2$. It remains to consider $e(x)$, according to its label:
\begin{itemize}
\item T or M: $e(x)$ is an entry candidate point in $\relint\Bot(Y)\subset\partial Y$.
\item E: $e(x)=r(t_S)$ lies on $\partial Y$ by (2).
\item D: $e(x)=r(t_D)$ with $t_D\le t_S$. If $t_D<t_S$, it lies in no square by (1) (note $t_D=\delta-g>0$). If $t_D=t_S$, it is the contact point, on $\partial Y$ by (2).
\item W: $e(x)$ has abscissa $0$ or $k$, so it is not in $(0,k)^2$.
\item H inside a ray step: $e(x)=r(t)$ with $t<t_S$. If $t>0$, it lies in no square by (1). If $t=0$, it is the start of the step, i.e.\ $(x,0)$ or an exit point.
\item H at an exit point: $e(x)$ is an exit point. \qedhere
\end{itemize}
\end{proof}

\subsection{The exceptional set}\label{ssec:P4-exc}

By Lemma~\ref{lem:P4-traj}(iii), on a piece $I$ of $F$ the event points of $\Gamma_x$ are given by finitely many affine functions $x\mapsto e(x)$ on $I$. We call them the \emph{event-point functions} on $I$.

\begin{definition}[exceptional set]\label{def:P4-exc}
For a flow $F$ let $\Exc_F$ be the set of $y\in(0,h_F)$ such that, on some piece $I$ of $F$, some event-point function has the constant height $y$ on $I$. Put
\[
\Exc:=\Exc_{\Fz}.
\]
These are the sets $\Exc_F$ and $\Exc$ of Section~\ref{sec:flows}.
For $y\in(0,h_F)$ let $\Xi^F_y$ be the set of $x\in(0,k)\setminus\Sigma_F$ such that $\Gamma_x$ has an event point at height $y$.
\end{definition}

The set $\Exc_F$ does not depend on the choice of $\Sigma_F$. Replace $\Sigma_F$ by a finite superset: the new pieces are subintervals of the old ones, and the event-point functions are restricted. An affine function is constant on a non-empty open subinterval if and only if it is constant on the whole interval, so $\Exc_F$ does not change. Two admissible choices of $\Sigma_F$ have their union as a common finite superset.

\begin{lemma}\label{lem:P4-exc}
\begin{enumerate}[(i)]
\item $\Exc_F$ is finite.
\item If $y\in(0,h_F)\setminus\Exc_F$, then $\Xi^F_y$ is finite.
\item For every $s\in[y_0,y_1]$, $\Exc_{\Fs}\subseteq\Exc\cap(0,s+2)$. Thus the single finite set $\Exc$ serves as exceptional set for $\Fz$ and for all $\Fs$.
\item Every element of $\Exc$ is either the height of the bottom side or of the top side of a square $Z$ with $\varphi_Z=0$, or the constant height of the end points of a ray step that ends with rule D, on some piece of $\Fz$.
\end{enumerate}
\end{lemma}

\begin{proof}
(i) There are finitely many pieces and finitely many event-point functions on each. Each function with constant height contributes one value.

(ii) Let $I$ be a piece and $e$ an event-point function on $I$. Its height $e_y$ is affine and, since $y\notin\Exc_F$, not identically equal to $y$. Hence $e_y(x)=y$ for at most one $x\in I$. There are finitely many pairs $(I,e)$.

(iii) By definition $\Exc_{\Fs}\subseteq(0,s+2)$. Let $I'$ be a piece of $\Fs$; it lies in a piece $I$ of $\Fz$. By Lemma~\ref{lem:P4-refine}, every event point of the $\Fs$ trajectory on $I'$ is of one of two kinds. Either it is an event point of the $\Fz$ trajectory, i.e.\ the restriction to $I'$ of an event-point function $e$ of $\Fz$ on $I$; this covers the start, the entry and exit points, the end points of steps that are not cut, a $T_s$ point (the end point of a step of $\Fz$), an $H_s$ point at an exit point, and an inherited termination point. Or it is the end point of a step cut by $H_s$ inside the step, whose height is $s+2\notin(0,s+2)$. So let $y\in\Exc_{\Fs}$ come from an $\Fz$ event-point function $e$ with $e_y\equiv y$ on $I'$. Since $e_y$ is affine on $I$, $e_y\equiv y$ on $I$, and $y\in(0,s+2)\subseteq(0,\hmax)$. Hence $y\in\Exc$.

(iv) Let $I$ be a piece of $\Fz$, and let $e$ be an event-point function on $I$ whose height is a constant $y\in(0,\hmax)$. By Lemma~\ref{lem:P4-traj}(iii), the kind of the event point is the same for all $x\in I$. Let $p(x)$ be the start of the relevant ray step and $X$ its source, so that $p'=\lambda u_X$ with $|\lambda|\ge1$ (Lemma~\ref{lem:P4-recall}(iii)). We go through the kinds.
\begin{itemize}
\item \emph{The start $(x,0)$.} Its height is $0\notin(0,\hmax)$.
\item \emph{End point of a step ending with rule H.} Its height is $\hmax\notin(0,\hmax)$.
\item \emph{End point of a step ending with rule D.} This is the second alternative of (iv).
\item \emph{End point of a step ending at a contact.} By Lemma~\ref{lem:P4-recall}(ii) the contacted square $Y$ and the contacted side or vertex are the same for all $x\in I$. (This also follows from the affinity of $e$ on $I$. The contact points lie on the boundaries of the disjoint compact squares, so by continuity a single square $Y$ occurs. Moreover $e(I)$ is a non-degenerate segment, because $e'\cdot u_X=\lambda\ne0$ as shown below, and a non-degenerate segment contained in $\partial Y$ lies in one side of $Y$.) It is not a vertex. Indeed, the transversal coordinate $p(x)\cdot u_X$ has derivative $\lambda\ne0$, so the ray lines $p(x)+\R n_X$, $x\in I$, are pairwise distinct parallel lines, and at most one of them contains a given vertex. So the end point $e(x)=p(x)+t(x)n_X$ lies on the line of a fixed side $\sigma$ of $Y$ for all $x\in I$. Let $m$ be the outer unit normal of $\sigma$ and $v_\sigma$ a unit vector along $\sigma$. Differentiating $(e(x)-c_Y)\cdot m=\frac12$ gives $e'\cdot m=0$, so $e'=\kappa v_\sigma$ for some $\kappa$. Also $e'=\lambda u_X+t'n_X$, so $e'\cdot u_X=\lambda\ne0$ and therefore $\kappa\ne0$.
 \begin{itemize}
 \item If $\sigma$ is a lateral side, then $v_\sigma=\pm n_Y$ and the height has derivative $\pm\kappa\cos\varphi_Y\ne0$, because $\cos\varphi_Y\ge\cos\frac\pi4$. This case does not occur.
 \item If $\sigma$ is the bottom side, then $v_\sigma=\pm u_Y$ and the height has derivative $\pm\kappa\sin\varphi_Y$. It vanishes only if $\varphi_Y=0$, and then the height is the height $c_{Y,y}-\frac12$ of the bottom side. This covers the entry points, including the termination points with labels T and M. If $\sigma$ is the top side, the same argument shows that the height is constant only if $\varphi_Y=0$, and then it is the height $c_{Y,y}+\frac12$ of the top side.
 \end{itemize}
\item \emph{Exit points.} The exit point is $q(x)+n_Y$, where $q$ is the entry point; it has the same derivative as $q$. By the previous item, its height is constant only if $\varphi_Y=0$, and then it equals the height of $\Top(Y)$.
\item \emph{End point of a step ending with rule W.} The end point has abscissa $0$ or $k$; being affine on $I$, it is the same wall $\{z_x=w\}$, $w\in\{0,k\}$, for all $x\in I$. If $\varphi_X=0$, the ray is vertical with abscissa $p_x\in(0,k)$ (Lemma~\ref{lem:P4-position}(0)) and never meets a wall. So $\sin\varphi_X\ne0$. The end point is $p+t\,n_X$ with $p_x-t\sin\varphi_X=w$. Differentiating gives $\lambda\cos\varphi_X=t'\sin\varphi_X$. The height $p_y+t\cos\varphi_X$ therefore has derivative
\[
\lambda\sin\varphi_X+\lambda\frac{\cos^2\varphi_X}{\sin\varphi_X}=\frac{\lambda}{\sin\varphi_X}\ne0 ,
\]
and this case does not occur.
\end{itemize}
Every event point is of one of these kinds, which proves (iv).
\end{proof}

\begin{lemma}[vertices]\label{lem:P4-vertex}
Let $v$ be a vertex of a square $Y$ and $x\in(0,k)$. If $v\in\Gamma_x$, then $v=e(x)$, and the label of $x$ is one of the following:
\begin{itemize}
\item E: the last ray step reaches $v$ at $t_S<\min(t_W,t_D)$;
\item D: $t_D=t_S$ at $v$;
\item W: $t_W=t_S<t_D$ at $v$; then $v$ lies on a wall.
\end{itemize}
The set $\Upsilon_F$ of $x\in(0,k)$ such that $\Gamma_x$ contains a vertex of some square is finite.
\end{lemma}

\begin{proof}
A point of a pass segment of a square $Z$ has the form $c_Z+\xi u_Z+\zeta n_Z$ with $|\xi|<\frac12$. So it is not a vertex of $Z$, and it lies in no other square. In particular entry points and exit points are not vertices. The open part of a gap segment lies in the waste (Lemma~\ref{lem:P4-position}(i)). Hence $v=(x,0)$ or $v=e(x)$, and in both cases $v$ is the end point of a gap segment. Indeed, if $v=(x,0)$, then $t_S=0$ for the first ray step because $v\in Y$, so the first gap segment is $\{v\}$. If $v=e(x)$, the last segment is a gap segment, since otherwise $e(x)$ would be an exit point.

Let $r(t)=p+tn_X$ be the ray step whose gap segment ends at $v$, and let $t^*$ be its end parameter in $F$. If the source $X$ is a square, then $Y\ne X$: the point $r(0)=p\in\relint\Top(X)$ is not a vertex of $X$, and $r(t)\notin X$ for $t>0$. By the definition of $t_S$, $r(t)\notin Y$ for $t<t_S$, and $v\in Y$; hence $t^*\ge t_S$. As $t^*\le\min(t_S,t_W,t_D)$ (proof of Lemma~\ref{lem:P4-position}(i)), we get $t^*=t_S$. Rule H requires $t^*<\min(t_S,t_W,t_D)$ and cannot apply. So the step ends with D (if $t_D=t_S$), with W (if $t_W=t_S<t_D$; then $v$ has abscissa $0$ or $k$), or at a contact at $v$. A contact at $v$ is not an entry candidate because $v\notin\relint\Bot(Y)$; it is an E-termination. In each case the path terminates at $v$, so $v=e(x)$. If $v=(x,0)$, then $t_D=\delta>0=t_S$ and $t_W=+\infty$ for the vertical ray, so the label is E.

For finiteness, fix a piece $I$ of $F$ and a step $j$. For $x\in I$, the gap segment of step $j$ lies on the line $p_j(x)+\R n_{X_j}$, whose transversal coordinate $p_j(x)\cdot u_{X_j}$ is affine with derivative $\lambda_j\ne0$ (Lemma~\ref{lem:P4-recall}(iii); for $\Fs$ the same functions are used on the sub-pieces). A vertex $v$ lies on this line if and only if $p_j(x)\cdot u_{X_j}=v\cdot u_{X_j}$, which holds for at most one $x\in I$. There are finitely many pieces, steps and vertices, and $\Sigma_F$ is finite.
\end{proof}

\subsection{Statements}\label{ssec:P4-statements}

\begin{theorem}[P4$^*$]\label{thm:P4}
Let $k\ge4$ be an integer, let $\cP$ be a closed packing in $[0,k]^2$, and let the parameters be those of Section~\ref{sec:flows}. Let $F\in\{\Fz\}\cup\{\Fs:\ s\in[y_0,y_1]\}$, and let $\Exc$ be the finite set of Definition~\ref{def:P4-exc}; it is the same for all these flows.
\begin{enumerate}[(a)]
\item \emph{Crossings.} Let $y\in(0,h_F)$. For every $x$ with $\tau(x)\ge y$, the trajectory $\Gamma_x$ meets $\ell_y$ in exactly one point $P_y(x)$. If $y\notin\Exc$, there is a finite set of floor points outside of which $\tau(x)\ne y$, and $\tau(x)>y$ implies that $P_y(x)\in\mathrm{Wa}_y$ or $P_y(x)\in\operatorname{int}Z\cap\ell_y$ for a square $Z$. The set $\operatorname{int}Z\cap\ell_y$ is empty or the relative interior of the chord $Z\cap\ell_y$.
\item \emph{Bookkeeping.} For every $y\in[0,k]$,
\[
k=\sum_{Z\in\cP}c_Z(y)+\omega(y),
\]
and $|\mathrm{Wa}_y|=\omega(y)$ if $0<y<k$. For every $y\in(0,h_F)$,
\[
k=L(y)+\sum_{Z\in\cP}\mu_Z(y)+\Fgap(y)+\mathrm{Bd}(y),\qquad \mathrm{Bd}(y)\ge0,
\]
and $\mathrm{Bd}(y)=0$ if $y\notin\Exc$. Here $\mathrm{Bd}(y)$ is the measure of the set of $x$ that either terminate at height exactly $y$ or are alive at $y$ and cross $\ell_y$ at an entry point or an exit point (Remark~\ref{rem:P4-Bd}).
\item \emph{Shadow identity.} For every $y\in(0,h_F)$,
\[
\Sh(y)-\Ov(y)=L(y)+\Fgap(y)-\omega(y)+\mathrm{Bd}(y).
\]
In particular $\Sh(y)=L(y)+\Fgap(y)-\omega(y)+\Ov(y)$ if $y\notin\Exc$.
\item \emph{Square crossers.} For every $y\in(0,h_F)$ and every square $Z$, $\mu_Z(y)\le|\operatorname{int}Z\cap\ell_y|\le c_Z(y)$. Hence $\Ov\equiv0$ on $(0,h_F)$.
\item \emph{Gap crossers.} For every $y\in(0,h_F)$ the measure $\nu_y$ is absolutely continuous, its density $\rho_y$ vanishes almost everywhere outside $\mathrm{Wa}_y$, and
\[
\Fgap(y)-\omega(y)=\Ovgap(y)-U(y)\le\Ovgap(y).
\]
\item \emph{Shadow inequality.} For every $y\in(0,h_F)\setminus\Exc$,
\[
\Sh^F(y)=L^F_T(y)+\Lam^F(y)-U^F(y)\ \le\ L^F_T(y)+\Lam^F(y).
\]
If $F=\Fs$, then $L^{\Fs}_T(y)\le N(s)$ for every $y$ (this is also Corollary~\ref{cor:flows-LT}), and for $0<y<s+2$ the term $\Lam^{\Fs}(y)$ is bounded as in Proposition~\ref{prop:wall}.
\item \emph{Order of M and T.} The order ``M before T'' of rule R2 is not used in (a)--(f). Under the order ``T before M'' (Definition~\ref{def:P4-TM}), the trajectories, termination points and termination heights of all paths of $\Fz$ and of every $\Fs$ are unchanged; only the labels T and M change. Statements (a)--(f) hold verbatim with the labels of that order.
\end{enumerate}
\end{theorem}

\begin{proposition}[wall term]\label{prop:wall}
In the setting of Theorem~\ref{thm:P4}, let $s\in[y_0,y_1]$ and $0<y<s+2$. Then
\[
L^{\Fs}_D(y)\le|\cD|,\quad L^{\Fs}_M(y)\le|\cM|,\quad L^{\Fs}_E(y)\le L_E,\quad L^{\Fs}_W(y)\le2y\tan\alpha(s),\quad \Ovgap^{\Fs}(y)\le\Ovgap^{\Fz}(y),
\]
and consequently
\[
\Lam^{\Fs}(y)\ \le\ |\cD|+|\cM|+L_E+\Ovgap^{\Fz}(y)+2y\tan\alpha(s)\ \le\ \Lam_0+2(s+2)\tan\alpha(s)=\Lam^*(s).
\]
\end{proposition}

\begin{corollary}[the form used in Section~\ref{sec:P7}]\label{cor:P4-scale}
In the setting of Theorem~\ref{thm:P4}, for every $s\in[y_0,y_1]$ and every $y\in(0,s+2)\setminus\Exc$,
\[
\Sh^{\Fs}(y)\ \le\ N(s)+\Lam^*(s).
\]
\end{corollary}

\begin{remark}[what is used later]\label{rem:P4-used}
What Section~\ref{sec:P7} needs from this section is the content of Corollary~\ref{cor:P4-scale} in the following weak form: for each $s$ there is a finite set, here $\Exc\cap(0,s+2)$, outside of which the inequality holds for $y\in(0,s+2)$. Two further facts established in this section are not needed there: that the exceptional set can be chosen independent of $s$, and that the functions of $y$ involved are Borel (Remark~\ref{rem:P4-borel}). The size of $\Lam_0$ is estimated in Section~\ref{sec:P5} (Corollary~\ref{cor:Lambda0}); this section does not bound it.
\end{remark}

\begin{remark}[hypotheses used]\label{rem:P4-hyp}
Besides the definitions of the flows and Lemma~\ref{lem:P4-recall}, the proofs in this section use only \eqref{eq:P4-param}. They do not use the smallness of $\delta$ and $\amax$, nor any result of~[R]. The uniqueness of the R1 winner (N3) is used in part (d) and nowhere else in (a)--(f).

From rule R2, the proofs of (a)--(e) and of the identity in (f) use only that every square traversed in $F$ has $a(Y)<\alpha_F$ (Lemma~\ref{lem:P4-traj}(i)). These statements depend on the labels only through $L=\sum_\star L_\star$ and through $L_T+\Lam=L+\Ovgap$. So the clause of R2 by which tilt terminations occur only at bottom-side entries (side and vertex contacts give E, whatever the inclination) is not used in them; it enters only through the definitions of $N(s)$ and of the losses in (f) and in Proposition~\ref{prop:wall}.

Squares with $a(Y)=\pi/4$ need no separate treatment: the bottom side is fixed by the convention of Section~\ref{sec:flows}, and the arguments use only $|\varphi_Y|\le\frac\pi4$.
\end{remark}

\subsection{Crossings: proof of (a)}\label{ssec:P4-a}

\begin{lemma}[crossings]\label{lem:P4-crossing}
Let $y\in(0,h_F)$ and $x\in(0,k)$.
\begin{enumerate}[(i)]
\item If $\tau(x)\ge y$, then $\Gamma_x\cap\ell_y$ consists of one point $P_y(x)$. This holds for every $x$, including $x\in\Sigma_F$.
\item Suppose $\tau(x)>y$ and $P_y(x)$ is not an event point. Then exactly one of the following holds: $P_y(x)$ lies in the open part of a gap segment and $P_y(x)\in\mathrm{Wa}_y$; or $P_y(x)$ lies in the open part of the pass segment of a square $Y$ traversed in $F$, and $P_y(x)\in\operatorname{int}Y$.
\item If $\tau(x)=y$, then $P_y(x)=e(x)$; in particular $P_y(x)$ is an event point.
\item If $\tau(x)>y$ and $P_y(x)$ is an event point, then $P_y(x)$ is an entry point or an exit point of a traversed square $Y$, and $P_y(x)\in\partial Y$. In particular $P_y(x)\notin\mathrm{Wa}_y\cup\bigcup_Z\operatorname{int}Z$.
\item If $y\notin\Exc_F$ and $x\notin\Sigma_F\cup\Xi^F_y$, then $\tau(x)\ne y$; and if $\tau(x)>y$, then $P_y(x)\in\mathrm{Wa}_y\cup\bigcup_Z\operatorname{int}Z$.
\end{enumerate}
\end{lemma}

\begin{proof}
(i) By Lemma~\ref{lem:P4-traj}(ii), $\Gamma_x$ is a compact connected polygonal line from $(x,0)$ to $e(x)$, and by Lemma~\ref{lem:P4-traj}(i) the height is strictly increasing along it. The height of $e(x)$ is $\tau(x)\ge y$ if the label is not H. If the label is H, the height of $e(x)$ is $h_F$ (completion inside a ray step) or at least $h_F$ (completion at an exit point), and $h_F>y$. Since $0<y\le$ the height of $e(x)$, the intermediate value theorem gives a point of $\Gamma_x$ at height $y$. Strict monotonicity makes it unique. The argument does not use the pieces.

(ii) Since $P_y(x)$ is not an event point, it lies in the open part of a segment. By Lemma~\ref{lem:P4-position}(i), the open part of a gap segment lies in the waste, so then $P_y(x)\in\ell_y\cap\text{waste}=\mathrm{Wa}_y$. By Lemma~\ref{lem:P4-position}(ii), the open part of the pass segment of $Y$ lies in $\operatorname{int}Y$. The waste and the interiors of the squares are disjoint, so exactly one case occurs.

(iii) The label is not H, since $\tau=h_F>y$ for label H. So $e(x)$ has height $\tau(x)=y$, and $P_y(x)=e(x)$ by uniqueness.

(iv) The point $P_y(x)$ has height $y$. So it is not $(x,0)$, whose height is $0$. It is not $e(x)$ either: the height of $e(x)$ is $\tau(x)>y$ if the label is not H, and at least $h_F>y$ if it is. By Lemma~\ref{lem:P4-position}(iii), $P_y(x)$ is an entry point or an exit point of a traversed square $Y$, and it lies on $\partial Y$. A point of $\partial Y$ is not in the waste, not in $\operatorname{int}Y$, and not in any other square.

(v) If $\tau(x)=y$, or if $\tau(x)>y$ and $P_y(x)$ is an event point, then $\Gamma_x$ has an event point at height $y$ (by (iii)), so $x\in\Xi^F_y$. Otherwise (ii) applies.
\end{proof}

\begin{proof}[Proof of Theorem~\ref{thm:P4}(a)]
The first sentence is Lemma~\ref{lem:P4-crossing}(i). Let $y\in(0,h_F)\setminus\Exc$. Then $y\notin\Exc_F$ by Lemma~\ref{lem:P4-exc}(iii), and the set $\Sigma_F\cup\Xi^F_y$ is finite by Lemma~\ref{lem:P4-exc}(ii). Now Lemma~\ref{lem:P4-crossing}(v) gives the second sentence.

For the last sentence, let $Z$ be a square. The chord $Z\cap\ell_y$ is a compact segment, and $\operatorname{int}Z\cap\ell_y$ is a relatively open convex subset of it. Suppose $\operatorname{int}Z\cap\ell_y$ contains a point $z_0$, and let $z$ be a point of the relative interior of the chord. If $z\ne z_0$, choose a point $z_1$ of the chord such that $z$ lies strictly between $z_0$ and $z_1$; this is possible because $z$ is a relative interior point. The open segment between the interior point $z_0$ and the point $z_1$ of the convex set $Z$ lies in $\operatorname{int}Z$. So $z\in\operatorname{int}Z$. Hence $\operatorname{int}Z\cap\ell_y$ is either empty or the relative interior of the chord. It is empty, for example, when $\ell_y$ contains a horizontal side of $Z$.
\end{proof}

\begin{lemma}[measurability for fixed $y$]\label{lem:P4-meas}
Let $y\in(0,h_F)$. Each of the sets
\[
\mathcal L(y):=\{x:\tau(x)<y\},\quad \{x:\ \text{label }\star,\ \tau(x)<y\}\ (\star\in\{T,D,E,M,W\}),\quad A_Z(y),\quad A_{\rm gap}(y),
\]
\[
\mathcal B(y):=\{x:\tau(x)\ge y\}\setminus\Big(\bigcup_{Z\in\cP}A_Z(y)\cup A_{\rm gap}(y)\Big)
\]
is a finite union of intervals and points.
\end{lemma}

\begin{proof}
Let $I$ be a piece of $F$. On $I$ the label is constant, and $\tau$ is affine: it is the height of the termination point, an affine function by Lemma~\ref{lem:P4-traj}(iii), or the constant $h_F$ for label H. Hence $\{x\in I:\ \text{label }\star,\ \tau(x)<y\}$ is empty or an interval, and so is $\mathcal L(y)\cap I$.

Let $x\in I$ with $\tau(x)>y$. By Lemma~\ref{lem:P4-crossing}(ii) and (iv), $P_y(x)\in\operatorname{int}Z$ holds if and only if $P_y(x)$ lies in the open part of the pass segment $[q_Z(x),q_Z(x)+n_Z]$ of $Z$. Since the height is strictly increasing along this segment, this means that the paths of $I$ traverse $Z$ (a property of the type, the same for all $x\in I$) and
\[
q_{Z,y}(x)<y<q_{Z,y}(x)+\cos\varphi_Z .
\]
Similarly, $P_y(x)\in\mathrm{Wa}_y$ holds if and only if, for some gap segment of $\Gamma_x$ with start $p_j(x)$, length $t^*_j(x)$ and direction $d_j$,
\[
p_{j,y}(x)<y<p_{j,y}(x)+t^*_j(x)\,d_{j,y}.
\]
These are finitely many strict affine inequalities, together with $\tau(x)>y$. So $A_Z(y)\cap I$ and $A_{\rm gap}(y)\cap I$ are finite unions of open intervals, and $\mathcal B(y)\cap I$ is a finite union of intervals and points. Adding the finitely many points of $\Sigma_F$ proves the lemma.
\end{proof}

\subsection{The bookkeeping identities: proofs of (b) and (c)}\label{ssec:P4-bc}

\begin{proof}[Proof of Theorem~\ref{thm:P4}(b)]
\emph{First identity.} Fix $y\in(0,h_F)$. The sets $\mathcal L(y)$, $A_Z(y)$ ($Z\in\cP$), $A_{\rm gap}(y)$ and $\mathcal B(y)$ of Lemma~\ref{lem:P4-meas} partition $(0,k)$. Indeed, every $x$ has $\tau(x)<y$ or $\tau(x)\ge y$, and in the second case $P_y(x)$ is defined. The sets $A_Z(y)$ and $A_{\rm gap}(y)$ are subsets of $\{\tau>y\}$. They are pairwise disjoint, because the interiors of distinct squares are disjoint and the waste meets no square. Finally $\mathcal B(y)$ is the rest of $\{\tau\ge y\}$. The set $\mathcal L(y)$ is the disjoint union of the sets $\{\text{label }\star,\ \tau<y\}$, $\star\in\{T,D,E,M,W\}$; label H does not occur in it, since then $\tau=h_F>y$. Hence $|\mathcal L(y)|=L(y)$, and
\[
k=L(y)+\sum_Z\mu_Z(y)+\Fgap(y)+|\mathcal B(y)| .
\]
Comparing with \eqref{eq:P4-BdU} gives $\mathrm{Bd}(y)=|\mathcal B(y)|\ge0$. If $y\notin\Exc$, then $y\notin\Exc_F$ by Lemma~\ref{lem:P4-exc}(iii). By Lemma~\ref{lem:P4-crossing}(v), $\mathcal B(y)\subseteq\Sigma_F\cup\Xi^F_y$, which is finite by Lemma~\ref{lem:P4-exc}(ii). Hence $\mathrm{Bd}(y)=0$.

\emph{Second identity.} Let $y\in[0,k]$. Since $Z\subset[0,k]^2$, the closed chords $Z\cap\ell_y$ lie in $[0,k]\times\{y\}$. They are pairwise disjoint compact intervals, possibly empty or single points, because the squares are disjoint. By additivity of Lebesgue measure on $[0,k]\times\{y\}$,
\[
k=\sum_Zc_Z(y)+\Big|([0,k]\times\{y\})\setminus\bigcup_ZZ\Big|=\sum_Zc_Z(y)+\omega(y).
\]
If $0<y<k$, then $\mathrm{Wa}_y=((0,k)\times\{y\})\setminus\bigcup_ZZ$. It differs from $([0,k]\times\{y\})\setminus\bigcup_ZZ$ at most in the points $(0,y)$ and $(k,y)$, so $|\mathrm{Wa}_y|=\omega(y)$.
\end{proof}

\begin{proof}[Proof of Theorem~\ref{thm:P4}(c)]
For each $Z$, $(c_Z-\mu_Z)_+-(\mu_Z-c_Z)_+=c_Z-\mu_Z$. Summing over the finitely many squares and using (b),
\[
\Sh(y)-\Ov(y)=\sum_Zc_Z(y)-\sum_Z\mu_Z(y)=\big(k-\omega(y)\big)-\big(k-L(y)-\Fgap(y)-\mathrm{Bd}(y)\big),
\]
which is the identity. If $y\notin\Exc$, then $\mathrm{Bd}(y)=0$ by (b).
\end{proof}

\begin{remark}[the term $\mathrm{Bd}$ and degenerate configurations]\label{rem:P4-Bd}
\leavevmode
\begin{enumerate}[(1)]
\item \emph{The set $\mathcal B(y)$.} By Lemma~\ref{lem:P4-crossing}(iii) and (iv),
\[
\mathcal B(y)=\{x:\tau(x)=y\}\ \cup\ \{x:\ \tau(x)>y,\ P_y(x)\text{ is an entry point or an exit point}\}.
\]
A termination point at height $y\in(0,h_F)$ has one of the following labels. With T or M it lies on $\partial Y$. With E it lies on $\partial Y$. With D it lies in the waste if $t_D<\min(t_S,t_W)$, on $\partial Y$ if $t_D=t_S$, and on a wall if $t_D=t_W$. With W it lies on a wall. Label H is impossible, since then $\tau=h_F>y$. So $\mathrm{Bd}(y)$ counts the paths that terminate at height $y$ (in the waste, on a wall, or on the boundary of a square), and the alive paths that cross $\ell_y$ at an entry or exit point. By (b), $\mathrm{Bd}(y)=0$ for $y\notin\Exc$. At heights of $\Exc$ it can be positive (Proposition~\ref{prop:P4-example}).
\item \emph{Boundaries of squares.} A trajectory meets the boundary of a square only at event points (Lemma~\ref{lem:P4-position}). An alive path crosses $\ell_y$ at a boundary point only at an entry or exit point (Lemma~\ref{lem:P4-crossing}(iv)). The chord $c_Z(y)$ is the length of the closed chord. For example, at the height of the bottom side of a square with $\varphi_Z=0$ we have $c_Z=1$, while $\operatorname{int}Z\cap\ell_y=\emptyset$ and hence $\mu_Z=0$.
\item \emph{Vertices.} A vertex of a square lies on a trajectory only as its termination point, with label E, D or W, and only for the finitely many $x\in\Upsilon_F$ (Lemma~\ref{lem:P4-vertex}). The case W occurs when the ray reaches a wall exactly at a vertex of a square touching the wall ($t_W=t_S<t_D$); the priority of W over contact decides the label.
\item \emph{Squares on the floor.} Let $Y$ be a square with $\varphi_Y=0$ whose bottom side lies on the floor, and let $x\in(0,k)$ with $(x,0)\in\relint\Bot(Y)$. The first ray step of $x$ has $t_S=0<t_D=\delta$, $t_W=+\infty$ and $t_H=h_F>0$, so $(Y,(x,0))$ is an entry candidate with cumulative gap $0$. Every point reached by a ray step with a square source has positive height (Lemma~\ref{lem:P4-position}(0)). A ray step with the floor as source and start $(x',0)$ runs along the vertical half-line above $x'$. Hence $\mathcal A_Y((x,0))=\{x\}$, and $x$ is the R1 winner. Since $a(Y)=0<\alpha_F$, the path traverses $Y$ (in $\Fz$ and in every $\Fs$). Its first gap segment is the single point $(x,0)$. For $0<y<1$ (note $h_F\ge1$) the path is alive at $y$, and $P_y(x)=(x,y)\in\operatorname{int}Y$, so $x\in A_Y(y)$. If the end points of $\Bot(Y)$ have abscissae in $(0,k)$, these two floor points are vertex contacts, with label E and $\tau=0$. A square with $a(Y)>0$ meets the floor at most in its lowest vertex $(x_v,0)$; if $x_v\in(0,k)$, then $x_v$ has label E and $\tau=0$. These finitely many floor points are counted in $L(y)$ for all $y>0$ and have measure zero.
\item \emph{Walls.} Floor points lie in $(0,k)$, so no path starts on a wall. The open parts of gap segments have abscissae in $(0,k)$ (Lemma~\ref{lem:P4-position}(i)), so no trajectory runs along a wall. Points with label W lie on a wall and therefore not in the waste.
\end{enumerate}
\end{remark}

\begin{proposition}[exceptional heights are needed]\label{prop:P4-example}
Let $k\ge4$ be an integer, let the parameters be as in Section~\ref{sec:flows}, and let $0<b<\min(\delta,1)$. Let $\cP$ consist of the three squares
\[
Z_i:=\Big[\,i+\tfrac{i+1}{1000},\ i+1+\tfrac{i+1}{1000}\,\Big]\times[\,b,\ b+1\,],\qquad i=0,1,2 .
\]
Then $b\in\Exc$, and for $F=\Fz$ and for every $F=\Fs$, $s\in[y_0,y_1]$,
\[
\Sh^F(b)=3,\qquad L^F_T(b)+\Lam^F(b)=0,\qquad \mathrm{Bd}^F(b)=3 .
\]
So the inequality in Theorem~\ref{thm:P4}(f) fails at $y=b$. On the other hand $\Sh^F(y)=0$ for $b<y<\min(b+1,h_F)$. The example concerns $\mu_Z$ as defined in (N8), with crossing points in $\operatorname{int}Z$. If crossing points in the closed square $Z$ were counted instead, one would get $\mu_{Z_i}(b)=1$ and $\Sh^F(b)=0$.
\end{proposition}

\begin{proof}
Consecutive squares are at distance $\frac1{1000}$, and all three lie in $[0,4]^2\subseteq[0,k]^2$, so $\cP$ is a closed packing. Note $h_F\ge1>b$, since $h_{\Fz}=\frac k2-1\ge1$ and $h_{\Fs}=s+2>2$. Let $J\subset(0,k)$ be the union of the three open intervals $\big(i+\frac{i+1}{1000},\,i+1+\frac{i+1}{1000}\big)$, so $|J|=3$. No square meets the strip $\R\times[0,b)$, so every path starts with a vertical ray step.
\begin{itemize}
\item Let $x\in J$, with $x$ in the $i$-th interval. The first ray step meets $Z_i$ at $t_S=b$, with $t_D=\delta>b$, $t_W=+\infty$ and $t_H=h_F>b$. So $(Z_i,(x,b))$ is an entry candidate. Only the vertical ray of $x$ reaches $(x,b)$, so $x$ is the R1 winner. Since $a(Z_i)=0<\alpha_F$, the path traverses $Z_i$, and $\tau(x)>b$. Its crossing with $\ell_b$ is the entry point $(x,b)\in\partial Z_i$, so $x\in\mathcal B(b)$.
\item The six end points of the three intervals are vertex contacts at height $b$, with label E and $\tau=b$.
\item Let $x\in(0,k)$ lie outside the closure of $J$. The vertical line through $x$ meets no square, so $t_S=t_W=+\infty$. The path terminates with D at height $\delta$ or is completed at height $h_F$; both heights exceed $b$. So $P_b(x)=(x,b)\in\mathrm{Wa}_b$. These $x$ form a set of measure $k-3$.
\end{itemize}
Hence at $y=b$: $L(b)=0$; $\mu_{Z_i}(b)=0$ because $\operatorname{int}Z_i\cap\ell_b=\emptyset$; $\Fgap(b)=k-3$ and $\mathrm{Bd}(b)=3$. Also $c_{Z_i}(b)=1$ and $\omega(b)=k-3$. The gap crossers have $X_b(x)=x$, so $\rho_b=1$ almost everywhere on $\mathrm{Wa}_b$, and $\Ovgap(b)=U(b)=0$. Therefore $\Sh(b)=3$, while $L_T(b)=0$ and $\Lam(b)=0$.

On every piece of $F$ contained in $J$ (such pieces exist because $\Sigma_F$ is finite), the entry point $(x,b)$ has the constant height $b\in(0,h_F)$. So $b\in\Exc_F$; for $F=\Fz$ this gives $b\in\Exc$.

For $b<y<\min(b+1,h_F)$, every $x\in J$ is alive at $y$ and $P_y(x)=(x,y)\in\operatorname{int}Z_i$. So $\mu_{Z_i}(y)\ge1=c_{Z_i}(y)$, and $\Sh(y)=0$.

If crossing points in the closed square were counted, then at $y=b$ the floor points of $J$ would be counted, while the six vertex points are not alive at $b$. This would give $\mu_{Z_i}(b)=1=c_{Z_i}(b)$ and $\Sh(b)=0$.
\end{proof}

\subsection{Square crossers: proof of (d)}\label{ssec:P4-d}

\begin{proof}[Proof of Theorem~\ref{thm:P4}(d)]
Fix $y\in(0,h_F)$ and a square $Z$.

\emph{Step 1: the paths of $A_Z(y)$ traverse $Z$.} Let $x\in A_Z(y)$. The point $P_y(x)\in\operatorname{int}Z$ is not an event point (Lemma~\ref{lem:P4-position}(iii)), and it is not in the waste. By Lemma~\ref{lem:P4-crossing}(ii), it therefore lies in the open part of the pass segment of a square $Y$ traversed by $\pi_x$ in $F$, and $P_y(x)\in\operatorname{int}Y$. Since $\operatorname{int}Y\cap\operatorname{int}Z=\emptyset$ for $Y\ne Z$, we get $Y=Z$. So $\pi_x$ traverses $Z$ in $F$. A traversal in $\Fs$ is a traversal in $\Fz$, and in $\Fz$ only the R1 winner of an entry candidate traverses. Hence $x$ is the R1 winner at an entry candidate $(Z,q_Z(x))$ of $\Fz$, where $q_Z(x)\in\relint\Bot(Z)$. This entry point is unique, because a path has at most one entry candidate on $Z$ (Lemma~\ref{lem:P4-recall}(i)). Finally $P_y(x)=q_Z(x)+t\,n_Z$ with $0<t<1$.

\emph{Step 2: injectivity.} Let $x,x'\in A_Z(y)$ with $P_y(x)=P_y(x')$. Write $q_Z(x)=c_Z-\frac12n_Z+\xi u_Z$, $q_Z(x')=c_Z-\frac12n_Z+\xi'u_Z$, and $P_y(x)=q_Z(x)+tn_Z=q_Z(x')+t'n_Z$. Taking scalar products with $u_Z$ and with $n_Z$ gives $\xi=\xi'$ and $t=t'$. So $q_Z(x)=q_Z(x')=:q$, and both $x$ and $x'$ are R1 winners at $(Z,q)$. An entry candidate has at most one R1 winner (N3), so $x=x'$. Thus $X_y$ is injective on $A_Z(y)$.

\emph{Step 3: expansion.} By the proof of Lemma~\ref{lem:P4-meas}, $A_Z(y)\setminus\Sigma_F$ is a finite union of pairwise disjoint open intervals $J$. Each $J$ is contained in a piece of $F$ on which the paths traverse $Z$ at the same step $j$. On $J$ the entry point $q_Z(x)$ is the end point $q_j(x)$ of step $j$, and Lemma~\ref{lem:P4-recall}(iii) gives $q_Z'=\mu u_Z$ with $|\mu|\ge1$. For $\Fs$, $J$ lies in a piece of $\Fz$ and the same functions are used. Since $P_y(x)=q_Z(x)+t(x)n_Z$ has height $y$, we have $t(x)=(y-q_{Z,y}(x))/\cos\varphi_Z$ and
\[
X_y(x)=q_{Z,x}(x)-\big(y-q_{Z,y}(x)\big)\tan\varphi_Z .
\]
Hence
\[
X_y'=\mu\cos\varphi_Z+\mu\sin\varphi_Z\tan\varphi_Z=\frac{\mu}{\cos\varphi_Z},\qquad |X_y'|\ge|\mu|\ge1 .
\]
So $X_y$ is affine on $J$ with slope of absolute value at least $1$, and $|X_y(J)|=|X_y'|\,|J|\ge|J|$.

\emph{Step 4: summation.} By Step 2 the intervals $X_y(J)$ are pairwise disjoint, and by Step 1 they are contained in $\operatorname{int}Z\cap\ell_y$, viewed as a subset of $\R$. The set $\Sigma_F$ is finite. Therefore
\[
\mu_Z(y)=\sum_J|J|\le\sum_J|X_y(J)|\le|\operatorname{int}Z\cap\ell_y|\le|Z\cap\ell_y|=c_Z(y).
\]
Hence $(\mu_Z(y)-c_Z(y))_+=0$ for every $Z$, and $\Ov(y)=0$.
\end{proof}

Step 2 is the only place in (a)--(f) where the uniqueness of the R1 winner is used.

\subsection{Gap crossers: proof of (e)}\label{ssec:P4-e}

\begin{proof}[Proof of Theorem~\ref{thm:P4}(e)]
Fix $y\in(0,h_F)$.

\emph{Step 1: decomposition.} For a piece $I$ of $F$ and a ray step $j$ of the paths of $I$, let
\[
J_{I,j}:=\{x\in I:\ \tau(x)>y,\ p_{j,y}(x)<y<p_{j,y}(x)+t^*_j(x)\,d_{j,y}\}.
\]
This is an open interval, possibly empty. By the proof of Lemma~\ref{lem:P4-meas}, $A_{\rm gap}(y)\setminus\Sigma_F$ is the union of the sets $J_{I,j}$. These sets are pairwise disjoint. Different pieces are disjoint. For a fixed $x$, the point $P_y(x)$ lies in the open part of at most one segment, because the height is strictly increasing along $\Gamma_x$ and the open parts of different segments therefore have disjoint ranges of heights. Let $c$ run over the pairs $(I,j)$ with $J_c:=J_{I,j}\ne\emptyset$.

\emph{Step 2: slopes.} Let $c=(I,j)$, let $X$ be the source of step $j$, and recall $p_j'=\lambda u_X$ with $|\lambda|\ge1$. For $x\in J_c$ we have $P_y(x)=p_j(x)+t\,n_X$ with $t=(y-p_{j,y}(x))/\cos\varphi_X$. So on $J_c$
\[
X_c(x):=X_y(x)=p_{j,x}(x)-\big(y-p_{j,y}(x)\big)\tan\varphi_X ,\qquad X_c'=\lambda\cos\varphi_X+\lambda\sin\varphi_X\tan\varphi_X=\frac{\lambda}{\cos\varphi_X},
\]
and $|X_c'|\ge1$.

\emph{Step 3: density.} Under an affine map of slope $a\ne0$, Lebesgue measure on an interval $J$ is mapped to $|a|^{-1}$ times Lebesgue measure on the image interval. Since $\Sigma_F$ is finite,
\[
\nu_y=\sum_c\,(X_c)_\#\big(|\cdot|\,\big|_{J_c}\big)=\Big(\sum_c|X_c'|^{-1}\mathbf 1_{X_c(J_c)}\Big)\dd q .
\]
So $\nu_y$ is absolutely continuous, with density $\rho_y=\sum_c|X_c'|^{-1}\mathbf 1_{X_c(J_c)}$ almost everywhere. By the definition of $A_{\rm gap}(y)$, $X_c(J_c)\subseteq\mathrm{Wa}_y$. Hence $\rho_y=0$ almost everywhere outside $\mathrm{Wa}_y$. Further, $\rho_y$ is bounded, and
\[
\Fgap(y)=|A_{\rm gap}(y)|=\nu_y(\R)=\int_{\mathrm{Wa}_y}\rho_y\dd q .
\]

\emph{Step 4: conclusion.} By (b), $\omega(y)=|\mathrm{Wa}_y|=\int_{\mathrm{Wa}_y}1\dd q$. All integrals are finite, because $\rho_y$ is bounded and $|\mathrm{Wa}_y|\le k$. Writing $\rho_y-1=(\rho_y-1)_+-(1-\rho_y)_+$,
\[
\Fgap(y)-\omega(y)=\int_{\mathrm{Wa}_y}(\rho_y-1)\dd q=\Ovgap(y)-U(y)\le\Ovgap(y).
\]
Another version of $\rho_y$, equal to this one almost everywhere, gives the same $\Ovgap(y)$ and $U(y)$.
\end{proof}

Part (e) uses no bound on the density $\rho_y$; the size of $\Ovgap$ is estimated in Section~\ref{sec:P5}.

\subsection{The shadow inequality and the wall term}\label{ssec:P4-f}

\begin{proof}[Proof of Theorem~\ref{thm:P4}(f)]
Let $y\in(0,h_F)\setminus\Exc$. By (b), $\mathrm{Bd}(y)=0$, and by (d), $\Ov(y)=0$. So (c) gives $\Sh(y)=L(y)+\Fgap(y)-\omega(y)$, and (e) turns this into $\Sh(y)=L(y)+\Ovgap(y)-U(y)$. Since $L=L_T+L_D+L_E+L_M+L_W$, this is $\Sh(y)=L_T(y)+\Lam(y)-U(y)$, and $U(y)\ge0$.

Let $F=\Fs$. By (N4), a path of $\Fs$ terminates with label T only at a $T_s$ point. This is either a cut of type $T_s$, or an inherited $\Fz$ termination with label T, which is itself a $T_s$ point. In both cases $x\in\cT_s$. Hence $\{x:\ \text{label T in }\Fs,\ \tau_{\Fs}(x)<y\}\subseteq\cT_s$ and $L^{\Fs}_T(y)\le|\cT_s|=N(s)$ for every $y$. The bound for $\Lam^{\Fs}$ is Proposition~\ref{prop:wall}, whose proof follows.
\end{proof}

\begin{lemma}[$\Fs$ and $\Fz$]\label{lem:P4-FsF0}
Let $s\in[y_0,y_1]$ and $x\in(0,k)$.
\begin{enumerate}[(i)]
\item $\Gamma^{\Fs}_x$ is an initial part of $\Gamma^{\Fz}_x$. Its pass segments are pass segments of $\Gamma^{\Fz}_x$, and its gap segments are those of $\Gamma^{\Fz}_x$, except that the last one may be an initial part of the corresponding gap segment of $\Gamma^{\Fz}_x$.
\item $\tau_{\Fs}(x)\le s+2$ and $\tau_{\Fs}(x)\le\tau_{\Fz}(x)$.
\item If $x$ has label D, M, E or W in $\Fs$, then it has the same label and the same termination point in $\Fz$.
\item Every square traversed in $\Fs$ has $a(Y)<\alpha(s)$.
\end{enumerate}
\end{lemma}

\begin{proof}
(i) and (iii) follow from (N4). The $\Fs$ path is the $\Fz$ path cut at an entry point ($T_s$), at a point inside a ray step or at an exit point ($H_s$), or at the termination of the $\Fz$ path. The cuts carry the labels T and H, so the labels D, M, E and W of $\Fs$ come from terminations of $\Fz$, at the same point. Part (iv) is Lemma~\ref{lem:P4-traj}(i).

(ii) Every ray step of $\Fs$ ends at height at most $s+2$ (proof of Lemma~\ref{lem:P4-position}(i)), and every entry point of $\Fs$ has height at most $s+2$ (proof of Lemma~\ref{lem:P4-traj}(ii)). A termination with a label other than H is the end point of a ray step or an entry point, so it has height $\le s+2$. With label H, $\tau_{\Fs}=s+2$ by definition. This proves $\tau_{\Fs}(x)\le s+2$.

If $x$ has label H in $\Fz$, then $\tau_{\Fz}(x)=\hmax\ge s+2\ge\tau_{\Fs}(x)$. Otherwise $\tau_{\Fz}(x)$ is the height of the last point of $\Gamma^{\Fz}_x$. By (i) and Lemma~\ref{lem:P4-traj}(i), this is at least the height of the last point $e^{\Fs}(x)$ of $\Gamma^{\Fs}_x$. That height equals $\tau_{\Fs}(x)$ if the label in $\Fs$ is not H, and it is at least $s+2=\tau_{\Fs}(x)$ if the label is H, the cut being at height $s+2$ inside a ray step or at an exit point of height $\ge s+2$.
\end{proof}

\begin{proof}[Proof of Proposition~\ref{prop:wall}]
Let $s\in[y_0,y_1]$ and $0<y<s+2$. Note that $y<\hmax$, so the quantities of $\Fz$ at $y$ are defined.

\emph{Step 1: deaths, merges, end collisions.} By Lemma~\ref{lem:P4-FsF0}(iii), $\{x:\ \text{label D in }\Fs,\ \tau_{\Fs}(x)<y\}\subseteq\cD$, so $L^{\Fs}_D(y)\le|\cD|$. In the same way $L^{\Fs}_M(y)\le|\cM|$ and $L^{\Fs}_E(y)\le L_E$.

\emph{Step 2: gap overlap.} Let $x\in A^{\Fs}_{\rm gap}(y)$, so $\tau_{\Fs}(x)>y$ and $P^{\Fs}_y(x)\in\mathrm{Wa}_y$. By Lemma~\ref{lem:P4-FsF0}(ii), $\tau_{\Fz}(x)\ge\tau_{\Fs}(x)>y$. By Lemma~\ref{lem:P4-FsF0}(i), $\Gamma^{\Fs}_x\subseteq\Gamma^{\Fz}_x$, so the unique point of $\Gamma^{\Fz}_x$ on $\ell_y$ is $P^{\Fz}_y(x)=P^{\Fs}_y(x)$. Hence $x\in A^{\Fz}_{\rm gap}(y)$, with the same crossing abscissa. So $\nu^{\Fs}_y(B)\le\nu^{\Fz}_y(B)$ for every Borel set $B\subseteq\R$. Both densities are bounded (proof of (e)). Taking $B=\{\rho^{\Fs}_y>\rho^{\Fz}_y\}$ gives $\int_B(\rho^{\Fs}_y-\rho^{\Fz}_y)\le0$ with a positive integrand, so $|B|=0$, i.e.\ $\rho^{\Fs}_y\le\rho^{\Fz}_y$ almost everywhere. Therefore
\[
\Ovgap^{\Fs}(y)=\int_{\mathrm{Wa}_y}(\rho^{\Fs}_y-1)_+\le\int_{\mathrm{Wa}_y}(\rho^{\Fz}_y-1)_+=\Ovgap^{\Fz}(y)\le\sup_{0<y'<\hmax}\Ovgap^{\Fz}(y').
\]

\emph{Step 3: walls.} Let $x$ have label W in $\Fs$ with $\tau(x):=\tau_{\Fs}(x)<y$. By Lemma~\ref{lem:P4-FsF0}(iv), every segment of $\Gamma^{\Fs}_x$ has direction $e_2$ or $n_Y$ with $|\varphi_Y|<\alpha(s)$, so its direction $d$ satisfies $|d_x|\le d_y\tan\alpha(s)$. Summing over the segments of $\Gamma^{\Fs}_x$, which goes from $(x,0)$ to $e(x)$ and gains the height $\tau(x)$, we get $|e(x)_x-x|\le\tau(x)\tan\alpha(s)<y\tan\alpha(s)$. Since $e(x)_x\in\{0,k\}$, this gives $x<y\tan\alpha(s)$ or $k-x<y\tan\alpha(s)$. Hence $L^{\Fs}_W(y)\le2y\tan\alpha(s)$.

\emph{Step 4: summation.} Adding Steps 1--3,
\[
\Lam^{\Fs}(y)=\big[L_D+L_E+L_M+L_W\big]^{\Fs}(y)+\Ovgap^{\Fs}(y)\le|\cD|+|\cM|+L_E+\Ovgap^{\Fz}(y)+2y\tan\alpha(s),
\]
and the right-hand side is at most $\Lam_0+2y\tan\alpha(s)\le\Lam_0+2(s+2)\tan\alpha(s)$, because $0<y<s+2$ and $\tan\alpha(s)>0$.
\end{proof}

\begin{proof}[Proof of Corollary~\ref{cor:P4-scale}]
Let $s\in[y_0,y_1]$ and $y\in(0,s+2)\setminus\Exc$. Since $h_{\Fs}=s+2$, Theorem~\ref{thm:P4}(f) applies to $F=\Fs$ and gives $\Sh^{\Fs}(y)\le L^{\Fs}_T(y)+\Lam^{\Fs}(y)$. By (f), $L^{\Fs}_T(y)\le N(s)$, and by Proposition~\ref{prop:wall}, $\Lam^{\Fs}(y)\le\Lam^*(s)$.
\end{proof}

\subsection{The order of M and T: proof of (g)}\label{ssec:P4-g}

\begin{definition}[order ``T before M'']\label{def:P4-TM}
The master flow $F^{0,\mathrm{TM}}$ is defined like $\Fz$ in (N2)--(N3), with one change at entry candidates $(Y,q)$ with $a(Y)\ge\amax$: there every path that reaches $(Y,q)$ alive terminates at $q$ with label T, and R1 is not applied. At entry candidates with $a(Y)<\amax$, R1 and the traversal rule are applied as before, by the same induction over the squares in increasing order of centre height. For $s\in[y_0,y_1]$, the scale flow $F^{\mathrm{TM}}_s$ is the path of $F^{0,\mathrm{TM}}$ cut at the first of the following points: an entry candidate reached alive with $a(Y)\ge\alpha(s)$ (label T, whether or not the path would win R1 there); an $H_s$ point as in (N4); or the termination of the path of $F^{0,\mathrm{TM}}$, with its label.
\end{definition}

The induction is well defined for $F^{0,\mathrm{TM}}$ because Lemma~\ref{lem:P4-recall}(i) is a statement about the geometry of a single path that traverses only squares with $a(Y)<\amax$, and such paths are the same in both orders.

\begin{proof}[Proof of Theorem~\ref{thm:P4}(g)]
\emph{Step 1: the master flows.} We show by induction over the squares, in increasing order of centre height (ties in the fixed order used for R1), the following claim. For every square $Y$ and every $q\in\relint\Bot(Y)$, the set $\mathcal A_Y(q)$ of floor points that reach $(Y,q)$ alive is the same in $\Fz$ and in $F^{0,\mathrm{TM}}$. Moreover, each of these paths either terminates at $q$ in both flows or traverses $Y$ in the same state in both flows.

Assume the claim for all squares before $Y$. Apart from R1 and R2, every rule (D, W, the classification of contacts, E, H, traversal) depends only on the current state $(p,d,X,g)$ of the path. Let $x$ reach $(Y,q)$ alive in one of the two flows. Its path up to $q$ consists of ray steps, which are determined by the state, and of traversals of squares $Y'$ at entry candidates $(Y',q')$. By Lemma~\ref{lem:P4-recall}(i) these squares $Y'$ have smaller centre height than $Y$, so by the induction hypothesis $x$ traverses each of them in the same state in the other flow as well. Hence $x$ reaches $(Y,q)$ alive in the other flow, and $\mathcal A_Y(q)$ is the same in both flows. If $a(Y)\ge\amax$, every path of $\mathcal A_Y(q)$ terminates at $q$ in both flows: in $\Fz$ the R1 winner with label T and the others with label M, in $F^{0,\mathrm{TM}}$ all with label T. If $a(Y)<\amax$, both flows apply R1 to the same set with the same cumulative gaps $g_x(q)$ (the paths agree up to $q$), so they have the same winner. The winner traverses $Y$ in the same state, and the other paths terminate at $q$ with label M. This proves the claim.

Consequently, every floor point has the same trajectory and the same termination point in $\Fz$ and in $F^{0,\mathrm{TM}}$. The labels agree, except that at entry candidates with $a(Y)\ge\amax$ some paths with label M in $\Fz$ have label T in $F^{0,\mathrm{TM}}$.

\emph{Step 2: the scale flows.} Fix $s$ and $x$. By Step 1 the master path of $x$ is the same in both orders: the same trajectory, the same entry candidates, reached alive or not in the same way. Let $C$ be the first entry candidate on it that is reached alive and has $a(Y_C)\ge\alpha(s)$, if there is one. The $H_s$ points and the master termination point are the same in both orders, since they depend only on the trajectory. Before $C$ there is no $T_s$ point in $\Fs$, because a $T_s$ point is such an entry candidate. Hence, in both orders, the path is cut or terminates at the first of $C$, the first $H_s$ point and the master termination. At $C$, $F^{\mathrm{TM}}_s$ gives label T. In $\Fs$, the path gives label T if it is the R1 winner at $C$ ($T_s$). Otherwise $C$ is an entry candidate where the path loses R1, so the $\Fz$ path terminates at $C$ with label M, and $\Fs$ inherits this. Elsewhere the labels agree. So every $x$ has the same trajectory, termination point and termination height in $\Fs$ and in $F^{\mathrm{TM}}_s$, and only labels M of $\Fs$ can become labels T.

\emph{Step 3: consequences.} The pieces $\Sigma_0$, $\Sigma_s$ serve for both orders. Trajectories coincide, and on a piece the labels of the order T before M are constant, because they are determined by the type. The following quantities are defined from trajectories, termination points, termination heights and aliveness only, so they coincide for the two orders: $P_y$, $A_Z$, $\mu_Z$, $A_{\rm gap}$, $\Fgap$, $c_Z$, $\omega$, $\nu_y$, $\rho_y$, $\Sh$, $\Ov$, $\Ovgap$, $U$, $\mathrm{Bd}$ and $\Exc$. So does $L=\sum_\star L_\star$, because every $x$ with $\tau(x)<y$ is counted once whatever its label among T, D, E, M, W. The labels D, E and W are the same in both orders. Only $L_T$ and $L_M$ change, and $L^{\mathrm{TM}}_T-L_T=L_M-L^{\mathrm{TM}}_M\ge0$.

Hence (a)--(e) hold for the order T before M: every quantity occurring in them is the same for the two orders. So does the identity in (f), because $L_T+\Lam=L+\Ovgap$ is the same for the two orders. The bound $L^{\mathrm{TM}}_T(y)\le N^{\mathrm{TM}}(s):=|\{x:\ \text{label T in }F^{\mathrm{TM}}_s\}|$ holds trivially. In Proposition~\ref{prop:wall}, the cuts of $F^{\mathrm{TM}}_s$ carry the labels T and H, so the labels M of $F^{\mathrm{TM}}_s$ are inherited from $F^{0,\mathrm{TM}}$. Thus $L^{F^{\mathrm{TM}}_s}_M(y)\le|\cM^{\mathrm{TM}}|$, where $\cM^{\mathrm{TM}}\subseteq\cM$ is the set of points with label M in $F^{0,\mathrm{TM}}$. Steps 1--3 of that proof are otherwise unchanged, and $\Lam_0^{\mathrm{TM}}\le\Lam_0$.
\end{proof}

\begin{remark}[where the order is used]\label{rem:P4-order-used}
The order M before T is used later, in Section~\ref{sec:P8}, through the inclusion $\cT_s\subseteq\{T^*\le s\}$. The function $T^*$ is defined from the entry candidates at which the path is the R1 winner (Section~\ref{sec:flows}). So the inclusion needs that the terminal entry of every $x\in\cT_s$ be an entry at which $x$ is the R1 winner. Under the order M before T this holds by (N4). Under the order T before M, the set of points with label T in $F^{\mathrm{TM}}_s$ may also contain R1 losers. The order M before T is sufficient for the inclusion. We make no claim about whether the later arguments could be carried out with the order T before M.
\end{remark}

\subsection{Ceiling flows and measurability in \texorpdfstring{$y$}{y}}\label{ssec:P4-ceiling}

\begin{remark}[ceiling flows]\label{rem:P4-ceiling}
Theorem~\ref{thm:P4}, Proposition~\ref{prop:wall} and Corollary~\ref{cor:P4-scale} hold for every closed packing. Let $\varrho(x,y):=(x,k-y)$ be the reflection of Lemma~\ref{lem:flows-refl} (denoted $\rho$ there and in the other sections); in this remark we write $\varrho$ for it, to keep it apart from the densities $\rho_y$. The reflected packing $\varrho(\cP)$ is a closed packing in $[0,k]^2$, and the ceiling flows of Section~\ref{sec:flows} are $\tilde F=\varrho\circ F[\varrho(\cP)]\circ\varrho$. Translated back to the original coordinates, a ceiling path moves downwards. Its termination height $\tilde\tau(x)$ is the original height of its termination point, with $\tilde\tau(x):=k-h_F$ for completed paths. It is alive at $y$ if $\tilde\tau(x)<y$. The ceiling losses are $\tilde L_\star(y):=|\{x:\ \text{label }\star,\ \tilde\tau(x)>y\}|$, and $c_Z(y)=c_{\varrho Z}(k-y)$, $\omega(y)=\omega_{\varrho(\cP)}(k-y)$. With these conventions every quantity of a ceiling flow at height $y$ equals the corresponding quantity of the floor flow of $\varrho(\cP)$ at height $k-y$. Hence, with $\widetilde{\Exc}:=k-\Exc[\varrho(\cP)]$:
\begin{itemize}
\item $\widetilde{\Sh}(y)=\tilde L_T(y)+\tilde\Lam(y)-\tilde U(y)\le\tilde L_T(y)+\tilde\Lam(y)$ for $y\in(k-h_F,k)\setminus\widetilde{\Exc}$;
\item $\tilde\Lam^{\tilde F_s}(y)\le\tilde\Lam_0+2(k-y)\tan\alpha(s)\le\tilde\Lam_0+2(s+2)\tan\alpha(s)$ for $k-s-2<y<k$.
\end{itemize}
Here $\tilde\Lam_0$ is $\Lam_0$ computed for $\varrho(\cP)$.
\end{remark}

\begin{remark}[Borel measurability in $y$]\label{rem:P4-borel}
For a fixed flow $F$, the functions
\[
y\mapsto L_\star(y),\ \mu_Z(y),\ \Fgap(y),\ c_Z(y),\ \omega(y),\ \Sh(y),\ \Ov(y),\ \Ovgap(y),\ U(y),\ \mathrm{Bd}(y)
\]
are Borel on $(0,h_F)$. This is not used in the later sections (Remark~\ref{rem:P4-used}).

\emph{Proof.} (1) Let $A_\bullet(y)$ be one of the sets $\{\text{label }\star,\ \tau<y\}$, $A_Z(y)$, $A_{\rm gap}(y)$, and put $\mathbf A_\bullet:=\{(x,y):\ 0<y<h_F,\ x\in A_\bullet(y)\}$. For $x$ in a piece $I$, the conditions that define $x\in A_\bullet(y)$ in the proof of Lemma~\ref{lem:P4-meas} are finitely many strict inequalities between $y$ and affine functions of $x$. So $\mathbf A_\bullet\cap(I\times(0,h_F))$ is a finite union of open sets. For each of the finitely many $x\in\Sigma_F$, the set $\{y:\ x\in A_\bullet(y)\}$ is a finite union of intervals, determined by the height ranges of the finitely many segments of $\Gamma_x$. Hence $\mathbf A_\bullet$ is a Borel set, and by Tonelli's theorem $y\mapsto|A_\bullet(y)|=\int\mathbf 1_{\mathbf A_\bullet}(x,y)\dd x$ is Borel. This covers $L_\star$, $\mu_Z$ and $\Fgap$; then $\mathrm{Bd}=k-L-\sum_Z\mu_Z-\Fgap$.

(2) $c_Z(y)=\int\mathbf 1_Z(q,y)\dd q$ is Borel by Tonelli's theorem, and $\omega=k-\sum_Zc_Z$ by (b). $\Sh$ and $\Ov$ are finite sums of positive parts of Borel functions.

(3) For a piece $I$ and a ray step $j$ let $\mathbf J_{I,j}:=\{(x,y):\ x\in I,\ 0<y<h_F,\ \tau(x)>y,\ p_{j,y}(x)<y<p_{j,y}(x)+t^*_j(x)d_{j,y}\}$, an open set. The map $\Psi_{I,j}(x,y):=\big(p_{j,x}(x)-(y-p_{j,y}(x))\tan\varphi_{X_j},\,y\big)$ is affine in $(x,y)$. The derivative of its first component in $x$ is $\lambda_j/\cos\varphi_{X_j}\ne0$, so $\Psi_{I,j}$ is an affine bijection of $\R^2$, and $\mathbf Q_{I,j}:=\Psi_{I,j}(\mathbf J_{I,j})$ is open. Put
\[
\bar\rho(q,y):=\sum_{I,j}\Big|\frac{\cos\varphi_{X_j}}{\lambda_j}\Big|\,\mathbf 1_{\mathbf Q_{I,j}}(q,y).
\]
For fixed $y$, the slice of $\mathbf Q_{I,j}$ at height $y$ is $X_c(J_c)$ with $c=(I,j)$ as in the proof of (e). So $\bar\rho(\cdot,y)$ is a version of $\rho_y$. The functions $(q,y)\mapsto(\bar\rho(q,y)-1)_+\mathbf 1_{\rm waste}(q,y)$ and $(q,y)\mapsto(1-\bar\rho(q,y))_+\mathbf 1_{\rm waste}(q,y)$ are Borel, and Tonelli's theorem shows that $\Ovgap$ and $U$ are Borel. \qed
\end{remark}

\section{Deaths, merges and gap overlap (P5)}\label{sec:P5}

We keep the standing assumptions and the notation of Section~\ref{ssec:P3-setting}, and we use sources and delays (Definition~\ref{def:P3-reach}) and Lemmas~\ref{lem:P3-reach}--\ref{lem:P3-mult}.

Let $\cD$ and $\cM\subset(0,k)$ be the sets of $x$ such that $\pi_x$ terminates with~D, respectively with~M, and let $\tilde\cD$, $\tilde\cM$ be their ceiling versions (Definitions~\ref{def:flows-losses} and~\ref{def:flows-ceiling}). For $y\in(0,\hmax)$, $\mathrm{Wa}_y=\ell_y\cap{\rm waste}$, $A_{\rm gap}(y)$, $X_y$, the density $\rho_y$ of the gap crossers and the gap overlap $\Ovgap^{\Fz}(y)=\int_{\mathrm{Wa}_y}(\rho_y-1)_+$ are those of Definitions~\ref{def:flows-line} and~\ref{def:flows-ovgap} and Lemma~\ref{lem:flows-density}; a path is \emph{alive at $y$} if its termination height exceeds $y$ (Definition~\ref{def:flows-traj}). Finally (Definition~\ref{def:flows-ovgap})
\[
\Lam_0:=|\cD|+|\cM|+L_E+\sup_{0<y<\hmax}\Ovgap^{\Fz}(y),
\]
and $\tilde\Lam_0:=\Lam_0[\rho\cP]$. We put $Q_*:=0.32$, the constant of [R, Proposition~\ref{imp:P4.5}].

\begin{theorem}[P5]\label{thm:P5}
Assume the standing assumptions of Section~\ref{ssec:P3-setting}. Let $K$ be the constant of Lemma~\ref{lem:P3-mult} and
\[
S_K:=(1+10^{-11})\,\frac{4K}{Q_*\,(2-3\cdot10^{-6})},\qquad\text{so that}\qquad S_9\le56.2501,\quad S_{13}\le81.2502
\]
(rounded up). Then:
\begin{enumerate}
\item[(D)] $|\cD|+|\tilde\cD|\ \le\ (1+S_K\delta^2)\,\dfrac W\delta$;
\item[(M)] $|\cM|+|\tilde\cM|\ \le\ \dfrac{4}{\cos(\pi/4+\amax)}\,S_K\,\delta W$;
\item[(G)] $\sup_y\Ovgap^{\Fz}(y)+\sup_y\Ovgap^{\tilde F^0}(y)\ \le\ \dfrac{S_K}{\cos\amax}\,\delta W$, the suprema being over all heights at which the gap overlaps are defined.
\end{enumerate}
Numerically (all values rounded up): with $K=9$ (using [R, Lemma~\ref{imp:L4.10}]),
\[
|\cD|+|\tilde\cD|\le C_D\frac W\delta,\quad |\cM|+|\tilde\cM|\le C_M\,\delta W,\quad \sup_y\Ovgap^{\Fz}+\sup_y\Ovgap^{\tilde F^0}\le C_G\,\delta W,
\]
with $C_D=1+56.2501\,\delta^2$, $C_M=318.20$, $C_G=56.2501$; with $K=13$ (using only [R, Lemma~\ref{imp:L4.9}]) the same holds with $C_D=1+81.26\,\delta^2$, $C_M=459.7$, $C_G=81.26$.
\end{theorem}

\begin{corollary}\label{cor:Lambda0}
Under the standing assumptions,
\[
\Lam_0+\tilde\Lam_0\ \le\ C_\Lam\,\frac W\delta ,
\]
where $C_\Lam=1+1185.6\,\delta^2<1.0000002$ \textup{(}Definition~\ref{def:flows-consts}\textup{)} if $K=9$ is used; if $K=13$ is used, the same holds with $C^{(13)}_\Lam=1+1712.6\,\delta^2<1.0000002$ in place of $C_\Lam$.
\end{corollary}

\begin{remark}[the hypothesis $\delta\le\frac12\cos2\amax$]\label{rem:P5-C6}
Under the standing assumptions, $\delta\le10^{-5}<0.4999999<\frac12\cos(3\cdot10^{-6})\le\frac12\cos2\amax$; this is constraint~(C6) of Definition~\ref{def:flows-constraints}, which follows from~(C1) (Lemma~\ref{lem:flows-constraints}(a)). It is a hypothesis of the geometric Lemmas~\ref{lem:P5-two} and~\ref{lem:P5-refined} below; Lemma~\ref{lem:P5-atmost2}, Lemma~\ref{lem:P5-pairsum} and Theorem~\ref{thm:P5} depend on it through these two lemmas. Section~\ref{sec:P3} does not use it.
\end{remark}

The proof of Theorem~\ref{thm:P5} rests on one geometric fact (Lemmas~\ref{lem:P5-atmost2} and~\ref{lem:P5-refined}): a point is reached by at most two paths; if it is reached by two, their sources are two squares $X_1,X_2$ with $\varphi_{X_1}>\varphi_{X_2}$, the point lies in a rectangle of width $\delta\tan\theta/\cos\theta$ and height $\delta$ above the upper left vertex of $X_1$ ($\theta=\varphi_{X_1}-\varphi_{X_2}$), and the two squares are at distance less than $\delta\tan\theta$. The rest consists of summing the areas and lengths of these rectangles with the pair waste of~[R].

\subsection{Points with two sources}\label{ssec:P5-two}

\begin{lemma}[floor and square sources]\label{lem:P5-floor}
No point is reached both from the floor $X_0$ and from a square.
\end{lemma}

\begin{proof}
Let $q$ be reached from $X_0$ with delay $t_0$ and from a square $X$ with delay $t$. By parts (a) and (e) of Lemma~\ref{lem:P3-reach}, $q=(x,t_0)$ with $t_0<\delta$, and $t>0$. The start point of the second step is $p:=q-t\,n_X$, a point of $\relint\Top(X)$. The point $p-n_X$ lies in $\Bot(X)\subset X$, and its height is
\[
q_y-(t+1)\cos\varphi_X\ <\ \delta-\cos\amax\ <\ 0,
\]
since $|\varphi_X|<\amax$ (Lemma~\ref{lem:P3-reach}(b)) and $\delta\le10^{-5}<\cos\amax$. This contradicts $X\subset[0,k]^2$.
\end{proof}

The next lemma is purely geometric; it is stated with its own parameters $\amax$ and $\delta$.

\begin{lemma}[two sources]\label{lem:P5-two}
Let $0<\amax<\pi/4$ and $\delta>0$ satisfy
\[
\delta\ \le\ \tfrac12\cos2\amax .
\]
Let $X_1,X_2$ be two disjoint closed unit squares with $|\varphi_{X_1}|,|\varphi_{X_2}|<\amax$, and let $q$ be a point such that $q-t_i\,n_{X_i}\in\relint\Top(X_i)$ with $0<t_i<\delta$ for $i=1,2$. Then $\varphi_{X_1}\ne\varphi_{X_2}$. Label the squares so that $\varphi_{X_1}>\varphi_{X_2}$, and put $\theta:=\varphi_{X_1}-\varphi_{X_2}\in(0,2\amax)$ and $\tau_i:=(q-c_{X_i})\cdot u_{X_i}$. Then
\[
\xi_1:=\tau_1+\tfrac12\ <\ \tfrac12\sin\theta,\qquad \xi_2:=\tfrac12-\tau_2\ <\ \tfrac12\sin\theta .
\]
In particular, under the standing assumptions this applies to any two distinct square sources $X_1,X_2$ of a point $q$, with their delays $t_1,t_2$ (Lemma~\ref{lem:P3-reach}(a),(b),(e)).
\end{lemma}

So $q$ lies just above the left end of $\Top(X_1)$ and just above the right end of $\Top(X_2)$.

\begin{proof}
Write $u_i,n_i,c_i,\varphi_i$ for $u_{X_i},n_{X_i},c_{X_i},\varphi_{X_i}$. Since $q-t_in_i\in\relint\Top(X_i)$, we have $(q-c_i)\cdot n_i=\frac12+t_i$ and $\tau_i\in(-\frac12,\frac12)$. Note $|\varphi_1-\varphi_2|<2\amax<\pi/2$, so $\cos(\varphi_1-\varphi_2)\ge\cos2\amax$.

Put $z:=q-\frac12n_1$. In the frame of $X_1$, $(z-c_1)\cdot u_1=\tau_1\in(-\frac12,\frac12)$ and $(z-c_1)\cdot n_1=t_1\in(0,\frac12]$; hence $z\in X_1$. In the frame of $X_2$, by~\eqref{eq:P3-frames},
\[
(z-c_2)\cdot n_2=\tfrac12+t_2-\tfrac12\cos(\varphi_1-\varphi_2)\in\big[t_2,\ t_2+\tfrac12(1-\cos2\amax)\big],
\qquad
(z-c_2)\cdot u_2=\tau_2-\tfrac12\sin(\varphi_2-\varphi_1).
\]
Because $0<t_2<\delta\le\frac12\cos2\amax$, the $n_2$-coordinate $(z-c_2)\cdot n_2$ lies in $(0,\frac12]$. \emph{This is where the hypothesis $\delta\le\frac12\cos2\amax$ enters the proof} (and in the same step with the roles of $X_1$ and $X_2$ exchanged). Since $X_1\cap X_2=\emptyset$, $z\notin X_2$, and therefore
\begin{equation}\label{eq:P5-tau2}
\big|\tau_2-\tfrac12\sin(\varphi_2-\varphi_1)\big|>\tfrac12 .
\end{equation}
Exchanging the roles of $X_1$ and $X_2$ (with $z':=q-\frac12n_2\in X_2$) gives
\begin{equation}\label{eq:P5-tau1}
\big|\tau_1-\tfrac12\sin(\varphi_1-\varphi_2)\big|>\tfrac12 .
\end{equation}
If $\varphi_1=\varphi_2$, \eqref{eq:P5-tau2} says $|\tau_2|>\frac12$, which is false. So $\varphi_1\ne\varphi_2$; let $\varphi_1>\varphi_2$ and $s:=\sin\theta\in(0,1)$. Since $\tau_2>-\frac12$, we have $\tau_2+\frac s2>-\frac12$, so~\eqref{eq:P5-tau2} gives $\tau_2+\frac s2>\frac12$, that is, $\xi_2<\frac s2$. Since $\tau_1<\frac12$, we have $\tau_1-\frac s2<\frac12$, so~\eqref{eq:P5-tau1} gives $\tau_1-\frac s2<-\frac12$, that is, $\xi_1<\frac s2$.
\end{proof}

\begin{remark}[the hypothesis cannot be weakened]\label{rem:P5-sharp}
The hypothesis $\delta\le\frac12\cos2\amax$ cannot be replaced by a weaker bound on $\delta$ in Lemma~\ref{lem:P5-two}. Let $0<\amax<\pi/4$ and $\delta>\frac12\cos2\amax$. Choose $\theta\in(0,2\amax)$ with $\frac12\cos\theta<\delta$ (possible by continuity), then $t\in(\frac12\cos\theta,\delta)$, put $v:=\sin\theta\,(t-\frac12\cos\theta)>0$, and choose $\eta\in(0,\min(v,1))$. In coordinates with $q=(0,0)$ let
\[
X_2:=[\eta-1,\eta]\times[-1-t,-t],\qquad X_1:=\{v_1+a\,u-b\,n:\ a,b\in[0,1]\},
\]
where $u=(\cos\theta,\sin\theta)$, $n=(-\sin\theta,\cos\theta)$ and $v_1:=-\frac12\sin\theta\,u-t\,n$. Then $q-t\,e_2=(0,-t)\in\relint\Top(X_2)$, since $0\in(\eta-1,\eta)$, and $q-t\,n=v_1+\frac12\sin\theta\,u\in\relint\Top(X_1)$, since $\frac12\sin\theta\in(0,1)$; both delays equal $t\in(0,\delta)$. The first coordinate of $v_1$ is $-\frac12\sin\theta\cos\theta+t\sin\theta=v$, and every point of $X_1$ has first coordinate $v+a\cos\theta+b\sin\theta\ge v>\eta$, while $X_2\subset\{x\le\eta\}$; so $X_1\cap X_2=\emptyset$. (Moreover the segment from $(0,-t)$ to $q$ has first coordinate $0<v$ and misses $X_1$, and the points $q-s\,n=(s\sin\theta,-s\cos\theta)$, $0\le s\le t$, have second coordinate $>-t$ and miss $X_2$.) Rotating the configuration about $q$ by $-\theta/2$ gives phases $\theta/2$ and $-\theta/2$, of absolute value $<\amax$, and the frames used above become the frames of these phases. All hypotheses of Lemma~\ref{lem:P5-two} hold except $\delta\le\frac12\cos2\amax$, but $\xi_1=\frac12\sin\theta$, so the conclusion $\xi_1<\frac12\sin\theta$ fails. Under the standing assumptions, $\delta$ is far below the threshold (Remark~\ref{rem:P5-C6}).
\end{remark}

\begin{lemma}[at most two paths through a point]\label{lem:P5-atmost2}
For every point $q$, at most two $x\in(0,k)$ have the property that $\pi_x$ reaches $q$ (in the gap case or in the contact case). If there are two, their sources are two distinct squares $X_1,X_2$, and $q$, $X_1,X_2$ and the two delays satisfy the hypotheses of Lemma~\ref{lem:P5-two}. In particular $M_g(q)\le2$ for every $q$.
\end{lemma}

\begin{proof}
Let $x_1,\dots,x_m$ be distinct and let $\pi_{x_i}$ reach $q$ from the source $X_i$ with delay $t_i$ (unique by Lemma~\ref{lem:P3-reach}(f)). By Lemma~\ref{lem:P3-samesource}(b) the sources $X_i$ are pairwise distinct. So at most one of them is the floor, and if one of them is the floor, then none is a square by Lemma~\ref{lem:P5-floor}, and $m=1$. Otherwise all $X_i$ are squares with $|\varphi_{X_i}|<\amax$ and $0<t_i<\delta$ (Lemma~\ref{lem:P3-reach}(a),(b)), and Lemma~\ref{lem:P5-two} applies to each pair (Remark~\ref{rem:P5-C6}); in particular the phases are pairwise distinct. Suppose $m\ge3$ and let $A,B,C$ be three of the sources with $\varphi_A>\varphi_B>\varphi_C$. Applied to $(A,B)$, where $B$ has the smaller phase, Lemma~\ref{lem:P5-two} gives $\tau_B>\frac12-\frac12\sin(\varphi_A-\varphi_B)>0$, because $0<\varphi_A-\varphi_B<2\amax<\pi/2$. Applied to $(B,C)$, where $B$ has the larger phase, it gives $\tau_B<-\frac12+\frac12\sin(\varphi_B-\varphi_C)<0$. This is a contradiction, so $m\le2$, and if $m=2$ both sources are squares.
\end{proof}

\begin{lemma}[two sources: refined form]\label{lem:P5-refined}
In the situation of Lemma~\ref{lem:P5-two}, with $\varphi_{X_1}>\varphi_{X_2}$ and the notation there, let $v_1$ be the upper left vertex of $X_1$ and $v_2$ the upper right vertex of $X_2$:
\[
v_1:=c_{X_1}-\tfrac12u_{X_1}+\tfrac12n_{X_1},\qquad v_2:=c_{X_2}+\tfrac12u_{X_2}+\tfrac12n_{X_2}.
\]
Then $\xi_1=(q-v_1)\cdot u_{X_1}$, $t_1=(q-v_1)\cdot n_{X_1}$, $\xi_2=(v_2-q)\cdot u_{X_2}$, $t_2=(q-v_2)\cdot n_{X_2}$, and
\begin{equation}\label{eq:P5-star}
\cos\theta\ \xi_1+\xi_2\ \le\ \sin\theta\,\max\big(t_1,\,t_2/\cos\theta\big)\ <\ \delta\tan\theta .
\end{equation}
Consequently:
\begin{enumerate}[label=(\roman*)]
\item $0<\xi_1<\delta\tan\theta/\cos\theta$ and $0<\xi_2<\delta\tan\theta$;
\item $\dist(X_1,X_2)\le\sin\theta\,\max(t_1,t_2/\cos\theta)<\delta\tan\theta$;
\item $q=v_1+a\,u_{X_1}+b\,n_{X_1}$ with $0<a<\delta\tan\theta/\cos\theta$ and $0<b<\delta$.
\end{enumerate}
\end{lemma}

\begin{proof}
Write $u_i,n_i,c_i$ as before and $\kappa:=\cos\theta$, $\varsigma:=\sin\theta$. The four identities follow from $(v_1-c_1)\cdot u_1=-\frac12$, $(v_1-c_1)\cdot n_1=\frac12$, $(v_2-c_2)\cdot u_2=\frac12$, $(v_2-c_2)\cdot n_2=\frac12$ and $(q-c_i)\cdot n_i=\frac12+t_i$. Since $\varphi_2=\varphi_1-\theta$,
\[
u_2=\kappa\,u_1-\varsigma\,n_1,\qquad n_2=\varsigma\,u_1+\kappa\,n_1,
\]
so $u_1\cdot u_2=\kappa$, $n_1\cdot u_2=-\varsigma$, $u_1\cdot n_2=\varsigma$, $n_1\cdot n_2=\kappa$. Moreover
\[
X_1=\{v_1+a\,u_1-b\,n_1:\ a,b\in[0,1]\},\qquad X_2=\{v_2-a\,u_2-b\,n_2:\ a,b\in[0,1]\}.
\]
Let $\bar t:=\max(t_1,t_2/\kappa)$. Note $\kappa>\cos2\amax>0$ and $\bar t<\delta/\kappa\le\frac12\cos2\amax/\cos\theta<\frac12$.

\emph{A family of points of $X_1$.} For $\varepsilon\in(0,1]$ and $\eta>0$ put
\[
\hat t:=\max\big(t_1+\eta,\ (t_2+\varsigma\varepsilon)/\kappa\big),\qquad h:=\hat t-t_1\ (\ge\eta),\qquad z:=v_1+\varepsilon\,u_1-h\,n_1 .
\]
(The letters $\bar t$, $\hat t$, $h$, $\kappa$, $\varsigma$, $a_z$, $b_z$ are local to this proof.) As $\varepsilon,\eta\to0^+$, $\hat t\to\bar t<\frac12$; so for small $\varepsilon,\eta$ we have $h\le1$ and $z\in X_1$. Using $v_2-q=\xi_2u_2-t_2n_2$ and $q-z=(\xi_1-\varepsilon)u_1+\hat t\,n_1$, the coordinates of $z$ in the frame of $X_2$ are
\[
a_z:=(v_2-z)\cdot u_2=(\kappa\,\xi_1+\xi_2)-\kappa\,\varepsilon-\varsigma\,\hat t,\qquad
b_z:=(v_2-z)\cdot n_2=-t_2+\varsigma(\xi_1-\varepsilon)+\kappa\,\hat t .
\]
Since $\kappa\hat t\ge t_2+\varsigma\varepsilon$ and $\xi_1>0$, we get $b_z\ge\varsigma\xi_1\ge0$. Also $b_z\le\varsigma\xi_1+\kappa\hat t$, which tends to $\varsigma\xi_1+\kappa\bar t<\frac12\varsigma^2+\frac12<1$ by Lemma~\ref{lem:P5-two} and $\kappa\bar t<\delta\le\frac12$; so $b_z\in[0,1]$ for small $\varepsilon,\eta$. Finally $a_z\le\kappa\xi_1+\xi_2<\frac12\varsigma(\kappa+1)<1$ by Lemma~\ref{lem:P5-two}.

\emph{Proof of~\eqref{eq:P5-star}.} Suppose $\kappa\xi_1+\xi_2>\varsigma\bar t$. As $\varepsilon,\eta\to0^+$, $\kappa\varepsilon+\varsigma\hat t\to\varsigma\bar t$, so $a_z>0$ for small $\varepsilon,\eta$. Then $a_z,b_z\in[0,1]$, so $z\in X_2$; but $z\in X_1$, contradicting $X_1\cap X_2=\emptyset$. Hence $\kappa\xi_1+\xi_2\le\varsigma\bar t$. Since $t_1<\delta\le\delta/\kappa$ and $t_2/\kappa<\delta/\kappa$, $\varsigma\bar t<\varsigma\delta/\kappa=\delta\tan\theta$.

(i) $\xi_1=\tau_1+\frac12>0$ and $\xi_2=\frac12-\tau_2>0$ because $\tau_i\in(-\frac12,\frac12)$. Then~\eqref{eq:P5-star} gives $\kappa\xi_1<\delta\tan\theta$ and $\xi_2<\delta\tan\theta$.

(ii) Now~\eqref{eq:P5-star} holds. For small $\varepsilon,\eta$ take $z\in X_1$ as above. Since $\hat t\ge\bar t$, $a_z\le\varsigma\bar t-\kappa\varepsilon-\varsigma\hat t\le-\kappa\varepsilon<0$, and $b_z\in[0,1]$. Put $z':=z+a_z\,u_2$. Then $(v_2-z')\cdot u_2=a_z-a_z=0$ and $(v_2-z')\cdot n_2=b_z$, so $z'\in X_2$, and $|z-z'|=|a_z|=\kappa\varepsilon+\varsigma\hat t-(\kappa\xi_1+\xi_2)\le\kappa\varepsilon+\varsigma\hat t$. Letting $\varepsilon,\eta\to0^+$ gives $\dist(X_1,X_2)\le\varsigma\bar t<\delta\tan\theta$.

(iii) Take $a=\xi_1$ and $b=t_1$, and use (i) and $0<t_1<\delta$.
\end{proof}

\begin{lemma}[small discs]\label{lem:P5-disc}
Every closed disc of radius $r\le10^{-4}$ meets at most $6$ squares of $\cP$.
\end{lemma}

\begin{proof}
Let $S$ be a square and $p\in S$, $p=c_S+a\,u_S+b\,n_S$ with $|a|,|b|\le\frac12$. Put $\operatorname{sgn}0:=1$ and
\[
Z_p:=\{c_S+a'u_S+b'n_S:\ a'\ \text{between}\ a\ \text{and}\ a-\tfrac12\operatorname{sgn}a,\ \ b'\ \text{between}\ b\ \text{and}\ b-\tfrac12\operatorname{sgn}b\}.
\]
If $a\ge0$, the values $a'$ lie in $[a-\frac12,a]\subset[-\frac12,\frac12]$, and if $a<0$, in $[a,a+\frac12]\subset[-\frac12,\frac12]$; likewise for $b'$. So $Z_p\subset S$ is a square of side $\frac12$ and area $\frac14$ with vertex $p$, and $Z_p\subset\bar B(p,\sqrt2/2)$. Let $\bar B(o,r)$ be the disc, and for each square $S$ meeting it choose $p_S\in S\cap\bar B(o,r)$. The squares $Z_{p_S}$ lie in pairwise disjoint squares, so they are pairwise disjoint, and they lie in $\bar B(o,r+\sqrt2/2)$. Hence their number is at most $4\pi(r+\sqrt2/2)^2\le4\pi\cdot0.70721^2<6.29$.
\end{proof}

\subsection{Valley pairs}\label{ssec:P5-valley}

\begin{definition}[valley pairs]\label{def:P5-valley}
A point reached by two paths of $\Fz$ is a \emph{double point}. By Lemma~\ref{lem:P5-atmost2} its two sources are squares $X_1,X_2$; we label them so that $\varphi_{X_1}>\varphi_{X_2}$ and call $e=\{X_1,X_2\}$ the \emph{valley pair} of the double point. For a valley pair $e$ let $\operatorname{Dbl}_e$ be the set of its double points, and put
\[
\theta_e:=\varphi_{X_1}-\varphi_{X_2},\qquad w_e:=\frac{\delta\tan\theta_e}{\cos\theta_e},\qquad
\operatorname{Rect}_e:=\{v_1+a\,u_{X_1}+b\,n_{X_1}:\ 0<a<w_e,\ 0<b<\delta\},
\]
where $v_1$ is the upper left vertex of $X_1$. Let $\mathcal Q_b$ be the set of valley pairs of $\Fz$, and let $\mathcal Q_c$ be the set of valley pairs of the master flow of $\rho\cP$, transported to $\cP$ by $\rho$ (with $\operatorname{Dbl}_e$ and $\operatorname{Rect}_e$ transported by $\rho$; $\theta_e$ and $w_e$ are unchanged).
\end{definition}

By Lemma~\ref{lem:P5-two}, $0<\theta_e<2\amax\le3\cdot10^{-6}$, so $\theta_e=\theta(X_1,X_2)$. By Lemma~\ref{lem:P5-refined}(iii), $\operatorname{Dbl}_e\subset\operatorname{Rect}_e$; and $|\operatorname{Rect}_e|=\delta\,w_e$. Both sets $\mathcal Q_b,\mathcal Q_c$ are finite.

\begin{lemma}[sums over valley pairs]\label{lem:P5-pairsum}
Let $K$ be the constant of Lemma~\ref{lem:P3-mult}. Then:
\begin{enumerate}[label=(\alph*)]
\item $\mathcal Q_b\cap\mathcal Q_c=\emptyset$;
\item $\displaystyle\sum_{e\in\mathcal Q_b\cup\mathcal Q_c}\theta_e\ \le\ \frac{4K}{Q_*(2-3\cdot10^{-6})}\,W$;
\item $\displaystyle\sum_{e\in\mathcal Q_b\cup\mathcal Q_c}w_e\ \le\ S_K\,\delta W$, where $S_9\le56.2501$ and $S_{13}\le81.2502$ (rounded up).
\end{enumerate}
\end{lemma}

\begin{proof}
(a) The squares of a pair in $\mathcal Q_b$ are sources of points reached by paths of $\Fz$, hence floor squares; those of a pair in $\mathcal Q_c$ are ceiling squares. By Lemma~\ref{lem:P3-separation} no square is both.

(b) Let $e=\{X_1,X_2\}\in\mathcal Q_b$. By Lemma~\ref{lem:P5-refined}(ii), $\dist(X_1,X_2)<\delta\tan\theta_e<10^{-5}$, and $\theta_e<3\cdot10^{-6}\le10^{-5}$. By [R, Proposition~\ref{imp:P4.5}] with $d=\theta_1=10^{-5}$, the waste in $R_e=R(X_1,X_2)$ has area
\[
f_e\ \ge\ \frac{Q_*\,\theta_e(2-\theta_e)}{4}\ \ge\ \frac{Q_*(2-3\cdot10^{-6})}{4}\,\theta_e .
\]
For $e\in\mathcal Q_c$ the same holds: apply the argument to $\rho\cP$ and use that $\rho$ preserves distances, $\theta$, the rectangles $R(\cdot,\cdot)$ and the waste. By (a), $\mathcal Q_b\cup\mathcal Q_c$ is a set of pairs of squares at distance $<10^{-5}\le10^{-4}$, so Lemma~\ref{lem:P3-mult} gives $\sum_ef_e\le KW$, and (b) follows.

(c) For $0<\theta\le3\cdot10^{-6}$,
\[
\frac{\tan\theta}{\theta\cos\theta}=\frac{\sin\theta}{\theta}\cdot\frac1{\cos^2\theta}\ \le\ \frac1{\cos^2(3\cdot10^{-6})}\ \le\ 1+10^{-11},
\]
so $w_e\le(1+10^{-11})\,\delta\,\theta_e$, and (c) follows from (b). Numerically, $4\cdot9/(0.32\,(2-3\cdot10^{-6}))=56.2500843\ldots$ and $4\cdot13/(0.32\,(2-3\cdot10^{-6}))=81.2501218\ldots$; multiplied by $1+10^{-11}$ and rounded up they give $56.2501$ and $81.2502$.
\end{proof}

\subsection{Deaths}\label{ssec:P5-D}

\begin{proof}[Proof of Theorem~\ref{thm:P5}\,\textup{(D)}]
Consider first the floor flow.

\emph{Cells.} By~\ref{it:P3-S2}, $\cD\setminus\Sigma$ is the union of the pieces $I$ that it contains. Fix such an $I$. For $x\in I$ the path $\pi_x$ has the ray steps $j=0,1,\dots,J$, where $J$ and the sources $X_j$ do not depend on $x\in I$, with start points $p_j(x)$, cumulative gaps $g_j(x)$ at the start and lengths $t^*_j(x)$, all affine in $x$. The steps $j<J$ end with contacts that are entry candidates won by $\pi_x$ and followed by a pass, and step $J$ ends with~D. Passes do not change the cumulative gap, so $g_0=0$ and $g_{j+1}=g_j+t^*_j$, and $g_J+t^*_J=\delta$ because step $J$ ends at $t_D=\delta-g_J$. For $j=0,\dots,J$ let
\[
C_{I,j}:=\{(x,t):\ x\in I,\ g_j(x)<t<g_j(x)+t^*_j(x)\},\qquad
\gamma_{I,j}(x,t):=p_j(x)+\big(t-g_j(x)\big)\,n_{X_j}.
\]
For fixed $x\in I$ the intervals $(g_j(x),g_j(x)+t^*_j(x))$, $j=0,\dots,J$, are pairwise disjoint and cover $(0,\delta)$ up to the points $g_1(x),\dots,g_J(x)$. So the sets $C_{I,j}$ are pairwise disjoint, and their union is $I\times(0,\delta)$ up to the graphs of the affine functions $g_1,\dots,g_J$, which are null sets. Summing over $I$,
\begin{equation}\label{eq:P5-cells}
\sum_{I,j}|C_{I,j}|=\delta\,|\cD| .
\end{equation}

\emph{Jacobian.} The map $\gamma_{I,j}$ is affine, with $\partial_x\gamma_{I,j}=p_j'-g_j'\,n_{X_j}=\lambda_j\,u_{X_j}-g_j'\,n_{X_j}$ and $\partial_t\gamma_{I,j}=n_{X_j}$, where $\lambda_j\ge1$ by~\ref{it:P3-S3}. Hence $\det D\gamma_{I,j}=\lambda_j\det[u_{X_j},n_{X_j}]=\lambda_j$, so $\gamma_{I,j}$ is injective and
\[
|\gamma_{I,j}(C_{I,j})|=\lambda_j\,|C_{I,j}|\ \ge\ |C_{I,j}| .
\]

\emph{Images.} For $(x,t)\in C_{I,j}$ the point $\gamma_{I,j}(x,t)$ is reached by $\pi_x$ from the source $X_j$ with delay $t-g_j(x)\in(0,t^*_j(x))$, that is, in the gap case. By Lemma~\ref{lem:P3-reach}(d) it is a waste point of the strip $S_b:=(0,k)\times(0,\hmax)$.

\emph{Multiplicity.} For a point $z$ let $\operatorname{mult}(z):=\sum_{I,j}\mathbf 1_{\gamma_{I,j}(C_{I,j})}(z)$; each term equal to $1$ corresponds to exactly one preimage $(x,t)\in C_{I,j}$ of $z$. Two such preimages $(x,t)\in C_{I,j}$ and $(x',t')\in C_{I',j'}$ with $(I,j)\ne(I',j')$ have $x\ne x'$: if $x=x'$, then $I=I'$, and $\pi_x$ would reach $z$ in the steps $j$ and $j'$, so $j=j'$ by Lemma~\ref{lem:P3-reach}(f). All these $x$ reach $z$ in the gap case, so $\operatorname{mult}(z)\le M_g(z)$. By Lemma~\ref{lem:P5-atmost2}, $M_g(z)\le2$, and if $M_g(z)=2$, then $z$ is a double point and $z\in\operatorname{Dbl}_e\subset\operatorname{Rect}_e$ for its valley pair $e\in\mathcal Q_b$. Thus $\operatorname{mult}(z)\le2$; if $\operatorname{mult}(z)\ge1$, then $z$ is a waste point of $S_b$ by the previous step; and if $\operatorname{mult}(z)=2$, then $z\in\operatorname{Rect}_e$ for some $e\in\mathcal Q_b$. Hence, for every point $z$,
\[
\operatorname{mult}(z)\ \le\ \mathbf 1_{{\rm waste}\cap S_b}(z)+\sum_{e\in\mathcal Q_b}\mathbf 1_{\operatorname{Rect}_e}(z).
\]

\emph{Integration.} By~\eqref{eq:P5-cells} and the two previous steps,
\begin{align*}
\delta\,|\cD|=\sum_{I,j}|C_{I,j}|&\le\sum_{I,j}|\gamma_{I,j}(C_{I,j})|=\int \operatorname{mult}(z)\,dz\\
&\le|{\rm waste}\cap S_b|+\sum_{e\in\mathcal Q_b}|\operatorname{Rect}_e|=|{\rm waste}\cap S_b|+\delta\sum_{e\in\mathcal Q_b}w_e .
\end{align*}

\emph{Ceiling flow and conclusion.} Applying this to $\rho\cP$ and reflecting, $\delta|\tilde\cD|\le|{\rm waste}\cap S_c|+\delta\sum_{e\in\mathcal Q_c}w_e$ with $S_c:=\rho(S_b)=(0,k)\times(k-\hmax,k)$. Since $\hmax<k-\hmax$, the strips $S_b$ and $S_c$ are disjoint, so $|{\rm waste}\cap S_b|+|{\rm waste}\cap S_c|\le W$. With Lemma~\ref{lem:P5-pairsum}(a),(c),
\[
|\cD|+|\tilde\cD|\ \le\ \frac W\delta+\sum_{e\in\mathcal Q_b\cup\mathcal Q_c}w_e\ \le\ \frac W\delta+S_K\,\delta W=(1+S_K\delta^2)\,\frac W\delta . \qedhere
\]
\end{proof}

\begin{remark}\label{rem:P5-CD2}
Using only $\operatorname{mult}\le M_g\le2$ (Lemma~\ref{lem:P5-atmost2}), the same computation gives $|\cD|+|\tilde\cD|\le2W/\delta$ without any result of~[R]. The bound $(1+O(\delta^2))W/\delta$ uses that double points lie in the rectangles $\operatorname{Rect}_e$, together with Lemma~\ref{lem:P5-pairsum}.
\end{remark}

\subsection{Merges}\label{ssec:P5-M}

\begin{proof}[Proof of Theorem~\ref{thm:P5}\,\textup{(M)}]
Consider the floor flow. For a square $Y$ let $\cM_Y$ be the set of $x\in\cM$ whose path stops with~M at an entry candidate on $Y$; then $\cM$ is the disjoint union of the sets $\cM_Y$. For $x\in\cM_Y$ let $q(x)\in\relint\Bot(Y)$ be that entry candidate.

\emph{One loser per entry point.} Let $x\in\cM_Y$. The paths that reach the entry candidate $(Y,q(x))$ alive are the paths that reach $q(x)$ in the contact case; they include $x$ and the winner of R1 there, which is different from $x$. By Lemma~\ref{lem:P5-atmost2} there are exactly two of them, their sources form a valley pair $e\in\mathcal Q_b$, and $q(x)\in\operatorname{Dbl}_e$. Since only one of the two paths loses, the map $x\mapsto q(x)$ is injective on $\cM_Y$.

\emph{Expansion.} By~\ref{it:P3-S2} and~\ref{it:P3-S3}, $\cM_Y\setminus\Sigma$ is a finite union of pieces $I$, and on each of them $q$ is affine along the line of $\Bot(Y)$ with $q'=\mu u_Y$, $\mu\ge1$. With injectivity,
\[
|\cM_Y|=\sum_I|I|\le\sum_I\mathcal H^1(q(I))=\mathcal H^1\big(q(\cM_Y\setminus\Sigma)\big)\ \le\ \sum_{e}\mathcal H^1\big(\Bot(Y)\cap\operatorname{Rect}_e\big),
\]
the last sum over the $e\in\mathcal Q_b$ with $\operatorname{Dbl}_e\cap\relint\Bot(Y)\ne\emptyset$, because $q(x)\in\operatorname{Dbl}_e\subset\operatorname{Rect}_e$.

\emph{Length of one intersection.} Let $e=\{X_1,X_2\}$. The rectangle $\operatorname{Rect}_e$ lies in the strip $\{z:0<(z-v_1)\cdot u_{X_1}<w_e\}$. Along the line of $\Bot(Y)$, which has the direction $u_Y$, the function $(z-v_1)\cdot u_{X_1}$ changes at the rate $u_Y\cdot u_{X_1}=\cos(\varphi_Y-\varphi_{X_1})$. Since $|\varphi_{X_1}|<\amax$ and $|\varphi_Y|\le\pi/4$, we have $|\varphi_Y-\varphi_{X_1}|<\pi/4+\amax$, so this rate is $>\cos(\pi/4+\amax)>0$, and
\[
\mathcal H^1\big(\Bot(Y)\cap\operatorname{Rect}_e\big)\ \le\ \frac{w_e}{\cos(\pi/4+\amax)} .
\]

\emph{At most four squares per pair.} Fix $e=\{X_1,X_2\}\in\mathcal Q_b$, and let $Y$ be a square with a point $q\in\operatorname{Dbl}_e\cap\relint\Bot(Y)$. As $q\in\operatorname{Rect}_e$ and $w_e<3.1\cdot10^{-6}\delta$, $|q-v_1|<(w_e^2+\delta^2)^{1/2}<1.01\,\delta$. By Lemma~\ref{lem:P5-refined}(i), $|q-v_2|=(\xi_2^2+t_2^2)^{1/2}<\delta(\tan^2\theta_e+1)^{1/2}<1.01\,\delta$, with $v_2$ the upper right vertex of $X_2$. So the closed disc $\bar B(v_1,3\delta)$ contains $v_1\in X_1$, $v_2\in X_2$ and $q\in Y$. Moreover $Y\notin\{X_1,X_2\}$, because $(q-c_{X_i})\cdot n_{X_i}=\frac12+t_i>\frac12$ shows $q\notin X_i$. Hence $X_1$, $X_2$ and all such squares $Y$ are distinct squares meeting $\bar B(v_1,3\delta)$, a disc of radius $3\delta\le3\cdot10^{-5}$. By Lemma~\ref{lem:P5-disc} there are at most $6$ of them, so at most $4$ such squares $Y$.

\emph{Summation.} Summing over $Y$ and exchanging the order of summation,
\begin{align*}
|\cM|=\sum_Y|\cM_Y|&\le\frac{1}{\cos(\pi/4+\amax)}\sum_{e\in\mathcal Q_b}w_e\cdot\#\{Y:\ \operatorname{Dbl}_e\cap\relint\Bot(Y)\ne\emptyset\}\\
&\le\frac{4}{\cos(\pi/4+\amax)}\sum_{e\in\mathcal Q_b}w_e .
\end{align*}
Applying this to $\rho\cP$ gives the same bound for $|\tilde\cM|$ with $\mathcal Q_c$. By Lemma~\ref{lem:P5-pairsum}(a),(c),
\[
|\cM|+|\tilde\cM|\ \le\ \frac{4}{\cos(\pi/4+\amax)}\sum_{e\in\mathcal Q_b\cup\mathcal Q_c}w_e\ \le\ \frac{4}{\cos(\pi/4+\amax)}\,S_K\,\delta W .
\]
Numerically, $4/\cos(\pi/4+1.5\cdot10^{-6})=5.65686273\ldots\le5.656863$ (rounded up), and $5.656863\cdot56.2501=318.1991\ldots\le318.20$, $5.656863\cdot81.2502=459.6212\ldots\le459.7$.
\end{proof}

The proof shows more precisely that, for each square $Y$, the losers at entry candidates on $Y$ whose two sources form the pair $e$ have measure at most $w_e/\cos(\varphi_Y-\varphi_{X_1})$; and the two last gap directions of two merging paths are always different ($\theta_e>0$, Lemma~\ref{lem:P5-two}).

\subsection{Gap overlap}\label{ssec:P5-G}

\begin{proof}[Proof of Theorem~\ref{thm:P5}\,\textup{(G)}]
Fix $y\in(0,\hmax)$ and consider the floor flow. Only paths alive at $y$, that is, with termination height $>y$, enter $A_{\rm gap}(y)$. In particular a path whose cumulative gap reaches $\delta$ exactly at a point of height $y$ (a death point at height $y$) is not counted at $y$; this agrees with the gap case of Definition~\ref{def:P3-reach}, whose delays are $<t^*$, and with the definition of $M_g$, which counts only points of the open parts of gap segments.

\emph{Density.} By Lemma~\ref{lem:flows-density}(c), $A_{\rm gap}(y)$ is the union of finitely many points and finitely many disjoint open intervals $J_c$; on each $J_c$ the point $P_y(x)$ lies on the open part of the gap segment of a fixed ray step, $X_y$ is affine with slope $\kappa_c\ge1$, and
\[
\rho_y=\sum_c\kappa_c^{-1}\,\mathbf 1_{X_y(J_c)}\ \le\ \sum_c\mathbf 1_{X_y(J_c)}\qquad\text{almost everywhere.}
\]
For $z\in\mathrm{Wa}_y$, $\sum_c\mathbf 1_{X_y(J_c)}(z)$ is the number of $x\in\bigcup_cJ_c$ with $P_y(x)=z$, because $X_y$ is injective on each $J_c$ and the intervals $J_c$ are disjoint. For each such $x$ the point $z$ lies on the open part of a gap segment of the trajectory of $\pi_x$, that is, $\pi_x$ reaches $z$ in the gap case; so the sum is at most $M_g(z)\le2$ (Lemma~\ref{lem:P5-atmost2}). If it equals $2$, then $z$ is a double point and lies in $\operatorname{Rect}_e$ for its valley pair $e\in\mathcal Q_b$. Hence, almost everywhere on $\mathrm{Wa}_y$,
\[
(\rho_y-1)_+\ \le\ \sum_{e\in\mathcal Q_b}\mathbf 1_{\operatorname{Rect}_e}.
\]

\emph{Length of one intersection.} Let $e=\{X_1,X_2\}\in\mathcal Q_b$, with $v_1$ the upper left vertex of $X_1$. Along $\ell_y$ (direction $e_1$) the function $(z-v_1)\cdot u_{X_1}$ changes at the rate $e_1\cdot u_{X_1}=\cos\varphi_{X_1}\ge\cos\amax$. So $\ell_y$ meets the strip $\{0<(z-v_1)\cdot u_{X_1}<w_e\}\supset\operatorname{Rect}_e$ in a segment of length at most $w_e/\cos\amax$, and
\[
\Ovgap^{\Fz}(y)=\int_{\mathrm{Wa}_y}(\rho_y-1)_+\ \le\ \sum_{e\in\mathcal Q_b}\mathcal H^1(\ell_y\cap\operatorname{Rect}_e)\ \le\ \frac1{\cos\amax}\sum_{e\in\mathcal Q_b}w_e .
\]
This holds for every $y\in(0,\hmax)$, including the heights of horizontal sides of squares and the exceptional heights of Definition~\ref{def:flows-exc}: the waste does not contain the points of the closed squares, and the argument uses no property of $y$. The bound does not depend on $y$, so it bounds $\sup_y\Ovgap^{\Fz}(y)$. Applying this to $\rho\cP$ gives
\[
\sup_y\Ovgap^{\tilde F^0}(y)\ \le\ \frac1{\cos\amax}\sum_{e\in\mathcal Q_c}w_e ,
\]
and by Lemma~\ref{lem:P5-pairsum}(a),(c),
\[
\sup_y\Ovgap^{\Fz}(y)+\sup_y\Ovgap^{\tilde F^0}(y)\ \le\ \frac1{\cos\amax}\sum_{e\in\mathcal Q_b\cup\mathcal Q_c}w_e\ \le\ \frac{S_K}{\cos\amax}\,\delta W .
\]
Numerically (rounded up),
\[
S_9/\cos(1.5\cdot10^{-6})<56.2501,\qquad S_{13}/\cos(1.5\cdot10^{-6})<81.2502\le81.26 . \qedhere
\]
\end{proof}

\subsection{The bound on \texorpdfstring{$\Lam_0$}{Lambda0} and the wall term}\label{ssec:P5-Lambda}

\begin{proof}[Proof of Corollary~\ref{cor:Lambda0}]
By definition,
\[
\Lam_0+\tilde\Lam_0=\big(|\cD|+|\tilde\cD|\big)+\big(|\cM|+|\tilde\cM|\big)+\big(L_E+\tilde L_E\big)+\Big(\sup_y\Ovgap^{\Fz}(y)+\sup_y\Ovgap^{\tilde F^0}(y)\Big).
\]
With $K=9$, Theorem~\ref{thm:P5} and Proposition~\ref{prop:P3} give
\begin{align*}
\Lam_0+\tilde\Lam_0&\le(1+56.2501\,\delta^2)\frac W\delta+(318.20+754.8+56.2501)\,\delta W\\
&=\big(1+1185.5002\,\delta^2\big)\frac W\delta\ \le\ \big(1+1185.6\,\delta^2\big)\frac W\delta .
\end{align*}
With $K=13$ the same computation gives $81.26+459.7+1090.3+81.26=1712.52\le1712.6$. Since $\delta\le10^{-5}$, $1185.6\,\delta^2\le1.1856\cdot10^{-7}$ and $1712.6\,\delta^2\le1.7126\cdot10^{-7}$, both less than $2\cdot10^{-7}$.
\end{proof}

\begin{remark}[the wall term]\label{rem:P5-wall}
The quantity $\Lam_0$ collects the losses of the master flow that are bounded in terms of $W$. The loss term $\Lam^{F_s}(y)$ of a scale flow $F_s$ (Section~\ref{sec:flows}) also contains the wall loss of $F_s$, which is not bounded in terms of $W$. Proposition~\ref{prop:wall} bounds $\Lam^{F_s}(y)$, for $s\in[y_0,y_1]$ and $0<y<s+2$, by $\Lam_0+2(s+2)\tan\alpha(s)$; the wall term $2(s+2)\tan\alpha(s)$ is not part of $\Lam_0$, and it is treated separately in Section~\ref{sec:P8}.
\end{remark}

\begin{remark}[dependence on {[R]}]\label{rem:P5-dependence}
The bounds of Theorem~\ref{thm:P5} with $K=9$, and hence the constant $C_\Lam=1+1185.6\,\delta^2$ of Corollary~\ref{cor:Lambda0}, use [R, Proposition~\ref{imp:P4.5}] and [R, Lemma~\ref{imp:L4.10}] (through Lemma~\ref{lem:P3-mult}), together with [R, Lemma~\ref{imp:L4.8}(a)] for $L_E$ (Proposition~\ref{prop:P3}); [R, Lemma~\ref{imp:L4.10}] is computer-assisted in~[R] (Remark~\ref{rem:imp-dep}). The bounds with $K=13$, and $C^{(13)}_\Lam=1+1712.6\,\delta^2$, use [R, Lemma~\ref{imp:L4.9}] in place of [R, Lemma~\ref{imp:L4.10}] and do not depend on the computer-assisted lemma.
\end{remark}

\begin{remark}[degenerate cases]\label{rem:P5-degenerate}
The lemmas of Section~\ref{ssec:P5-two} are statements about single points, so contacts at vertices and rays running along a side of a square (which occur for finitely many $x$, by~\ref{it:P3-S2}) do not affect them. Ties between the event parameters, which may occur for sets of $x$ of positive measure, are decided by the priority rule, and the delay bound of Lemma~\ref{lem:P3-reach}(a) holds also at ties. A square resting on the floor gives contacts with delay $0$ from the floor, which Lemma~\ref{lem:P5-floor} allows (it only uses delays $<\delta$); a square source always has a positive delay. Contacts at height exactly $\hmax$ are processed, and Lemma~\ref{lem:P3-separation} uses only $q_y\le\hmax$. Ties of the cumulative gap at an entry candidate are decided by R1 through the floor abscissa; they do not enter the bounds above, which only use that each entry candidate has exactly one winner. Squares with equal centre heights cause no difficulty (Lemma~\ref{lem:DAG}). Squares of phase $\pi/4$ are never sources ($\amax<\pi/4$). Two distinct squares with the same phase are never the two sources of one point (Lemma~\ref{lem:P5-two}). No ray step starts on a side wall of the container (Lemma~\ref{lem:flows-traj}(b)). The bound (G) holds at every height, including heights of horizontal sides of squares.
\end{remark}

\section{P6: the path version of \texorpdfstring{[R, Lemma~\ref{imp:L4.14}]}{[R, Lemma 4.14]}}\label{sec:P6}

This section establishes Theorem~\ref{thm:P6}, the analogue of [R, Lemma~\ref{imp:L4.14}] for the master flows of
Section~\ref{sec:flows}, and its scale form, Corollary~\ref{cor:P6scale}, which is the form used in
Section~\ref{sec:P8}. In [R, Lemma~\ref{imp:L4.14}] a column is a vertical segment with vertical waste less than
$\delta$, and two squares that are consecutive on it are consecutive on the whole vertical line. The paths of $\Fz$
follow instead the directions $n_X$ of the squares they pass, so two consecutively entered squares need not be
consecutive on vertical lines. The transfer to vertical lines is Lemma~\ref{lem:P6V} (the vertical comparison P6.V).
It rests on Lemma~\ref{lem:P6-hull}: for a pair of consecutively entered squares, the whole region swept by the gap
segments over the convex hull of the exit points is waste. The rest of the argument follows the proof of
[R, Lemma~\ref{imp:L4.14}]: a triangle inequality along each path, the slit inequality, the waste in the rectangles
of [R, Proposition~\ref{imp:P4.5}] and [R, Lemma~\ref{imp:L4.8}(a)], the multiplicity bound of
Proposition~\ref{imp:pairs}, and the Cauchy--Schwarz inequality.

\subsection{Setting, notation, and facts used from Section~\ref{sec:flows}}\label{ssec:P6-setting}

Throughout this section $k\ge4$ is an integer, $\cP$ is a closed packing of $[0,k]^2$ with waste area $W=k^2-N$, the
parameters satisfy \textup{(C0)} and \textup{(C1)} of Definition~\ref{def:flows-constraints}, so that
\begin{equation}\label{eq:P6-params}
\amax=\alpha(y_0)\ \le\ \bbar=1.5\cdot10^{-6},\qquad \hmax=\tfrac k2-1,
\end{equation}
$0<\delta\le10^{-5}$, and $\Fz=F(\amax,\delta,\hmax)$ is the master flow of Definition~\ref{def:F0}. For $x\in(0,k)$,
$\pi^0_x$ denotes the path of $\Fz$ that starts at $(x,0)$. The constants $A$, $A'$, $A_{13}$, $A'_{13}$ are those of
Definition~\ref{def:flows-consts}:
\[
A=\frac{4\sqrt{K_w^*/Q_*}}{2-3\cdot10^{-6}},\qquad A'=\frac{A}{\cos\bbar}+10^{-8},\qquad
A_{13}=\frac{4\sqrt{K_w/Q_*}}{2-3\cdot10^{-6}},\qquad A'_{13}=\frac{A_{13}}{\cos\bbar}+10^{-8},
\]
with $Q_*=0.32$, $K_w^*=9$, $K_w=13$; $A$ is the constant of [R, Lemma~\ref{imp:L4.13}].

\emph{Ceiling versions.} The ceiling master flow is $\tilde F^0=\rho\circ\Fz[\rho\cP]\circ\rho$, $\rho(x,y)=(x,k-y)$
(Definition~\ref{def:flows-ceiling}). By Lemma~\ref{lem:flows-refl}, $\rho\cP$ is a closed packing with the same
waste area $W$, and $\rho$ preserves inclinations, the function $\theta$ below, distances, vertical lines and lengths,
and maps the waste of $\rho\cP$ onto the waste of $\cP$. Every statement below that is established for $\Fz$ and an
arbitrary closed packing therefore applies to $\rho\cP$, and transported by $\rho$ it gives the corresponding
statement for $\tilde F^0$, its \emph{ceiling version}. Ceiling objects carry a tilde.

\emph{Notation.} We use the phases $\varphi_S\in(-\frac\pi4,\frac\pi4]$ (with the $45^\circ$ convention), the
frames $u_S,n_S$, the centres $c_S$, the sides $\Top(S)$, $\Bot(S)$ and their relative interiors, and the floor
$X_0$ (with $\varphi_{X_0}=0$, $u_{X_0}=e_1$, $n_{X_0}=e_2$, ``top side'' $[0,k]\times\{0\}$) of
Definition~\ref{def:flows-frames}, and the identities \eqref{eq:flows-frame}. Thus
$\Top(S)=\{c_S+\frac12n_S+\xi u_S:|\xi|\le\frac12\}$, $\Bot(S)=\{c_S-\frac12n_S+\xi u_S:|\xi|\le\frac12\}$, the
relative interiors are given by $|\xi|<\frac12$, and $a(S)=|\varphi_S|$. For points $\xi\in\R^2$ we call
$\xi\cdot e_1$ the abscissa and $\xi\cdot e_2$ the height of $\xi$. For squares $S_1,S_2$, $\theta(S_1,S_2)$ is the
distance of $\varphi_{S_1}$ and $\varphi_{S_2}$ in $\R/\frac\pi2\Z$, and, as in [R], $\theta(X_0,S):=a(S)=\theta(S,S_0)$,
which is the same distance with $\varphi_{X_0}=0$. Hence $\theta$ satisfies the triangle inequality on
$\cP\cup\{X_0\}$, and $\theta(S_1,S_2)\le|\varphi_{S_2}-\varphi_{S_1}|$.

Let $S$ be a square with inclination $a=a(S)$ and lowest point at height $v_S\ge0$. Its vertical extent is
$[v_S,\,v_S+\cos a+\sin a]\subset[0,k]$ by [R, Lemma~\ref{imp:L3.1}(a)], and
$c_S\cdot e_2=v_S+\frac12(\cos a+\sin a)$ because $S$ is symmetric about its centre. By \eqref{eq:flows-frame}, for
$|\xi|\le\frac12$,
\[
(c_S\pm\tfrac12n_S+\xi u_S)\cdot e_2=c_S\cdot e_2\pm\tfrac12\cos\varphi_S+\xi\sin\varphi_S,\qquad
|\xi\sin\varphi_S|\le\tfrac12\sin a ;
\]
hence
\begin{equation}\label{eq:P6-heights}
\xi\cdot e_2\ \ge\ v_S+\cos a>0\quad(\xi\in\Top(S)),\qquad
\xi\cdot e_2\ \le\ v_S+\sin a\ \le\ k-\cos a<k\quad(\xi\in\Bot(S)).
\end{equation}

\emph{Entries.} Following Definition~\ref{def:F0}, the \emph{R1-passed entries} of $\pi^0_x$ are the entry candidates
$(Y,q)$ of $\pi^0_x$ (Definition~\ref{def:flow}: $q\in\relint\Bot(Y)$ is a contact point of a processed ray step)
at which $x$ is the winner (Definition~\ref{def:R1}). We list them as $(Z_1,q_1),\dots,(Z_J,q_J)$, $J=J(x)\ge0$, call
them simply the \emph{entries} of $\pi^0_x$, and say that $\pi^0_x$ \emph{enters} $Z_j$ at $q_j$. For $j<J$ the path
passes $Z_j$ (Definition~\ref{def:R2}): it traverses the pass segment $[q_j,q_j+n_{Z_j}]$ and leaves $Z_j$ at the
\emph{exit point} $q_j+n_{Z_j}\in\relint\Top(Z_j)$; $Z_J$ is passed or T-terminates $\pi^0_x$. Contacts with a left
or right side, or with a vertex, are end collisions and never entries.

We use the following facts from Section~\ref{sec:flows}.
\begin{enumerate}[label=(F\arabic*)]
\item\label{it:P6-F1} (Lemma~\ref{lem:flows-R1}(c).) At most one path of $\Fz$ passes, or is T-terminated at, a
given entry point $(Y,q)$.
\item\label{it:P6-F2} (Lemma~\ref{lem:DAG}.) The squares $Z_1,\dots,Z_J$ entered by one path have strictly increasing
centre heights; in particular they are pairwise distinct.
\item\label{it:P6-F3} (Lemmas~\ref{lem:structure}(a) and~\ref{lem:expansion}(a).) There is a finite set
$\Sigma\subset[0,k]$ such that on each component (\emph{piece}) of $[0,k]\setminus\Sigma$ the combinatorial type of
$\pi^0_x$ is constant; in particular its sequence of entries is constant, and so is, for each entry, whether the
path passes or is T-terminated. On a piece, the start point $p_j(x)$ of the $j$-th ray step is an affine function of
$x$ with $p_j'(x)=\lambda_ju_{X_j}$, where $X_j$ is the source of that step ($X_0$ the floor) and
$\lambda_j=\prod_{i=1}^j\sec(\varphi_{X_i}-\varphi_{X_{i-1}})\ge1$. Each factor is positive, because
$|\varphi_{X_i}-\varphi_{X_{i-1}}|\le\frac\pi4+|\varphi_{X_{i-1}}|<\frac\pi2$; so the slope $\lambda_j$ is positive.
\end{enumerate}
The next lemma recalls what one ray step between two consecutive entries gives.

\begin{lemma}[one step]\label{lem:P6-step}
Let $x\in(0,k)$, let $Z_1,\dots,Z_J$ be the squares entered by $\pi^0_x$, in order, and put $Z_0:=X_0$. Fix
$1\le i\le J$ and write $X:=Z_{i-1}$, $Y:=Z_i$. Let $p_x$ be the point where $\pi^0_x$ leaves $X$ \textup{(}its exit
point from $X$ if $i\ge2$, and $p_x=(x,0)$ if $i=1$\textup{)}, and let $q_x:=q_i$ be its entry point into $Y$. Then:
\begin{enumerate}[label=(\alph*)]
\item $q_x=p_x+t_xn_X$ with $0\le t_x<\delta$, and $t_x>0$ if $X$ is a square;
\item $q_x\cdot e_2\le\hmax$; if $X$ is a square, its entry point $q_{i-1}$ also has height $\le\hmax$;
\item the open segment $\{p_x+s\,n_X:\ 0<s<t_x\}$ lies in the waste; if $X=X_0$, then the point $p_x=(x,0)$ lies in
no square other than $Y$;
\item if $X$ is a square, then $a(X)<\amax\le\bbar$; in all cases $|\varphi_X|<\bbar$.
\end{enumerate}
\end{lemma}

\begin{proof}
If $i\ge2$, then $\pi^0_x$ has an entry after the one into $X$, so it passes $X$, and $a(X)<\amax$ by
Definition~\ref{def:R2}; this gives (d), as $\varphi_{X_0}=0<\bbar$. After leaving $X$ at $p_x$, the path performs
a ray step $r(t)=p_x+tn_X$ with source $X$; since it has a further entry, this step ends with a processed contact at
an entry candidate at which $x$ wins, and this is the entry into $Y$. So $q_x=r(t_S)$, and we put $t_x:=t_S=t^*$. At
a processed contact the cumulative gap is $<\delta$ (Lemma~\ref{lem:flows-traj}(f)), so $t_x<\delta$; and $t_S>0$
if the source is a square (Lemma~\ref{lem:flows-traj}(b)). This gives (a). A contact is processed only if
$t_S=t^*\le t_H$ (Definition~\ref{def:flow}); as the height increases along the ray at the rate
$n_X\cdot e_2=\cos\varphi_X>0$, the contact point then has height at most
$p_x\cdot e_2+t_H\cos\varphi_X=h_F=\hmax$. The entry into $X$ (if $X$ is a square) is also a processed contact, so (b)
follows. The open part of the gap segment
$[p_x,q_x]$ lies in the waste by Lemma~\ref{lem:flows-traj}(c). If $X=X_0$ and $(x,0)$ lies in a square $Z$, then
$t_S=0$, so $q_x=(x,0)$, and $Z=Y$ because $q_x$ lies in exactly one square (Lemma~\ref{lem:flows-traj}(d)). This
gives (c).
\end{proof}

\begin{definition}[consecutive entry pairs]\label{def:P6-pairs}
In the notation of Lemma~\ref{lem:P6-step}, the ordered pairs $(Z_{i-1},Z_i)$, $1\le i\le J$, are the
\emph{consecutive entry pairs} of $\pi^0_x$; $Z_{i-1}$ is the \emph{first} and $Z_i$ the \emph{second member}. For
an ordered pair $e=(X,Y)$ with $X\in\cP\cup\{X_0\}$ and $Y\in\cP$ put
\[
P_e:=\{x\in(0,k):\ e\ \text{is a consecutive entry pair of}\ \pi^0_x\},\qquad m_e:=|P_e|,
\]
and let $\mathcal C:=\{e:\ P_e\neq\emptyset\}$. For $e=(X,Y)$ put $\theta_e:=\theta(X,Y)$ (so $\theta_e=a(Y)$ if
$X=X_0$) and $\theta'_e:=\min(\theta_e,\bbar)$. The ceiling objects $\tilde{\mathcal C}$ and, for
$e\in\tilde{\mathcal C}$, $P_e$, $m_e$, $\theta_e$, $\theta'_e$ are defined in the same way from $\tilde F^0$; the
first member of $e\in\tilde{\mathcal C}$ is a square or the ceiling pseudo-square $\tilde X_0=\rho X_0$, whose
``side'' is $[0,k]\times\{k\}$, and $\theta(\tilde X_0,Y):=a(Y)$.
\end{definition}

Since $\cP$ is finite, $\mathcal C$ and $\tilde{\mathcal C}$ are finite. By~\ref{it:P6-F2} the consecutive entry
pairs of one path are pairwise distinct, as their second members are.

\begin{theorem}[path version of {[R, Lemma~\ref{imp:L4.14}]}]\label{thm:P6}
Let $k\ge4$ be an integer, $\cP$ a closed packing of $[0,k]^2$ with waste area $W$, $0<\delta\le10^{-5}$, and
assume \textup{(C0)} and \textup{(C1)}. Let $\cX\subset(0,k)$ be a Lebesgue measurable set and $\zeta:\cX\to[0,\bbar]$
a measurable function such that for every $x\in\cX$ the path $\pi^0_x$ enters a square $Z_x$ \textup{(}that is, one
of its R1-passed entries is an entry into $Z_x$; $Z_x$ may be passed or T-terminate the path\textup{)} with
$a(Z_x)\ge\zeta(x)$. Let $\cX'\subset(0,k)$ and $\zeta':\cX'\to[0,\bbar]$ satisfy the same with $\tilde F^0$ in place
of $\Fz$.
\begin{enumerate}[label=(\alph*)]
\item $\displaystyle\int_{\cX}\zeta(x)\,dx+\int_{\cX'}\zeta'(x)\,dx\ \le\ A'W$, where $A'=\dfrac{A}{\cos\bbar}+10^{-8}<10.6068$.
\item More strongly, independently of $\cX,\zeta,\cX',\zeta'$,
\[
\sum_{e\in\mathcal C\cup\tilde{\mathcal C}}\theta'_e\,m_e\ \le\ \Big(\frac{A}{\cos\bbar}+8.5\cdot10^{-10}\Big)W .
\]
\item Without [R, Lemma~\ref{imp:L4.10}], using [R, Lemma~\ref{imp:L4.9}] in its place, \textup{(a)} and \textup{(b)}
hold with $A$ replaced by $A_{13}<12.7476$, $A'$ by $A'_{13}<12.7476$, and $8.5\cdot10^{-10}$ by $1.22\cdot10^{-9}$.
\end{enumerate}
All numerical bounds are rounded up; $A<10.6067$.
\end{theorem}

\begin{remark}[dependence on {[R, Lemma~\ref{imp:L4.10}]}]\label{rem:P6-dependency}
Parts (a) and (b) use the multiplicity bound $K_w^*=9$ of Proposition~\ref{imp:pairs}, which rests on the proof of
[R, Lemma~\ref{imp:L4.10}] and [R, Remark~\ref{imp:R4.11}]; the proof of [R, Lemma~\ref{imp:L4.10}] is
computer-assisted (Remark~\ref{rem:imp-dep}). Part (c) uses instead the bound $K_w=13$ of
Proposition~\ref{imp:pairs}, which rests on the analytic [R, Lemma~\ref{imp:L4.9}], and does not use
[R, Lemma~\ref{imp:L4.10}]. Apart from this, the proofs of (a)--(c) are the same. Through
Corollary~\ref{cor:P6scale}, part (a) enters the proof of Theorem~\ref{thm:A} and part (c) the proof of
Theorem~\ref{thm:A13} (Section~\ref{sec:assembly}).
\end{remark}

The rest of the section contains the density lemma (\S\ref{ssec:P6-density}), the hull slit (\S\ref{ssec:P6-hull}),
the vertical comparison (\S\ref{ssec:P6-vertical}), the disjointness of the vertical slits (\S\ref{ssec:P6-disjoint}),
the slit inequality, the pair waste and the multiplicity (\S\ref{ssec:P6-pairs}), the proof of Theorem~\ref{thm:P6}
(\S\ref{ssec:P6-proof}), the scale form (\S\ref{ssec:P6-scale}), and remarks on the rules and on degenerate cases
(\S\ref{ssec:P6-remarks}).

\subsection{Density}\label{ssec:P6-density}

For $e=(X,Y)\in\mathcal C$ put $o_X:=c_X+\frac12n_X$ if $X$ is a square and $o_{X_0}:=(0,0)$. For $x\in P_e$ let
$p_x$ be the point where $\pi^0_x$ leaves $X$ (Lemma~\ref{lem:P6-step}), and define the \emph{exit coordinate}
$w_X(x)$ by $p_x=o_X+w_X(x)\,u_X$. Thus $|w_X(x)|<\frac12$ if $X$ is a square, and $w_{X_0}(x)=x$. Put
$I_e:=w_X(P_e)$.

\begin{lemma}[density]\label{lem:P6-density}
Let $e=(X,Y)\in\mathcal C$. Then $P_e$ and $I_e$ are Borel sets, $w_X$ is injective on $P_e$, and $|I_e|\ge m_e$.
\end{lemma}

\begin{proof}
\emph{Pieces.} By~\ref{it:P6-F3} the sequence of entries of $\pi^0_x$ is constant on each piece, so each piece is
either contained in $P_e$ or disjoint from it. Hence $P_e$ is the union of the pieces $I\subset P_e$ and of a subset
of the finite set $\Sigma$; it is a Borel set, and $m_e=\sum_{I\subset P_e}|I|$. If $X=X_0$, then $w_X$ is the
identity and $I_e=P_e$, and there is nothing more to show. Let $X$ be a square and $I\subset P_e$ a piece. On $I$,
$X$ is the source $X_j$ of a fixed ray step $j$, and $p_x=p_j(x)$; by~\ref{it:P6-F3},
$\frac{d}{dx}w_X(x)=p_j'(x)\cdot u_X=\lambda_j\ge1$. So $w_X$ is affine on $I$ with positive slope $\lambda_j$, and
$w_X(I)$ is an interval of length $\lambda_j|I|\ge|I|$.

\emph{Injectivity.} Let $x\neq x'$ be points of $P_e$ with $w_X(x)=w_X(x')$. Both paths pass $X$ and leave it at the
same exit point; since the exit point is the entry point plus $n_X$, they entered $X$ at the same point
$q:=c_X-\frac12n_X+w_X(x)u_X$. Two paths would then pass $X$ from the same entry point $(X,q)$, contradicting
\ref{it:P6-F1}. So $w_X$ is injective on $P_e$.

\emph{Measure.} By injectivity, the intervals $w_X(I)$, $I\subset P_e$ a piece, are pairwise disjoint, and $I_e$ is
their union together with the finitely many points $w_X(P_e\cap\Sigma)$. Hence $I_e$ is a Borel set and
$|I_e|=\sum_I\lambda_{j(I)}|I|\ge\sum_I|I|=m_e$.
\end{proof}

\subsection{The frame of a pair and the hull slit}\label{ssec:P6-hull}

Fix $e=(X,Y)\in\mathcal C$ and write $u:=u_X$, $n:=n_X$, $\varphi:=\varphi_X$, $o:=o_X$. Every point of the plane is
$o+wu+sn$ for a unique $(w,s)\in\R^2$, and $(w,s)\mapsto o+wu+sn$ is an isometry. (In this subsection and in Lemma~\ref{lem:P6-step}(c), $s$ denotes a coordinate or a length
parameter, not a scale; the lines $\ell_w$ below are lines in the direction $n_X$, not horizontal lines $\ell_y$.) Put $\ell_w:=\{o+wu+sn:\ s\in\R\}$
and let $L_Y:=\{\xi:\ (\xi-c_Y)\cdot n_Y=-\frac12\}$ be the line of $\Bot(Y)$. By Lemma~\ref{lem:P6-step}(d),
$|\varphi|<\bbar$, so $|\varphi_Y-\varphi|<\frac\pi4+\bbar<\frac\pi2$ and, by \eqref{eq:flows-frame},
$n\cdot n_Y=\cos(\varphi_Y-\varphi)>0$. Hence $\ell_w$ meets $L_Y$ in exactly one point $o+wu+t(w)\,n$, where
\begin{equation}\label{eq:P6-t}
t(w)=t_0+\kappa_e w,\qquad \kappa_e:=-\frac{u\cdot n_Y}{n\cdot n_Y}=\tan(\varphi_Y-\varphi),\qquad
|\kappa_e|<\tan\big(\tfrac\pi4+\bbar\big)<1.00001,
\end{equation}
for a constant $t_0$ (the last bound is rounded up). Along $\ell_w$, in the direction $n$, the function
$\xi\mapsto(\xi-c_Y)\cdot n_Y$ increases at the rate $n\cdot n_Y>0$; since
$Y\subset\{\xi:\ (\xi-c_Y)\cdot n_Y\ge-\frac12\}$,
\begin{equation}\label{eq:P6-belowY}
\text{no point $o+wu+sn$ with $s<t(w)$ belongs to $Y$.}
\end{equation}
For $x\in P_e$, Lemma~\ref{lem:P6-step}(a) gives $q_x=o+w_X(x)u+t_xn\in\Bot(Y)\subset L_Y$; hence
\begin{equation}\label{eq:P6-tx}
t(w_X(x))=t_x\in[0,\delta),\qquad\text{and}\quad t(w_X(x))>0\ \text{ if $X$ is a square.}
\end{equation}
Let $w_1:=\inf I_e$ and $w_2:=\sup I_e$. Then $[w_1,w_2]\subset[-\frac12,\frac12]$ if $X$ is a square, and
$[w_1,w_2]\subset[0,k]$ if $X=X_0$.

\begin{lemma}[the hull slit is waste]\label{lem:P6-hull}
Let $e=(X,Y)\in\mathcal C$.
\begin{enumerate}[label=(\roman*)]
\item For every $w\in[w_1,w_2]$: $o+wu\in\Top(X)$ \textup{(}$o+wu=(w,0)\in[0,k]\times\{0\}$ if $X=X_0$\textup{)},
$o+wu+t(w)n\in\Bot(Y)$, and $0\le t(w)\le\delta$. For every $w\in(w_1,w_2)$ moreover $o+wu\in\relint\Top(X)$
\textup{(}$0<w<k$ if $X=X_0$\textup{)}, $o+wu+t(w)n\in\relint\Bot(Y)$, $t(w)<\delta$, and $t(w)>0$ if $X$ is a square.
\item The set
\[
\mathcal O_e:=\{o+wu+sn:\ w_1<w<w_2,\ 0<s<t(w)\}
\]
is open and contained in the waste.
\end{enumerate}
\end{lemma}

\begin{proof}
(i) Let $w\in(w_1,w_2)$. By the definition of $w_1$ and $w_2$ there are $p_1,p_2\in I_e$ with $p_1<w<p_2$; write
$w=(1-\mu)p_1+\mu p_2$ with $0<\mu<1$, and $p_j=w_X(x_j)$ with $x_j\in P_e$. By Lemma~\ref{lem:P6-step}(a) and
\eqref{eq:P6-tx}, $o+p_ju=p_{x_j}\in\relint\Top(X)$ (respectively $0<p_j<k$ if $X=X_0$),
$o+p_ju+t(p_j)n=q_{x_j}\in\relint\Bot(Y)$, and $0\le t(p_j)<\delta$, with $t(p_j)>0$ if $X$ is a square. The maps
$w'\mapsto o+w'u$ and $w'\mapsto o+w'u+t(w')n$ are affine, and $\relint\Top(X)$, $\relint\Bot(Y)$ and $(0,k)$ are
convex. So the images of $w$ lie in $\relint\Top(X)$ (respectively $0<w<k$) and in $\relint\Bot(Y)$, and
$t(w)=(1-\mu)t(p_1)+\mu t(p_2)$ lies in $[0,\delta)$, and in $(0,\delta)$ if $X$ is a square. For $w\in I_e$ the claims
of the first sentence hold by Lemma~\ref{lem:P6-step}(a) and \eqref{eq:P6-tx}. If $w_1=w_2$, then $I_e=\{w_1\}$ and
we are done. If $w_1<w_2$, each endpoint $w_j$ is a limit of points of $(w_1,w_2)$; since $\Top(X)$,
$[0,k]\times\{0\}$ and $\Bot(Y)$ are closed and $t$ is continuous, we obtain $o+w_ju\in\Top(X)$ (respectively
$\in[0,k]\times\{0\}$), $o+w_ju+t(w_j)n\in\Bot(Y)$ and $0\le t(w_j)\le\delta$.

(ii) The set $\{(w,s):\ w_1<w<w_2,\ 0<s<t(w)\}$ is open since $t$ is continuous, and $\mathcal O_e$ is its image under
an isometry, so $\mathcal O_e$ is open. Each point $o+wu+sn$ of $\mathcal O_e$ lies on the segment from
$o+wu\in[0,k]^2$ to $o+wu+t(w)n\in\Bot(Y)\subset[0,k]^2$ (by (i)), so $\mathcal O_e\subset[0,k]^2$; being open,
$\mathcal O_e\subset(0,k)^2$. If $X$ is a square, then $X\subset\{\xi:\ (\xi-c_X)\cdot n\le\frac12\}$ while
$(o+wu+sn-c_X)\cdot n=\frac12+s>\frac12$, so $\mathcal O_e\cap X=\emptyset$. By \eqref{eq:P6-belowY},
$\mathcal O_e\cap Y=\emptyset$.

Suppose now that a square $Z$, with $Z\neq Y$ and $Z\neq X$, contains a point $v=o+wu+sn\in\mathcal O_e$, so
$w_1<w<w_2$ and $0<s<t(w)$.

\emph{Claim 1. For every $p\in[w_1,w_2]$ and every point $z=o+pu+\sigma n$ of $Z$ we have $\sigma<t(p)$; moreover
$\sigma>0$ if $X$ is a square, and $\sigma\ge0$ if $X=X_0$.}
Let $X$ be a square. The disjoint compact convex sets $X$ and $Z$ are strictly separated by a line: there are
$\nu\in\R^2$ and $c\in\R$ with $\nu\cdot\xi<c$ for $\xi\in X$ and $\nu\cdot\xi>c$ for $\xi\in Z$. By (i), $o+wu\in X$,
and $v\in Z$, so $s\,(\nu\cdot n)=\nu\cdot v-\nu\cdot(o+wu)>0$, whence $\nu\cdot n>0$. By (i), $o+pu\in X$ for every
$p\in[w_1,w_2]$, including the endpoints (this is where the closed hull is needed), so
$\sigma\,(\nu\cdot n)=\nu\cdot z-\nu\cdot(o+pu)>0$, and $\sigma>0$. If $X=X_0$, then $o=(0,0)$, $n=e_2$, and
$\sigma=z\cdot e_2\ge0$ because $Z\subset[0,k]^2$. For the upper bound, separate $Z$ and $Y$ in the same way:
$\nu'\cdot\xi<c'$ on $Z$ and $\nu'\cdot\xi>c'$ on $Y$. By (i), $o+wu+t(w)n\in Y$; as $(o+wu+t(w)n)-v=(t(w)-s)\,n$
with $t(w)-s>0$, we get $\nu'\cdot n>0$. By (i), $o+pu+t(p)n\in Y$ for every $p\in[w_1,w_2]$, so
$(t(p)-\sigma)(\nu'\cdot n)=\nu'\cdot(o+pu+t(p)n)-\nu'\cdot z>0$, and $\sigma<t(p)$.

\emph{Claim 2. $Z\cap\ell_p=\emptyset$ for every $p\in I_e$; consequently, if $p_1,p_2\in I_e$ and $p_1<w<p_2$, then
every point of $Z$ has the form $o+w'u+\sigma n$ with $p_1<w'<p_2$.}
Let $p=w_X(x)$ with $x\in P_e$, and let $z=o+pu+\sigma n\in Z$. By Claim~1 and \eqref{eq:P6-tx}, $\sigma<t(p)=t_x$. If
$\sigma>0$, then $z=p_x+\sigma n$ lies on the open segment of Lemma~\ref{lem:P6-step}(c), which lies in the waste;
this is impossible for a point of $Z$. If $\sigma\le0$, then by Claim~1 $X=X_0$ and $\sigma=0$, so $z=(p,0)=p_x$ lies
in the square $Z\neq Y$, which contradicts Lemma~\ref{lem:P6-step}(c). This proves the first assertion. The function
$\xi\mapsto(\xi-o)\cdot u$ is continuous on the connected set $Z$, so its image is an interval; it contains $w$ (the
value at $v$) and, by the first assertion, neither $p_1$ nor $p_2$. Hence it lies in $(p_1,p_2)$.

\emph{Claim 3: a contradiction.} Choose $p_1,p_2\in I_e$ with $p_1<w<p_2$ (possible since $w_1<w<w_2$). By Claim~2,
every $z\in Z$ is $o+w'u+\sigma n$ with $w'\in(p_1,p_2)\subset(w_1,w_2)$, and then Claim~1 and (i) give
$0\le\sigma<t(w')<\delta$. So the set $\{(z-o)\cdot n:\ z\in Z\}$ is contained in $[0,\delta)$, an interval of length
$\delta<1$. But for every unit vector $\mathbf m$ the set $\{\mathbf m\cdot z:\ z\in Z\}$ is an interval of length
$|\mathbf m\cdot u_Z|+|\mathbf m\cdot n_Z|\ge(\mathbf m\cdot u_Z)^2+(\mathbf m\cdot n_Z)^2=1$: every orthogonal
projection of a unit square to a line has length at least $1$. With $\mathbf m=n$ this is a contradiction. Hence no
square meets $\mathcal O_e$, and $\mathcal O_e$ lies in the waste.
\end{proof}

The lemma concerns the geometry of the hull only: points of $(w_1,w_2)\setminus I_e$ need not be exit coordinates of
any path of $P_e$, and paths leaving $X$ there may have been terminated later. Claim~3 plays the role, for the
direction $n_X$, of the middle-chord argument in the proof of [R, Lemma~\ref{imp:L4.12}(a)].

\subsection{The vertical comparison}\label{ssec:P6-vertical}

For $e=(X,Y)\in\mathcal C$ let $A_e:=X$ if $X$ is a square and $A_e:=[0,k]\times\{0\}$ if $X=X_0$, and let
$V_{x'}:=\{x'\}\times\R$. Let $E^v_e$ be the set of $x'\in[0,k]$ such that $V_{x'}$ meets $A_e$ and $Y$, every point
of $A_e\cap V_{x'}$ has height at most the height of every point of $Y\cap V_{x'}$, and the open segment between the
highest point of $A_e\cap V_{x'}$ and the lowest point of $Y\cap V_{x'}$ meets no square; for $x'\in E^v_e$ let
$g^v_e(x')$ be the length of this segment. This is the set $E_{AB}$ and the gap function of the vertical version of
[R, Lemma~\ref{imp:L4.12}] with $A=A_e$ below and $B=Y$ above. Put
\[
\chi_e:=\cos\varphi_X-\kappa_e\sin\varphi_X,\qquad \Delta_e:=\frac{\delta\,|\sin\varphi_X|}{\chi_e}.
\]
($\Delta_e$ is a length; it is unrelated to the phase difference $\Delta$ used in Section~\ref{sec:P3}.)

\begin{lemma}[vertical comparison P6.V]\label{lem:P6V}
Let $e=(X,Y)\in\mathcal C$. Then $\chi_e>0.99999849$ \textup{(}rounded down\textup{)}, and there is an open interval
$J'_e\subset E^v_e\cap(0,k)$, possibly empty, with the following properties.
\begin{enumerate}[label=(\alph*)]
\item For $x'\in J'_e$, the gap between $A_e$ and $Y$ on $V_{x'}$ \textup{(}as in the definition of $E^v_e$\textup{)}
is the open segment from a point $b(x')$ to the point $c(x')=b(x')+g^v_e(x')\,e_2$, where $b(x')\in\relint\Top(X)$
is the highest point of $X\cap V_{x'}$ \textup{(}$b(x')=(x',0)$ if $X=X_0$\textup{)} and $c(x')\in\relint\Bot(Y)$ is
the lowest point of $Y\cap V_{x'}$; this open segment lies in $\mathcal O_e$, and
$0\le g^v_e(x')<\delta/\chi_e<1.0000016\,\delta$ \textup{(}rounded up\textup{)}.
\item On $J'_e$ the function $g^v_e$ is affine with slope $\tan\varphi_Y-\tan\varphi_X$. Either $g^v_e>0$ on $J'_e$,
or $g^v_e\equiv0$ on $J'_e$; the latter happens only if $X=X_0$, $\varphi_Y=0$ and $\Bot(Y)\subset[0,k]\times\{0\}$,
and then $\theta_e=0$.
\item $|J'_e|\ \ge\ \cos\varphi_X\,(m_e-\Delta_e)\ \ge\ \cos\bbar\cdot m_e-1.0000016\,\delta\bbar$
\textup{(}rounded up\textup{)}.
\end{enumerate}
The ceiling version holds for $e\in\tilde{\mathcal C}$, with ``highest'' and ``lowest'' exchanged: there the gap is the
open vertical segment between the highest point of the second member \textup{(}below\textup{)} and the lowest point of
the first member \textup{(}above\textup{)}, or the point $(x',k)$ if the first member is $\tilde X_0$.
\end{lemma}

\begin{proof}
Write $\varphi:=\varphi_X$, $\kappa:=\kappa_e$, $\chi:=\chi_e$, and use the frame of \S\ref{ssec:P6-hull}.

\emph{The constants.} By Lemma~\ref{lem:P6-step}(d), $|\varphi|<\bbar=1.5\cdot10^{-6}$, and $|\kappa|<1.00001$ by
\eqref{eq:P6-t}. Hence
$\chi\ge\cos\varphi-|\kappa|\,|\sin\varphi|>\cos(1.5\cdot10^{-6})-1.00001\cdot1.5\cdot10^{-6}>0.99999849$, and
$1/\chi<1/0.99999849<1.0000016$.

\emph{Construction.} For $w\in(w_1,w_2)$ put $b_w:=o+wu$, $\lambda_w:=t(w)/\chi\ge0$ and $c_w:=b_w+\lambda_w e_2$.
By \eqref{eq:flows-frame}, $c_w=o+\hat w(w)\,u+\lambda_w\cos\varphi\,n$ with $\hat w(w):=w+\lambda_w\sin\varphi$. Since
$t$ is affine, $t(\hat w(w))=t(w)+\kappa\lambda_w\sin\varphi=\lambda_w(\chi+\kappa\sin\varphi)=\lambda_w\cos\varphi$, so
$c_w=o+\hat w(w)u+t(\hat w(w))n\in L_Y$. Moreover
\[
\hat w(w)=w+\frac{(t_0+\kappa w)\sin\varphi}{\chi}=\frac{\cos\varphi}{\chi}\,w+\frac{t_0\sin\varphi}{\chi}
\]
is an increasing affine function of $w$. Let $G_e:=\{w\in(w_1,w_2):\ \hat w(w)\in(w_1,w_2)\}$, an open interval
(possibly empty) as the intersection of two open intervals, and
\[
J'_e:=\{b_w\cdot e_1:\ w\in G_e\}=\{o\cdot e_1+w\cos\varphi:\ w\in G_e\},
\]
an open interval with $|J'_e|=\cos\varphi\,|G_e|$. For $x'\in J'_e$ let $w(x')\in G_e$ be the unique $w$ with
$b_w\cdot e_1=x'$, and put $b(x'):=b_{w(x')}$, $c(x'):=c_{w(x')}$.

\emph{(a).} Fix $w\in G_e$ and $0<\mu<\lambda_w$. By \eqref{eq:flows-frame},
$b_w+\mu e_2=o+(w+\mu\sin\varphi)u+\mu\cos\varphi\,n$. The first coordinate lies between $w$ and $\hat w(w)$, both
in $(w_1,w_2)$, hence in $(w_1,w_2)$. The second satisfies $\mu\cos\varphi>0$, and
$\mu\cos\varphi<t(w+\mu\sin\varphi)=t(w)+\kappa\mu\sin\varphi$, because this is equivalent to $\mu\chi<t(w)$, that is,
to $\mu<\lambda_w$. So $b_w+\mu e_2\in\mathcal O_e$, and the open vertical segment from $b_w$ to $c_w$ lies in
$\mathcal O_e$, hence in the waste (Lemma~\ref{lem:P6-hull}(ii)).

Let $x':=b_w\cdot e_1$. If $X$ is a square, then $b_w\in\relint\Top(X)\subset X$ by Lemma~\ref{lem:P6-hull}(i), and
for $\mu>0$ we have $(b_w+\mu e_2-c_X)\cdot n=\frac12+\mu\cos\varphi>\frac12$, so $b_w+\mu e_2\notin X$; thus $b_w$ is
the highest point of $X\cap V_{x'}$. If $X=X_0$, then $b_w=(w,0)=(x',0)$ is the point of $A_e\cap V_{x'}$. Since
$\hat w(w)\in(w_1,w_2)$, Lemma~\ref{lem:P6-hull}(i) gives $c_w=o+\hat w(w)u+t(\hat w(w))n\in\relint\Bot(Y)\subset Y$.
For $\mu>0$, $(c_w-\mu e_2-c_Y)\cdot n_Y=-\frac12-\mu\cos\varphi_Y<-\frac12$, because
$\cos\varphi_Y\ge\cos\frac\pi4>0$ (also when $\varphi_Y=\frac\pi4$); so $c_w-\mu e_2\notin Y$, and $c_w$ is the lowest
point of $Y\cap V_{x'}$. Consequently every point of $A_e\cap V_{x'}$ has height $\le b_w\cdot e_2$, every point of
$Y\cap V_{x'}$ has height $\ge c_w\cdot e_2\ge b_w\cdot e_2$, and the open segment between $b_w$ and $c_w$ meets no
square. So $x'\in E^v_e$ and $g^v_e(x')=|c_w-b_w|=\lambda_w=t(w)/\chi<\delta/\chi$ by Lemma~\ref{lem:P6-hull}(i).
Finally $x'\in(0,k)$: if $\lambda_w>0$, the segment lies in $\mathcal O_e\subset(0,k)^2$; if $\lambda_w=0$, then
$t(w)=0$, so $X=X_0$ by Lemma~\ref{lem:P6-hull}(i), and $x'=w\in(w_1,w_2)\subset(0,k)$.

\emph{(b).} For $x'\in J'_e$, $b(x')$ lies on the line of $\Top(X)$ (on $\R\times\{0\}$ if $X=X_0$) and $c(x')$ on
$L_Y$. These lines have directions $u_X$ and $u_Y$ and are not vertical, as $\cos\varphi_X,\cos\varphi_Y>0$; their
heights above the abscissa $x'$ are therefore affine functions of $x'$ with slopes $\tan\varphi_X$ and
$\tan\varphi_Y$, and $g^v_e(x')=c(x')\cdot e_2-b(x')\cdot e_2$ is affine with slope $\tan\varphi_Y-\tan\varphi_X$. Next,
$t$ is affine and, by Lemma~\ref{lem:P6-hull}(i), nonnegative on $[w_1,w_2]$; if $t(w)=0$ for some $w\in(w_1,w_2)$,
then $t\equiv0$ on $[w_1,w_2]$. For a square $X$ this is excluded by Lemma~\ref{lem:P6-hull}(i). For $X=X_0$ it means
$\kappa_e=0$ and $t_0=0$, that is, $\varphi_Y=0$ and $L_Y=\R\times\{0\}$; then $\theta_e=a(Y)=0$. Since
$g^v_e(x')=t(w(x'))/\chi$, this proves (b).

\emph{(c).} For $w\in(w_1,w_2)$ we have $t(w)<\delta$, hence
$|\hat w(w)-w|=\lambda_w|\sin\varphi|\le\delta|\sin\varphi|/\chi=\Delta_e$. If $\sin\varphi\ge0$, then
$w\le\hat w(w)\le w+\Delta_e$, so $(w_1,w_2-\Delta_e)\subset G_e$; if $\sin\varphi<0$, then
$w-\Delta_e\le\hat w(w)\le w$, so $(w_1+\Delta_e,w_2)\subset G_e$. In both cases $|G_e|\ge(w_2-w_1-\Delta_e)_+$.
Since $I_e\subset[w_1,w_2]$, Lemma~\ref{lem:P6-density} gives $w_2-w_1\ge|I_e|\ge m_e$. Therefore
\[
|J'_e|=\cos\varphi\,|G_e|\ \ge\ \cos\varphi\,(w_2-w_1-\Delta_e)_+\ \ge\ \cos\varphi\,(m_e-\Delta_e).
\]
Finally $\cos\varphi\ge\cos\bbar$, and
$\cos\varphi\cdot\Delta_e\le\delta\,|\sin\varphi|/\chi<1.0000016\,\delta\sin\bbar\le1.0000016\,\delta\bbar$, so
$|J'_e|\ge\cos\bbar\cdot m_e-1.0000016\,\delta\bbar$. (If $w_1=w_2$, then $G_e=J'_e=\emptyset$ and
$m_e\le|I_e|=0$, and (c) holds trivially.)

The ceiling version follows by applying the above to $\rho\cP$ and transporting by $\rho$, which preserves abscissae
and vertical lengths and exchanges ``highest'' and ``lowest''.
\end{proof}

\begin{remark}[the exceptional set]\label{rem:P6-exceptional}
If $w_1<w_2$, the abscissae $o\cdot e_1+w\cos\varphi_X$, $w\in(w_1,w_2)$, form an open interval of length
$\cos\varphi_X(w_2-w_1)$. By the proof of Lemma~\ref{lem:P6V}(c), $J'_e$ contains all of it except one end interval
of length at most $\cos\varphi_X\,\Delta_e<1.0000016\,\delta\bbar$, at the end towards which the vertical direction
drifts in the frame of $X$; only in this end interval can the construction of Lemma~\ref{lem:P6V} fail. Points of the
hull that are not exit coordinates (holes of $I_e$) cause no loss, since Lemma~\ref{lem:P6-hull} concerns the whole
hull.
\end{remark}

\subsection{Disjointness of the vertical slits}\label{ssec:P6-disjoint}

For $e\in\mathcal C$ put
\[
\mathcal S_e:=\{(x',y):\ x'\in J'_e,\ b(x')\cdot e_2<y<c(x')\cdot e_2\},\qquad s_e:=\int_{J'_e}g^v_e(x')\,dx' ,
\]
with $b,c$ as in Lemma~\ref{lem:P6V}(a); for $e\in\tilde{\mathcal C}$ let $J'_e$ and $s_e$ be defined through $\rho\cP$,
and let $\mathcal S_e$ be the image under $\rho$ of the corresponding set for $\rho\cP$. Let $\mathcal K$ be the set of
squares that are members of some $e\in\mathcal C$, and $\tilde{\mathcal K}$ the set of squares that are members of some
$e\in\tilde{\mathcal C}$.

\begin{lemma}[disjoint vertical slits]\label{lem:P6-disjoint}
\begin{enumerate}[label=(\alph*)]
\item $\mathcal K\cap\tilde{\mathcal K}=\emptyset$.
\item Each $\mathcal S_e$, $e\in\mathcal C\cup\tilde{\mathcal C}$, is an open subset of the waste with
$|\mathcal S_e|=s_e$.
\item The sets $\mathcal S_e$, $e\in\mathcal C\cup\tilde{\mathcal C}$, are pairwise disjoint. Consequently
$\sum_{e\in\mathcal C\cup\tilde{\mathcal C}}s_e\le W$.
\end{enumerate}
\end{lemma}

\begin{proof}
(a) Let $S\in\mathcal K$, a member of $e\in\mathcal C$, and let $x\in P_e$. If $S$ is the second member, it contains
the entry point $q_x$, of height $\le\hmax=\frac k2-1$; if it is the first member, it contains its entry point, also
of height $\le\hmax$ (Lemma~\ref{lem:P6-step}(b)). By the ceiling version, every $S\in\tilde{\mathcal K}$ contains a
point of height $\ge k-\hmax=\frac k2+1$. A square of $\mathcal K\cap\tilde{\mathcal K}$ would have vertical extent at
least $2$, whereas its vertical extent is $\cos a+\sin a\le\sqrt2$ by [R, Lemma~\ref{imp:L3.1}(a)].

(b) Let $e\in\mathcal C$. On the open interval $J'_e$ the functions $x'\mapsto b(x')\cdot e_2$ and
$x'\mapsto c(x')\cdot e_2$ are affine (proof of Lemma~\ref{lem:P6V}(b)), so $\mathcal S_e$ is open and, by Tonelli's
theorem, $|\mathcal S_e|=\int_{J'_e}(c(x')-b(x'))\cdot e_2\,dx'=s_e$. By Lemma~\ref{lem:P6V}(a),
$\mathcal S_e\subset\mathcal O_e$, which lies in the waste. For $e\in\tilde{\mathcal C}$ the same follows by
reflection, since $\rho$ maps the waste of $\rho\cP$ onto the waste of $\cP$ and preserves areas.

(c) For $x'\in(0,k)$ let $\mathcal U_{x'}:=\{y\in(0,k):\ (x',y)\ \text{lies in no square}\}$, an open subset of $\R$.
For $e\in\mathcal C\cup\tilde{\mathcal C}$ and $x'\in J'_e$, the section $\mathcal S_e^{x'}:=\{y:\ (x',y)\in\mathcal S_e\}$
is an open interval (empty if $g^v_e(x')=0$). Let $e=(X,Y)\in\mathcal C$ and $\mathcal S_e^{x'}\neq\emptyset$. Then
$\mathcal S_e^{x'}\subset\mathcal U_{x'}$ by (b); its upper endpoint is $c(x')\cdot e_2$ with $c(x')\in Y$, and its lower
endpoint is $b(x')\cdot e_2$ with $b(x')\in X$ if $X$ is a square and $b(x')\cdot e_2=0$ if $X=X_0$. So neither
endpoint lies in $\mathcal U_{x'}$, and $\mathcal S_e^{x'}$ is a connected component of $\mathcal U_{x'}$. By reflection,
for $e=(X,Y)\in\tilde{\mathcal C}$ every nonempty section $\mathcal S_e^{x'}$ is a component of $\mathcal U_{x'}$ whose
lower endpoint is the height of a point of $Y$ and whose upper endpoint is the height of a point of $X$ if $X$ is a
square, and equals $k$ if $X=\tilde X_0$.

Let $e\neq e'$ in $\mathcal C\cup\tilde{\mathcal C}$ and suppose that $(x',y)\in\mathcal S_e\cap\mathcal S_{e'}$. Then
$\mathcal S_e^{x'}$ and $\mathcal S_{e'}^{x'}$ are the component of $\mathcal U_{x'}$ that contains $y$, so they have the
same endpoints; the corresponding points on $V_{x'}$ coincide.

\emph{Case $e=(X,Y)$, $e'=(X',Y')\in\mathcal C$.} The common upper endpoint $c(x')$ lies in $Y$ and in $Y'$, so
$Y=Y'$, since distinct squares are disjoint. If $X$ is a square, then $b(x')\in\relint\Top(X)$ has height $>0$ by
\eqref{eq:P6-heights}; so the common lower endpoint has height $0$ if and only if $X=X_0$, and likewise for $X'$.
Hence either $X=X'=X_0$, or $b(x')\in X\cap X'$ and $X=X'$. In both cases $e=e'$, a contradiction.

\emph{Case $e,e'\in\tilde{\mathcal C}$.} This is the previous case for $\rho\cP$.

\emph{Case $e=(X,Y)\in\mathcal C$, $e'=(X',Y')\in\tilde{\mathcal C}$.} The upper endpoint of $\mathcal S_e^{x'}$ is the
height of $c(x')\in\relint\Bot(Y)$, which is $<k$ by \eqref{eq:P6-heights}. So the upper endpoint of
$\mathcal S_{e'}^{x'}$ is not $k$; hence $X'$ is a square and contains $c(x')$. Then $X'=Y\in\mathcal K\cap\tilde{\mathcal K}$,
contradicting (a).

So the sets $\mathcal S_e$ are pairwise disjoint subsets of the waste, and
$\sum_es_e=\big|\bigcup_e\mathcal S_e\big|\le W$; the family is finite.
\end{proof}

\subsection{Slit inequality, pair waste and multiplicity}\label{ssec:P6-pairs}

\begin{lemma}[slit inequality]\label{lem:P6-slit}
For every $e\in\mathcal C\cup\tilde{\mathcal C}$,
\begin{equation}\label{eq:P6-slit}
\theta'_e\,(2-\bbar)\,|J'_e|^2\ \le\ 4\,s_e .
\end{equation}
\end{lemma}

\begin{proof}
Let $e=(X,Y)\in\mathcal C$, and assume that the interval $J:=J'_e$ is not empty (otherwise there is nothing to show).
By parts (a) and (b) of Lemma~\ref{lem:P6V}, the function $g:=g^v_e$ is affine and nonnegative on $J$, with
\[
|g'|=|\tan\varphi_Y-\tan\varphi_X|\ \ge\ |\varphi_Y-\varphi_X|\ \ge\ \theta_e ,
\]
by the mean value theorem ($\tan'=\sec^2\ge1$) and since $\theta_e\le|\varphi_Y-\varphi_X|$; for $X=X_0$ this reads
$|\tan\varphi_Y|\ge|\varphi_Y|=a(Y)=\theta_e$. Let $x_*$ be an endpoint of $\bar J$ at which the continuous extension
of $g$ to $\bar J$ is smallest; it is $\ge0$. Then $g(x')\ge|g'|\,|x'-x_*|\ge\theta_e|x'-x_*|$ on $J$, so
$s_e=\int_Jg\ge\theta_e|J|^2/2$. As $\theta'_e\le\theta_e$ and $2-\bbar\le2$, we get
$\theta'_e(2-\bbar)|J|^2\le2\theta_e|J|^2\le4s_e$. (This is the monotone case of [R, Lemma~\ref{imp:L4.1}].) For
$e\in\tilde{\mathcal C}$ the claim follows by reflection.
\end{proof}

For two squares $S_1,S_2$, and for a square $S_1$ and $S_2=S_0=\partial[0,k]^2$, let $R(S_1,S_2)$ be the rectangle
defined before [R, Proposition~\ref{imp:P4.3}]; for $R(S_1,S_0)$ the point of $S_0$ nearest to the centre of $S_1$
used there is fixed once and for all for each square $S_1$. For $e=(X,Y)\in\mathcal C\cup\tilde{\mathcal C}$ put
\[
R_e:=R(X,Y)\ \text{if $X$ is a square},\qquad R_e:=R(Y,S_0)\ \text{if $X\in\{X_0,\tilde X_0\}$},\qquad
f_e:=|R_e\cap\text{waste}| .
\]

\begin{lemma}[pair waste]\label{lem:P6-pairwaste}
Let $e=(X,Y)\in\mathcal C\cup\tilde{\mathcal C}$. If $X$ is a square, then $\dist(X,Y)<\delta$; if $X=X_0$
\textup{(}resp.\ $X=\tilde X_0$\textup{)}, then $Y$ is at distance less than $\delta$ from the bottom side
$[0,k]\times\{0\}$ \textup{(}resp.\ the top side $[0,k]\times\{k\}$\textup{)}, and hence from $S_0$. Moreover
\begin{equation}\label{eq:P6-pairwaste}
f_e\ \ge\ \frac{Q_*\,\theta'_e\,(2-\bbar)}{4} .
\end{equation}
\end{lemma}

\begin{proof}
Let $e\in\mathcal C$ and $x\in P_e$. By Lemma~\ref{lem:P6-step}(a), $|q_x-p_x|=t_x<\delta$, where $q_x\in Y$, and
$p_x\in X$ if $X$ is a square, while $p_x=(x,0)$ lies on the bottom side if $X=X_0$. This gives the distance bounds;
for $e\in\tilde{\mathcal C}$ they follow by reflection, which preserves distances and maps the bottom side onto the
top side. Note that $\theta(Y,S_0)=a(Y)=\theta_e$ if $X\in\{X_0,\tilde X_0\}$.

If $\theta_e\le10^{-5}$, [R, Proposition~\ref{imp:P4.5}], applied with $S_1=X$, $S_2=Y$ (respectively $S_1=Y$,
$S_2=S_0$), $d=\delta\le10^{-5}$ and $\theta_1=10^{-5}$, gives $f_e\ge Q_*\theta_e(2-\theta_e)/4\ge Q_*\theta'_e(2-\bbar)/4$;
the last step holds because $\theta'_e\le\theta_e$, $\theta'_e\le\bbar$ and $x\mapsto x(2-x)$ is increasing on
$[0,1]$, so that $\theta'_e(2-\bbar)\le\theta'_e(2-\theta'_e)\le\theta_e(2-\theta_e)$. If $\theta_e>10^{-5}$,
[R, Lemma~\ref{imp:L4.8}(a)], with $d=\delta\le10^{-4}$, gives
$f_e\ge\theta_e(2-\theta_e)/40\ge10^{-5}(2-10^{-5})/40>4.99\cdot10^{-7}$ (as $\theta_e\le\frac\pi4<1$), while
$Q_*\theta'_e(2-\bbar)/4\le0.32\cdot1.5\cdot10^{-6}\cdot2/4=2.4\cdot10^{-7}$.

If $X\in\{X_0,\tilde X_0\}$ and $Y$ is near a corner of $[0,k]^2$, the rectangle $R(Y,S_0)$ may stand on the left
or the right side of $[0,k]^2$ rather than on the bottom or the top side. This does not matter here:
[R, Proposition~\ref{imp:P4.5}] and [R, Lemma~\ref{imp:L4.8}(a)] concern $R(S_1,S_0)$ for any square $S_1$ at distance
less than $d$ from $S_0$, with $\theta(S_1,S_0)=a(S_1)$ (as noted in the proof of [R, Lemma~\ref{imp:L4.13}]), and
Proposition~\ref{imp:pairs} allows this situation (see the remark following it).
\end{proof}

For $e=(X,Y)\in\mathcal C\cup\tilde{\mathcal C}$ let $\bar e:=\{\hat X,Y\}$ be the unordered pair with $\hat X:=X$ if
$X$ is a square, $\hat X:=[0,k]\times\{0\}$ (the bottom side) if $X=X_0$, and $\hat X:=[0,k]\times\{k\}$ (the top side)
if $X=\tilde X_0$.

\begin{lemma}[multiplicity]\label{lem:P6-overlap}
The unordered pairs $\bar e$, $e\in\mathcal C\cup\tilde{\mathcal C}$, are pairwise distinct, and
\begin{equation}\label{eq:P6-overlap}
\sum_{e\in\mathcal C\cup\tilde{\mathcal C}}f_e\ \le\ K_w^*\,W=9\,W .
\end{equation}
Without [R, Lemma~\ref{imp:L4.10}], $\sum_{e\in\mathcal C\cup\tilde{\mathcal C}}f_e\le K_wW=13\,W$.
\end{lemma}

\begin{proof}
\emph{Distinctness.} Let $e\neq e'$ with $\bar e=\bar e'$. If $e=(X,Y)$ and $e'$ are both in $\mathcal C$, then, as
$X_0$ is never a second member, $e'=(Y,X)$ with $X$ and $Y$ squares; by~\ref{it:P6-F2}, $c_Y\cdot e_2>c_X\cdot e_2$
(from $e$) and $c_X\cdot e_2>c_Y\cdot e_2$ (from $e'$), which is impossible. Both in $\tilde{\mathcal C}$: by
reflection. If $e\in\mathcal C$ and $e'\in\tilde{\mathcal C}$, the second member of $e$ is a square of
$\bar e=\bar e'$, so it lies in $\mathcal K\cap\tilde{\mathcal K}$, contradicting Lemma~\ref{lem:P6-disjoint}(a).

\emph{The bound.} Let $\mathfrak F:=\{\bar e:\ e\in\mathcal C\cup\tilde{\mathcal C}\}$, a family of distinct unordered
pairs. Each pair consists of two squares, or of a square and one of the two opposite sides $[0,k]\times\{0\}$,
$[0,k]\times\{k\}$, and its members are at distance less than $\delta\le10^{-4}$ (Lemma~\ref{lem:P6-pairwaste});
the rectangle that Proposition~\ref{imp:pairs} attaches to $\bar e$ is $R_e$. Proposition~\ref{imp:pairs}, with
$d:=\delta$ and $\sigma,\sigma'$ the bottom and the top side, gives $\sum_{e}f_e\le K_w^*W=9W$, and
$\sum_ef_e\le K_wW=13W$ by the analytic argument. That proposition is the form of the multiplicity bound given by
the proof of [R, Lemma~\ref{imp:L4.10}], as stated in [R, Remark~\ref{imp:R4.11}] (``The proof bounds the number of
\emph{all} pairs of squares (or of a square and one of two opposite sides) at distance less than $d$ whose rectangles
contain $p$, without using the graph $G$''), and of [R, Lemma~\ref{imp:L4.9}] by the paragraph following its proof;
the formal statements of [R, Lemmas~\ref{imp:L4.9} and~\ref{imp:L4.10}] concern edges of the visibility graphs, and
the pairs $\bar e$ need not be such edges.
\end{proof}

\subsection{Proof of Theorem~\ref{thm:P6}}\label{ssec:P6-proof}

\emph{Step 1: the chain.} Let $x\in\cX$, let $Z_1,\dots,Z_J$ be the squares entered by $\pi^0_x$, in order, and
$Z_0:=X_0$. By~\ref{it:P6-F2}, $Z_x=Z_j$ for exactly one $j\in\{1,\dots,J\}$. By the triangle inequality for $\theta$,
\[
\zeta(x)\ \le\ a(Z_j)=\theta(Z_0,Z_j)\ \le\ \sum_{i=1}^j\theta(Z_{i-1},Z_i).
\]
With $\theta':=\min(\theta,\bbar)$ also $\sum_{i=1}^j\theta'(Z_{i-1},Z_i)\ge\zeta(x)$: if some term has
$\theta(Z_{i-1},Z_i)\ge\bbar$, its capped value is $\bbar\ge\zeta(x)$; otherwise no term is capped. The pairs
$(Z_{i-1},Z_i)$, $1\le i\le J$, are pairwise distinct and are exactly the $e\in\mathcal C$ with $x\in P_e$; the terms
with $i>j$ are nonnegative. Hence
\[
\zeta(x)\ \le\ \vartheta(x):=\sum_{e\in\mathcal C}\theta'_e\,\mathbf 1_{P_e}(x).
\]
$\vartheta$ is a nonnegative Borel function on $(0,k)$ (Lemma~\ref{lem:P6-density}; the sum is finite), so
\begin{equation}\label{eq:P6-chain}
\int_{\cX}\zeta\ \le\ \int_0^k\vartheta\ =\ \sum_{e\in\mathcal C}\theta'_e\,m_e ,\qquad\text{and likewise}\qquad
\int_{\cX'}\zeta'\ \le\ \sum_{e\in\tilde{\mathcal C}}\theta'_e\,m_e
\end{equation}
by the ceiling version.

\emph{Step 2: vertical comparison.} Let $e=(X,Y)\in\mathcal C\cup\tilde{\mathcal C}$. By Lemma~\ref{lem:P6V}(c),
$m_e\le|J'_e|/\cos\varphi_X+\Delta_e$, where $\cos\varphi_X\ge\cos\bbar$ and
$\Delta_e=\delta|\sin\varphi_X|/\chi_e<1.0000016\,\delta\bbar$. Hence
\begin{equation}\label{eq:P6-comparison}
\sum_{e\in\mathcal C\cup\tilde{\mathcal C}}\theta'_e\,m_e\ \le\ \frac1{\cos\bbar}\sum_{e\in\mathcal C\cup\tilde{\mathcal C}}\theta'_e\,|J'_e|
\ +\ 1.0000016\,\delta\bbar\sum_{e\in\mathcal C\cup\tilde{\mathcal C}}\theta'_e .
\end{equation}

\emph{Step 3: Cauchy--Schwarz.} For $e\in\mathcal C\cup\tilde{\mathcal C}$, multiplying \eqref{eq:P6-slit} and
\eqref{eq:P6-pairwaste} gives $Q_*\theta_e'^2(2-\bbar)^2|J'_e|^2\le16\,s_ef_e$, that is,
$\theta'_e|J'_e|\le\frac{4}{(2-\bbar)\sqrt{Q_*}}\sqrt{s_ef_e}$. By the Cauchy--Schwarz inequality,
Lemma~\ref{lem:P6-disjoint}(c) and \eqref{eq:P6-overlap},
\[
\begin{aligned}
\sum_{e}\theta'_e|J'_e|\ &\le\ \frac{4}{(2-\bbar)\sqrt{Q_*}}\Big(\sum_es_e\Big)^{1/2}\Big(\sum_ef_e\Big)^{1/2}\\
&\le\ \frac{4}{(2-\bbar)\sqrt{Q_*}}\sqrt{W\cdot9W}\ =\ \frac{4\sqrt{9/Q_*}}{2-\bbar}\,W\ \le\ A\,W,
\end{aligned}
\]
the sums being over $\mathcal C\cup\tilde{\mathcal C}$ and the last step using $2-\bbar\ge2-3\cdot10^{-6}$.

\emph{Step 4: the error term.} By \eqref{eq:P6-pairwaste} and \eqref{eq:P6-overlap},
\[
\sum_{e\in\mathcal C\cup\tilde{\mathcal C}}\theta'_e\ \le\ \frac{4}{Q_*(2-\bbar)}\sum_ef_e\ \le\ \frac{36}{0.32\,(2-1.5\cdot10^{-6})}\,W\ <\ 56.26\,W ,
\]
so the last term of \eqref{eq:P6-comparison} is at most
$1.0000016\cdot10^{-5}\cdot1.5\cdot10^{-6}\cdot56.26\,W<8.5\cdot10^{-10}\,W$ (rounded up).

\emph{Step 5: conclusion.} Steps 2--4 give (b). Adding the two inequalities of \eqref{eq:P6-chain} and using (b)
gives $\int_{\cX}\zeta+\int_{\cX'}\zeta'\le(A/\cos\bbar+8.5\cdot10^{-10})W\le A'W$, which is (a). For the numerical
bounds: $A=4\sqrt{9/0.32}/(2-3\cdot10^{-6})<10.6067$ and $1/\cos\bbar<1+1.2\cdot10^{-12}$, hence
$A'<10.6067\,(1+1.2\cdot10^{-12})+10^{-8}<10.60671<10.6068$ (all rounded up).

\emph{Step 6: part (c).} Replace \eqref{eq:P6-overlap} by $\sum_ef_e\le13W$ (Lemma~\ref{lem:P6-overlap}). Step~3
gives $\sum_e\theta'_e|J'_e|\le4\sqrt{13/Q_*}\,W/(2-\bbar)\le A_{13}W$, and Step~4 gives
$\sum_e\theta'_e\le52\,W/(0.32\,(2-1.5\cdot10^{-6}))<81.26\,W$ and an error term at most
$1.0000016\cdot10^{-5}\cdot1.5\cdot10^{-6}\cdot81.26\,W<1.22\cdot10^{-9}\,W$. Hence
$\sum_e\theta'_em_e\le(A_{13}/\cos\bbar+1.22\cdot10^{-9})W$ and $\int_{\cX}\zeta+\int_{\cX'}\zeta'\le A'_{13}W$.
Numerically $A_{13}=4\sqrt{13/0.32}/(2-3\cdot10^{-6})<12.74757$ and
$A'_{13}<12.74757\,(1+1.2\cdot10^{-12})+10^{-8}<12.7476$ (rounded up). \qed

\subsection{The scale form}\label{ssec:P6-scale}

Recall from Definition~\ref{def:flows-Tstar} the tilt time $\Tstar$: with $(Z_j,q_j)$, $1\le j\le J(x)$, the
R1-passed entries of $\pi^0_x$ and $\eta_j:=q_j\cdot e_2$,
\[
s_j:=\max\big(y_0,\ \eta_j-2,\ (c_0/a(Z_j))^2\big)\ \text{ if } a(Z_j)>0,\qquad s_j:=+\infty\ \text{ if } a(Z_j)=0,
\]
and $\Tstar(x):=\min_js_j$ if this minimum is at most $y_1=\frac k2-3$; otherwise (in particular if $J(x)=0$)
$x\notin\operatorname{dom}\Tstar$. Here $\alpha(s)=c_0s^{-1/2}$ and $\amax=\alpha(y_0)$
(Definition~\ref{def:flows-params}). Only R1-passed entries enter; deaths, merges and end collisions do not. The
ceiling tilt time $\tilde T^*=\Tstar[\rho\cP]$ is defined from $\tilde F^0$ (Definition~\ref{def:flows-ceiling}).

\begin{corollary}[scale form]\label{cor:P6scale}
Let $k\ge4$ be an integer, $\cP$ a closed packing of $[0,k]^2$ with waste area $W$, $0<\delta\le10^{-5}$, and assume
\textup{(C0)} and \textup{(C1)}. Then $\operatorname{dom}\Tstar$ and $\operatorname{dom}\tilde T^*$ are Borel subsets of
$(0,k)$, $\Tstar$ and $\tilde T^*$ are Borel functions on them with values in $[y_0,y_1]$, and
\[
\int_{\operatorname{dom}\Tstar}\alpha(\Tstar(x))\,dx+\int_{\operatorname{dom}\tilde T^*}\alpha(\tilde T^*(x))\,dx\ \le\ A'W,
\qquad A'<10.6068 .
\]
Without [R, Lemma~\ref{imp:L4.10}], using [R, Lemma~\ref{imp:L4.9}] in its place, the same holds with $A'_{13}<12.7476$
in place of $A'$.
\end{corollary}

\begin{proof}
\emph{Borel measurability.} By Lemma~\ref{lem:flows-meas}(a),(b), $\operatorname{dom}\Tstar$ is a finite union of
intervals and points, and $\Tstar$ is a Borel function on it (on each piece of $\Fz$ the entries $Z_j$ are fixed and
the $\eta_j$ are affine, so $\Tstar$ is continuous on the intersection of $\operatorname{dom}\Tstar$ with each piece).
Its values lie in $[y_0,y_1]$ because $s_j\ge y_0$. The ceiling tilt time is treated by reflection.

\emph{The hypotheses of Theorem~\ref{thm:P6}.} Let $x\in\operatorname{dom}\Tstar$ and choose $j$ with $\Tstar(x)=s_j$.
As $s_j\le y_1<\infty$, we have $a(Z_j)>0$. Since $\alpha$ is decreasing and
$s_j\ge\max\big(y_0,(c_0/a(Z_j))^2\big)$,
\[
\begin{aligned}
0<\alpha(\Tstar(x))=c_0\,s_j^{-1/2}\ &\le\ \min\Big(c_0\,y_0^{-1/2},\ c_0\big((c_0/a(Z_j))^2\big)^{-1/2}\Big)\\
&=\ \min\big(\amax,a(Z_j)\big)\ \le\ \min\big(\bbar,a(Z_j)\big),
\end{aligned}
\]
by (C1). So $\zeta:=\alpha\circ\Tstar$ is a Borel function from $\cX:=\operatorname{dom}\Tstar$ to $[0,\bbar]$ ($\alpha$
is continuous on $(0,\infty)$), and $\pi^0_x$ enters $Z_x:=Z_j$ (an R1-passed entry at $q_j\in\relint\Bot(Z_j)$) with
$a(Z_x)\ge\zeta(x)$. The same holds for $\cX':=\operatorname{dom}\tilde T^*$ and $\zeta':=\alpha\circ\tilde T^*$ with
$\tilde F^0$. Theorem~\ref{thm:P6}(a), respectively (c), gives the claim.
\end{proof}

All scale flows $\Fs$ are truncations of the single master flow $\Fz$, in which R1 is decided once
(Definition~\ref{def:Fs} and Remark~\ref{rem:flows-joint}); accordingly Corollary~\ref{cor:P6scale} involves the paths
of $\Fz$ only, and the injectivity of Lemma~\ref{lem:P6-density} is that of this one family of paths.

\subsection{Remarks}\label{ssec:P6-remarks}

\begin{remark}[where the rules R1 and R2 enter]\label{rem:P6-rules}
(1) R1 is used only in Lemma~\ref{lem:P6-density}, through~\ref{it:P6-F1} applied at the first member $X$ of a pair:
two floor points with the same exit point from $X$ entered $X$ at the same point, and at most one path passes $X$ from
there. A path that passes $X$ is the R1 winner at its entry point into $X$; as $a(X)<\amax$, this does not depend on
the order in which M and T are decided at an entry candidate (compare Remark~\ref{rem:flows-order}). Nothing about the
entry into the second member is used in Lemma~\ref{lem:P6-density}.

(2) R2 enters in two ways. First, the first member $X$ of every pair with $X\neq X_0$ is a passed square, so
$|\varphi_X|<\amax\le\bbar$ (Lemma~\ref{lem:P6-step}(d)); this makes $\Delta_e$ of order $\delta\bbar$ in
Lemma~\ref{lem:P6V}(c). Secondly, the second member $Y$ of every pair, in particular $Z_x$ in the last pair of the
chain in Step~1 of \S\ref{ssec:P6-proof}, is entered through $\relint\Bot(Y)$, since side and vertex contacts are
never entries. Hence, over the hull, the boundary of $Y$ facing $X$ is the single line $L_Y$
(Lemma~\ref{lem:P6-hull}), and $c(x')$ is the lowest point of $Y$ on $V_{x'}$ (Lemma~\ref{lem:P6V}(a)).
\end{remark}

\begin{remark}[degenerate cases]\label{rem:P6-degenerate}
The general conventions are those of Remark~\ref{rem:flows-degenerate}. \emph{Vertices and grazing contacts}: a first contact at a vertex
terminates the path (with E, or with D or W at a tie) and is never an entry, so it produces no pairs; only contacts with
$\relint\Bot(Y)$ are used. \emph{Ties} between
events are resolved by the priority D $>$ W $>$ contact $>$ H; the proofs use only the consequences $g+t_S<\delta$,
$t_S<t_W$ and $t_S\le t_H$ at processed contacts (Definition~\ref{def:flow} and Lemma~\ref{lem:flows-traj}(f); they enter
through Lemma~\ref{lem:P6-step}(a),(b)), and a contact at height exactly $\hmax$ is
processed and covered by Lemma~\ref{lem:P6-step}(b). \emph{Holes} of $I_e$, and exit coordinates of paths that are
later D-, M- or E-terminated, do not matter, because Lemma~\ref{lem:P6-hull} concerns the whole hull. \emph{A square
resting with a side on the floor} gives the case $g^v_e\equiv0$ of Lemma~\ref{lem:P6V}(b), with $\theta_e=0$ and
$\mathcal S_e=\emptyset$; empty sections are excluded in the proof of Lemma~\ref{lem:P6-disjoint}(c).
\emph{Inclination $\frac\pi4$:} for the second member only $\cos\varphi_Y>0$, $|\varphi_Y-\varphi_X|<\frac\pi2$,
$|\kappa_e|<1.00001$ and $\theta_e\le|\varphi_Y-\varphi_X|$ are used, all valid under the $45^\circ$ convention; the
first member always has $|\varphi_X|<\bbar$. \emph{Walls:} a wall event terminates the path (W), and by the priority W $>$ contact of Definition~\ref{def:flow} a
contact with $t_S=t_W$ is not processed; so only contacts reached before a wall event form pairs, and $\mathcal O_e\subset(0,k)^2$ by Lemma~\ref{lem:P6-hull}(ii).
\end{remark}

\section{P7: the per-scale count}\label{sec:P7}

In this section we fix a scale $s\in[y_0,y_1]$. We compare the measure $N(s)$ of the set of floor points whose path is $T_s$-terminated in the scale flow $\Fs$ with the number of squares of a family $\cZ(s)$ (Theorem~\ref{thm:P7}), and we bound $|\cZ(s)|$ from below by an integral over lines of kind~(iii) (Theorem~\ref{thm:P7lower}). Both theorems have ceiling versions, and both are used in Section~\ref{sec:P8}. Besides the definitions and lemmas of Section~\ref{sec:flows}, the proofs use Theorem~\ref{thm:P2}, Theorem~\ref{thm:P4}, Proposition~\ref{prop:wall} and [R, Lemmas~\ref{imp:L3.1}, \ref{imp:L3.4} and~\ref{imp:L3.5}].

\subsection{Setting, notation and statements}\label{ssec:P7-setting}

\emph{Standing assumptions.} Throughout this section $k\ge4$ is an integer, $\cP$ is a closed packing of $[0,k]^2$, $\delta=10^{-5}$ (Definition~\ref{def:flows-params}), and the parameters $c_0,c_1,y_0,\eps,\omega_0$ satisfy the constraints (C0)--(C7) of Definition~\ref{def:flows-constraints}. We write $y_1=k/2-3$, $\hmax=k/2-1$, $\alpha(t)=c_0t^{-1/2}$ and $\hat\beta(t)=c_1t^{-3/4}$ for $t>0$, and $\amax=\alpha(y_0)$. The proofs of the results of this section use directly only the constraints
\begin{itemize}
\item[(C0)] $c_0,c_1>0$, $y_0$ is a positive integer, $0<\eps<1$ and $\omega_0>0$;
\item[(C1)] $\amax=c_0y_0^{-1/2}\le1.5\cdot10^{-6}$;
\item[(C2)] $\hat\beta(y_0)=c_1y_0^{-3/4}\le1.5\cdot10^{-6}$,
\end{itemize}
together with the hypotheses of the results quoted from Sections~\ref{sec:flows}, \ref{sec:P12} and~\ref{sec:P4}, which are among the standing assumptions. Constraint (C3) is not used (Remark~\ref{rem:P7-C3}).

\emph{Lines.} We use the following notation of Section~\ref{sec:flows} (Definitions~\ref{def:flows-params} and~\ref{def:flows-lines}), repeated for convenience. For $y\in\R$ put $\ell_y:=\R\times\{y\}$; a square $S$ \emph{meets the line $y$} if $S\cap\ell_y\ne\emptyset$ (since $S\subset[0,k]^2$, this is the same as meeting $[0,k]\times\{y\}$). For a square $S$, $c_S(y)$ is the length of the closed chord $S\cap\ell_y$, $\Phi(S)=\{y:0<c_S(y)<1\}$, and $\cR(S)$ is the ramp set of [R, Lemma~\ref{imp:L3.1}]; $\omega(y)$ is the line waste and
\[
E(y)=\sum_{S:\ c_S(y)\ge1}\big(c_S(y)-1\big).
\]
Since $c_S(y)$ is the length of the section of $S$ at height $y$, $\int_\R c_S(y)\dd y=\operatorname{area}(S)=1$ (Cavalieri's principle). Further, $d(y)=\min(y,k-y)$ and
\[
\beta(y):=\hat\beta\big(d(y)\big)=c_1d(y)^{-3/4}\qquad(0<y<k),
\]
so that $\beta(t)=\hat\beta(t)$ for $0<t\le k/2$. Finally $\dist(y,\Z)=\min(\frc y,1-\frc y)$, so that for every $w>0$
\[
\frc y\in[w,1-w]\iff\dist(y,\Z)\ge w .
\]
Put
\[
w_0:=0.50001\,c_0^2\Big(1+\frac{1.0001}{y_0}\Big)+\delta+1.0001\,c_1\big((1-\eps)y_0\big)^{-3/4}+10^{-12},
\]
\[
H:=\big\{y\in(0,k):\ y_0\le d(y)\le(1-\eps)y_1,\ \frc y\in[w_0,1-w_0]\big\}.
\]
Every $y\in H$ is of exactly one of the following kinds:
\begin{itemize}
\item[(i)] $\omega(y)\ge\omega_0$;
\item[(ii)] $\omega(y)<\omega_0$, and the line $y$ meets a square $S$ with $a(S)\ge\beta(y)$;
\item[(iii)] $\omega(y)<\omega_0$, and every square $S$ meeting the line $y$ has $a(S)<\beta(y)$.
\end{itemize}
The corresponding subsets of $H$ are $Y_i$, $Y_{ii}$ and $Y_{iii}$, and $Y_b:=Y_{iii}\cap(0,k/2)$, $Y_t:=Y_{iii}\cap(k/2,k)$. For $s\in[y_0,y_1]$ put
\[
\mathcal W_s:=\big[(1-\eps)s,\ s\big],\qquad V_s:=\big[(1-\eps)s-1.0001,\ s+1.0001\big],\qquad |V_s|=\eps s+b_W,\quad b_W:=2.0002,
\]
\[
\cZ(s):=\big\{Z\in\cP:\ \Phi(Z)\cap Y_b\cap\mathcal W_s\ne\emptyset\big\},\qquad
\tilde\cZ(s):=\big\{Z\in\cP:\ \Phi(Z)\cap Y_t\cap(k-\mathcal W_s)\ne\emptyset\big\},
\]
where $k-Y:=\{k-y:y\in Y\}$ for $Y\subset\R$. For every $Z\in\cZ(s)$ we fix a \emph{witness} $y_Z\in\Phi(Z)\cap Y_b\cap\mathcal W_s$.

\emph{Flows.} The master flow is $\Fz=F(\amax,\delta,\hmax)$ (Definition~\ref{def:F0}); R1 is decided in $\Fz$ only. The \emph{R1-passed entries} of $\pi^0_x$ are the entry candidates of $\pi^0_x$ at which $x$ is the winner of R1 (Definition~\ref{def:F0}). For $s\in[y_0,y_1]$ the path $\pi_x$ of the scale flow $\Fs$ (written $\pi^s_x$ in Definition~\ref{def:Fs}) is the path $\pi^0_x$ of $\Fz$, truncated at the first, in path order, of the following events:
\begin{itemize}
\item[(a)] $T_s$: an R1-passed entry $(Y,q)$ of $\pi^0_x$ with $a(Y)\ge\alpha(s)$; the path ends at $q$;
\item[(b)] $H_s$: a gap point at height $s+2$, reached in a ray step with $t_H^{(s)}:=(s+2-p_y)/d_y<\min(t_S,t_W,t_D)$, or an exit point of height $\ge s+2$;
\item[(c)] a termination of $\pi^0_x$ in $\Fz$, which is then inherited.
\end{itemize}
Here $p$, $d$, $t_S$, $t_W$, $t_D$ are the start point, the direction and the parameters of the ray step (Definition~\ref{def:flow}). A T-termination of $\pi^0_x$ in $\Fz$ happens at an R1-passed entry $(Y,q)$ with $a(Y)\ge\amax\ge\alpha(s)$; by Definition~\ref{def:Fs}, if it is reached in $\Fs$ it is a $T_s$ event, so that ``T-terminated in $\Fs$'' and ``$T_s$-terminated in $\Fs$'' mean the same. We put
\[
\cT_s:=\{x\in(0,k):\ \pi_x\text{ is $T_s$-terminated in }\Fs\},\qquad N(s):=|\cT_s| .
\]
For $F\in\{\Fz,\Fs\}$ the \emph{trajectory} $\Gamma_x$ of $\pi_x$ is the union of the segments $\{p+td:0\le t\le t^*\}$ of its ray steps, where $t^*$ is the length of the step in $F$, and of its pass segments $[q,q+n_Y]$; its \emph{termination point} $e(x)$ is its last point in path order, and $\tau_F(x)$ is the height of $e(x)$, except that $\tau_F(x):=h_F$ if $\pi_x$ is H-completed; here $h_{\Fz}=\hmax$ and $h_{\Fs}=s+2$ (Definition~\ref{def:flows-traj}). For $y\in(0,h_F)$ and $\tau_F(x)\ge y$, $\Gamma_x\cap\ell_y$ is a single point $P_y(x)$ (Definition~\ref{def:flows-line}). A path of $F$ \emph{makes an alive contact with $\relint\Bot(Y)$} if one of its ray steps ends in the contact branch of the priority rule at a point $q\in\relint\Bot(Y)$, that is, if $(Y,q)$ is an entry candidate reached by the path; this refers to the moment before R1 is applied.

As in Definitions~\ref{def:flows-losses}, \ref{def:flows-line} and~\ref{def:flows-ovgap}: for $\star\in\{T,D,E,M,W\}$ and $y\in\R$ let $L^F_\star(y):=|\{x:\pi_x\text{ is $\star$-terminated in }F,\ \tau_F(x)<y\}|$, and for $y\in(0,h_F)$
\[
\mu^F_Z(y):=\big|\{x:\ \tau_F(x)>y,\ P_y(x)\in\operatorname{int}Z\}\big|,\qquad \Sh^F(y):=\sum_{S\in\cP}\big(c_S(y)-\mu^F_S(y)\big)_+,
\]
and $\Lambda^F(y):=[L^F_D+L^F_E+L^F_M+L^F_W](y)+\Ovgap^F(y)$. With $\cD$, $\cM$, $\cE$ the sets of floor points that are D-, M- and E-terminated in $\Fz$ and $L_E:=|\cE|$,
\[
\Lambda_0:=|\cD|+|\cM|+L_E+\sup_{0<y<\hmax}\Ovgap^{\Fz}(y),\qquad \Lambda^*(s):=\Lambda_0+2(s+2)\tan\alpha(s).
\]
As noted in Section~\ref{ssec:flows-quantities}, $\Lambda_0\le2k<\infty$.

\emph{Ceiling.} With $\rho(x,y)=(x,k-y)$, the ceiling flows are $\tilde F^0=\rho\circ\Fz[\rho\cP]\circ\rho$ and $\tilde F_s=\rho\circ\Fs[\rho\cP]\circ\rho$, and $\tilde N(s)$, $\tilde\Lambda_0$, $\tilde T^*$ are the quantities $N(s)$, $\Lambda_0$, $\Tstar$ computed for the packing $\rho\cP$ (Definition~\ref{def:flows-ceiling}); $\tilde\Lambda^*(s):=\tilde\Lambda_0+2(s+2)\tan\alpha(s)$.

\begin{theorem}[per-scale count]\label{thm:P7}
Assume the standing assumptions of Section~\ref{ssec:P7-setting}, and let $s\in[y_0,y_1]$ and $b_W=2.0002$.
\begin{enumerate}[label=(\alph*)]
\item $\displaystyle N(s)\ \ge\ \frac{|\cZ(s)|}{\eps s+b_W}-\Lambda^*(s)$, where $\Lambda^*(s)=\Lambda_0+2(s+2)\tan\alpha(s)$.
\item $\displaystyle \tilde N(s)\ \ge\ \frac{|\tilde\cZ(s)|}{\eps s+b_W}-\tilde\Lambda^*(s)$, where $\tilde\Lambda^*(s)=\tilde\Lambda_0+2(s+2)\tan\alpha(s)$.
\end{enumerate}
\end{theorem}

\begin{theorem}[lower bound for $|\cZ(s)|$]\label{thm:P7lower}
Assume the standing assumptions of Section~\ref{ssec:P7-setting}, and let $s\in[y_0,y_1]$.
\begin{enumerate}[label=(\alph*)]
\item We have
\[
\int_{Y_b\cap\mathcal W_s}\big(1-\omega_0-E(y)\big)_+\dd y\ \le\ \sum_{Z\in\cZ(s)}\int_{\Phi(Z)\cap Y_b\cap\mathcal W_s}c_Z(y)\dd y\ \le\ |\cZ(s)|\,\hat\beta\big((1-\eps)s\big),
\]
and consequently
\[
|\cZ(s)|\ \ge\ \frac{((1-\eps)s)^{3/4}}{c_1}\int_{Y_b\cap\mathcal W_s}\big(1-\omega_0-E(y)\big)_+\dd y .
\]
\item $\displaystyle |\tilde\cZ(s)|\ \ge\ \frac{((1-\eps)s)^{3/4}}{c_1}\int_{Y_t\cap(k-\mathcal W_s)}\big(1-\omega_0-E(y)\big)_+\dd y$.
\end{enumerate}
\end{theorem}

All sets of heights occurring here are finite unions of intervals and all integrands are Borel functions (Lemma~\ref{lem:P7-meas}); $\Lambda_0,\tilde\Lambda_0\le2k$. The proofs of the floor parts (a) are given in Sections~\ref{ssec:P7-proofA} and~\ref{ssec:P7-proofB}, and the ceiling parts (b) are deduced by reflection in Section~\ref{ssec:P7-ceiling}.

\subsection{Parameter facts}\label{ssec:P7-params}

\begin{lemma}\label{lem:P7-params}
Let $s\in[y_0,y_1]$.
\begin{enumerate}[label=(\alph*)]
\item If $y\in H$, then $d(y)\ge y_0$ and $0<\beta(y)\le\hat\beta(y_0)\le1.5\cdot10^{-6}$.
\item If $y\in H$, then $d(y)\le(1-\eps)y_1<k/2-1$. Moreover $k/2\notin H$, and $d(y)=y$ for every $y\in H\cap(0,k/2)$.
\item If $Z\in\cZ(s)$, then
\[
0<a(Z)<\beta(y_Z)=\hat\beta(y_Z)\le\min\big\{\hat\beta((1-\eps)s),\ \hat\beta(y_0)\big\}\le1.5\cdot10^{-6}.
\]
\item If $Z\in\cZ(s)$, then $\cos a(Z)+\sin a(Z)\le1+a(Z)<1.0001$.
\item If $Z\in\cZ(s)$ and $0<q_y\le s+1.0001$, then
\[
w(s,Z):=0.50001\,\alpha(s)^2q_y+\delta+\sin a(Z)+\big(1-\cos a(Z)\big)\ <\ w_0-10^{-12}\ <\ w_0 .
\]
\item $\alpha(s)\le\amax$, $s+2\le\hmax$ and $s<k/2$.
\end{enumerate}
\end{lemma}

The quantity $w(s,Z)$ in (e) is the width $w(s,Z;q)$ of Theorem~\ref{thm:P2} at a point $q$ of height $q_y$.

\begin{proof}
(a) By the definition of $H$, $d(y)\ge y_0>0$, so $\beta(y)=c_1d(y)^{-3/4}>0$. As $\hat\beta$ is decreasing, $\beta(y)=\hat\beta(d(y))\le\hat\beta(y_0)$, and $\hat\beta(y_0)\le1.5\cdot10^{-6}$ is (C2).

(b) Since $y_1\ge s\ge y_0>0$ and $\eps>0$, we have $(1-\eps)y_1<y_1=k/2-3<k/2-1$. As $d(k/2)=k/2>(1-\eps)y_1$, the point $k/2$ is not in $H$. For $0<y<k/2$ we have $d(y)=y$.

(c) Since $y_Z\in\Phi(Z)$, $c_Z(y_Z)>0$; thus $\Phi(Z)\ne\emptyset$ and $Z$ meets the line $y_Z$. By [R, Lemma~\ref{imp:L3.4}], $\Phi(Z)=\emptyset$ if $a(Z)=0$; hence $a(Z)>0$. As $y_Z\in Y_{iii}$ and $Z$ meets the line $y_Z$, the definition of kind (iii) gives $a(Z)<\beta(y_Z)$. As $y_Z\in Y_b\subset H\cap(0,k/2)$, part (b) gives $d(y_Z)=y_Z$; hence $\beta(y_Z)=\hat\beta(y_Z)$, and $y_Z=d(y_Z)\ge y_0$ by the definition of $H$. Moreover $y_Z\ge(1-\eps)s>0$ because $y_Z\in\mathcal W_s$. Since $\hat\beta$ is decreasing, $\hat\beta(y_Z)\le\min\{\hat\beta((1-\eps)s),\hat\beta(y_0)\}$, and $\hat\beta(y_0)\le1.5\cdot10^{-6}$ by (C2).

(d) For $a\ge0$ we have $\cos a\le1$ and $\sin a\le a$; and $a(Z)<1.5\cdot10^{-6}$ by (c).

(e) Since $\alpha(s)^2=c_0^2/s$, $q_y\le s+1.0001$ and $s\ge y_0$,
\[
0.50001\,\alpha(s)^2q_y\ \le\ 0.50001\,c_0^2\,\frac{s+1.0001}{s}\ \le\ 0.50001\,c_0^2\Big(1+\frac{1.0001}{y_0}\Big).
\]
Put $a:=a(Z)$. By (c), $0<a<1.5\cdot10^{-6}<2\cdot10^{-4}$. Using $\sin a\le a$ and $1-\cos a\le a^2/2\le10^{-4}a$, then (c) and $s\ge y_0$,
\[
\sin a+(1-\cos a)\ \le\ 1.0001\,a\ <\ 1.0001\,\hat\beta\big((1-\eps)s\big)\ \le\ 1.0001\,\hat\beta\big((1-\eps)y_0\big)=1.0001\,c_1\big((1-\eps)y_0\big)^{-3/4}.
\]
Adding $\delta$ and these two bounds gives $w(s,Z)<w_0-10^{-12}$.

(f) $\alpha$ is decreasing and $s\ge y_0$; and $s\le y_1=k/2-3$.
\end{proof}

\begin{remark}\label{rem:P7-sin}
The inequality $\sin a+(1-\cos a)\le1.0001\,a$ used in Lemma~\ref{lem:P7-params}(e) holds for $0\le a\le2\cdot10^{-4}$, by $\sin a\le a$ and $1-\cos a\le a^2/2$, but not for all inclinations: for $a=\pi/4$ the left side equals $1$, while $1.0001\cdot\pi/4=0.7854767\ldots$. There it is applied only with $0<a<1.5\cdot10^{-6}$.
\end{remark}

\begin{remark}\label{rem:P7-R}
For later use we note that Lemma~\ref{lem:P7-params}(a),(b) give the hypotheses of two lemmas of [R] on the lines of $H$; these lemmas are not used in this section. Let $y\in Y_{ii}$. Then $0\le\beta(y)\le1.5\cdot10^{-6}$, $\omega(y)<\omega_0<\frac12$ by (C4), and the line $y$ meets a square of inclination $\ge\beta(y)$; these are the hypotheses of [R, Lemma~\ref{imp:L4.13}] with $\bar\beta=1.5\cdot10^{-6}$, the set $Y_{ii}$ and the function $\beta$. Let $y\in Y_{iii}$ and $\alpha:=\beta(y)$. Then $0<\alpha\le1.5\cdot10^{-6}$, $y\in(0,k)$, $d(y)\le k/2-1$, and every square meeting the line $y$ has inclination $<\alpha$; as $\delta=10^{-5}$, these are the hypotheses of [R, Lemma~\ref{imp:L4.15}], which gives $E(y)\le0.50001\,\beta(y)\big(A+\beta(y)/\delta\big)W$ with the constant $A$ of [R, Lemma~\ref{imp:L4.13}].
\end{remark}

\subsection{Tilt terminations in the scale flows}\label{ssec:P7-T}

Fix $x\in(0,k)$. At an R1-passed entry $(Y,q)$, $\pi^0_x$ passes $Y$ if $a(Y)<\amax$, and is T-terminated at $q$ otherwise (rule R2, Definition~\ref{def:R2}); R1 is applied before R2 (Definitions~\ref{def:R1} and~\ref{def:R2}). As in Definition~\ref{def:F0}, let
\[
(Z_1,q_1),\ \dots,\ (Z_J,q_J),\qquad J=J(x)\ge0,
\]
be the R1-passed entries of $\pi^0_x$ in path order, and $\eta_j:=(q_j)_y$; only the last one can be a T-termination. These are the entries in the definition of $\Tstar$ (Definition~\ref{def:flows-Tstar}): $s_j:=\max\{y_0,\ \eta_j-2,\ (c_0/a(Z_j))^2\}$ if $a(Z_j)>0$, $s_j:=+\infty$ if $a(Z_j)=0$, and $\Tstar(x):=\min_js_j$ if this minimum is at most $y_1$; otherwise $\Tstar(x)$ is undefined. We write $\{\Tstar\le s\}$ for the set of $x$ for which $\Tstar(x)$ is defined and $\Tstar(x)\le s$.

We use two simple facts about trajectories. First, every segment of a trajectory of $\Fz$ or $\Fs$ is parallel to $n_X$, where $X$ is the floor or a square passed by the path; a passed square has $a(X)<\amax$ by R2, so $n_X\cdot e_2=\cos\varphi_X\ge\cos\amax>0$. Hence the height is nondecreasing along a trajectory and strictly increasing along every segment of positive length (Lemma~\ref{lem:flows-traj}(a)). Second, a ray step ends with D, W, H, E, or at an entry candidate, and the path continues after the step only if the step ends at an R1-passed entry at which the path passes the square and is not completed at its exit point (Definitions~\ref{def:flow}--\ref{def:R2}). Hence, before the first R1-passed entry and between two consecutive ones, $\pi^0_x$ performs exactly one ray step: the step from $(x,0)$ ends at $q_1$, and for $j\ge2$ the step ending at $q_j$ starts at the exit point $q_{j-1}+n_{Z_{j-1}}$ of $Z_{j-1}$.

Lemmas~\ref{lem:P7-Tchar} and~\ref{lem:P7-LT} below restate, in the form used in this section, Lemma~\ref{lem:flows-Ts} (with Lemma~\ref{lem:expansion} for the scale flows) and Corollary~\ref{cor:flows-LT}, under the same conventions; we include the proofs so that this section can be read on its own. Precisely, Lemma~\ref{lem:P7-Tchar}(a), (b) are Lemma~\ref{lem:flows-Ts}(a), (b); Lemma~\ref{lem:P7-Tchar}(c) contains Lemma~\ref{lem:flows-Ts}(d), (e); Lemma~\ref{lem:P7-Tchar}(d) restates Lemma~\ref{lem:flows-Ts}(f), the finite set $\Sigma_s$ constructed there possibly differing by finitely many points from the one of Lemma~\ref{lem:flows-Ts}(f) (only its finiteness and the stated properties are used); and Lemma~\ref{lem:P7-LT} is Corollary~\ref{cor:flows-LT}.

\begin{lemma}\label{lem:P7-Tchar}
Let $s\in[y_0,y_1]$.
\begin{enumerate}[label=(\alph*)]
\item $\eta_1<\eta_2<\dots<\eta_J$. For $j\ge2$, every point of the trajectory of $\pi^0_x$ that precedes $q_j$ has height $<\eta_j$; for $j=1$ such points have height $\le\eta_1$.
\item $\pi_x$ is T-terminated in $\Fs$ if and only if there is $j\le J(x)$ with $a(Z_j)\ge\alpha(s)$ and $\eta_j\le s+2$. In this case the termination point is $q_{j^*}$, where $j^*$ is the least such index. In particular, every T-termination in $\Fs$, including an inherited T-termination of $\Fz$, is the $T_s$-truncation at $q_{j^*}$.
\item $\cT_s=\{x:\ \exists j\le J(x):\ a(Z_j)\ge\alpha(s),\ \eta_j\le s+2\}$, and this set is a finite union of intervals and points. If $y_0\le s\le s'\le y_1$, then $\cT_s\subset\cT_{s'}$. Moreover $\{\Tstar\le s\}=\cT_s$. Consequently $s\mapsto N(s)$ is nondecreasing on $[y_0,y_1]$ and $N(s)=|\{\Tstar\le s\}|$.
\item There is a finite set $\Sigma_s\subset[0,k]$, containing the finite set $\Sigma_0:=\Sigma_{\Fz}$ of Lemma~\ref{lem:structure} for $\Fz$, such that on each component $I$ of $[0,k]\setminus\Sigma_s$ the following holds: the number of ray steps and pass segments of $\pi_x$ in $\Fs$ and its termination kind are constant; for each ray step the source $X$ is constant, and its start point $p(x)$ is an affine function of $x$ with $p'(x)=\lambda u_X$, $|\lambda|\ge1$; and the termination point $e(x)$ and the termination height $\tau_{\Fs}(x)$ are affine functions of $x$. In particular, for $\star\in\{T,D,E,M,W\}$ and $y\in\R$, the set $\{x:\pi_x\text{ is $\star$-terminated in }\Fs,\ \tau_{\Fs}(x)<y\}$ is a finite union of intervals and points.
\end{enumerate}
\end{lemma}

\begin{proof}
(a) Let $j\ge2$. The ray step ending at $q_j$ starts at the exit point $p=q_{j-1}+n_{Z_{j-1}}$ of $Z_{j-1}$. The pass segment $[q_{j-1},p]$ has length $1$, so $p_y>\eta_{j-1}$. The source of the ray step is the square $Z_{j-1}$, which contains $p$ and is disjoint from $Z_j\ni q_j$; hence the segment $[p,q_j]$ has positive length and $\eta_j>p_y$. Points preceding $p$ have height $\le p_y$, and the points of the segment $[p,q_j)$ have height $<\eta_j$. Hence every point preceding $q_j$ has height $<\eta_j$; in particular $\eta_{j-1}<\eta_j$. For $j=1$, the points preceding $q_1$ lie on the vertical segment from $(x,0)$ to $q_1$ and have height $\le\eta_1$.

(b) For a ray step with start point $p$ and direction $d$ we use $t_H^{(s)}=(s+2-p_y)/d_y$. If such a step of $\pi^0_x$ ends in $\Fz$ at a processed contact $q$, then $t_S=t^*$ while neither D nor W applies, so $t_S<t_D$ and $t_S<t_W$; thus $\min(t_S,t_W,t_D)=t_S$, and
\begin{equation}\label{eq:P7-Hs}
\text{the $H_s$ condition $t_H^{(s)}<\min(t_S,t_W,t_D)$ holds in this step}\iff q_y>s+2 .
\end{equation}

($\Leftarrow$) Let $j$ be the least index with $a(Z_j)\ge\alpha(s)$ and $\eta_j\le s+2$. We show that none of the three truncation events of $\Fs$ occurs before $q_j$, and that a $T_s$ event occurs at $q_j$.
\begin{itemize}
\item $\pi^0_x$ is not terminated in $\Fz$ before $q_j$, because it reaches the entry candidate $(Z_j,q_j)$ and wins R1 there.
\item No $T_s$ event occurs before $q_j$: the R1-passed entries before $q_j$ are the $(Z_i,q_i)$ with $i<j$; by (a), $\eta_i<\eta_j\le s+2$, so the minimality of $j$ gives $a(Z_i)<\alpha(s)$.
\item No $H_s$ event occurs before $q_j$. By (a), the exit points preceding $q_j$ (there are none if $j=1$) have height $<\eta_j\le s+2$. The ray steps up to $q_j$ are the steps ending at $q_1,\dots,q_j$; each ends at a processed contact of height $\eta_i\le\eta_j\le s+2$, so by \eqref{eq:P7-Hs} the $H_s$ condition fails in each of them. (If $\eta_i=s+2$, then $t_S=t_H^{(s)}$ and the contact is processed, since contact precedes H.)
\end{itemize}
Hence $\pi_x$ coincides in $\Fs$ with $\pi^0_x$ up to $q_j$, with the same cumulative gap, and reaches the R1-passed entry $(Z_j,q_j)$; the decision of R1 is the one taken in $\Fz$ (Definition~\ref{def:Fs}). As $a(Z_j)\ge\alpha(s)$, $\pi_x$ is $T_s$-terminated at $q_j$.

($\Rightarrow$) Let $\pi_x$ be T-terminated in $\Fs$. By the definition of $\Fs$, its termination point is either a $T_s$-truncation point or a T-termination point of $\Fz$; in both cases it is an R1-passed entry $(Y,q)=(Z_i,q_i)$ of $\pi^0_x$ with $a(Y)\ge\alpha(s)$ (in the second case $a(Y)\ge\amax\ge\alpha(s)$ by Lemma~\ref{lem:P7-params}(f)). The path is not truncated by $H_s$ before reaching $q_i$; in particular the $H_s$ condition fails in the ray step ending at $q_i$, and \eqref{eq:P7-Hs} gives $\eta_i\le s+2$. So the index $i$ has the property in (b). By ($\Leftarrow$), $\pi_x$ is $T_s$-terminated at $q_{j^*}$, $j^*$ the least index with this property. Since the termination point is unique and the $q_j$ are distinct by (a), $i=j^*$. This proves (b), including the last sentence.

(c) The description of $\cT_s$ is (b). On each component $I$ of $[0,k]\setminus\Sigma_0$, the combinatorial type of $\pi^0_x$ is constant, including the outcome of R1 at each entry candidate, and the contact points are affine in $x$ (Lemma~\ref{lem:structure}). Hence $J$ and the squares $Z_j$ are constant on $I$, and each $\eta_j$ is an affine function on $I$. Therefore
\[
\cT_s\cap I=\bigcup_{j:\ a(Z_j)\ge\alpha(s)}\{x\in I:\ \eta_j(x)\le s+2\}
\]
is a finite union of intervals; together with the finite set $\Sigma_0$ this shows that $\cT_s$ is a finite union of intervals and points. If $s\le s'$, then $\alpha(s')\le\alpha(s)$ and $s+2\le s'+2$, so the condition in (b) for $s$ implies the condition for $s'$; thus $\cT_s\subset\cT_{s'}$. Finally let $s\in[y_0,y_1]$. If $\min_js_j\le s$, then $\min_js_j\le y_1$, so $\Tstar(x)$ is defined and equals $\min_js_j$; conversely $\Tstar(x)\le s$ means $\min_js_j\le s$. Hence $x\in\{\Tstar\le s\}$ if and only if $s_j\le s$ for some $j$. Since $y_0\le s$, and since for $a(Z_j)>0$
\[
\Big(\frac{c_0}{a(Z_j)}\Big)^2\le s\iff a(Z_j)\ge c_0s^{-1/2}=\alpha(s),
\]
the condition $s_j\le s$ holds if and only if $a(Z_j)>0$, $\eta_j\le s+2$ and $a(Z_j)\ge\alpha(s)$. Indices with $a(Z_j)=0$ satisfy neither $s_j\le s$ nor $a(Z_j)\ge\alpha(s)$, because $s_j=+\infty$ and $\alpha(s)>0$. So $\{\Tstar\le s\}=\cT_s$ by (b). The last sentence of (c) follows.

(d) Fix a component $I$ of $[0,k]\setminus\Sigma_0$. On $I$, the combinatorial type of $\pi^0_x$ is constant; the start points, the cumulative gaps, the step lengths $t^*_i(x)$, the contact points and the termination point of $\pi^0_x$ are affine in $x$; and a ray step with source $X$ has $p'(x)=\lambda u_X$ with $|\lambda|\ge1$ (Lemmas~\ref{lem:structure} and~\ref{lem:expansion}). By definition, $\pi_x$ is $\pi^0_x$ truncated at the first, in path order, of the following events:
\begin{itemize}
\item a $T_s$ event at an R1-passed entry $(Z_j,q_j)$ with $a(Z_j)\ge\alpha(s)$; the set of such indices $j$ is constant on $I$;
\item an $H_s$ event in the $i$-th ray step. If the event of this step in $\Fz$ is D, W or a contact, then $t^*_i=\min(t_S,t_W,t_D)$ (in $\Fz$ the rule H applies only if $t_H<\min(t_S,t_W,t_D)$, where $t_H=(\hmax-p_y)/d_y$), and the $H_s$ condition is the affine strict inequality $t_H^{(s)}(x)<t^*_i(x)$. If the event of this step in $\Fz$ is H, then $t_H^{(s)}\le t_H<\min(t_S,t_W,t_D)$ because $s+2\le\hmax$ (Lemma~\ref{lem:P7-params}(f)), and the $H_s$ condition holds;
\item an $H_s$ event at the exit point of the $i$-th pass, whose height $q_{i,y}(x)+\cos\varphi_{Y_i}$ is affine; the condition is that this height is $\ge s+2$.
\end{itemize}
If none of these events occurs, $\pi_x$ inherits the termination of $\pi^0_x$, which is then of kind D, W, E or M: an H-termination of $\pi^0_x$ (at a gap point of height $\hmax$, or at an exit point of height $\ge\hmax\ge s+2$) satisfies the condition of the second or the third item, and a T-termination of $\pi^0_x$ occurs at an R1-passed entry $(Z_j,q_j)$ with $a(Z_j)\ge\amax\ge\alpha(s)$, which is an event of the first item. Let $\Sigma_s$ be the union of $\Sigma_0$ and of the zeros, in the components of $[0,k]\setminus\Sigma_0$, of the finitely many affine functions $t_H^{(s)}-t^*_i$ and $q_{i,y}+\cos\varphi_{Y_i}-(s+2)$ that do not vanish identically. On each component of $[0,k]\setminus\Sigma_s$ all the conditions above have constant truth values, so the position (the step) and the kind of the first truncation, or its absence, are constant. The ray steps of $\pi_x$ in $\Fs$ are ray steps of $\pi^0_x$, the last one possibly shortened by $H_s$ at a gap point, so their sources are constant and $p'=\lambda u_X$ with $|\lambda|\ge1$. The termination point is affine: it is $q_j(x)$ for $T_s$, $p_i(x)+t_H^{(s)}(x)\,d_i$ for $H_s$ at a gap point, the exit point $q_i(x)+n_{Y_i}$ for $H_s$ at an exit, and the termination point of $\pi^0_x$ otherwise. The termination height is the height of $e(x)$, or the constant $s+2$ for $H_s$; it is affine. Therefore $\{x:\pi_x\text{ is $\star$-terminated in }\Fs,\ \tau_{\Fs}(x)<y\}$ meets each component of $[0,k]\setminus\Sigma_s$ in an interval or in the empty set, and the last assertion follows since $\Sigma_s$ is finite.
\end{proof}

\begin{lemma}\label{lem:P7-LT}
Let $s\in[y_0,y_1]$. Then $L^{\Fs}_T(y)\le N(s)$ for every $y\in\R$.
\end{lemma}

\begin{proof}
By Lemma~\ref{lem:P7-Tchar}(b), the set of $x$ for which $\pi_x$ is T-terminated in $\Fs$ is $\cT_s$. Hence $L^{\Fs}_T(y)=|\{x\in\cT_s:\tau_{\Fs}(x)<y\}|\le|\cT_s|=N(s)$.
\end{proof}

\begin{remark}\label{rem:P7-Tmax}
Lemma~\ref{lem:P7-LT} uses that an inherited T-termination of $\Fz$ is a $T_s$ event of $\Fs$ (Definition~\ref{def:Fs}) and is therefore counted in $N(s)$. Lemma~\ref{lem:P7-Tchar} uses that R1 is applied before R2, so that T-terminations occur only at R1-passed entries. Lemma~\ref{lem:P7-Tchar}(c) gives the equality $\{\Tstar\le s\}=\cT_s$, and not only an inclusion, together with the monotonicity of $N(s)$; both rest on the fact that R1 is decided once, in $\Fz$, for all scales.
\end{remark}

\subsection{Squares without alive bottom contacts}\label{ssec:P7-avoid}

\begin{lemma}\label{lem:P7-avoid}
Let $s\in[y_0,y_1]$, and let $Y\in\cP$ be a square such that no path of $\Fs$ makes an alive contact with $\relint\Bot(Y)$.
\begin{enumerate}[label=(\alph*)]
\item For every $x\in(0,k)$, the trajectory $\Gamma_x$ of $\pi_x$ in $\Fs$ meets $Y$ in at most one point. If it does, this point is the termination point $e(x)$; it lies in $\partial Y$, it is the end point $p+t_Sd$ of the last ray step of $\pi_x$ and the first point of $Y$ on that ray, and $\pi_x$ is D-, W- or E-terminated there.
\item If $a(Y)>0$, then for every $y\in\R$ the set $\{x:\ e(x)\in\partial Y\cap\ell_y\}$ is finite.
\item $\mu^{\Fs}_Y(y)=0$ for every $y\in(0,s+2)$; more precisely, there is no $x$ with $\tau_{\Fs}(x)>y$ and $P_y(x)\in Y$. If $a(Y)>0$, then also $|\{x:\ \tau_{\Fs}(x)\ge y,\ P_y(x)\in Y\}|=0$ for every $y\in(0,s+2)$.
\end{enumerate}
\end{lemma}

\begin{proof}
(a) Consider a ray step of $\pi_x$ in $\Fs$, with start point $p$, direction $d=n_X$, where $X$ is the floor or the square last passed, parameters $t_S,t_W,t_D$, and $t_H:=t_H^{(s)}$. Its length in $\Fs$ is $t^*=\min(t_S,t_W,t_D,t_H)$: if the $H_s$ condition holds, the step is cut at $t_H<\min(t_S,t_W,t_D)$; otherwise $t_H\ge\min(t_S,t_W,t_D)$, the event of the step in $\Fz$ is not H (as in the proof of Lemma~\ref{lem:P7-Tchar}(d)), and the step ends at $\min(t_S,t_W,t_D)$. In particular $t^*\le t_S$. Write $r(t)=p+td$.

\emph{If $X$ is a square, then $t^*>0$} (this is also part of Lemma~\ref{lem:flows-traj}(b)). Indeed $p\in\relint\Top(X)\subset X$, and:
$t_S>0$, because $X$ is disjoint from each of the finitely many other squares, which are closed, so $r(t)$ lies in no other square for small $t\ge0$;
$t_D=\delta-g>0$, because the contact at which the path entered $X$ was processed, so its cumulative gap $g$ is $<\delta$ (D precedes contact), and a pass does not change $g$;
$t_W>0$, because the abscissa of $p$ lies in $(0,k)$: along $\Top(X)=\{c_X+\frac12n_X+\xi u_X:|\xi|\le\frac12\}$ the abscissa $c_{X,x}-\frac12\sin\varphi_X+\xi\cos\varphi_X$ is strictly increasing in $\xi$ and lies in $[0,k]$ for $\xi=\pm\frac12$, since $X\subset[0,k]^2$, so it lies in $(0,k)$ for $|\xi|<\frac12$;
$t_H>0$, because a ray step starts at an exit point $p$ only if the path was not H-completed there, that is, only if $p_y<s+2$.

We now go through the points of $\Gamma_x$.
\begin{enumerate}[label=(\roman*)]
\item \emph{Points $r(t)$ with $0<t<t^*$.} Since $t<t_S$, $r(t)$ lies in no square other than $X$. If $X$ is a square, then $(r(t)-c_X)\cdot n_X=\frac12+t>\frac12$, so $r(t)\notin X$. Thus $r(t)$ lies in no square.
\item \emph{The start point $p$.} If $X$ is a square, then $p\in X$ and $X\ne Y$: the path passed $X$, which requires an alive contact with $\relint\Bot(X)$, and there is none with $\relint\Bot(Y)$. As $X\cap Y=\emptyset$, $p\notin Y$. If $X$ is the floor, then $p=(x,0)$, and if $p\in Y$, then $t_S=0=t^*$, so that $p$ is also the end point of the step and is treated in (iii).
\item \emph{The end point $r(t^*)$.} If $t^*<t_S$, then $r(t^*)$ lies in no square: by the argument of (i), using $t^*>0$ when $X$ is a square. Let $t^*=t_S$ and $r(t_S)\in Y$. Then $Y$ is the only square containing $r(t_S)$, and $r(t)\notin Y$ for $0\le t<t_S$. Hence $r(t_S)\in\partial Y$: for $t_S>0$ it is a limit of points $r(t)\notin Y$, and for $t_S=0$ it is the point $(x,0)$ of height $0$, which lies on $\partial Y$ because $Y\subset[0,k]^2$. At the parameter $t^*=t_S$ the rules are applied in the order D, W, contact (the rule H and the $H_s$ condition do not apply, since $t^*=t_S$). If the contact branch is reached, then $r(t_S)\notin\relint\Bot(Y)$ by hypothesis, so the event is E. In every case $\pi_x$ terminates at $r(t_S)$, by D, W or E.
\item \emph{Pass segments $[q,q+n_{X'}]$.} They lie in the passed square $X'$, and $X'\ne Y$ as in (ii); so they do not meet $Y$.
\end{enumerate}
The trajectory $\Gamma_x$ is the union of the points considered in (i)--(iv); the truncation points of $\Fs$ are among them ($T_s$ points are entry points into squares $X'\ne Y$, and $H_s$ points are gap points as in (iii) with $t^*<t_S$, or exit points, which lie on pass segments as in (iv)). Hence the only point of $\Gamma_x$ that can lie in $Y$ is a point $r(t_S)$ as in the second case of (iii), and $\pi_x$ terminates there. This proves (a).

(b) Let $a(Y)>0$ and $y\in\R$. Let $I$ be a component of $[0,k]\setminus\Sigma_s$ (Lemma~\ref{lem:P7-Tchar}(d)) and $I_Y:=\{x\in I: e(x)\in Y\}$; assume $I_Y\ne\emptyset$. By (a), for $x\in I_Y$ the path $\pi_x$ is D-, W- or E-terminated at the end point of its last ray step. As the termination kind and the index of the last ray step are constant on $I$, for every $x\in I$ the termination point is the end point $e(x)=p(x)+t^*(x)\,d$ of the last ray step, whose source $X$ and direction $d=n_X$ are constant on $I$; $p$ is affine with $p'=\lambda u_X$, $|\lambda|\ge1$, and $e$ is affine. Since $d\perp u_X$, $e(x)\cdot u_X=p(x)\cdot u_X$ has derivative $\lambda\ne0$; hence $e$ is injective on $I$. For $x\in I_Y$, by (a), $e(x)=p(x)+t_S(x)\,d$ is the first point of $Y$ on the ray $\{p(x)+td:t\ge0\}$.
\begin{enumerate}[label=(\arabic*)]
\item \emph{Vertices.} As $e$ is injective on $I$, at most four $x\in I$ have $e(x)$ at a vertex of $Y$. (This suffices here. In fact no $x\in I$ has $e(x)$ at a vertex: $I$ lies in a piece of $\Fz$, the ray steps of $\pi_x$ in $\Fs$ are initial parts of those of $\pi^0_x$, and by Lemma~\ref{lem:structure}(c) for $\Fz$ the lines carrying them contain no vertex of a square.) Every other $x\in I_Y$ has $e(x)\in\partial Y$ (by (a)) not a vertex, so $e(x)$ lies in the relative interior of one of the four sides of $Y$. For a side $\sigma$ of $Y$ let $I_\sigma:=\{x\in I_Y: e(x)\in\relint\sigma\}$.
\item \emph{If $I_\sigma\ne\emptyset$, then $d$ is not parallel to $\sigma$.} Let $x\in I_\sigma$. First, $p(x)\notin Y$: if $X$ is a square, then $p(x)\in X$ and $X\cap Y=\emptyset$ (as in (a)(ii)); if $X$ is the floor and $p(x)=(x,0)\in Y$, then $t_S=0$ and $e(x)=(x,0)$ is a point of $Y$ at height $0$, i.e.\ a lowest point of $Y$; since $a(Y)>0$, the lowest point of $Y$ is a vertex, contradicting $e(x)\in\relint\sigma$. If $d$ were parallel to $\sigma$, the line of the ray would contain $\sigma$, because it meets $\sigma$ at $e(x)$; that line meets $Y$ exactly in $\sigma$, and a ray starting at $p(x)\notin\sigma$ first meets $\sigma$ at one of its end points, which is a vertex. This contradicts $e(x)\in\relint\sigma$.
\item \emph{An affine parametrization.} Let $I_\sigma\ne\emptyset$, and let $L_\sigma$ be the line containing $\sigma$. For $x\in I$ let $h(x)$ be the intersection point of the line $\{p(x)+td:t\in\R\}$ with $L_\sigma$; it exists and is unique by (2). If $\nu$ is a unit normal of $L_\sigma$ and $L_\sigma=\{z:\nu\cdot z=c\}$, then $h(x)=p(x)+\frac{c-\nu\cdot p(x)}{\nu\cdot d}\,d$, an affine function of $x$. Since $d=n_X\perp u_X$, we have $h(x)\cdot u_X=p(x)\cdot u_X$, whose derivative is $\lambda\ne0$; hence $h'\ne0$, and $h'$ is parallel to $\sigma$ because $h(x)\in L_\sigma$. For $x\in I_\sigma$, $e(x)$ lies on the ray and on $L_\sigma$, so $e(x)=h(x)$.
\item \emph{Monotone height.} The side $\sigma$ is parallel to $u_Y$ or to $n_Y$, and the second components of these vectors are $\sin\varphi_Y$ and $\cos\varphi_Y$. Since $0<a(Y)=|\varphi_Y|\le\pi/4$, both are nonzero. Hence the height of $h(x)$ is a strictly monotone affine function on $I$, and at most one $x\in I_\sigma$ has $e(x)\in\ell_y$.
\end{enumerate}
Thus each component of $[0,k]\setminus\Sigma_s$ contains at most eight points $x$ with $e(x)\in\partial Y\cap\ell_y$: at most four with $e(x)$ at a vertex, and at most one in each $I_\sigma$. There are finitely many components, and $\Sigma_s$ is finite; this proves (b).

(c) Let $y\in(0,s+2)$, $\tau_{\Fs}(x)>y$ and $P_y(x)\in Y$. By (a), $P_y(x)\in\Gamma_x\cap Y$ is the termination point $e(x)$, and $\pi_x$ is D-, W- or E-terminated there, so $\tau_{\Fs}(x)$ is the height of $e(x)$, namely $y$; this contradicts $\tau_{\Fs}(x)>y$. Hence no such $x$ exists, and in particular $\mu^{\Fs}_Y(y)=0$. If $a(Y)>0$, $\tau_{\Fs}(x)\ge y$ and $P_y(x)\in Y$, then in the same way $e(x)=P_y(x)\in\partial Y\cap\ell_y$, and by (b) there are only finitely many such $x$.
\end{proof}

\begin{remark}\label{rem:P7-count}
Part (b) can also be obtained without Lemma~\ref{lem:structure}, with the explicit bound $2(|\cP|+1)$. Since no side of $Y$ is horizontal when $a(Y)>0$, the line $\ell_y$ meets $\partial Y$ in at most two points. By (a), each $x$ in the set of (b) ends a ray step at a point $z\in\partial Y\cap\ell_y$, so $z=p+t\,n_X$ with $p$ the start point and $X$ the source of that step. For fixed $z$ and $X$ there is at most one such $x$. If $X$ is the floor, then $x=z_x$. If $X$ is a square, then $p$ is the unique point of the line of $\Top(X)$ with $z-p$ parallel to $n_X$; so all such paths leave $X$ at the same exit point $p$, hence enter $X$ at the same point $p-n_X\in\relint\Bot(X)$, where only one path is the winner of R1. As there are at most $|\cP|+1$ sources, the set in (b) has at most $2(|\cP|+1)$ elements.
\end{remark}

\begin{remark}\label{rem:P7-a0}
The hypothesis $a(Y)>0$ in (b) and in the last assertion of (c) cannot be omitted. Under the standing assumptions (so that $y_0\le y_1$), let $\cP$ consist of the two squares $X=[1,2]\times[0,1]$ and $Y=[1,2]\times[1+\delta,2+\delta]$, and let $s\in[y_0,y_1]$. For $x\in(1,2)$, the first ray step of $\pi_x$ has $t_S=0$ at the point $(x,0)\in\relint\Bot(X)$; this contact is processed ($t_D=\delta>0$ and $t_W=+\infty$), $x$ is the only path reaching it, and since $a(X)=0<\alpha(s)\le\amax$ the path passes $X$ and leaves it at $(x,1)$ with cumulative gap $0$. In the next ray step $t_S=t_D=\delta$, $t_W=+\infty$ and $t_H^{(s)}=s+1>\delta$, so D applies and $\pi_x$ terminates at $(x,1+\delta)\in\Bot(Y)$. For $x\notin[1,2]$ the vertical ray from $(x,0)$ meets neither square, and for $x\in\{1,2\}$ the path is E-terminated at the vertex $(x,0)$ of $X$. So no path of $\Fs$ makes an alive contact with $\relint\Bot(Y)$, but every $x\in(1,2)$ has $e(x)\in\partial Y\cap\ell_{1+\delta}$, and $|\{x:\tau_{\Fs}(x)\ge1+\delta,\ P_{1+\delta}(x)\in Y\}|=1$. In this example $\mu^{\Fs}_Y\equiv0$, in accordance with the first assertion of (c), which needs no assumption on $a(Y)$. The squares of $\cZ(s)$, to which Lemma~\ref{lem:P7-avoid} is applied below, satisfy $a(Z)>0$ (Lemma~\ref{lem:P7-params}(c)).
\end{remark}

\subsection{Measurability}\label{ssec:P7-meas}

\begin{lemma}\label{lem:P7-meas}
Let $s\in[y_0,y_1]$.
\begin{enumerate}[label=(\alph*)]
\item The set $\cT_s$ and the sets $\{x:\pi_x\text{ is $\star$-terminated in }\Fs,\ \tau_{\Fs}(x)<y\}$, for $\star\in\{T,D,E,M,W\}$ and $y\in\R$, are finite unions of intervals and points; in particular $N(s)\in[0,k]$ and $L^{\Fs}_\star(y)$ are well defined. Moreover $\Lambda_0<\infty$.
\item For every square $S$, $c_S$ is piecewise linear with finitely many pieces. The functions $\omega$ (on $[0,k]$) and $E$ are piecewise linear with finitely many pieces, and the sets $H$, $Y_i$, $Y_{ii}$, $Y_{iii}$, $Y_b$, $Y_t$, $\mathcal W_s$, $V_s$ and $\Phi(S)$ are finite unions of intervals. In particular $y\mapsto(1-\omega_0-E(y))_+$ and $y\mapsto c_S(y)\mathbf 1_{\Phi(S)}(y)$ are Borel functions.
\item On $[y_0,y_1]$ the function $s\mapsto N(s)$ is nondecreasing, hence Borel; the function $s\mapsto|\cZ(s)|$ is a finite sum of indicator functions of finite unions of intervals, hence Borel; and $(y,s)\mapsto\mathbf 1[y\in Y_b\cap\mathcal W_s]$ is a Borel function on $\R\times[y_0,y_1]$.
\end{enumerate}
The same statements hold for the ceiling flows and for $\tilde\cZ(s)$, $Y_t$ and $k-\mathcal W_s$.
\end{lemma}

\begin{proof}
(a) The first assertion is Lemma~\ref{lem:P7-Tchar}(c),(d). The sets $\cD$, $\cM$, $\cE$ are finite unions of intervals and points by Lemma~\ref{lem:structure} (see also Lemma~\ref{lem:flows-meas}(a)), and $\Lambda_0\le2k<\infty$ (Section~\ref{ssec:flows-quantities}).

(b) If $a(S)=0$, then $c_S$ is the indicator function of an interval of length $1$. If $a(S)>0$, then $c_S$ vanishes outside the vertical extent of $S$ and is affine between consecutive heights of vertices of $S$, because each end point of the chord moves along a fixed side between two such heights. Since $\cP$ is finite and the chords $S\cap\ell_y$, $S\in\cP$, are pairwise disjoint subsets of $[0,k]\times\{y\}$, the functions $\omega=k-\sum_Sc_S$ (on $[0,k]$) and $E=\sum_S(c_S-1)_+$ are piecewise linear with finitely many pieces; hence $\{\omega<\omega_0\}$, $\{\omega\ge\omega_0\}$ and $\Phi(S)=\{0<c_S<1\}$ are finite unions of intervals. For a square $S$ let
\[
A_S:=\{y\in(0,k):\ \text{$S$ does not meet the line $y$, or }a(S)<\beta(y)\}.
\]
If $a(S)=0$, then $A_S=(0,k)$, since $\beta>0$. If $a(S)>0$, then $a(S)<\beta(y)$ if and only if $d(y)<(c_1/a(S))^{4/3}$, so
\[
A_S=\big((0,k)\setminus[v_S,\ v_S+\cos a(S)+\sin a(S)]\big)\cup\big\{y\in(0,k):\ d(y)<(c_1/a(S))^{4/3}\big\},
\]
where $v_S$ is the height of the lowest point of $S$; this is a finite union of intervals. The set $\{y\in(0,k):y_0\le d(y)\le(1-\eps)y_1\}$ consists of at most two intervals, and $\{y\in(0,k):\dist(y,\Z)\ge w_0\}$ is a finite union of intervals (one condition in each of the finitely many intervals $[m,m+1]$, $m\in\Z$); so $H$ is a finite union of intervals. Then $Y_{iii}=H\cap\{\omega<\omega_0\}\cap\bigcap_{S\in\cP}A_S$, $Y_{ii}=(H\cap\{\omega<\omega_0\})\setminus Y_{iii}$, $Y_i=H\cap\{\omega\ge\omega_0\}$, $Y_b$ and $Y_t$ are finite unions of intervals. The Borel property of the two functions follows.

(c) The monotonicity of $N$ is Lemma~\ref{lem:P7-Tchar}(c). For $Z\in\cP$ the set $G_Z:=\Phi(Z)\cap Y_b$ is a finite union of intervals by (b). For an interval $J$, the set of $s$ with $J\cap\mathcal W_s\ne\emptyset$ is the intersection of $\{s:(1-\eps)s\le\sup J\}$ and $\{s:s\ge\inf J\}$, with strict inequality where the corresponding end point of $J$ does not belong to $J$; as both end points of $\mathcal W_s$ increase with $s$, this is an interval. Hence $\{s:Z\in\cZ(s)\}$ is a finite union of intervals, and $|\cZ(s)|=\sum_{Z\in\cP}\mathbf 1[Z\in\cZ(s)]$ is Borel. Finally $\mathbf 1[y\in Y_b\cap\mathcal W_s]=\mathbf 1_{Y_b}(y)\,\mathbf 1_C(y,s)$ with the closed set $C=\{(y,s):(1-\eps)s\le y\le s\}$.

The ceiling statements follow in the same way, or by applying (a)--(c) to the packing $\rho\cP$ (Section~\ref{ssec:P7-ceiling}).
\end{proof}

\subsection{The per-scale count}\label{ssec:P7-proofA}

\begin{proof}[Proof of Theorem~\ref{thm:P7}(a)]
Fix $s\in[y_0,y_1]$. If $\cZ(s)=\emptyset$, the claim reads $N(s)\ge-\Lambda^*(s)$, which holds because $N(s)\ge0$ and $\Lambda^*(s)\ge0$. Let $\cZ(s)\ne\emptyset$.

\emph{Step 1: $Z\subset\R\times V_s$ for every $Z\in\cZ(s)$.} Let $a:=a(Z)$, and let $v$ be the height of the lowest point of $Z$. By [R, Lemma~\ref{imp:L3.1}(a)], the vertical extent of $Z$ is $[v,v+\cos a+\sin a]$; it contains $y_Z$ because $c_Z(y_Z)>0$. By Lemma~\ref{lem:P7-params}(d), its length $\cos a+\sin a$ is less than $1.0001$. Hence, as $y_Z\in\mathcal W_s$,
\[
v\ >\ y_Z-1.0001\ \ge\ (1-\eps)s-1.0001,\qquad v+\cos a+\sin a\ <\ y_Z+1.0001\ \le\ s+1.0001 .
\]
Thus $Z\subset\R\times V_s$, and $c_Z(y)=0$ for $y\notin V_s$.

\emph{Step 2: for $Z\in\cZ(s)$, no path of $\Fs$ makes an alive contact with $\relint\Bot(Z)$, and $\mu^{\Fs}_Z(y)=0$ for all $y\in(0,s+2)$.} These assertions are also contained in Theorem~\ref{thm:P2Z}(b),(d), and Step~3 below is Theorem~\ref{thm:P2Z}(e); we derive them here from Theorem~\ref{thm:P2}, using only (C1) and (C2). Suppose that a path of $\Fs$ makes an alive contact with $\relint\Bot(Z)$ at a point $q$.
\begin{itemize}
\item Since $q\in Z$, Step~1 gives $q_y\le s+1.0001\le s+2$.
\item $q_y>0$. Indeed $a(Z)>0$ by Lemma~\ref{lem:P7-params}(c), so $\sin\varphi_Z\ne0$, and along $\Bot(Z)=\{c_Z-\frac12n_Z+\xi u_Z:|\xi|\le\frac12\}$ the height $c_{Z,y}-\frac12\cos\varphi_Z+\xi\sin\varphi_Z$ is a strictly monotone function of $\xi$. Hence a point of $\relint\Bot(Z)$ lies strictly above one of the two end points of $\Bot(Z)$, whose height is $\ge0$ because $Z\subset[0,k]^2$.
\item By Theorem~\ref{thm:P2}(a), which applies because (Q1) follows from (C1), $s\in[y_0,y_1]$, and $q\in\Bot(Z)$ lies on the trajectory of a path of $\Fs$, every $y\in\cR(Z)$ satisfies $\dist(y,\Z)<w(s,Z)$, with $w(s,Z)$ as in Lemma~\ref{lem:P7-params}(e). By the two previous items and Lemma~\ref{lem:P7-params}(e), $w(s,Z)<w_0$.
\item On the other hand, $a(Z)>0$, so $y_Z\in\Phi(Z)\subset\cR(Z)$ by [R, Lemma~\ref{imp:L3.4}]; and $y_Z\in H$, so $\dist(y_Z,\Z)\ge w_0$. This contradicts the previous item.
\end{itemize}
Hence no such contact exists; in particular no path of $\Fs$ enters $Z$. Lemma~\ref{lem:P7-avoid}, applied with $Y=Z$, gives $\mu^{\Fs}_Z(y)=0$ for every $y\in(0,s+2)$; as $a(Z)>0$, the same holds if in the definition of $\mu$ the interior of $Z$ is replaced by $Z$ or the condition $\tau_{\Fs}(x)>y$ by $\tau_{\Fs}(x)\ge y$.

\emph{Step 3: a lower bound for the shadow.} For every $y\in(0,s+2)$, by Step~2 and since all terms are nonnegative,
\[
\Sh^{\Fs}(y)=\sum_{S\in\cP}\big(c_S(y)-\mu^{\Fs}_S(y)\big)_+\ \ge\ \sum_{Z\in\cZ(s)}\big(c_Z(y)-0\big)_+=\sum_{Z\in\cZ(s)}c_Z(y).
\]

\emph{Step 4: the shadow inequality.} Let $\Omega_s:=\Exc\cap(0,s+2)$, where $\Exc=\Exc_{\Fz}$ is the finite set of exceptional heights (Definition~\ref{def:flows-exc} and Lemma~\ref{lem:flows-exc}) off which Theorem~\ref{thm:P4}(f) holds. For every $y\in(0,s+2)\setminus\Omega_s$,
\[
\sum_{Z\in\cZ(s)}c_Z(y)\ \le\ \Sh^{\Fs}(y)\ \le\ L^{\Fs}_T(y)+\Lambda^{\Fs}(y)\ \le\ N(s)+\Lambda_0+2(s+2)\tan\alpha(s)=N(s)+\Lambda^*(s).
\]
The first inequality is Step~3; the second is Theorem~\ref{thm:P4}(f) for $F=\Fs$, whose height bound is $h_{\Fs}=s+2$; the third follows from Lemma~\ref{lem:P7-LT} and Proposition~\ref{prop:wall}, which applies since $y<s+2$.

\emph{Step 5: integration.} Let $c_{\cZ(s)}(y):=\sum_{Z\in\cZ(s)}c_Z(y)$ (the total length of the chords of the squares of $\cZ(s)$), a finite sum of bounded nonnegative Borel functions. By Step~1 and since $Z\subset[0,k]^2$, $c_{\cZ(s)}$ vanishes outside $V_s\cap[0,k]$. As $s+1.0001<k$, we have $V_s\cap[0,k]\subset[0,s+1.0001]\subset[0,s+2)$. Hence, by Step~4, $c_{\cZ(s)}(y)\le N(s)+\Lambda^*(s)$ for every $y\in V_s\cap[0,k]$ outside the finite set $\Omega_s\cup\{0\}$; for $y\in V_s$ with $y<0$ this holds as well, because $c_{\cZ(s)}(y)=0$. Since $\int_\R c_Z=1$ for every square,
\[
|\cZ(s)|=\sum_{Z\in\cZ(s)}\int_\R c_Z(y)\dd y=\int_{V_s}c_{\cZ(s)}(y)\dd y\ \le\ |V_s|\,\big(N(s)+\Lambda^*(s)\big)=(\eps s+b_W)\big(N(s)+\Lambda^*(s)\big).
\]
Dividing by $\eps s+b_W>0$ gives Theorem~\ref{thm:P7}(a).
\end{proof}

\begin{remark}\label{rem:P7-proof}
(1) The proof does not use any measurability of $y\mapsto\Sh^{\Fs}(y)$: it uses the pointwise bound of Step~4 outside a finite set and integrates only the piecewise linear function $c_{\cZ(s)}$.

(2) In Step~1 the number $1.0001$ may be replaced by $\max_{Z\in\cZ(s)}(\cos a(Z)+\sin a(Z))<1+1.5\cdot10^{-6}$. Hence Theorem~\ref{thm:P7}(a) also holds with $b_W=2+3\cdot10^{-6}$; we keep $b_W=2.0002$.

(3) Of the definition of $\cZ(s)$, the proof of Theorem~\ref{thm:P7}(a) uses only the following properties of each $Z\in\cZ(s)$: there is a point $y_Z\in\Phi(Z)\cap H\cap\mathcal W_s$ with $y_Z<k/2$, and $a(Z)<\beta(y_Z)$. The condition $y_Z<k/2$ gives $d(y_Z)=y_Z$, hence $\beta(y_Z)=\hat\beta(y_Z)$ and $y_Z\ge y_0$ (Lemma~\ref{lem:P7-params}(c)); these properties enter the height bound of Step~1 through Lemma~\ref{lem:P7-params}(d) and the inequality $w(s,Z)<w_0$ of Step~2 through Lemma~\ref{lem:P7-params}(e). The condition $\omega(y_Z)<\omega_0$ of kind (iii) is not used here; it is used in Theorem~\ref{thm:P7lower}.
\end{remark}

\begin{remark}\label{rem:P7-C3}
(1) Constraint (C3), which says that $\hat\beta((1-\eps)s)<\alpha(s)$ for all $s\ge y_0$, is not used in this section. Theorem~\ref{thm:P2} applies to every square $Z$ and does not require $a(Z)<\alpha(s)$, and Step~2 shows that a square of $\cZ(s)$ receives no alive contact from $\Fs$ at all; so it does not matter whether R2 would stop a path at such a square.

(2) Constraint (Q2) is not used either: Step~2 uses (C1), through (Q1) in Theorem~\ref{thm:P2}, and (C2), through Lemma~\ref{lem:P7-params}(c)--(e), for every $\eps\in(0,1)$. Only the statement about alive contacts with $\relint\Bot(Z)$ is needed in this paper.
\end{remark}

\subsection{The lower bound for \texorpdfstring{$|\cZ(s)|$}{|Z(s)|}}\label{ssec:P7-proofB}

\begin{proof}[Proof of Theorem~\ref{thm:P7lower}(a)]
Fix $s\in[y_0,y_1]$.

\emph{Step 1: one line.} Let $y\in Y_b\cap\mathcal W_s$. Then $y\in H\subset(0,k)$, so [R, Lemma~\ref{imp:L3.5}] applies ($k\ge2$). Since $\omega(y)<\omega_0$ for $y\in Y_{iii}$,
\[
\sum_{S:\ 0<c_S(y)<1}c_S(y)\ \ge\ 1-\omega(y)-E(y)\ >\ 1-\omega_0-E(y).
\]
The left side is nonnegative, so it is $\ge(1-\omega_0-E(y))_+$.

\emph{Step 2: the squares in the sum.} Every square $S$ in the sum of Step~1 satisfies $y\in\Phi(S)\cap Y_b\cap\mathcal W_s$, so $S\in\cZ(s)$ by definition. Hence, for every $y\in Y_b\cap\mathcal W_s$,
\[
\big(1-\omega_0-E(y)\big)_+\ \le\ \sum_{Z\in\cZ(s)}c_Z(y)\,\mathbf 1_{\Phi(Z)}(y).
\]

\emph{Step 3: integration.} The set $Y_b\cap\mathcal W_s$ is a finite union of intervals and both sides are Borel functions (Lemma~\ref{lem:P7-meas}). Integrating over $Y_b\cap\mathcal W_s$ and interchanging the integral with the finite sum,
\[
\int_{Y_b\cap\mathcal W_s}\big(1-\omega_0-E(y)\big)_+\dd y\ \le\ \sum_{Z\in\cZ(s)}\int_{\Phi(Z)\cap Y_b\cap\mathcal W_s}c_Z(y)\dd y\ \le\ \sum_{Z\in\cZ(s)}\int_{\Phi(Z)}c_Z(y)\dd y .
\]
Each square is counted once here. A square may have short chords on many lines of $Y_b\cap\mathcal W_s$; its contribution is the integral of these chords over all such lines, which is at most $\int_{\Phi(Z)}c_Z$, while $|\cZ(s)|$ counts the square once.

\emph{Step 4: the area of the short chords.} Let $Z\in\cZ(s)$. By Lemma~\ref{lem:P7-params}(c), $0<a(Z)<\hat\beta((1-\eps)s)$, and by [R, Lemma~\ref{imp:L3.4}], $\int_{\Phi(Z)}c_Z=\sin a(Z)\cos a(Z)\le a(Z)$. Hence the right side of Step~3 is at most $|\cZ(s)|\,\hat\beta((1-\eps)s)$; if $\cZ(s)=\emptyset$, all three quantities in Theorem~\ref{thm:P7lower}(a) are $0$. This proves the chain of inequalities. Since $\hat\beta((1-\eps)s)=c_1((1-\eps)s)^{-3/4}>0$, dividing by it gives the lower bound for $|\cZ(s)|$.
\end{proof}

\begin{remark}\label{rem:P7-E}
The term $E(y)$ is kept inside the positive part. In Section~\ref{sec:P8} it is bounded by means of [R, Lemma~\ref{imp:L4.15}], whose hypotheses on the lines of kind (iii) are recorded in Remark~\ref{rem:P7-R}.
\end{remark}

\subsection{The ceiling versions}\label{ssec:P7-ceiling}

\begin{proof}[Proof of Theorems~\ref{thm:P7}(b) and~\ref{thm:P7lower}(b)]
Let $\rho(x,y)=(x,k-y)$ and $\cP':=\rho\cP$. Quantities computed for $\cP'$ are marked with a prime.
\begin{enumerate}[label=(\arabic*)]
\item $\rho$ is an isometry of the plane that maps $[0,k]^2$ onto itself. Hence $\cP'$ is a closed packing of $[0,k]^2$ with the same number of squares and the same waste area.
\item For every square $S$, $\rho S$ is a square with $a(\rho S)=a(S)$, since $\rho$ maps horizontal and vertical directions to horizontal and vertical directions. Since $\rho S\cap\ell_y=\rho(S\cap\ell_{k-y})$, we have $c'_{\rho S}(y)=c_S(k-y)$, and therefore $\Phi(\rho S)=k-\Phi(S)$.
\item The square $\rho S$ meets the line $y$ if and only if $S$ meets the line $k-y$. Moreover $\omega'(y)=\omega(k-y)$, $E'(y)=E(k-y)$, $d(k-y)=d(y)$ and $\beta(k-y)=\beta(y)$.
\item $H$ is invariant under $y\mapsto k-y$: $d$ is invariant, and since $k$ is an integer, $\Z$ is invariant under $t\mapsto k-t$, so $\dist(k-y,\Z)=\dist(y,\Z)$. This is where the integrality of $k$ is used.
\item By (2)--(4), the kind of $y\in(0,k)$ with respect to $\cP'$ is the kind of $k-y$ with respect to $\cP$. Hence $Y'_{iii}=k-Y_{iii}$, and $Y'_b=Y'_{iii}\cap(0,k/2)=k-Y_t$.
\item For $Z\in\cP$: $\rho Z\in\cZ'(s)$ if and only if $(k-\Phi(Z))\cap(k-Y_t)\cap\mathcal W_s\ne\emptyset$, that is, if and only if $\Phi(Z)\cap Y_t\cap(k-\mathcal W_s)\ne\emptyset$. Hence $\cZ'(s)=\rho\,\tilde\cZ(s)$ and $|\cZ'(s)|=|\tilde\cZ(s)|$.
\item By the definition of the ceiling flows, $\tilde N(s)=N'(s)$ and $\tilde\Lambda_0=\Lambda'_0$, so $\tilde\Lambda^*(s)=\Lambda'^*(s)$; the floor coordinate $x$ is the same, as $\rho(x,0)=(x,k)$.
\item All results used in Sections~\ref{ssec:P7-params}--\ref{ssec:P7-proofB} (the lemmas of Section~\ref{sec:flows}, Theorems~\ref{thm:P2} and~\ref{thm:P4}, Proposition~\ref{prop:wall} and the results of [R]) hold for every closed packing of $[0,k]^2$, and the constraints on the parameters do not depend on the packing. Hence Theorems~\ref{thm:P7}(a) and~\ref{thm:P7lower}(a) hold for $\cP'$. Theorem~\ref{thm:P7}(b) follows from (6) and (7). For Theorem~\ref{thm:P7lower}(b), the substitution $y\mapsto k-y$, together with (3) and (5), gives
\[
\int_{Y'_b\cap\mathcal W_s}\big(1-\omega_0-E'(y)\big)_+\dd y=\int_{Y_t\cap(k-\mathcal W_s)}\big(1-\omega_0-E(y)\big)_+\dd y,
\]
and (6) gives $|\cZ'(s)|=|\tilde\cZ(s)|$.\qedhere
\end{enumerate}
\end{proof}

\begin{corollary}\label{cor:P7-ceiling}
Lemmas~\ref{lem:P7-Tchar}, \ref{lem:P7-LT}, \ref{lem:P7-avoid} and~\ref{lem:P7-meas} hold for the ceiling flows. In particular, for $s\in[y_0,y_1]$, the set of $x$ for which $\tilde T^*(x)$ is defined and $\tilde T^*(x)\le s$ is the set of $x$ whose ceiling path is $T_s$-terminated in $\tilde F_s$, the function $s\mapsto\tilde N(s)$ is nondecreasing on $[y_0,y_1]$, and $\tilde N(s)$ is the measure of $\{x:\tilde T^*(x)\le s\}$.
\end{corollary}

\begin{proof}
By the definition of the ceiling flows, these statements are the statements of Lemmas~\ref{lem:P7-Tchar}--\ref{lem:P7-meas} for the closed packing $\rho\cP$, to which the lemmas apply as in step (8) of the preceding proof.
\end{proof}

\section{The integration inequality}\label{sec:P8}

In this section we prove the inequality $(\bigstar)$ of Theorem~\ref{thm:star}. It bounds the waste $W$ from
below by an explicit expression in $k$ and the parameters. Horizontal lines are weighted with the line measure
$\ell$ of Definition~\ref{def:flows-lines}; we write
\[
\ell(Y)=\int_Y d(y)^{-3/4}\dd y,\qquad
\int_Y g\dd\ell:=\int_Y g(y)\,d(y)^{-3/4}\dd y ,
\]
where $Y\subset H$ is a Borel set and $g\ge0$ is a Borel function. The heights $y\in H$ are of three kinds. Lines of
the first kind carry line waste. Lines of the second kind meet a square of large inclination; they are paid for by
[R, Lemma~\ref{imp:L4.13}]. On lines of the third kind the excess of the long chords is paid for by
[R, Lemma~\ref{imp:L4.15}], and the short chords are counted by the per-scale bounds of Section~\ref{sec:P7}; these
bounds are integrated over the scales $s$ against $|\alpha'(s)|\dd s$ and compared with the scale form of
Section~\ref{sec:P6}.

\subsection{Standing assumptions and notation}\label{subsec:P8-notation}

Throughout Sections~\ref{sec:P8}--\ref{sec:constants}, $k\ge4$ is an integer, $\cP$ is a closed packing of
$[0,k]^2$, $N=|\cP|$ and $W=k^2-N$; $\delta=10^{-5}$, $\bbar=1.5\cdot10^{-6}$ and $b_W=2.0002$
(Definition~\ref{def:flows-params}); and the parameters $\pi=(c_0,c_1,y_0,\eps,\omega_0)$ satisfy the constraints
(C0)--(C7) of Definition~\ref{def:flows-constraints} for the given $k$. We use the notation of
Section~\ref{sec:flows} and recall the parts that are needed.

\begin{enumerate}[label=(N\arabic*),leftmargin=*]
\item (Definition~\ref{def:flows-params}.) $\alpha(s)=c_0s^{-1/2}$ for $s>0$ and $\amax=\alpha(y_0)$;
  $\hat\beta(t)=c_1t^{-3/4}$ for $t>0$; $d(y)=\min(y,k-y)$ and $\beta(y)=\hat\beta(d(y))$ for $0<y<k$;
  $b^*=\hat\beta((1-\eps)y_0)$; $y_1=k/2-3$, $\hmax=k/2-1$;
  \[
  w_0=0.50001\,c_0^2\Big(1+\frac{1.0001}{y_0}\Big)+\delta+1.0001\,b^*+10^{-12},\qquad h_0=1-2w_0 .
  \]
\item (Definition~\ref{def:flows-lines}.) With $\frc{y}$ the fractional part of $y$,
  \[
  H=\big\{y\in(0,k):\ y_0\le d(y)\le(1-\eps)y_1,\ \frc{y}\in[w_0,1-w_0]\big\},
  \]
  $H_b=H\cap(0,\tfrac k2)$ and $H_t=H\cap(\tfrac k2,k)$. A height $y\in H$ is of kind (i) if $\omega(y)\ge\omega_0$;
  of kind (ii) if $\omega(y)<\omega_0$ and the line $y$ meets a square $S$ with $a(S)\ge\beta(y)$; of kind (iii) if
  $\omega(y)<\omega_0$ and every square meeting the line $y$ has $a(S)<\beta(y)$. The sets of heights of the three
  kinds are $Y_i,Y_{ii},Y_{iii}$, and $Y_b=Y_{iii}\cap(0,\frac k2)$, $Y_t=Y_{iii}\cap(\frac k2,k)$. For
  $s\in[y_0,y_1]$, $\mathcal W_s=[(1-\eps)s,s]$,
  \[
  \cZ(s)=\{Z\in\cP:\ \Phi(Z)\cap Y_b\cap\mathcal W_s\ne\emptyset\},\qquad
  \tilde\cZ(s)=\{Z\in\cP:\ \Phi(Z)\cap Y_t\cap(k-\mathcal W_s)\ne\emptyset\},
  \]
  where $k-Y=\{k-y:y\in Y\}$ for $Y\subset\R$. Here $\omega(y)$ is the line waste, $c_S(y)$ the length of the chord
  $S\cap(\R\times\{y\})$, $\Phi(S)=\{y:0<c_S(y)<1\}$ and $E(y)=\sum_{S\in\cP:\,c_S(y)\ge1}(c_S(y)-1)$
  ([R, Lemmas~\ref{imp:L3.4} and~\ref{imp:L3.5}]).
\item (Definitions~\ref{def:F0}, \ref{def:Fs} and~\ref{def:flows-ceiling}.) $\Fz$ is the master flow, with paths
  $\pi^0_x$ and R1-passed entries $(Z_j,q_j)$, $j\le J(x)$, of heights $\eta_j=(q_j)_y$. For $s\in[y_0,y_1]$, $\Fs$ is
  the scale flow, with paths $\pi^s_x$; $\cT_s$ is the set of $x$ for which $\pi^s_x$ is T-terminated (equivalently,
  $T_s$-terminated), and $N(s)=|\cT_s|$. The ceiling versions carry a tilde.
\item (Definition~\ref{def:flows-Tstar}.) $s_j=\max(y_0,\eta_j-2,(c_0/a(Z_j))^2)$ if $a(Z_j)>0$ and $s_j=+\infty$ if
  $a(Z_j)=0$; $\Tstar(x)=\min_js_j$ on $\operatorname{dom}\Tstar=\{x:\min_js_j\le y_1\}$; $\tilde T^*$ is its ceiling
  version.
\item (Definition~\ref{def:flows-ovgap}.)
  $\Lam_0=|\cD|+|\cM|+L_E+\sup_{0<y<\hmax}\Ovgap^{\Fz}(y)$ and $\Lam^*(s)=\Lam_0+2(s+2)\tan\alpha(s)$, with ceiling
  versions $\tilde\Lam_0$, $\tilde\Lam^*(s)=\tilde\Lam_0+2(s+2)\tan\alpha(s)$.
\end{enumerate}

\begin{lemma}\label{lem:P8-cons}
Under the standing assumptions:
\begin{enumerate}[label=(\alph*)]
\item $0<\alpha(s)\le\alpha(y_0)\le1.5\cdot10^{-6}$ for every $s\ge y_0$; \textup{(Q1)} and \textup{(C6)} hold;
\item $0<w_0<\frac12$ and $0<h_0<1$;
\item $y_0+1\le(1-\eps)y_1<y_1<\frac k2-1$, so $[y_0,y_1]$ is a nondegenerate interval;
\item for $y\in H$: $y_0\le d(y)\le(1-\eps)y_1<\frac k2-1$ and $0<\beta(y)\le\hat\beta(y_0)\le1.5\cdot10^{-6}$;
  moreover $\frac k2\notin H$, $d(y)=y$ on $H_b$ and $d(y)=k-y$ on $H_t$;
\item $H=Y_i\sqcup Y_{ii}\sqcup Y_{iii}$ and $Y_{iii}=Y_b\sqcup Y_t$ \textup{(}disjoint unions\textup{)}.
\end{enumerate}
\end{lemma}

\begin{proof}
(a) $\alpha$ is decreasing and positive, and $\alpha(y_0)\le1.5\cdot10^{-6}$ is (C1); (Q1) and (C6) follow from (C1)
(Lemma~\ref{lem:flows-constraints}(a)).
(b) $w_0>0$ is clear and $w_0<\frac12$ is (C5).
(c) The first inequality is (C7); then $y_1>0$ and $(1-\eps)y_1<y_1=\frac k2-3$.
(d) The bounds on $d(y)$ are the definition of $H$ and (c); $\beta(y)=c_1d(y)^{-3/4}\le c_1y_0^{-3/4}$, which is at
most $1.5\cdot10^{-6}$ by (C2). As $d(k/2)=k/2>(1-\eps)y_1$, $\frac k2\notin H$; the formulas for $d$ are its
definition.
(e) The three kinds are defined by mutually exclusive conditions which together cover $H$, and
$Y_{iii}\subset H\subset(0,k)\setminus\{\frac k2\}$.
\end{proof}

The constraints enter this section as follows: (C0) everywhere (the integrality of $y_0$ in Step~2 below); (C1) in
Lemma~\ref{lem:P8-cons}(a) and Step~11; (C2) in Steps~4 and~5 (the hypotheses of [R, Lemmas~\ref{imp:L4.13}
and~\ref{imp:L4.15}]); (C4) in Step~4; (C5) in Step~2; (C7) in Lemma~\ref{lem:P8-cons}(c) and Step~2. Constraint (C3)
is not used in Sections~\ref{sec:P8}--\ref{sec:constants}, and neither is (Q2). The results of
Sections~\ref{sec:flows}--\ref{sec:P7} quoted below are stated under hypotheses that are part of the standing
assumptions.

\subsection{The constants}\label{subsec:P8-constants}

We use the constants of Definition~\ref{def:flows-consts}: with $Q_*=0.32$ and $K_w^*=9$,
\begin{align*}
A&=\frac{4\sqrt{K_w^*/Q_*}}{2-3\cdot10^{-6}}&&\text{(the constant of [R, Lemma~\ref{imp:L4.13}])},\\
A'&=\frac{A}{\cos\bbar}+10^{-8}&&\text{(the constant of Corollary~\ref{cor:P6scale})},\\
C_\Lam&=1+1185.6\,\delta^2&&\text{(the constant of Corollary~\ref{cor:Lambda0})},
\end{align*}
and, for the given $\pi$ and $k$,
\begin{align*}
\invF&=\frac{4(1-\eps)^{3/4}\big(1-(1-\eps)^{3/4}\big)}{3\eps\,\big(1+b_W/(\eps y_0)\big)},\qquad
\Gamma=1.00002\Big(2c_1A\,y_0^{-1/2}+\frac45\,\frac{c_1^2}{\delta}\,y_0^{-5/4}\Big),\\
\Omega_W&=\int_{y_0}^{y_1}2(s+2)\tan\alpha(s)\,\frac{c_0}{2}s^{-3/2}\dd s,\qquad
K_W=\frac{4c_1}{c_0\,\invF}\cdot1.000001\,c_0^2,\\
B&=\frac1{\omega_0y_0^{3/4}}+\frac{A}{2c_1}+\frac{2c_1}{c_0\,\invF}\Big(A'+\frac{\alpha(y_0)\,C_\Lam}{\delta}\Big)+\Gamma .
\end{align*}
In addition we put
\begin{gather*}
C_{\rm wall}:=\frac{4c_1}{c_0\,\invF}\quad(\text{so that }K_W=1.000001\,c_0^2\,C_{\rm wall}),\\
\gamma(y):=0.50001\,\beta(y)\Big(A+\frac{\beta(y)}{\delta}\Big)\ \ (y\in H),\qquad
\Gamma_H:=\int_H\gamma\dd\ell .
\end{gather*}
Thus $\Gamma_H$ is an integral, which depends on $k$ through $H$, while $\Gamma$ is the closed-form expression above,
which depends only on $\pi$; Lemma~\ref{lem:P8-Gamma} shows $\Gamma_H\le\Gamma$, and it is $\Gamma$ that enters $B$.
(These are the numbers $\Gamma$ and $\Gamma_H$ of Section~\ref{ssec:flows-lines}; they are not trajectories, which
always carry a floor point as an index.)
$\Omega_W$ depends on $k$ through $y_1$; $\invF$, $\Gamma$, $B$, $C_{\rm wall}$ and $K_W$ depend only on $\pi$ and
the fixed constants.

\begin{lemma}\label{lem:P8-invF}
$0<\invF<\big(1+b_W/(\eps y_0)\big)^{-1}<1$.
\end{lemma}

\begin{proof}
Put $v:=(1-\eps)^{1/4}\in(0,1)$, so that $\eps=1-v^4$ and
$\invF\cdot\big(1+b_W/(\eps y_0)\big)=4v^3(1-v^3)/\big(3(1-v^4)\big)>0$. The polynomial identity
\[
3(1-v^4)-4v^3(1-v^3)=4v^6-3v^4-4v^3+3=(v-1)^2\,(4v^4+8v^3+9v^2+6v+3)
\]
(expand the right-hand side) shows that $3(1-v^4)-4v^3(1-v^3)>0$ for $0<v<1$. Hence
$\invF\cdot(1+b_W/(\eps y_0))<1$, and $1+b_W/(\eps y_0)>1$.
\end{proof}

\subsection{Results used}\label{subsec:P8-used}

Besides the definitions recalled above we use the following results; under the standing assumptions their
hypotheses are satisfied (Lemma~\ref{lem:P8-cons}).

\begin{enumerate}[label=(U\arabic*),leftmargin=*]
\item\label{it:P8-L413} [R, Lemma~\ref{imp:L4.13}]: if $0<\bar\beta_1\le1.5\cdot10^{-6}$, $Y\subset[0,k]$,
  $\beta_1:Y\to[0,\bar\beta_1]$, and for every $y\in Y$ we have $\omega(y)<\frac12$ and the line $y$ meets a square of
  inclination $\ge\beta_1(y)$, then $\int_Y2\beta_1(y)\dd y\le AW$.
\item\label{it:P8-L415} [R, Lemma~\ref{imp:L4.15}]: if $0<\delta\le10^{-5}$, $0<\alpha_1\le1.5\cdot10^{-6}$,
  $y\in(0,k)$, $d(y)\le\frac k2-1$, and every square meeting the line $y$ has inclination $<\alpha_1$, then
  $E(y)\le0.50001\,\alpha_1(A+\alpha_1/\delta)W$. (Its proof in [R] uses [R, Lemma~\ref{imp:L4.14}] with
  $\bar\beta:=\alpha_1$ and $\bar y_x:=d(y)$, and through it the pair estimates [R, Proposition~\ref{imp:P4.5} and
  Lemma~\ref{imp:L4.10}] with $d=\delta$; see Section~\ref{sec:imported}.)
\item\label{it:P8-Lam} Corollary~\ref{cor:Lambda0}: $\Lam_0+\tilde\Lam_0\le C_\Lam\,W/\delta$.
\item\label{it:P8-P6} Corollary~\ref{cor:P6scale}: $\operatorname{dom}\Tstar$ and $\operatorname{dom}\tilde T^*$ are
  Borel sets, $\Tstar$ and $\tilde T^*$ are Borel functions on them with values in $[y_0,y_1]$, and
  \[
  \int_{\operatorname{dom}\Tstar}\alpha(\Tstar(x))\dd x+\int_{\operatorname{dom}\tilde T^*}\alpha(\tilde T^*(x))\dd x\ \le\ A'W .
  \]
\item\label{it:P8-P7} Theorem~\ref{thm:P7}: for every $s\in[y_0,y_1]$,
  \[
  N(s)\ \ge\ \frac{|\cZ(s)|}{\eps s+b_W}-\Lam^*(s),\qquad \tilde N(s)\ \ge\ \frac{|\tilde\cZ(s)|}{\eps s+b_W}-\tilde\Lam^*(s).
  \]
\item\label{it:P8-P7l} Theorem~\ref{thm:P7lower}: for every $s\in[y_0,y_1]$, with $f(y):=(1-\omega_0-E(y))_+$,
  \begin{align*}
  |\cZ(s)|&\ \ge\ G_b(s):=\frac{((1-\eps)s)^{3/4}}{c_1}\int_{Y_b\cap\mathcal W_s}f(y)\dd y,\\
  |\tilde\cZ(s)|&\ \ge\ G_t(s):=\frac{((1-\eps)s)^{3/4}}{c_1}\int_{Y_t\cap(k-\mathcal W_s)}f(y)\dd y .
  \end{align*}
\item\label{it:P8-flows} Lemma~\ref{lem:flows-Ts}(c),(e): every first contact point on the trajectory of $\pi^s_x$
  has height $\le s+2$, and $\{x\in\operatorname{dom}\Tstar:\ \Tstar(x)\le s\}=\cT_s$ for $s\in[y_0,y_1]$; and
  Lemma~\ref{lem:flows-meas}(b),(c): the set $\{(x,s):x\in\operatorname{dom}\Tstar,\ \Tstar(x)\le s\}$ is a Borel subset
  of $\R^2$, and $s\mapsto N(s)$ is non-decreasing. The same holds for the ceiling versions.
\end{enumerate}
Here $|\cZ(s)|$ is the number of squares in $\cZ(s)$. The function $f$ should not be confused with the waste areas
$f_e$ of [R] or with the auxiliary function $f(a)$ of Lemma~\ref{lem:P12-aux}. In \ref{it:P8-L413} and \ref{it:P8-L415} the threshold functions of [R] are renamed $\beta_1$, $\alpha_1$,
since $\beta$ and $\alpha$ denote the thresholds of (N1).

\subsection{The inequality}\label{subsec:P8-thm}

\begin{theorem}[integration inequality]\label{thm:star}
Let $k\ge4$ be an integer, $\cP$ a closed packing of $[0,k]^2$ with waste $W=k^2-|\cP|$, and let
$\pi=(c_0,c_1,y_0,\eps,\omega_0)$ satisfy \textup{(C0)--(C7)}. Then
\begin{equation}\label{eq:P8-star}\tag{$\bigstar$}
(1-\omega_0)\,\ell(H)\ \le\ B\,W+C_{\rm wall}\,\Omega_W .
\end{equation}
Moreover
\begin{align}
\ell(H)&\ \ge\ 8h_0\Big(\big((1-\eps)y_1\big)^{1/4}-(y_0+1)^{1/4}\Big),\label{eq:P8-lH}\\
\Omega_W&\ \le\ 1.000001\,c_0^2\Big(\log\frac{y_1}{y_0}+\frac2{y_0}\Big),\label{eq:P8-Om}
\end{align}
and consequently, with $y_1=\frac k2-3$,
\begin{equation}\label{eq:P8-Psi}
W\ \ge\ \Psi_\pi(k):=\frac1B\Big[\,8h_0(1-\omega_0)\Big(\big((1-\eps)y_1\big)^{1/4}-(y_0+1)^{1/4}\Big)
-K_W\Big(\log\frac{y_1}{y_0}+\frac2{y_0}\Big)\Big].
\end{equation}
\end{theorem}

\begin{remark}\label{rem:P8-larger}
$B$ is increasing in each of $A$, $A'$, $C_\Lam$ (also through $\Gamma$), while $C_{\rm wall}$, $K_W$, $h_0$ and
$\Omega_W$ do not depend on them. Since $W\ge0$, \eqref{eq:P8-star} remains true when $B$ is computed with any
numbers $\bar A\ge A$, $\bar A'\ge A'$, $\bar C_\Lam\ge C_\Lam$ in place of $A,A',C_\Lam$; dividing by the
corresponding $\bar B>0$ shows that \eqref{eq:P8-Psi} also remains true with $\bar B$ in place of $B$.
\end{remark}

\subsection{Proof of Theorem~\ref{thm:star}}\label{subsec:P8-proof}

The proof has twelve steps. The standing assumptions are in force.

\medskip\noindent\textbf{Step 1 (measurability).}
For every square $S$, $c_S$ is a piecewise linear function of $y$ with finitely many pieces (it is continuous if
$a(S)>0$, and the indicator function of a closed interval of length $1$ if $a(S)=0$). The chords of the squares on a
line $y\in[0,k]$ are pairwise disjoint subsets of $[0,k]\times\{y\}$, so $\omega(y)=k-\sum_{S\in\cP}c_S(y)$ there, and
$E(y)=\sum_{S\in\cP}(c_S(y)-1)_+$. Hence $\omega$, $E$ and $f=(1-\omega_0-E)_+$ are Borel functions, and
$0\le f\le1-\omega_0<1$. The set $\{y:y_0\le d(y)\le(1-\eps)y_1\}$ is a union of two closed intervals and
$\{y\in(0,k):\frc{y}\in[w_0,1-w_0]\}$ is a finite union of intervals, so $H$ is a finite union of intervals. For a
square $S$, the set of $y$ such that the line $y$ meets $S$ is a closed interval (the vertical extent of $S$), and the
set of $y\in(0,k)$ with $a(S)\ge\beta(y)$ is empty if $a(S)=0$ and equals $\{y:d(y)\ge(c_1/a(S))^{4/3}\}$ if $a(S)>0$.
Therefore $Y_i,Y_{ii},Y_{iii},Y_b,Y_t$ are Borel sets. Since $d\ge y_0>0$ on $H$, $\ell(Y)\le y_0^{-3/4}k<\infty$ for
every Borel $Y\subset H$. By \ref{it:P8-flows}, the set $\{(x,s):x\in\operatorname{dom}\Tstar,\ \Tstar(x)\le s\}$ is a
Borel subset of $\R^2$, and $N$ is a non-decreasing, hence Borel, function on $[y_0,y_1]$ with values in $[0,k]$; the
same holds for the ceiling versions.

\medskip\noindent\textbf{Step 2 (lower bound for $\ell(H)$; proof of \eqref{eq:P8-lH}).}
\emph{Reflection.} The map $y\mapsto k-y$ maps $(0,\frac k2)$ onto $(\frac k2,k)$ and preserves $d$. If $y\notin\Z$,
then $\frc{k-y}=1-\frc{y}$ because $k$ is an integer, so $\frc{y}\in[w_0,1-w_0]$ if and only if
$\frc{k-y}\in[w_0,1-w_0]$; if $y\in\Z$, then $\frc{y}=\frc{k-y}=0\notin[w_0,1-w_0]$ because $w_0>0$. Hence
$H_t=k-H_b$, and since $d$ is preserved, $\ell(H_t)=\ell(H_b)$. As $\frac k2\notin H$
(Lemma~\ref{lem:P8-cons}(d)), $\ell(H)=2\,\ell(H_b)$.

\emph{Unit intervals.} Put $j_0:=\lceil y_0\rceil=y_0$ (an integer by (C0)) and $j_1:=\lfloor(1-\eps)y_1\rfloor$. Let
$m$ be an integer with $j_0\le m\le j_1-1$. Then $[m+w_0,\,m+1-w_0]\subset H_b$: indeed $m+w_0>j_0=y_0$;
$m+1-w_0\le j_1-w_0<j_1\le(1-\eps)y_1<\frac k2$, so $d(y)=y\in[y_0,(1-\eps)y_1]$ on this interval; and since
$0<w_0<\frac12$ (Lemma~\ref{lem:P8-cons}(b)), every $y$ in it satisfies $m<y<m+1$ and $\frc{y}=y-m\in[w_0,1-w_0]$.

\emph{Comparison.} Since $y^{-3/4}$ is decreasing and the interval has length $1-2w_0=h_0>0$,
\[
\int_{m+w_0}^{m+1-w_0}y^{-3/4}\dd y\ \ge\ h_0\,(m+1-w_0)^{-3/4}\ \ge\ h_0\,(m+1)^{-3/4}\ \ge\ h_0\int_{m+1}^{m+2}y^{-3/4}\dd y .
\]

\emph{Summation.} The intervals $[m+w_0,m+1-w_0]$, $j_0\le m\le j_1-1$, are pairwise disjoint, so
\[
\ell(H_b)\ \ge\ h_0\int_{j_0+1}^{j_1+1}y^{-3/4}\dd y\ =\ 4h_0\Big((j_1+1)^{1/4}-(j_0+1)^{1/4}\Big).
\]
If $j_1\le j_0$ the sum is empty, but then the right-hand side is $\le0$ and the inequality still holds. Finally
$j_1+1>(1-\eps)y_1$, $j_0+1=y_0+1$ and $h_0>0$, so $\ell(H)=2\ell(H_b)\ge8h_0\big(((1-\eps)y_1)^{1/4}-(y_0+1)^{1/4}\big)$,
which is \eqref{eq:P8-lH}.

\medskip\noindent\textbf{Step 3 (kind (i)).}
On $H$ we have $d(y)\ge y_0$, hence $d(y)^{-3/4}\le y_0^{-3/4}$. Therefore
\[
W=\int_0^k\omega(y)\dd y\ \ge\ \int_{Y_i}\omega(y)\dd y\ \ge\ \omega_0|Y_i|\ \ge\ \omega_0\,y_0^{3/4}\,\ell(Y_i),
\qquad\text{that is}\qquad \ell(Y_i)\le\frac{W}{\omega_0y_0^{3/4}} .
\]

\medskip\noindent\textbf{Step 4 (kind (ii)).}
Apply \ref{it:P8-L413} with $\bar\beta_1:=1.5\cdot10^{-6}$, $Y:=Y_{ii}$ and $\beta_1:=\beta|_{Y_{ii}}$. By
Lemma~\ref{lem:P8-cons}(d), $\beta_1$ takes values in $(0,1.5\cdot10^{-6}]$. For $y\in Y_{ii}$ we have
$\omega(y)<\omega_0<\frac12$ by (C4), and the line $y$ meets a square of inclination $\ge\beta(y)$ by definition of
kind (ii). Since $\beta(y)=c_1d(y)^{-3/4}$,
\[
2c_1\,\ell(Y_{ii})=\int_{Y_{ii}}2\beta(y)\dd y\ \le\ AW,\qquad\text{that is}\qquad \ell(Y_{ii})\le\frac{A}{2c_1}\,W .
\]

\medskip\noindent\textbf{Step 5 (kind (iii): the excess of the long chords).}

\begin{lemma}\label{lem:P8-Gamma}
$\Gamma_H\le\Gamma$, and $E(y)\le\gamma(y)\,W$ for every $y\in Y_{iii}$.
\end{lemma}

\begin{proof}
Let $y\in Y_{iii}$. Apply \ref{it:P8-L415} with $\alpha_1:=\beta(y)$: we have $\delta=10^{-5}$,
$0<\beta(y)\le1.5\cdot10^{-6}$ and $d(y)\le(1-\eps)y_1<\frac k2-1$ by Lemma~\ref{lem:P8-cons}(d), and every square meeting
the line $y$ has inclination $<\beta(y)$ by definition of kind (iii). This gives $E(y)\le\gamma(y)W$.

On $H_b$ we have $d(y)=y\ge y_0$, so
$\gamma(y)\,d(y)^{-3/4}=0.50001\,\big(Ac_1\,y^{-3/2}+c_1^2\delta^{-1}y^{-9/4}\big)$ and
\[
\int_{H_b}\gamma\dd\ell\ \le\ 0.50001\int_{y_0}^{\infty}\big(Ac_1\,y^{-3/2}+c_1^2\delta^{-1}y^{-9/4}\big)\dd y
=0.50001\Big(2Ac_1\,y_0^{-1/2}+\frac45\,\frac{c_1^2}{\delta}\,y_0^{-5/4}\Big).
\]
The substitution $y\mapsto k-y$ maps $H_t$ onto $H_b$ (Step~2) and preserves $d$, hence $\gamma$ and $\ell$; so
$\int_{H_t}\gamma\dd\ell=\int_{H_b}\gamma\dd\ell$. Since $\frac k2\notin H$ and $2\cdot0.50001=1.00002$,
$\Gamma_H=\int_{H_b}\gamma\dd\ell+\int_{H_t}\gamma\dd\ell\le\Gamma$.
\end{proof}

For $y\in Y_{iii}$ we have $1-\omega_0-E(y)\le f(y)$, so by Lemma~\ref{lem:P8-Gamma},
$1-\omega_0\le f(y)+\gamma(y)W$. Integrating over $Y_{iii}$ with respect to $\ell$ and using $\gamma\ge0$,
$Y_{iii}\subset H$ and Lemma~\ref{lem:P8-Gamma},
\begin{equation}\label{eq:P8-iii}
(1-\omega_0)\,\ell(Y_{iii})\ \le\ \int_{Y_{iii}}f\dd\ell+\Gamma_H\,W\ \le\ \int_{Y_{iii}}f\dd\ell+\Gamma\,W .
\end{equation}
This is the integral of a pointwise inequality in which the same $W$ appears for every $y$; no disjointness is used.

\medskip\noindent\textbf{Step 6 (the lower bounds for $|\cZ(s)|$ and $|\tilde\cZ(s)|$).}
By \ref{it:P8-P7l}, $|\cZ(s)|\ge G_b(s)$ and $|\tilde\cZ(s)|\ge G_t(s)$ for every $s\in[y_0,y_1]$ (these bounds rest on
[R, Lemmas~\ref{imp:L3.4} and~\ref{imp:L3.5}]). The functions $G_b$ and $G_t$ are continuous on $[y_0,y_1]$. Indeed,
since $0\le f\le1$, for $s,s'\in[y_0,y_1]$
\[
\Big|\int_{Y_b\cap\mathcal W_s}f-\int_{Y_b\cap\mathcal W_{s'}}f\Big|\ \le\ \big|\mathcal W_s\,\triangle\,\mathcal W_{s'}\big|
\ \le\ (1-\eps)|s-s'|+|s-s'|\ \le\ 2|s-s'| ,
\]
and the same holds for the windows $k-\mathcal W_s$; the factor $((1-\eps)s)^{3/4}/c_1$ is continuous. In particular
$G_b$ and $G_t$ are bounded Borel functions on $[y_0,y_1]$.

\medskip\noindent\textbf{Step 7 (the inclusion $\cT_s\subset\{\Tstar\le s\}$).}
By \ref{it:P8-flows}, $\{x\in\operatorname{dom}\Tstar:\Tstar(x)\le s\}=\cT_s$ for $s\in[y_0,y_1]$. Only the inclusion
$\cT_s\subset\{\Tstar\le s\}$ is used below; we recall why it holds. Let $x\in\cT_s$. The path $\pi^s_x$ coincides with
$\pi^0_x$ up to its termination, and it is T-terminated at an R1-passed entry $(Z_j,q_j)$ of $\pi^0_x$ (by the definition
of $T_s$, and, for a T-termination inherited from $\Fz$, because R1 is applied before R2,
Definitions~\ref{def:R1} and~\ref{def:R2}); moreover $a(Z_j)\ge\alpha(s)=c_0s^{-1/2}>0$, that is $(c_0/a(Z_j))^2\le s$.
The point $q_j$ is a first contact point on the trajectory of $\pi^s_x$, so $\eta_j\le s+2$ by \ref{it:P8-flows}. Since
also $y_0\le s$, we get $s_j=\max(y_0,\eta_j-2,(c_0/a(Z_j))^2)\le s\le y_1$; hence $x\in\operatorname{dom}\Tstar$ and
$\Tstar(x)\le s$. Consequently
\[
\big|\{x\in\operatorname{dom}\Tstar:\ \Tstar(x)\le s\}\big|\ \ge\ N(s),\qquad
\big|\{x\in\operatorname{dom}\tilde T^*:\ \tilde T^*(x)\le s\}\big|\ \ge\ \tilde N(s)
\]
(the second by the same argument for the reflected packing).

\medskip\noindent\textbf{Step 8 (first use of Tonelli's theorem).}
For $s>0$ we have $|\alpha'(s)|=\frac{c_0}{2}s^{-3/2}$. If $\tau\in[y_0,y_1]$, then
\[
\int_{y_0}^{y_1}\mathbf 1[\tau\le s]\,\frac{c_0}{2}s^{-3/2}\dd s=\int_{\tau}^{y_1}\frac{c_0}{2}s^{-3/2}\dd s
=c_0\big(\tau^{-1/2}-y_1^{-1/2}\big)\ \le\ c_0\tau^{-1/2}=\alpha(\tau).
\]
Since $\Tstar$ takes values in $[y_0,y_1]$, applying this with $\tau=\Tstar(x)$, integrating over
$x\in\operatorname{dom}\Tstar$ and using Tonelli's theorem for the nonnegative Borel function
$(x,s)\mapsto\mathbf 1[\Tstar(x)\le s]\,\frac{c_0}2s^{-3/2}$ on $\operatorname{dom}\Tstar\times[y_0,y_1]$ (Step~1), and then
Step~7, we get
\[
\int_{\operatorname{dom}\Tstar}\alpha(\Tstar(x))\dd x\ \ge\ \int_{y_0}^{y_1}\big|\{\Tstar\le s\}\big|\,\frac{c_0}2s^{-3/2}\dd s
\ \ge\ \int_{y_0}^{y_1}N(s)\,\frac{c_0}2s^{-3/2}\dd s .
\]
With the same argument for $\tilde T^*$ and \ref{it:P8-P6},
\begin{equation}\label{eq:P8-Ton1}
A'W\ \ge\ \int_{y_0}^{y_1}\big(N(s)+\tilde N(s)\big)\,\frac{c_0}{2}s^{-3/2}\dd s .
\end{equation}

\medskip\noindent\textbf{Step 9 (the per-scale bound; the loss and wall terms).}
Let $s\in[y_0,y_1]$. By \ref{it:P8-P7} and Step~6,
\[
N(s)+\tilde N(s)\ \ge\ \Theta(s):=\frac{G_b(s)+G_t(s)}{\eps s+b_W}-(\Lam_0+\tilde\Lam_0)-4(s+2)\tan\alpha(s) .
\]
The right-hand side may be negative; we do not take positive parts but integrate the lower bound as it stands. By
\ref{it:P8-Lam}, $\Lam_0+\tilde\Lam_0$ is finite, so $\Theta$ is continuous and bounded on $[y_0,y_1]$ (Step~6), and we may
multiply by $\frac{c_0}2s^{-3/2}$, integrate over $[y_0,y_1]$ and split the integral into three terms. The loss term
satisfies
\[
(\Lam_0+\tilde\Lam_0)\int_{y_0}^{y_1}\frac{c_0}{2}s^{-3/2}\dd s=(\Lam_0+\tilde\Lam_0)\,c_0\big(y_0^{-1/2}-y_1^{-1/2}\big)
\ \le\ \alpha(y_0)\,(\Lam_0+\tilde\Lam_0)\ \le\ \alpha(y_0)\,C_\Lam\,\frac{W}{\delta}
\]
by \ref{it:P8-Lam}, and the wall term is
$\int_{y_0}^{y_1}4(s+2)\tan\alpha(s)\,\frac{c_0}2s^{-3/2}\dd s=2\,\Omega_W$ (one $\Omega_W$ for each family). With
\eqref{eq:P8-Ton1} and
\[
\mathcal J:=\int_{y_0}^{y_1}\frac{G_b(s)+G_t(s)}{\eps s+b_W}\,\frac{c_0}{2}s^{-3/2}\dd s
\]
we obtain
\begin{equation}\label{eq:P8-J}
\mathcal J\ \le\ \Big(A'+\frac{\alpha(y_0)\,C_\Lam}{\delta}\Big)W+2\,\Omega_W .
\end{equation}

\medskip\noindent\textbf{Step 10 (second use of Tonelli's theorem; the constant $\invF$).}
By Lemma~\ref{lem:P8-cons}(d), every $y\in Y_b$ satisfies $y_0\le y=d(y)\le(1-\eps)y_1$, so
\[
\{s\in[y_0,y_1]:\ (1-\eps)s\le y\le s\}=\big[y,\ y/(1-\eps)\big].
\]
The function $(y,s)\mapsto\mathbf 1_{Y_b}(y)\,\mathbf 1[(1-\eps)s\le y\le s]\,f(y)\,s^{3/4}s^{-3/2}/(\eps s+b_W)$ is a
nonnegative Borel function on $\R\times[y_0,y_1]$ (the set $\{(y,s):(1-\eps)s\le y\le s\}$ is closed), so by Tonelli's
theorem
\[
\int_{y_0}^{y_1}\frac{G_b(s)}{\eps s+b_W}\,\frac{c_0}{2}s^{-3/2}\dd s
=\frac{c_0(1-\eps)^{3/4}}{2c_1}\int_{Y_b}f(y)\Big[\int_y^{y/(1-\eps)}\frac{s^{-3/4}}{\eps s+b_W}\dd s\Big]\dd y .
\]
For $s\ge y\ge y_0$ we have $\eps s+b_W=\eps s\,\big(1+b_W/(\eps s)\big)\le\eps s\,\big(1+b_W/(\eps y_0)\big)$, hence
\[
\int_y^{y/(1-\eps)}\frac{s^{-3/4}}{\eps s+b_W}\dd s\ \ge\ \frac{1}{\eps\,(1+b_W/(\eps y_0))}\int_y^{y/(1-\eps)}s^{-7/4}\dd s
=\frac{4\big(1-(1-\eps)^{3/4}\big)}{3\eps\,(1+b_W/(\eps y_0))}\;y^{-3/4} .
\]
Multiplying by $(1-\eps)^{3/4}$ gives $\invF\cdot y^{-3/4}=\invF\cdot d(y)^{-3/4}$. Hence the bottom part of $\mathcal J$ is at
least $\frac{c_0\,\invF}{2c_1}\int_{Y_b}f\dd\ell$. For the top part, write $u:=k-y=d(y)$ for $y\in Y_t$; then
$y\in k-\mathcal W_s$ if and only if $(1-\eps)s\le u\le s$, and $y_0\le u\le(1-\eps)y_1$; the same computation gives at
least $\frac{c_0\,\invF}{2c_1}\int_{Y_t}f\dd\ell$. Since $Y_{iii}=Y_b\sqcup Y_t$ (Lemma~\ref{lem:P8-cons}(e)),
\begin{equation}\label{eq:P8-Ton2}
\mathcal J\ \ge\ \frac{c_0\,\invF}{2c_1}\int_{Y_{iii}}f\dd\ell .
\end{equation}

\medskip\noindent\textbf{Step 11 (upper bound for $\Omega_W$; proof of \eqref{eq:P8-Om}).}
The function $\tan x/x$ is increasing on $(0,\pi/2)$ (the Taylor coefficients of $\tan$ are positive), and
$\tan(1.5\cdot10^{-6})\le1.000001\cdot1.5\cdot10^{-6}$; hence $\tan x\le1.000001\,x$ for $0\le x\le1.5\cdot10^{-6}$. For
$s\ge y_0$ we have $0<\alpha(s)\le1.5\cdot10^{-6}$ (Lemma~\ref{lem:P8-cons}(a)), so
\[
2(s+2)\tan\alpha(s)\,\frac{c_0}{2}s^{-3/2}\ \le\ 1.000001\,c_0^2\,(s+2)\,s^{-2}=1.000001\,c_0^2\Big(\frac1s+\frac2{s^2}\Big).
\]
Integrating over $[y_0,y_1]$ gives
\[
\Omega_W\ \le\ 1.000001\,c_0^2\Big(\log\frac{y_1}{y_0}+2\big(y_0^{-1}-y_1^{-1}\big)\Big)\ \le\ 1.000001\,c_0^2\Big(\log\frac{y_1}{y_0}+\frac{2}{y_0}\Big).
\]

\medskip\noindent\textbf{Step 12 (summation).}
From \eqref{eq:P8-J} and \eqref{eq:P8-Ton2}, and since $\invF>0$ (Lemma~\ref{lem:P8-invF}),
\begin{equation}\label{eq:P8-f}
\int_{Y_{iii}}f\dd\ell\ \le\ \frac{2c_1}{c_0\,\invF}\Big(A'+\frac{\alpha(y_0)\,C_\Lam}{\delta}\Big)W+C_{\rm wall}\,\Omega_W .
\end{equation}
By Lemma~\ref{lem:P8-cons}(e), $\ell(H)=\ell(Y_i)+\ell(Y_{ii})+\ell(Y_{iii})$; as $0<1-\omega_0\le1$, Steps~3 and~4,
\eqref{eq:P8-iii} and \eqref{eq:P8-f} give
\begin{align*}
(1-\omega_0)\,\ell(H)&\ \le\ \ell(Y_i)+\ell(Y_{ii})+(1-\omega_0)\,\ell(Y_{iii})\\
&\ \le\ \Big[\frac{1}{\omega_0y_0^{3/4}}+\frac{A}{2c_1}+\frac{2c_1}{c_0\,\invF}\Big(A'+\frac{\alpha(y_0)C_\Lam}{\delta}\Big)+\Gamma\Big]W
+C_{\rm wall}\,\Omega_W .
\end{align*}
The bracket is $B$, so \eqref{eq:P8-star} holds. Inserting \eqref{eq:P8-lH} (the coefficient $(1-\omega_0)$ of $\ell(H)$
is positive) and \eqref{eq:P8-Om} (the coefficient $C_{\rm wall}$ of $\Omega_W$ is positive) into \eqref{eq:P8-star} and
dividing by $B>0$ gives \eqref{eq:P8-Psi}, because $C_{\rm wall}\cdot1.000001\,c_0^2=K_W$. This completes the proof of
Theorem~\ref{thm:star}.\qed

\subsection{Remarks on the proof}\label{subsec:P8-remarks}

\begin{remark}[directions of the inequalities]\label{rem:P8-directions}
Step~2 gives a lower bound for $\ell(H)$, which enters \eqref{eq:P8-star} with the positive coefficient $1-\omega_0$.
Steps~3--5 give upper bounds for $\ell(Y_i)$, $\ell(Y_{ii})$ and $\int_{Y_{iii}}E\dd\ell$. Step~6 gives lower bounds for
$|\cZ(s)|$, $|\tilde\cZ(s)|$; Step~7 the inequality $|\{\Tstar\le s\}|\ge N(s)$; Step~8 bounds the left-hand side of
\ref{it:P8-P6} from below; in Step~9 the subtracted loss and wall terms are bounded from above; Step~10 bounds $\mathcal J$ from
below, using $1/(\eps s+b_W)\ge1/\big(\eps s(1+b_W/(\eps y_0))\big)$ for $s\ge y_0$; Step~11 bounds $\Omega_W$ from above
(using $\tan x\le1.000001x$ and dropping $-2/y_1$). The only positive part taken is $f=(1-\omega_0-E)_+$, in
\ref{it:P8-P7l}; in Step~5 it is used through $1-\omega_0-E\le f$.
\end{remark}

\begin{remark}[how the waste is used]\label{rem:P8-bookkeeping}
The upper bounds in terms of $W$ that enter $B$ are: the line waste in Step~3 ($\int_{Y_i}\omega\le\int_0^k\omega=W$),
[R, Lemma~\ref{imp:L4.13}] in Step~4, [R, Lemma~\ref{imp:L4.15}] applied separately to each $y\in Y_{iii}$ in Step~5,
and Corollaries~\ref{cor:P6scale} and~\ref{cor:Lambda0} in Steps~8--9 (these two together give the third term of $B$).
Each holds with the waste $W$ of the whole packing, and \eqref{eq:P8-star} is obtained by adding them. No disjointness
between the parts of the waste used by different bounds is claimed or needed; the price is that their contributions to
$B$ are added. The floor and ceiling families are combined within Corollaries~\ref{cor:Lambda0} and~\ref{cor:P6scale},
each of which bounds a sum over both families by one multiple of $W$, and each of these two corollaries is used once;
Theorems~\ref{thm:P7} and~\ref{thm:P7lower} are applied to each family separately. The wall term $C_{\rm wall}\Omega_W$
does not involve $W$.
\end{remark}

\begin{remark}[integrality of $k$]\label{rem:P8-integer}
In this section the integrality of $k$ is used in Step~2 (the reflection $y\mapsto k-y$ preserves the condition
$\frc{y}\in[w_0,1-w_0]$) and, for the ceiling family, through the results quoted in \ref{it:P8-P7} and \ref{it:P8-P7l}.
\end{remark}

\section{Proof of \texorpdfstring{Theorems~\ref{thm:A} and~\ref{thm:A13} and of Corollary~\ref{cor:A}}{Theorems A and A13 and of Corollary A}}\label{sec:assembly}

By Theorem~\ref{thm:star}, every closed packing of $[0,k]^2$ satisfies $W\ge\Psi_\pi(k)$ whenever $\pi$ satisfies
(C0)--(C7) for $k$. In this section we choose explicit parameters, show that the bound stays above $c_*k^{1/4}$ for
all $k$ beyond a threshold (Proposition~\ref{prop:monotone}), and deduce Theorems~\ref{thm:A} and~\ref{thm:A13} and
Corollary~\ref{cor:A}. The numerical values used are listed in Section~\ref{sec:constants}, together with the direction
in which each displayed value is rounded; they are recomputed by the script described in
Section~\ref{subsec:const-script}. The notation is that of Sections~\ref{sec:flows} and~\ref{sec:P8}. As in
Section~\ref{sec:intro}, $M(k)$ is the largest number of squares in a closed packing of $[0,k]^2$ (it exists, since
$N\le k^2$), and $W_{\min}(k)=k^2-M(k)$ is the least waste of a closed packing of $[0,k]^2$.

\subsection{Upper bounds for the constants}\label{subsec:assembly-upper}

Put
\[
\bar A:=10.6067,\qquad \bar A':=10.6068,\qquad \bar C_\Lam:=1.0000002 .
\]
By Table~\ref{tab:const-fixed}, $A<\bar A$, $A'<\bar A'$ and $C_\Lam<\bar C_\Lam$. Let $\bar B$ denote the number $B$
of Definition~\ref{def:flows-consts} computed with $\bar A,\bar A',\bar C_\Lam$ in place of $A,A',C_\Lam$ (also in
$\Gamma$), and let $\bar\Psi_\pi(k)$ be $\Psi_\pi(k)$ of \eqref{eq:P8-Psi} with $\bar B$ in place of $B$.

\begin{lemma}\label{lem:assembly-upper}
If $\pi$ satisfies \textup{(C0)--(C7)} for the integer $k\ge4$, then every closed packing of $[0,k]^2$ satisfies
$W\ge\bar\Psi_\pi(k)$.
\end{lemma}

\begin{proof}
This is Theorem~\ref{thm:star} together with Remark~\ref{rem:P8-larger}.
\end{proof}

All explicit bounds below are computed with $\bar B$. In the asymptotic statements (Section~\ref{subsec:assembly-asym})
the exact constants $A$, $A'$, $C_\Lam$ are used, which is allowed by Theorem~\ref{thm:star} itself.

\subsection{Monotonicity in \texorpdfstring{$k$}{k}}\label{subsec:assembly-mono}

For real $k$ with $y_0+1\le(1-\eps)(\frac k2-3)$ we use the right-hand side of \eqref{eq:P8-Psi} as the definition of
$\Psi_\pi(k)$ (with $y_1=y_1(k):=\frac k2-3$), and likewise for $\bar\Psi_\pi(k)$.

\begin{proposition}\label{prop:monotone}
Let $\pi$ satisfy \textup{(C0)--(C6)}, let $c_*>0$, and let $k_0>0$ be a real number with
$y_0+1\le(1-\eps)\big(\frac{k_0}2-3\big)$. Let $\hat B$ be either $B$ or $\bar B$, and $\hat\Psi_\pi$ the corresponding
function \eqref{eq:P8-Psi}. Put
\[
D_1:=\frac{h_0(1-\omega_0)(1-\eps)^{1/4}}{\hat B}-\frac{c_*}{4\cdot2^{3/4}} .
\]
Suppose that
\begin{enumerate}[label=\textup{(\roman*)}]
\item $\hat\Psi_\pi(k_0)\ge c_*k_0^{1/4}$;
\item $D_1>0$;
\item $y_1(k_0)^{1/4}D_1\ge K_W/(2\hat B)$.
\end{enumerate}
Then $\hat\Psi_\pi(k)\ge c_*k^{1/4}$ for every real $k\ge k_0$.
\end{proposition}

\begin{proof}
For $k\ge k_0$ we have $y_1(k)\ge y_1(k_0)>y_0>0$, so $\hat\Psi_\pi$ is differentiable on $[k_0,\infty)$. Let
$\phi(k):=\hat\Psi_\pi(k)-c_*k^{1/4}$ and write $y_1=y_1(k)$, so that $dy_1/dk=\frac12$ and $k=2y_1+6$. Then
\[
\phi'(k)=\frac{h_0(1-\omega_0)(1-\eps)^{1/4}}{\hat B}\,y_1^{-3/4}-\frac{K_W}{2\hat B\,y_1}-\frac{c_*}{4}\,(2y_1+6)^{-3/4}.
\]
Since $(2y_1+6)^{-3/4}<(2y_1)^{-3/4}=2^{-3/4}y_1^{-3/4}$,
\[
\phi'(k)\ >\ y_1^{-3/4}D_1-\frac{K_W}{2\hat B\,y_1}=y_1^{-1}\Big(y_1^{1/4}D_1-\frac{K_W}{2\hat B}\Big)
\ \ge\ y_1^{-1}\Big(y_1(k_0)^{1/4}D_1-\frac{K_W}{2\hat B}\Big)\ \ge\ 0,
\]
where we used $D_1>0$ and $y_1\ge y_1(k_0)$ in the second inequality and (iii) in the last. As $\phi(k_0)\ge0$ by (i),
the mean value theorem gives $\phi(k)\ge0$ for all $k\ge k_0$.
\end{proof}

\begin{remark}\label{rem:assembly-C7}
If (C7) holds for $k_0$, it holds for every $k\ge k_0$ (Lemma~\ref{lem:flows-constraints}(e)). Hence, if $\pi$
satisfies (C0)--(C6), (C7) holds for $k_0$ and (i)--(iii) hold with $\hat B=\bar B$, then by
Lemma~\ref{lem:assembly-upper} and Proposition~\ref{prop:monotone} every closed packing of $[0,k]^2$, for every integer
$k\ge\max(k_0,4)$, satisfies $W\ge\bar\Psi_\pi(k)\ge c_*k^{1/4}$. Note that the proof of
Proposition~\ref{prop:monotone} uses (i) only at the end: if (ii) holds and (iii) holds at some $k_1$ in place of
$k_0$, then $\phi$ is increasing on $[k_1,\infty)$.
\end{remark}

\subsection{The explicit part of Theorem~\ref{thm:A}}\label{subsec:assembly-explicit}

We prove: \emph{for every integer $k\ge k_0:=4.62\cdot10^{12}$, every closed packing of $[0,k]^2$ satisfies
$W\ge0.1\,k^{1/4}$; consequently $W\ge\lceil0.1\,k^{1/4}\rceil\ge147$.}

Let
\[
\pi^*=(c_0,c_1,y_0,\eps,\omega_0)=\big(0.3078,\ 0.2754,\ 42\,108\,000\,000,\ 5.8\cdot10^{-6},\ 1.7\cdot10^{-5}\big)
\]
as in Section~\ref{sec:flows}. By Lemma~\ref{lem:flows-pistar}(a) and Table~\ref{tab:const-main} (column
``Theorem~\ref{thm:A}''), $\pi^*$ satisfies (C0)--(C6), it satisfies (C7) for $k_0=4.62\cdot10^{12}$, and conditions
(i)--(iii) of Proposition~\ref{prop:monotone} hold with $\hat B=\bar B$ and $c_*=0.1$: indeed
$\bar\Psi_{\pi^*}(k_0)\ge146.62429>146.60896\ge0.1\,k_0^{1/4}$, $D_1\ge0.0086429>0$ and
$y_1(k_0)^{1/4}D_1\ge10.655>0.0044029\ge K_W/(2\bar B)$. By Remark~\ref{rem:assembly-C7}, $W\ge0.1\,k^{1/4}$ for every
integer $k\ge k_0$ and every closed packing of $[0,k]^2$. Since $W$ is an integer, $W\ge\lceil0.1k^{1/4}\rceil$, and
$\lceil0.1k^{1/4}\rceil\ge\lceil0.1k_0^{1/4}\rceil=147$ because $146<0.1\,k_0^{1/4}\le146.60896$.\qed

\subsection{The asymptotic constant}\label{subsec:assembly-asym}

Let $\lambda:=2\delta+2\cdot10^{-12}$ (as in Theorem~\ref{thm:A}), and put
\[
h_0^\infty(c):=1-\lambda-1.00002\,c^2,\qquad c_0^{\rm opt}:=\sqrt{\frac{1-\lambda}{5.0001}},
\]
\begin{equation}\label{eq:assembly-cinf}
c_\infty:=\frac{4\cdot2^{-1/4}\,h_0^\infty(c_0^{\rm opt})\,\sqrt{c_0^{\rm opt}}}{\sqrt{AA'}}
=\frac{16\cdot2^{-1/4}}{5\sqrt{AA'}}\,(1-\lambda)^{5/4}\,5.0001^{-1/4}.
\end{equation}
The right-hand side is the constant $c_\infty$ of Theorem~\ref{thm:A}. Here $A,A'$ are the exact constants of
Definition~\ref{def:flows-consts}; numerically $c_\infty=0.1696524112\ldots$ (Table~\ref{tab:const-fixed}). For $\pi$
satisfying (C0)--(C6) let
\[
\Psi_\infty(\pi):=\lim_{k\to\infty}\frac{\Psi_\pi(k)}{k^{1/4}}=\frac{8\cdot2^{-1/4}\,h_0(1-\omega_0)(1-\eps)^{1/4}}{B};
\]
the limit exists and has this value because $((1-\eps)(\frac k2-3))^{1/4}/k^{1/4}\to2^{-1/4}(1-\eps)^{1/4}$, while
$(y_0+1)^{1/4}/k^{1/4}\to0$ and $\log(y_1/y_0)/k^{1/4}\to0$.

\begin{lemma}\label{lem:assembly-c0}
On $(0,\infty)$ the function $c\mapsto h_0^\infty(c)\sqrt c$ attains its maximum exactly at $c=c_0^{\rm opt}$, and
$h_0^\infty(c_0^{\rm opt})=\frac45(1-\lambda)$. In particular $4\cdot2^{-1/4}h_0^\infty(c)\sqrt c/\sqrt{AA'}\le c_\infty$
for every $c>0$.
\end{lemma}

\begin{proof}
The derivative of $(1-\lambda-1.00002\,c^2)\sqrt c$ is $(1-\lambda-5.0001\,c^2)/(2\sqrt c)$, which is positive for
$c<c_0^{\rm opt}$ and negative for $c>c_0^{\rm opt}$. At $c=c_0^{\rm opt}$ we have
$1.00002\,c^2=(1-\lambda)/5$, so $h_0^\infty=\frac45(1-\lambda)$, and the two expressions in \eqref{eq:assembly-cinf}
agree.
\end{proof}

\begin{proposition}[the asymptotic part of Theorem~\ref{thm:A}]\label{prop:assembly-liminf}
$\displaystyle\liminf_{k\to\infty}\frac{W_{\min}(k)}{k^{1/4}}\ \ge\ c_\infty$.
\end{proposition}

\begin{proof}
Let $\pi$ satisfy (C0)--(C6). Then (C7) holds for all sufficiently large $k$, since $(1-\eps)(\frac k2-3)\to\infty$ as
$k\to\infty$ (and once it holds, it holds for all larger $k$, Remark~\ref{rem:assembly-C7}), and
Theorem~\ref{thm:star} gives $W_{\min}(k)\ge\Psi_\pi(k)$ for those $k$. Hence
$\liminf_kW_{\min}(k)/k^{1/4}\ge\Psi_\infty(\pi)$, and it suffices to find parameters with $\Psi_\infty(\pi)$
arbitrarily close to $c_\infty$.

For an integer $n\ge2$ let $\pi_n:=(c_0^{\rm opt},\,c_1^{\rm opt},\,n,\,n^{-1/2},\,n^{-1/4})$ with
$c_1^{\rm opt}:=\sqrt{A\,c_0^{\rm opt}/(4A')}$. As $n\to\infty$: $\alpha(n)=c_0^{\rm opt}n^{-1/2}\to0$;
$\hat\beta(n)=c_1^{\rm opt}n^{-3/4}\to0$ and $b^*=c_1^{\rm opt}((1-n^{-1/2})n)^{-3/4}\to0$; the left-hand side of (C3)
tends to $c_1^{\rm opt}$ while its right-hand side $c_0^{\rm opt}n^{1/4}$ tends to $\infty$; $n^{-1/4}<\frac12$ for
$n>16$; and $w_0\to0.50001(c_0^{\rm opt})^2+\delta+10^{-12}<0.10002<\frac12$, because
$(c_0^{\rm opt})^2=(1-\lambda)/5.0001<0.2$. Hence $\pi_n$ satisfies (C0)--(C6) for all large $n$ ((C6) follows from
(C1), Lemma~\ref{lem:flows-constraints}(a)).

Along this family $h_0\to1-1.00002\,(c_0^{\rm opt})^2-2\delta-2\cdot10^{-12}=h_0^\infty(c_0^{\rm opt})$ and
$(1-\omega_0)(1-\eps)^{1/4}\to1$. Since $\eps=n^{-1/2}\to0$ and
$(1-\eps)^{3/4}\big(1-(1-\eps)^{3/4}\big)/\eps\to\frac34$, while $b_W/(\eps n)=b_Wn^{-1/2}\to0$, we have $\invF\to1$.
Further $1/(\omega_0n^{3/4})=n^{-1/2}\to0$, $\alpha(n)C_\Lam/\delta\to0$ and $\Gamma\to0$. Therefore
\[
B\ \to\ \frac{A}{2c_1^{\rm opt}}+\frac{2c_1^{\rm opt}A'}{c_0^{\rm opt}}=2\sqrt{\frac{AA'}{c_0^{\rm opt}}},\qquad
\Psi_\infty(\pi_n)\ \to\ \frac{8\cdot2^{-1/4}h_0^\infty(c_0^{\rm opt})}{2\sqrt{AA'/c_0^{\rm opt}}}=c_\infty .
\]
This proves the proposition.
\end{proof}

The value $c_\infty$ is the best constant that the bounds \eqref{eq:P8-Psi} can give, if the fixed constants
$\delta$, $\bbar$, $b_W$, $C_\Lam$ and the numerical constants in $w_0$, $\Gamma$ and \eqref{eq:P8-Om} are kept as they
are; it is a supremum and is not attained.

\begin{proposition}\label{prop:assembly-sup}
Let $\Psi_\pi$, $\Psi_\infty$ and $c_\infty$ be computed with the same constants $A$, $A'$.
\begin{enumerate}[label=\textup{(\alph*)}]
\item If $\pi$ satisfies \textup{(C0)--(C7)} for $k$, then $\Psi_\pi(k)<c_\infty k^{1/4}$.
\item If $\pi$ satisfies \textup{(C0)--(C6)}, then $\Psi_\infty(\pi)<c_\infty$; and
  $\sup_\pi\Psi_\infty(\pi)=c_\infty$, the supremum being over all $\pi$ satisfying \textup{(C0)--(C6)}.
\end{enumerate}
The same holds with $\bar A,\bar A'$ (and $\bar B$) in place of $A,A'$ (and $B$), $c_\infty$ being replaced by the value
$c_\infty^+$ of \eqref{eq:assembly-cinf} computed with $\bar A,\bar A'$.
\end{proposition}

\begin{proof}
First, $h_0<h_0^\infty(c_0)$, because $w_0>0.50001\,c_0^2+\delta+10^{-12}$ (the terms $0.50001c_0^2\cdot1.0001/y_0$ and
$1.0001\,b^*$ are positive); and $h_0>0$ by (C5), so $h_0^\infty(c_0)>0$. Second, all four terms of $B$ are positive and
$0<\invF<1$ (Lemma~\ref{lem:P8-invF}), so by the inequality between arithmetic and geometric means
\[
B\ >\ \frac{A}{2c_1}+\frac{2c_1A'}{c_0\,\invF}\ \ge\ \frac{A}{2c_1}+\frac{2c_1A'}{c_0}\ \ge\ 2\sqrt{\frac{AA'}{c_0}} .
\]
(a) By (C7), $y_1>y_0$, so $\log(y_1/y_0)+2/y_0>0$ and the subtracted term in \eqref{eq:P8-Psi} is positive; also
$(y_0+1)^{1/4}>0$, $0<1-\omega_0<1$, $(1-\eps)^{1/4}<1$ and $y_1<k/2$. Hence, as $h_0>0$,
\[
\Psi_\pi(k)\ <\ \frac{8h_0\,(k/2)^{1/4}}{B}\ <\ \frac{8\cdot2^{-1/4}h_0^\infty(c_0)\,k^{1/4}}{2\sqrt{AA'/c_0}}
=\frac{4\cdot2^{-1/4}h_0^\infty(c_0)\sqrt{c_0}}{\sqrt{AA'}}\,k^{1/4}\ \le\ c_\infty k^{1/4}
\]
by Lemma~\ref{lem:assembly-c0}.
(b) In the same way
\[
\Psi_\infty(\pi)=\frac{8\cdot2^{-1/4}h_0(1-\omega_0)(1-\eps)^{1/4}}{B}\ <\ \frac{4\cdot2^{-1/4}h_0^\infty(c_0)\sqrt{c_0}}{\sqrt{AA'}}
\ \le\ c_\infty .
\]
The family $\pi_n$ in the proof of Proposition~\ref{prop:assembly-liminf} gives $\Psi_\infty(\pi_n)\to c_\infty$, so the
supremum is $c_\infty$, and by the strict inequality it is not attained.
The last sentence follows by the same argument, since only $A,A'>0$ and the form of $B$ were used.
\end{proof}

Theorem~\ref{thm:A} consists of the explicit statement of Section~\ref{subsec:assembly-explicit} and
Proposition~\ref{prop:assembly-liminf}; this completes its proof.

\subsection{\texorpdfstring{Theorem~\ref{thm:A13}}{Theorem A13}: the variant without the computer-assisted lemma}\label{subsec:assembly-A13}

By Remark~\ref{rem:imp-dep}, the computer-assisted [R, Lemma~\ref{imp:L4.10}] enters the proof of
Theorem~\ref{thm:star} only through the multiplicity bound $K_w^*=9$, namely through the constant $A$ of
[R, Lemmas~\ref{imp:L4.13}--\ref{imp:L4.15}] and, through Proposition~\ref{imp:pairs}, through the constants $C_\Lam$
of Corollary~\ref{cor:Lambda0} (Sections~\ref{sec:P3} and~\ref{sec:P5}) and $A'$ of Corollary~\ref{cor:P6scale}. With
the analytic [R, Lemma~\ref{imp:L4.9}] ($K_w=13$) in all these places, [R, Lemmas~\ref{imp:L4.13}--\ref{imp:L4.15}]
hold with $A$ replaced by $A_{13}$ (Remark~\ref{rem:imp-dep}, [R, Remark~7.5]), Corollary~\ref{cor:Lambda0} holds with
$C_\Lam^{(13)}=1+1712.6\,\delta^2$, and Corollary~\ref{cor:P6scale} holds with $A'_{13}$
(Definition~\ref{def:flows-consts} for $A_{13},A'_{13},C^{(13)}_\Lam$). Consequently Theorem~\ref{thm:star},
Remark~\ref{rem:P8-larger} and Sections~\ref{subsec:assembly-upper}--\ref{subsec:assembly-asym} hold verbatim with
$A_{13},A'_{13},C_\Lam^{(13)}$ in place of $A,A',C_\Lam$, and with the upper bounds
\[
\bar A_{13}:=12.7476,\qquad \bar A'_{13}:=12.7476,\qquad \bar C^{(13)}_\Lam:=1.0000002
\]
(Table~\ref{tab:const-fixed}) in place of $\bar A,\bar A',\bar C_\Lam$. Nothing else in the proof depends on
[R, Lemma~\ref{imp:L4.10}].

\begin{proof}[Proof of Theorem~\ref{thm:A13}]
\emph{Explicit part.} Let $k_0:=2.17\cdot10^{13}$ and
\[
\pi^*_{13}=(c_0,c_1,y_0,\eps,\omega_0)=\big(0.3674,\ 0.3012,\ 60\,000\,000\,000,\ 5\cdot10^{-6},\ 1.4\cdot10^{-5}\big)
\]
as in Section~\ref{sec:flows}. By Lemma~\ref{lem:flows-pistar}(b) and Table~\ref{tab:const-main} (column
``Theorem~\ref{thm:A13}''), $\pi^*_{13}$ satisfies (C0)--(C6), (C7) for $k_0$, and (i)--(iii) of
Proposition~\ref{prop:monotone} with $c_*=0.1$ and $\hat B=\bar B$ computed from
$\bar A_{13},\bar A'_{13},\bar C^{(13)}_\Lam$: $\bar\Psi_{\pi^*_{13}}(k_0)\ge215.83530>215.83156\ge0.1\,k_0^{1/4}$,
$D_1\ge0.0055790>0$ and $y_1(k_0)^{1/4}D_1\ge10.125>0.0052311\ge K_W/(2\bar B)$. As in
Section~\ref{subsec:assembly-explicit}, every closed packing of $[0,k]^2$ with an integer $k\ge k_0$ satisfies
$W\ge0.1\,k^{1/4}$, hence $W\ge\lceil0.1k^{1/4}\rceil\ge\lceil0.1k_0^{1/4}\rceil=216$, because
$215<0.1\,k_0^{1/4}\le215.83156$.

\emph{Asymptotic part.} Propositions~\ref{prop:assembly-liminf} and~\ref{prop:assembly-sup} with $A_{13},A'_{13}$ give
$\liminf_kW_{\min}(k)/k^{1/4}\ge c_\infty^{(13)}$, where $c^{(13)}_\infty$ is \eqref{eq:assembly-cinf} with
$A_{13},A'_{13}$, as in Theorem~\ref{thm:A13}; numerically $c^{(13)}_\infty=0.1411593387\ldots$
(Table~\ref{tab:const-fixed}).
\end{proof}

\subsection{Corollary~\ref{cor:A}}\label{subsec:assembly-cor}

We prove: \emph{for all integers $c\ge0$ and $k\ge\max(4.62\cdot10^{12},\,10^4c^4+1)$ one has $s(k^2-c)=k$.
Moreover $c^*(k)=W_{\min}(k)-1\ge\lceil0.1k^{1/4}\rceil-1$ for every integer $k\ge4.62\cdot10^{12}$; in particular
$c^*(k)\ge146$ there, and for $0\le c\le146$ the condition on $k$ is just $k\ge4.62\cdot10^{12}$. Without
[R, Lemma~\ref{imp:L4.10}], the same holds with $2.17\cdot10^{13}$ in place of $4.62\cdot10^{12}$, with $215$ in place
of $146$.} Here $s(n)$ and $c^*(k)$ are as in Section~\ref{sec:intro}.

\begin{proof}
By [R, Lemma~\ref{imp:L2.1}], for integers $k\ge2$ and $n\ge1$ we have $s(n)<k$ if and only if $n\le M(k)$, and
$c^*(k)=k^2-M(k)-1=W_{\min}(k)-1$. Moreover $s(n)\le k$ whenever $n\le k^2$: the squares $[i,i+1]\times[j,j+1]$
($0\le i,j<k$) have pairwise disjoint interiors, and any $n\le k^2$ of them lie in $[0,k]^2$.

Let $k\ge4.62\cdot10^{12}$ and $c\ge0$ be integers with $k\ge10^4c^4+1$. Then $k>10^4c^4=(c/0.1)^4$, so
$c<0.1\,k^{1/4}$. By Theorem~\ref{thm:A}, $W_{\min}(k)\ge0.1\,k^{1/4}>c$. Hence $n:=k^2-c$ satisfies
$n>k^2-W_{\min}(k)=M(k)$, and $n\ge1$ because $c<0.1\,k^{1/4}<k^2-1$. By [R, Lemma~\ref{imp:L2.1}], $s(n)\ge k$; and
$s(n)\le k$ since $n\le k^2$. So $s(k^2-c)=k$.

Since $W_{\min}(k)$ is an integer and $W_{\min}(k)\ge0.1k^{1/4}$, we get $W_{\min}(k)\ge\lceil0.1k^{1/4}\rceil$, so
$c^*(k)\ge\lceil0.1k^{1/4}\rceil-1\ge\lceil0.1(4.62\cdot10^{12})^{1/4}\rceil-1=146$
(Section~\ref{subsec:assembly-explicit}). Finally, $10^4c^4+1\le10^4\cdot146^4+1=4\,543\,718\,560\,001<4.62\cdot10^{12}$ for
$0\le c\le146$.

Without [R, Lemma~\ref{imp:L4.10}] the same argument applies with Theorem~\ref{thm:A13} in place of
Theorem~\ref{thm:A}: $\lceil0.1(2.17\cdot10^{13})^{1/4}\rceil-1=215$, and
$10^4\cdot215^4+1=21\,367\,506\,250\,001<2.17\cdot10^{13}$.
\end{proof}

\begin{remark}\label{rem:assembly-cor}
The condition $k\ge10^4c^4+1$ is equivalent to $c<0.1\,k^{1/4}$ for integers $c\ge0$ and $k\ge1$, because
$10^4c^4$ is an integer. For $c=147$ it reads $k\ge10^4\cdot147^4+1=4\,669\,488\,810\,001$, which exceeds
$4.62\cdot10^{12}$; so $146$ is the largest $c$ for which the threshold $4.62\cdot10^{12}$ alone suffices (and $215$ in
the variant, since $10^4\cdot216^4+1=21\,767\,823\,360\,001>2.17\cdot10^{13}$). For large $c$ the further thresholds of
Section~\ref{subsec:const-pairs} give better conditions.
\end{remark}

\section{Constants and thresholds}\label{sec:constants}

This section lists the numerical values used in Sections~\ref{sec:P8} and~\ref{sec:assembly} and in the statements
of Section~\ref{sec:intro}, with the direction in which every displayed number is rounded. All of them are
recomputed by the script described in Section~\ref{subsec:const-script}. The parameter values are fixed numbers; the
proofs use only the inequalities displayed below and do not depend on how the values were chosen. We do not claim that the thresholds $4.62\cdot10^{12}$ and $2.17\cdot10^{13}$ are the least ones
that \eqref{eq:P8-Psi} can give (for fixed parameters, see Section~\ref{subsec:const-sens}), and we do not prove bounds of
the form $W\ge c_*k^{1/4}$ for all $k$ below them; Table~\ref{tab:const-k} gives bounds at single values of $k$.

\subsection{Conventions}\label{subsec:const-conv}

\begin{itemize}
\item An equality sign between a quantity and a decimal number ($\eps y_0=244\,226.4$) means exact equality.
\item A decimal number followed by dots, $x=d\ldots$, is the truncation of $x$: $d\le x<d+u$, where $u$ is one unit
  of the last displayed digit of $d$ (all such quantities are positive).
\item An upper bound $x\le v$ is the value of $x$ rounded up at the last displayed digit ($v-u<x\le v$); a lower bound
  $x\ge v$ is the value rounded down ($v\le x<v+u$).
\item $x<v$ (as in $A<10.6067$) denotes a chosen bound which is not necessarily obtained by rounding at the last
  digit.
\end{itemize}
In every inequality of a proof the displayed values are used in the safe direction: quantities that are subtracted or
appear in a denominator are bounded from above, and the others from below. For example, the lower bound
$\bar\Psi_{\pi^*}(k_0)\ge146.62429$ and the upper bound $0.1\,k_0^{1/4}\le146.60896$ together give condition (i) of
Proposition~\ref{prop:monotone}.

\subsection{Fixed constants}\label{subsec:const-fixed}

The fixed constants are
\begin{gather*}
\delta=10^{-5},\qquad \bbar=1.5\cdot10^{-6},\qquad b_W=2.0002,\qquad Q_*=0.32,\\
\lambda=2\delta+2\cdot10^{-12}=2.0000002\cdot10^{-5},\qquad
c_0^{\rm opt}=\sqrt{\frac{1-\lambda}{5.0001}}=0.4472046513\ldots,
\end{gather*}
and the constants depending on the multiplicity bound $K$ ($K=K_w^*=9$ with [R, Lemma~\ref{imp:L4.10}], $K=K_w=13$ with
[R, Lemma~\ref{imp:L4.9}] only) are given in Table~\ref{tab:const-fixed}, where
\[
A=\frac{4\sqrt{K/Q_*}}{2-3\cdot10^{-6}},\qquad A'=\frac{A}{\cos\bbar}+10^{-8},\qquad
C_\Lam=\begin{cases}1+1185.6\,\delta^2&(K=9),\\ 1+1712.6\,\delta^2&(K=13).\end{cases}
\]
The bounds $\bar A,\bar A',\bar C_\Lam$ (and their analogues for $K=13$) are the numbers used in all explicit
computations (Sections~\ref{subsec:assembly-upper} and~\ref{subsec:assembly-A13}).

\begin{table}[ht]
\centering
\caption{Constants depending on the multiplicity bound $K$; $c_\infty^+$ denotes $c_\infty$ computed with the upper
bounds $\bar A,\bar A'$ (resp.\ $\bar A_{13},\bar A'_{13}$).}\label{tab:const-fixed}
\small
\begin{tabular}{|l|l|l|}
\hline
 & $K=9$ (with [R, Lemma~\ref{imp:L4.10}]) & $K=13$ ([R, Lemma~\ref{imp:L4.9}] only)\\
\hline
$A$ & $10.6066176277\ldots<10.6067=\bar A$ & $12.7475679053\ldots<12.7476=\bar A_{13}$\\
$A'$ & $10.6066176377\ldots<10.6068=\bar A'$ & $12.7475679153\ldots<12.7476=\bar A'_{13}$\\
$C_\Lam$ & $=1.00000011856$ & $=1.00000017126$\\
$\bar C_\Lam$, resp.\ $\bar C^{(13)}_\Lam$ & $1.0000002$ & $1.0000002$\\
$c_\infty$ (exact $A,A'$) & $0.1696524112\ldots$ & $0.1411593387\ldots$\\
$c_\infty^+$ & $0.1696502940\ldots$ & $0.1411589833\ldots$\\
\hline
\end{tabular}
\end{table}

The value $C_\Lam=1+1185.6\,\delta^2$ comes from the bound
$56.2501+318.20+56.2501+754.8\le1185.6$ for the coefficients of $\delta^2$ in Sections~\ref{sec:P3}
and~\ref{sec:P5}; for $K=13$ the corresponding coefficients give $81.26+459.7+81.26+1090.3\le1712.6$. The constant
$A'_{13}$ also bounds the expression $4\sqrt{13/Q_*}/\big((2-\bbar)\cos\bbar\big)+10^{-8}$, because
$2-\bbar\ge2-3\cdot10^{-6}$.

\subsection{The parameters of \texorpdfstring{Theorems~\ref{thm:A} and~\ref{thm:A13}}{Theorems A and A13}}\label{subsec:const-main}

Table~\ref{tab:const-main} lists the parameters $\pi^*$ (Theorem~\ref{thm:A}) and $\pi^*_{13}$ (Theorem~\ref{thm:A13}),
the values that establish the constraints (C1)--(C7), the constituents of $\bar B$, and the quantities in
conditions (i)--(iii) of Proposition~\ref{prop:monotone} with $c_*=0.1$. The column for Theorem~\ref{thm:A} uses
$\bar A,\bar A',\bar C_\Lam$, the column for Theorem~\ref{thm:A13} uses $\bar A_{13},\bar A'_{13},\bar C^{(13)}_\Lam$.
Constraint (C4) holds since $\omega_0<\frac12$; (C6) and (Q1) follow from (C1) (Lemma~\ref{lem:flows-constraints}(a));
(C7) holds because $y_0+1$ is smaller than the displayed lower bound for $(1-\eps)(k_0/2-3)$. The constraints (C3) and
(Q2), which are not used, hold as well (see also Lemma~\ref{lem:flows-pistar}).

\begin{table}[ht]
\centering
\caption{The parameters and the values at $k_0$ for Theorems~\ref{thm:A} and~\ref{thm:A13}.}\label{tab:const-main}
\small
\begin{tabular}{|l|l|l|}
\hline
 & Theorem~\ref{thm:A} & Theorem~\ref{thm:A13}\\
\hline
$c_0$ & $0.3078$ & $0.3674$\\
$c_1$ & $0.2754$ & $0.3012$\\
$y_0$ & $42\,108\,000\,000$ & $60\,000\,000\,000$\\
$\eps$ & $5.8\cdot10^{-6}$ & $5\cdot10^{-6}$\\
$\omega_0$ & $1.7\cdot10^{-5}$ & $1.4\cdot10^{-5}$\\
$k_0$ & $4.62\cdot10^{12}$ & $2.17\cdot10^{13}$\\
\hline
(C1): $\alpha(y_0)$ & $\le1.4999830\cdot10^{-6}$ & $\le1.4999043\cdot10^{-6}$\\
(C2): $\hat\beta(y_0)$ & $\le2.9628\cdot10^{-9}$ & $\le2.4846\cdot10^{-9}$\\
(C3): $c_1(1-\eps)^{-3/4}$ & $\le0.27541$ & $\le0.30121$\\
(C3): $c_0y_0^{1/4}$ & $\ge139.43$ & $\ge181.83$\\
(C5): $w_0$ & $\le0.04738138$ & $\le0.06750274$\\
(C7): $(1-\eps)(k_0/2-3)$ & $\ge2.3099866\cdot10^{12}$ & $\ge1.0849945\cdot10^{13}$\\
$b^*$ & $\le2.9628\cdot10^{-9}$ & $\le2.4846\cdot10^{-9}$\\
$h_0$ & $\ge0.9052372592$ & $\ge0.8649945353$\\
$\eps y_0$ & $=244\,226.4$ & $=300\,000$\\
$\invF$ & $\ge0.99998818$ & $\ge0.99999020$\\
\hline
$1/(\omega_0y_0^{3/4})$ & $\le0.00063282$ & $\le0.00058920$\\
$\bar A/(2c_1)$ & $\le19.256900$ & $\le21.161355$\\
$\frac{2c_1}{c_0\,\invF}\big(\bar A'+\alpha(y_0)\bar C_\Lam/\delta\big)$ & $\le19.249235$ & $\le21.147482$\\
$\Gamma$ & $\le2.8472\cdot10^{-5}$ & $\le3.1351\cdot10^{-5}$\\
$\bar B$ & $\le38.506796$ & $\le42.309457$\\
$C_{\rm wall}=4c_1/(c_0\,\invF)$ & $\le3.5789897$ & $\le3.2792918$\\
$K_W=1.000001\,c_0^2\,C_{\rm wall}$ & $\le0.3390769$ & $\le0.4426483$\\
\hline
bound \eqref{eq:P8-Om} for $\Omega_W$ at $k_0$ & $\le0.37941519$ & $\le0.70158383$\\
$C_{\rm wall}$ times this bound & $\le1.3579231$ & $\le2.3006981$\\
bound \eqref{eq:P8-lH} for $\ell(H)$ at $k_0$ & $\ge5647.4855$ & $\ge9134.3029$\\
$\bar\Psi_\pi(k_0)$ & $\ge146.62429$ & $\ge215.83530$\\
$0.1\,k_0^{1/4}$ & $\le146.60896$ & $\le215.83156$\\
$D_1$ & $\ge0.0086429$ & $\ge0.0055790$\\
$y_1(k_0)^{1/4}D_1$ & $\ge10.655$ & $\ge10.125$\\
$K_W/(2\bar B)$ & $\le0.0044029$ & $\le0.0052311$\\
$\lceil0.1\,k_0^{1/4}\rceil$ & $147$ & $216$\\
\hline
$\Psi_\infty(\pi)$ (computed with $\bar B$) & $0.1581428\ldots$ & $0.1375313\ldots$\\
\hline
\end{tabular}
\end{table}

The table shows directly that conditions (i)--(iii) of Proposition~\ref{prop:monotone} hold:
$146.62429>146.60896$, $0.0086429>0$, $10.655>0.0044029$ for Theorem~\ref{thm:A}, and
$215.83530>215.83156$, $0.0055790>0$, $10.125>0.0052311$ for Theorem~\ref{thm:A13}. For the exact values, the four
terms of $\bar B$ add up to $\bar B$, and $\bar\Psi_\pi(k_0)$ equals $(1-\omega_0)$ times the bound \eqref{eq:P8-lH},
minus $C_{\rm wall}$ times the bound \eqref{eq:P8-Om}, divided by $\bar B$. The displayed values are rounded
separately; in particular the four displayed terms of $\bar B$, each rounded up, add up to slightly more than the
displayed upper bound for $\bar B$.

\subsection{The asymptotic constant; the variant without \texorpdfstring{[R, Lemma~\ref{imp:L4.10}]}{[R, Lemma 4.10]}}\label{subsec:const-asym}

By Propositions~\ref{prop:assembly-liminf} and~\ref{prop:assembly-sup},
\[
c_\infty=\frac{4\cdot2^{-1/4}\,h_0^\infty(c_0^{\rm opt})\sqrt{c_0^{\rm opt}}}{\sqrt{AA'}}
=\frac{16\cdot2^{-1/4}}{5\sqrt{AA'}}\,(1-\lambda)^{5/4}\,5.0001^{-1/4}
\]
is the supremum, over all parameters satisfying (C0)--(C6), of the limits $\Psi_\infty(\pi)$ of $\Psi_\pi(k)/k^{1/4}$;
it is not attained. It depends on the fixed constants only through $A$, $A'$ and $\lambda$ (the numbers $5.0001$ and
$1.00002$ come from the factor $0.50001$ in $w_0$). The constants $C_\Lam$ and $b_W$, and the terms
$1/(\omega_0y_0^{3/4})$ and $\Gamma$, do not affect $c_\infty$, since along the family $\pi_n$ of
Proposition~\ref{prop:assembly-liminf} they enter only through quantities tending to $0$; they affect the thresholds. The
values of $c_\infty$ are in Table~\ref{tab:const-fixed}; with the exact constants, $c_\infty=0.1696524112\ldots$ for
$K=9$ (Theorem~\ref{thm:A}) and $c_\infty^{(13)}=0.1411593387\ldots$ for $K=13$ (Theorem~\ref{thm:A13}). With the
rounded constants $\bar A,\bar A'$ one gets the slightly smaller $c_\infty^+$. The parameters of
Table~\ref{tab:const-main} are tuned to the thresholds $k_0$, not to the asymptotics; their limit ratios
$\Psi_\infty(\pi)$ (last row) are smaller than $c_\infty$, as they must be by Proposition~\ref{prop:assembly-sup}.

Theorem~\ref{thm:A} (both parts) and Corollary~\ref{cor:A} with the threshold $4.62\cdot10^{12}$ depend on the
computer-assisted [R, Lemma~\ref{imp:L4.10}] through the constants $A$, $A'$, $C_\Lam$ (Remark~\ref{rem:imp-dep} and
Section~\ref{subsec:assembly-A13}).
Theorem~\ref{thm:A13} and the corresponding form of Corollary~\ref{cor:A} (threshold $2.17\cdot10^{13}$) use
[R, Lemma~\ref{imp:L4.9}] instead and do not depend on [R, Lemma~\ref{imp:L4.10}].

\subsection{Further thresholds and bounds at given \texorpdfstring{$k$}{k}}\label{subsec:const-pairs}

The proof of Section~\ref{subsec:assembly-explicit} applies verbatim to the pairs $(c_*,k_0)$ and parameters of
Table~\ref{tab:const-pairs}: for each of them $\pi$ satisfies (C0)--(C6) and (C7) for $k_0$, and conditions
(i)--(iii) of Proposition~\ref{prop:monotone} hold (the row $\bar\Psi_\pi(k_0)/k_0^{1/4}$ shows (i) in the form
$\bar\Psi_\pi(k_0)/k_0^{1/4}\ge c_*$). Hence every closed packing of $[0,k]^2$ with an integer $k\ge k_0$ satisfies
$W\ge c_*k^{1/4}$; the last column uses $\bar A_{13},\bar A'_{13},\bar C^{(13)}_\Lam$ and does not depend on
[R, Lemma~\ref{imp:L4.10}].

\begin{table}[ht]
\centering
\caption{Further thresholds: $W\ge c_*k^{1/4}$ for all integers $k\ge k_0$.}\label{tab:const-pairs}
\small
\begin{tabular}{|l|l|l|l|l|l|}
\hline
 & \multicolumn{4}{c|}{$K=9$} & $K=13$\\
\hline
$c_*$ & $0.12$ & $0.14$ & $0.15$ & $0.16$ & $0.12$\\
$k_0$ & $2.17\cdot10^{13}$ & $2.01\cdot10^{14}$ & $1.18\cdot10^{15}$ & $2.77\cdot10^{16}$ & $3.78\cdot10^{14}$\\
\hline
$c_0$ & $0.3674$ & $0.4086$ & $0.4243$ & $0.4376$ & $0.4153$\\
$c_1$ & $0.3009$ & $0.3173$ & $0.3234$ & $0.3284$ & $0.3203$\\
$y_0$ & $6\cdot10^{10}$ & $7.422\cdot10^{10}$ & $8.003\cdot10^{10}$ & $8.5109\cdot10^{10}$ & $7.667\cdot10^{10}$\\
$\eps$ & $5\cdot10^{-6}$ & $4.7\cdot10^{-6}$ & $4.6\cdot10^{-6}$ & $4.5\cdot10^{-6}$ & $4.6\cdot10^{-6}$\\
$\omega_0$ & $1.5\cdot10^{-5}$ & $1.5\cdot10^{-5}$ & $1.4\cdot10^{-5}$ & $1.4\cdot10^{-5}$ & $1.3\cdot10^{-5}$\\
\hline
$\bar\Psi_\pi(k_0)/k_0^{1/4}$ & $\ge0.1200448$ & $\ge0.1400019$ & $\ge0.1500058$ & $\ge0.1600075$ & $\ge0.1200086$\\
$D_1$ & $\ge0.0067036$ & $\ge0.0041135$ & $\ge0.0027026$ & $\ge0.0012500$ & $\ge0.0029558$\\
$y_1(k_0)^{1/4}D_1$ & $\ge12.166$ & $\ge13.024$ & $\ge13.320$ & $\ge13.560$ & $\ge10.959$\\
$K_W/(2\bar B)$ & $\le0.0062733$ & $\le0.0077586$ & $\le0.0083679$ & $\le0.0088999$ & $\le0.0066855$\\
\hline
\end{tabular}
\end{table}

For example, with $(c_*,k_0)=(0.16,\,2.77\cdot10^{16})$ the argument of Section~\ref{subsec:assembly-cor} gives
$s(k^2-c)=k$ for all integers $c\ge0$ and $k\ge2.77\cdot10^{16}$ with $k>(c/0.16)^4$; for $c\ge2065$ the second
condition implies the first, because $(2065/0.16)^4>2.77\cdot10^{16}$ (while $(2064/0.16)^4<2.77\cdot10^{16}$).

At a single value of $k$ no monotonicity is needed: if $\pi$ satisfies (C0)--(C7) for $k$, then
Lemma~\ref{lem:assembly-upper} gives $W\ge\bar\Psi_\pi(k)$, and since $W$ is an integer, $W\ge\lceil\bar\Psi_\pi(k)\rceil$.
Table~\ref{tab:const-k} lists such bounds; each parameter set satisfies (C0)--(C7) for the $k$ of its row.

\begin{table}[ht]
\centering
\caption{Lower bounds for $W$ at given $k$.}\label{tab:const-k}
\small
\begin{tabular}{|c|l|l|l|l|l|l|l|l|l|}
\hline
$K$ & $k$ & $c_0$ & $c_1$ & $y_0$ & $\eps$ & $\omega_0$ & $\bar B\le$ & $\bar\Psi_\pi(k)\ge$ & $W\ge$\\
\hline
9 & $10^{13}$ & $0.3421$ & $0.2903$ & $5.2015\cdot10^{10}$ & $5.3\cdot10^{-6}$ & $1.6\cdot10^{-5}$ & $36.52537$ & $196.77$ & $197$\\
9 & $10^{14}$ & $0.3993$ & $0.3136$ & $7.0863\cdot10^{10}$ & $4.8\cdot10^{-6}$ & $1.5\cdot10^{-5}$ & $33.80809$ & $426.16$ & $427$\\
9 & $10^{16}$ & $0.4345$ & $0.3273$ & $8.393\cdot10^{10}$ & $4.5\cdot10^{-6}$ & $1.4\cdot10^{-5}$ & $32.40967$ & $1575.76$ & $1576$\\
9 & $10^{20}$ & $0.4457$ & $0.3324$ & $2.458\cdot10^{11}$ & $2.7\cdot10^{-6}$ & $9.5\cdot10^{-6}$ & $31.91017$ & $16751.28$ & $16752$\\
9 & $10^{30}$ & $0.4472$ & $0.3343$ & $5.473\cdot10^{14}$ & $5.7\cdot10^{-8}$ & $5.3\cdot10^{-7}$ & $31.72495$ & $5363347.52$ & $5363348$\\
\hline
13 & $10^{13}$ & $0.3421$ & $0.2907$ & $5.2015\cdot10^{10}$ & $5.3\cdot10^{-6}$ & $1.4\cdot10^{-5}$ & $43.84613$ & $163.91$ & $164$\\
13 & $10^{14}$ & $0.3993$ & $0.314$ & $7.0863\cdot10^{10}$ & $4.8\cdot10^{-6}$ & $1.3\cdot10^{-5}$ & $40.58423$ & $355.01$ & $356$\\
13 & $10^{16}$ & $0.4345$ & $0.3276$ & $8.393\cdot10^{10}$ & $4.5\cdot10^{-6}$ & $1.3\cdot10^{-5}$ & $38.90551$ & $1312.66$ & $1313$\\
13 & $10^{20}$ & $0.4458$ & $0.3325$ & $1.928\cdot10^{11}$ & $3\cdot10^{-6}$ & $9.5\cdot10^{-6}$ & $38.33686$ & $13948.51$ & $13949$\\
13 & $10^{30}$ & $0.4472$ & $0.3343$ & $4.287\cdot10^{14}$ & $6.4\cdot10^{-8}$ & $5.3\cdot10^{-7}$ & $38.12805$ & $4462691.45$ & $4462692$\\
\hline
\end{tabular}
\end{table}

\subsection{Sensitivity of the threshold}\label{subsec:const-sens}

Condition (i) of Proposition~\ref{prop:monotone} for Theorem~\ref{thm:A} holds with a small margin. Keep $\pi^*$ and
$k_0=4.62\cdot10^{12}$, and keep the numerator of $\bar\Psi_{\pi^*}(k_0)$ fixed. Then (i) holds exactly when the
bracket does not exceed
\[
B_{\rm crit}:=\frac{(1-\omega_0)\cdot\big(\text{bound \eqref{eq:P8-lH} at }k_0\big)-C_{\rm wall}\cdot\big(\text{bound
\eqref{eq:P8-Om} at }k_0\big)}{0.1\,k_0^{1/4}}\ \ge\ 38.510824,
\]
and the relative margin is $(B_{\rm crit}-\bar B)/\bar B=1.0463\ldots\cdot10^{-4}$, less than $1.05\cdot10^{-4}$.

Consequently Theorem~\ref{thm:A}, with the same $\pi^*$ and $k_0$, remains valid if one of the constants entering
\eqref{eq:P8-star} is replaced by a larger value as follows (all other constants as in Table~\ref{tab:const-fixed}):
\begin{enumerate}[label=(\alph*)]
\item $A'$ by any value up to $10.60905$ (with $A\le10.6067$);
\item $A$ by any value up to $10.60891$ (with $A'\le10.6068$);
\item $A$ and $A'$ both, by values with $A\le10.60781$ and $A'\le A+10^{-4}$; in particular $A=A'\le10.60781$
  suffices;
\item $C_\Lam$ by any value up to $1.0150$;
\item $b_W$ in Theorem~\ref{thm:P7} by any value up to $53.09$;
\item the wall term $C_{\rm wall}\Omega_W$ by up to $1.434$ times the bound obtained from \eqref{eq:P8-Om}.
\end{enumerate}
An additional loss of the form $C'W/\delta$ in Theorem~\ref{thm:P7} acts as an increase of $C_\Lam$ by $C'$. In each
case it suffices to evaluate conditions (i)--(iii) at the extreme value: increasing any of these constants does not
decrease $\bar B$, $K_W$ or the wall term, and increases at least one of them; hence it decreases
$\bar\Psi_{\pi^*}(k_0)$ (whose numerator stays positive), does not increase $D_1$, and decreases
\[
y_1(k_0)^{1/4}D_1-\frac{K_W}{2\bar B}=\frac{y_1(k_0)^{1/4}h_0(1-\omega_0)(1-\eps)^{1/4}-K_W/2}{\bar B}
-\frac{y_1(k_0)^{1/4}c_*}{4\cdot2^{3/4}}
\]
as long as the numerator of the first fraction is positive, which holds at the extreme values. The listed values are
the largest possible ones for this $\pi^*$, rounded down: the largest values for which (i) still holds are
$10.6090514\ldots$ in (a), $10.6089191\ldots$ in (b), $10.6078176\ldots$ for $A$ in (c) with $A'=A+10^{-4}$,
$1.0150101\ldots$ in (d), $53.0942\ldots$ in (e) and $1.43499\ldots$ in (f).

If, on the other hand, $C_\Lam$ were $2$, then for the same $\pi^*$ conditions (ii) and (iii) would still hold at
$k=4.62\cdot10^{12}$, so by the proof of Proposition~\ref{prop:monotone} (Remark~\ref{rem:assembly-C7}) the function
$\bar\Psi_{\pi^*}(k)-0.1\,k^{1/4}$ would be increasing on $[4.62\cdot10^{12},\infty)$; it would be negative at
$4.62\cdot10^{12}$ and vanish at a single point $4.8449\ldots\cdot10^{12}$. Thus condition (i) would fail at
$4.62\cdot10^{12}$, and for this $\pi^*$ the threshold would move to $4.845\cdot10^{12}$ (rounded up), where conditions
(i)--(iii) and (C7) hold. We do not discuss other choices of $\pi$ for this case. The asymptotic constant $c_\infty$
depends only on $A$ and $A'$ (Section~\ref{subsec:const-asym}).

In terms of $k$ the margin is as follows. Keep $\pi^*$ and the bounds $\bar A,\bar A',\bar C_\Lam$. Conditions (ii)
and (iii) of Proposition~\ref{prop:monotone} hold already at the least integer $k$ for which (C7) holds, namely
$k=84\,216\,488\,464$, so by the proof of Proposition~\ref{prop:monotone} (Remark~\ref{rem:assembly-C7}) the function
$\bar\Psi_{\pi^*}(k)-0.1\,k^{1/4}$ is increasing wherever (C7) holds. Hence condition (i) holds exactly for $k\ge k_*$, where
\[
k_*=4.616673\ldots\cdot10^{12},\qquad 4\,616\,673\,963\,026<k_*\le4\,616\,673\,963\,027 ,
\]
and with this $\pi^*$ the argument gives $W\ge0.1\,k^{1/4}$ for every integer $k\ge4\,616\,673\,963\,027$, but not
for any smaller $k$ (this says nothing about $W$ itself for smaller $k$). The threshold $k_0=4.62\cdot10^{12}$ of Theorem~\ref{thm:A} is $k_*$ rounded up to three significant digits:
$k_0/k_*-1=7.20\ldots\cdot10^{-4}$.

For Theorem~\ref{thm:A13} the corresponding relative margin is $(B_{\rm crit}-\bar B)/\bar B=1.7385\ldots\cdot10^{-5}$
(with $\pi^*_{13}$, $k_0=2.17\cdot10^{13}$ and the constants for $K=13$). In terms of $k$: with $\pi^*_{13}$, conditions
(ii) and (iii) hold already at the least integer $k$ for which (C7) holds, $k=120\,000\,600\,012$, and condition (i)
holds exactly for $k\ge k^{(13)}_*$, where
\[
k^{(13)}_*=2.169597\ldots\cdot10^{13},\qquad 21\,695\,977\,343\,082<k^{(13)}_*\le21\,695\,977\,343\,083 ;
\]
the threshold $2.17\cdot10^{13}$ is $k^{(13)}_*$ rounded up to three significant digits, and
$2.17\cdot10^{13}/k^{(13)}_*-1=1.85\ldots\cdot10^{-4}$.

\subsection{The script}\label{subsec:const-script}

The script \texttt{k14\_constants\_check.py}, which accompanies this manuscript, recomputes every numerical constant
displayed in Sections~\ref{sec:intro}--\ref{sec:constants} (in particular the constants in Theorems~\ref{thm:A},
\ref{thm:A13} and Corollary~\ref{cor:A}, the numbers in the comparison with [R] in Section~\ref{sec:intro}, the
numerical constants of Sections~\ref{sec:imported}--\ref{sec:P7}, and all numbers of
Sections~\ref{sec:P8}--\ref{sec:constants}) and checks each displayed value in its stated rounding direction; it
consists of $555$ such checks. It requires Python~3 and the \texttt{mpmath} library and is run as
\begin{center}
\texttt{python k14\_constants\_check.py}
\end{center}
Every quantity is computed with $50$ significant digits and, in addition, enclosed in interval arithmetic with outward
rounding (\texttt{mpmath.iv}, $50$ digits). The endpoints of the enclosures are converted exactly to rational numbers
and compared with the displayed decimals according to the conventions of Section~\ref{subsec:const-conv}. The
constraints (C0)--(C4), (C7), (Q1), (Q2) are evaluated in exact rational arithmetic (after raising both sides to a
suitable power), displayed values that are rational functions of displayed decimals (such as $S_9$ in
Section~\ref{sec:P5}) are compared in exact rational arithmetic, and conditions (i)--(iii) of
Proposition~\ref{prop:monotone} and all other inequalities in interval arithmetic. It prints one line for each check
and a summary line at the end.

\section*{Verification status}
This paper was developed with AI assistance (see the section on the use of AI). Each part of the proof
(Sections~\ref{sec:P12}--\ref{sec:P8}) was first written as a separate draft, and each draft received one AI review and one
AI numerical test, made independently in separate sessions with their own code; the assembled argument received a review
and an end-to-end numerical test. The reviews judged the drafts correct, with minor and wording items only, and no test
found a counterexample; the resulting changes concern statements and proof details (for instance, a hypothesis that was
used but not stated is now constraint~(C6)) and are included here. The present text was written from the drafts, and the
script of Section~\ref{subsec:const-script} recomputes every displayed constant. A blank-slate AI review of the present
text, divided into twelve parts reviewed separately and accompanied by numerical searches for counterexamples, found no
critical or major issue and eighteen minor ones (wording, the use of the proofs in [R], notation, and two constants
displayed with a last digit larger than necessary); all of them have been addressed in this version. That review also
compared the statements quoted from [R] with [R] (version~1.2), recomputed the constants independently, and recomputed the
finite check in the proof of [R, Lemma~4.9] with a new program in interval arithmetic. All reviewers are AI systems of the
same family and may share blind spots, and the geometric tests use small analogues, not the actual parameter regime. The
proofs use [R] (version~1.2), itself not peer reviewed. Theorem~\ref{thm:A} also uses the computer-assisted
[R, Lemma~4.10], whose box computations were replayed by the Squares Project \cite{LevR} with the same programs but have
not been reproduced by an independent implementation; Theorem~\ref{thm:A13} does not use it. The programs, the numerical
tests and a summary of the reviews are available at \url{https://github.com/squarepacker/k2-minus-c-power}. No human has
checked the argument. This paper has not been peer reviewed.

\subsection*{Acknowledgements}
I thank J.~Levy and the Squares Project for an AI review of version~1.1 of [R] and for replaying the computation of
[R, Lemma~4.10], on which Theorem~\ref{thm:A} depends \cite{LevR}.

\subsection*{Use of AI}
Claude (Anthropic) was used extensively in developing the arguments, writing the text and the programs (including the
numerical tests described above), and in checking them. It is not an author; the author takes full responsibility for the
contents. The paper has not been refereed.

\subsection*{Licence}
This text is distributed under the Creative Commons Attribution 4.0 International licence (CC BY 4.0); the accompanying
programs are distributed under the MIT licence.

\begin{thebibliography}{McC}
\raggedright
\bibitem[AB]{AB} M.~Z.~Arslanov, H.~D.~Bui, Note on ``Efficient packings of unit squares in a large square'',
\emph{Discrete Comput.\ Geom.} \textbf{76} (2026), no.~2, 805--812 (published online 3 August 2025), doi:10.1007/s00454-025-00767-w.
\bibitem[Bui]{Bui} H.~D.~Bui, Square packing with $O(x^{0.6})$ wasted area, arXiv:2508.04603v2 (2025, revised 2026).
\bibitem[CG]{CG} F.~Chung, R.~Graham, Efficient packings of unit squares in a large square, \emph{Discrete Comput.\ Geom.}
\textbf{64} (2020), no.~3, 690--699, doi:10.1007/s00454-019-00088-9.
\bibitem[CG09]{CG09} F.~Chung, R.~Graham, Packing equal squares into a large square, \emph{J.\ Combin.\ Theory Ser.\ A}
\textbf{116} (2009), no.~6, 1167--1175, doi:10.1016/j.jcta.2009.02.005.
\bibitem[Dan]{Dan} E.~Daniel, \emph{square-packing}, GitHub repository, \url{https://github.com/evand/square-packing},
file \texttt{s12/search/FRIEDMAN.md}, commit \texttt{7ff3b21} (accessed 5 October 2026).
\bibitem[EG]{EG} P.~Erd\H{o}s, R.~L.~Graham, On packing squares with equal squares, \emph{J.\ Combin.\ Theory Ser.\ A}
\textbf{19} (1975), no.~1, 119--123, doi:10.1016/0097-3165(75)90099-0.
\bibitem[Fri]{Fri} E.~Friedman, Packing unit squares in squares: a survey and new results, \emph{Electron.\ J.\ Combin.},
Dynamic Survey DS7 (1998, revised 2009), doi:10.37236/28.
\bibitem[LevR]{LevR} J.~Levy, \emph{The Squares Project}, GitHub repository, \url{https://github.com/jlevy/squares}: AI review of
version~1.1 of [R], file \texttt{docs/project/reviews/review-2026-10-05-squarepacker-k2-minus-c.md}, commit \texttt{6953107},
and issue~\#368 (comments of 5 and 6 October 2026, UTC; accessed 7 October 2026).
\bibitem[McC]{McC} R.~McClenagan, Optimally packing a large square by unit squares, arXiv:2602.01484 (2026).
\bibitem[R]{R} S.~Ryu, Packing $k^2-c$ unit squares: $s(k^2-c)=k$ for all large $k$, preprint, version~1.2, 2026,
doi:10.5281/zenodo.23194104; programs and data: doi:10.5281/zenodo.23194031 and
\url{https://github.com/squarepacker/k2-minus-c} (release v1.2).
\bibitem[RV]{RV} K.~F.~Roth, R.~C.~Vaughan, Inefficiency in packing squares with unit squares, \emph{J.\ Combin.\ Theory Ser.\ A}
\textbf{24} (1978), no.~2, 170--186, doi:10.1016/0097-3165(78)90005-5.
\end{thebibliography}
\end{document}
